\documentclass[12pt]{article}

\usepackage[utf8]{inputenc}
\usepackage{CJKutf8}
\usepackage{amsmath,amssymb}
\usepackage{amsthm}
\usepackage{geometry}
\geometry{
  a4paper,
  top=25mm,
  bottom=25mm,
  left=20mm,
  right=20mm
}
\usepackage{xcolor}
\usepackage{url}
\usepackage[unicode,colorlinks=true,linkcolor=blue,urlcolor=blue]{hyperref}
\usepackage{xurl}
\Urlmuskip=0mu plus 2mu
\sloppy
\usepackage{tocloft}

% 目录中 section 条目使用点线填充到页码
\renewcommand{\cftsecleader}{\cftdotfill{\cftdotsep}}

% 基本定理环境（工作笔记：形式系统风格）
\newtheorem{definition}{定义}[section]
\newtheorem{axiom}{公理}[section]
\newtheorem{assumption}{假设}[section]
\newtheorem{theorem}{定理}[section]
\newtheorem{proposition}{命题}[section]
\newtheorem{corollary}{推论}[section]
\newtheorem{lemma}{引理}[section]
\newtheorem{remark}{备注}[section]
\newtheorem{Certificate}{证书}[section]

% 记号（仅用于本笔记自洽引用，避免依赖主稿宏定义）
\newcommand{\code}[1]{\path{#1}}
\newcommand{\GFramework}{\textsf{G-Framework}}
\newcommand{\NN}{\mathbb{N}}
\newcommand{\RR}{\mathbb{R}}

% 术语统一：全篇以 \KouJing 表示“表述约定”（并保证 PDF 书签可读）
\newcommand{\KouJing}{\texorpdfstring{表述约定}{表述约定}}

% 术语统一：本笔记中的“GZ富结构离散元”为离散统第一性基元；“莱布尼茨单子”主要在对照章中作为专名出现（少量前置章节仅作索引预告）
\newcommand{\GZMonFull}{\texorpdfstring{GZ富结构离散元}{GZ富结构离散元}}
\newcommand{\GZMon}{\texorpdfstring{GZ富结构离散元}{GZ富结构离散元}}
\newcommand{\GZMonization}{\texorpdfstring{GZ富结构离散元化}{GZ富结构离散元化}}
\newcommand{\GZMonState}{\texorpdfstring{GZ富结构离散元态}{GZ富结构离散元态}}
\newcommand{\GZMonMotherBundle}{\texorpdfstring{GZ富结构离散元母体主纤维丛}{GZ富结构离散元母体主纤维丛}}

% 避免少量不可控的换行溢出（尤其是混排 CJK 与等宽字体时）
\setlength{\emergencystretch}{6em}

\begin{document}
\begin{CJK*}{UTF8}{gkai}

\title{可修复障碍同伦类与广义分形：全集范畴与严格边界\\（工作笔记：形式系统版）}
\author{Gao Zheng（高政）}
\date{2026-02-21}
\maketitle

\section*{前言}
\addcontentsline{toc}{section}{前言}

本文的目标不是把“可修复障碍同伦类”与“广义分形”并列为两个相似概念，而是在
\GFramework{} 的障碍修复谱系中，严格区分全集范畴、严格子集、递归构造与从属
边界。若二者被混用，广义分形容易被过度提升为所有可修复障碍的总名称；本文的
核心价值正是给出边界：可修复障碍同伦类是更大的对象域，广义分形只是其中满足
特定递归、自相似或尺度结构条件的子类。

本文的第一层立场是对象域的防过度泛化。可修复障碍同伦类关注的是障碍能否通过
同伦、补边、修复塔或边缘算子平凡化被吸收；广义分形关注的是障碍生成是否具有
递归结构、层级自相似或尺度稳定形态。前者是修复范畴，后者是结构形态。二者
可以相交，也可以发生从属关系，但不能被无条件等同。

本文的第二层立场是修复内核的唯一性。障碍可修复并不意味着可任意命名；真正支撑
修复的是边缘算子平凡化、同伦类归并、修复塔终止或超限递归闭合等内核证书。
分形性若出现，只能作为某些障碍生成过程的附加形态，而不能替代修复证书本身。
因此，本文将“广义分形”降级为可修复障碍谱系中的从属描述。

本文的第三层立场是传统分形到广义分形的兼容边界。传统分形依赖较强的自相似、
尺度递归和几何生成条件；广义分形可以放宽这些条件，但不能放宽到吞并全部可修复
障碍。本文因此保留向下兼容：传统分形仍是广义分形的典型子类；同时保持向上边界：
广义分形仍只是修复障碍谱系的一部分。

本文的第四层立场是离散统下的递归生成边界。在离散统中，分形可以通过基元、障碍
驱动、超限递归与尺度嵌套产生；但这些构造仍必须服从修复塔与同伦类边界。本文
因此避免把“递归复杂”误写成“修复充分”，而是把递归形态、同伦归属与修复证书分层
处理。

本文的第五层立场是边缘算子平凡化的优先性。一个障碍是否可修复，首先取决于能否
通过适当边缘算子、补复形或同伦塔将非平凡障碍转化为边界或平凡类。分形结构若
存在，只能描述该障碍如何生成、如何重复、如何呈现层级形态；它不能替代“障碍是否
可被平凡化”的核心判据。

本文的第六层立场是全集范畴与子类范畴的治理意义。若把广义分形误当作全集范畴，
后续修复系统就会错误地要求所有障碍都具有分形性；若把可修复障碍同伦类误缩小为
分形子类，就会漏掉大量非分形但可修复的障碍。本文通过范畴边界校准，保证后续
统计命中和 PDE 残差低能集合可以接收更一般的障碍对象。

本文的第七层立场是为微观隧穿提供障碍分类底座。隧穿、命中和修复需要知道目标
障碍属于何种类型：是普通可修复障碍，还是具有递归分形形态的可修复障碍。本文
给出的分类边界，使后续“障碍穿透”不必假设所有障碍都具有同一种尺度结构。

因此，本文在整体 \GFramework{} 中承担的是“可修复障碍对象域的边界校准”作用。
它向前连接同伦修复与边缘算子，向中间明确广义分形的从属位置，向后为微观隧穿、
统计命中与 PDE 残差低能集合提供一个清晰的障碍分类底座。概括地说，本文把
“复杂递归障碍是否就是分形”这一问题，修正为“何种障碍属于可修复同伦类，何种
可修复障碍额外具有分形结构”的形式系统。

更具体地说，本文的价值在于防止概念膨胀。可修复障碍同伦类是一个由修复可能性
定义的范畴，而广义分形是一个由递归形态和尺度结构定义的子类。二者都重要，但
承担的任务不同：前者决定障碍是否能被修复，后者描述某些障碍如何以分形方式生成。

从方法论上看，本文把“形态描述”与“修复证书”分离。一个障碍可以很复杂但不可修复，
也可以不具备分形形态但完全可修复。只有当分形结构同时满足同伦归属、边缘平凡化
或修复塔吸收条件时，它才进入可修复障碍谱系中的相应子类。

从整体谱系看，本文为后续障碍穿透和低残差命中清理了对象边界。后续系统面对的
障碍并不必然是分形的，也不必然需要分形化处理；它首先需要被判定为可修复或不可
修复。这个边界校准使后续的隧穿、命中和 PDE 残差求解不会被过度泛化的“分形”语言
牵引偏离。

\setcounter{tocdepth}{3}
\setcounter{secnumdepth}{3}
\tableofcontents
\clearpage

%%%%%%%%%%%%%%%%%%%%%%%%%%%%%%%%%%%%%%%%%%%%%%%%%%%%%%%%%%%%
\section{可修复障碍同伦类：全集范畴；广义分形：严格子集}
\label{sec:tex29-ch1}
%%%%%%%%%%%%%%%%%%%%%%%%%%%%%%%%%%%%%%%%%%%%%%%%%%%%%%%%%%%%

\begin{remark}[核心陈述（边界条件）]
“可修复障碍同伦类”（亦称修复同伦类）是覆盖所有拓扑层级、所有障碍类型、所有开放/封闭系统场景的全集范畴；
其范畴覆盖范围等价于：
\[
  \{\text{一切可由边缘算子实现平凡化的非平凡拓扑障碍}\}/\simeq.
\]
而“广义分形”仅是该全集中满足“层间同伦等价递归约束”的一个子类（严格子集），
不等同于修复同伦类的全部。
\end{remark}

\begin{remark}[阅读导航：仓库内文档引用与外部教材参照]
本笔记在“自洽形式系统”\KouJing下推进推演，但为了便于复核与复用，这里登记两类引用参照点：
\begin{enumerate}
  \item \textbf{仓库内 TeX 文档：}更系统的体系化整理与术语表，可参照本仓库卷册主稿与发布稿（如 \cite{GaoZheng2025GFramework,GaoZhengAppVol1}）；

  与“边缘算子闭合/同伦修复塔/变分修复”等直接相关的发布稿链，可参照（如 \cite{RepoNB17EdgeOperator}）；
  \item \textbf{外部标准参照：}链复形/同调/同伦/障碍论/Postnikov 分层与形变收缩等标准概念，可参照代数拓扑教材（如 \cite{Hatcher2002AT,Switzer1975HomotopyHomology,
    Whitehead1978Elements}）；
  雅可比场/共轭点/比安基恒等式等几何侧概念，可参照黎曼几何教材（如 \cite{Lee2018RiemannianManifolds,Petersen2016RiemannianGeometry}）；
  “超限递归”之集合论基础\KouJing，可参照集合论教材（如 \cite{Jech2003SetTheory,Kunen2011SetTheory}）。
\end{enumerate}
以上引用仅用于给出外部参照点：本笔记的“跨分支同构/翻译”仍以工作\KouJing不变，不把外部教材结论当作本系统内的前提。
\end{remark}

\subsection{对象域、谓词与边缘算子：把“修复”写成可检验判定}
\label{subsec:tex29-domain-predicates}

\begin{definition}[系统、障碍与边缘算子（抽象链复形\KouJing）]
令 $X$ 表示系统（可为闭系统或开放系统的状态空间），$G$ 为系数群。
以 $\Omega^n(X;G)$ 表示第 $n$ 层的“障碍表达空间”（可理解为 $n$-上同调类的代表元集合或其等价表达层）。
称一个线性算子
\[
  \partial:\mathcal{D}^{n+1}(X;G)\to \Omega^n(X;G)
\]
为边缘算子（edge operator），其中 $\mathcal{D}^{n+1}(X;G)$ 表示“修复链空间”（用于承载修复证据链）。
若存在 $D\in\mathcal{D}^{n+1}(X;G)$ 使得 $\partial D=\omega$，则称 $\omega$ 被边缘算子平凡化。
\end{definition}

\begin{definition}[修复谓词与广义分形谓词（完全展开）]
\label{def:tex29-rep-fr}
对任意障碍代表元 $\omega\in\Omega^n(X;G)$，定义修复谓词
\[
  \mathsf{Rep}(\omega)\ :\Longleftrightarrow\ \exists D\in\mathcal{D}^{n+1}(X;G)\ \bigl(\partial D=\omega\bigr).
\]

进一步，把“层间同伦等价递归自相似约束”写成一个纯存在量词的可核验数据对象。定义广义分形谓词
$\mathsf{Fr}(\omega)$ 为：
\[
\begin{aligned}
  \mathsf{Fr}(\omega)\ :\Longleftrightarrow\
  &\exists \mathcal{C}\ \exists \{X_k\}_{k\in\NN}\ \exists \{\omega_k\}_{k\in\NN}\\
  &\exists \{D_k\}_{k\in\NN}\ \exists \{(h_k,k_k)\}_{k\in\NN}\ \text{s.t.}\\
  &\text{(F0)}\ \omega_0=\omega;\\
  &\text{(F1)}\ \forall k\in\NN,\ \omega_k\in\Omega^n(X_k;G);\\
  &\text{(F2)}\ \forall k\in\NN,\ D_k\in\mathcal{D}^{n+1}(X_k;G)\ \wedge\ \partial D_k=\omega_k;\\
  &\text{(F3)}\ \forall k\in\NN,\ \bigl(h_k:X_k\to X_{k+1},\ k_k:X_{k+1}\to X_k\bigr)\ \wedge\\
  &\qquad \bigl(h_k\circ k_k\simeq \mathrm{id}_{X_k}\bigr)\ \wedge\ \bigl(k_k\circ h_k\simeq \mathrm{id}_{X_{k+1}}\bigr);
    \\
  &\text{(F4)}\ \forall k\in\NN,\ (X_{k+1},\omega_{k+1},D_{k+1})=\mathcal{C}(X_k,\omega_k,D_k).
\end{aligned}
\]
其中 (F4) 中的 $\mathcal{C}$ 是“递归/自相似生成规则”的形式化占位符：它被视为\emph{定义性约束}的一部分，
而非本章内需要额外证明的结论。
\end{definition}

\begin{definition}[合成证据链、序列边缘算子与时间演化箭头的去歧义]
\label{def:tex29-composed-evidence}
为使“修复链”与“边缘算子 $\partial$”在形式系统中不产生符号歧义，给出如下完全展开。

\textbf{(1) 修复链（时间序列/层级顺序）.}
给定同一系统 $X$（或同伦等价塔的某一固定层）上的障碍序列
\[
  (\omega_0,\omega_1,\dots,\omega_N),\qquad (N\in\NN),\quad \omega_k\in\Omega^n(X;G).
\]
将“时间推进/扰动演化”记为符号 $\leadsto$，即
\[
  \omega_0 \leadsto \omega_1 \leadsto \cdots \leadsto \omega_N,
\]
其中 $\leadsto$ \emph{仅表示顺序}，不再与边缘算子 $\partial$ 混用。

\textbf{(2) 合成证据链（全局证书对象）.}
定义长度 $N$ 的合成证据链空间为
\[
  \widehat{\mathcal{D}}^{n+1}_N(X;G):=\prod_{k=0}^N \mathcal{D}^{n+1}(X;G),
\]
其元素记为
\[
  \widehat{D}:=(D_0,\dots,D_N).
\]

\textbf{(3) 序列边缘算子.}
定义
\[
  \widehat{\partial}_N:\widehat{\mathcal{D}}^{n+1}_N(X;G)\to \bigl(\Omega^n(X;G)\bigr)^{N+1},
  \qquad
  \widehat{\partial}_N(D_0,\dots,D_N):=(\partial D_0,\dots,\partial D_N).
\]

\textbf{(4) “整体修复成立”的精确定义.}
对给定障碍序列 $(\omega_0,\dots,\omega_N)$，称“整体修复成立”当且仅当
\[
  \exists \widehat{D}\in\widehat{\mathcal{D}}^{n+1}_N(X;G)\ \bigl(\widehat{\partial}_N(\widehat{D})=(\omega_0,\dots,
    \omega_N)\bigr),
\]
亦即逐项满足 $\partial D_k=\omega_k$。
\end{definition}

\begin{remark}[规则固定：修复的本质是障碍平凡化]
本章全程使用同一条底层规则：
\[
  \text{“边缘算子是修复核心；障碍是原生驱动力；修复的本质是障碍平凡化”。}
\]
因此后续所有范畴边界与分类，均以 $\mathsf{Rep}(\omega)$ 的真值为唯一入口；其余结构（递归、分形、自相似）
仅作为\emph{附加约束}进入讨论。
\end{remark}

\subsection{形式逻辑：公理降级为定理（数学归纳法证书）}
\label{subsec:tex29-axiom-downgraded}

\begin{theorem}[公理降级：修复闭包与“全集范畴”判定（数学归纳法证书）]
\label{thm:tex29-axiom-downgraded}
给定任意有限修复链（可对应开放系统中的时间序列扰动或多层同伦修复塔的有限截断）
\[
  \omega_0 \leadsto \omega_1 \leadsto \cdots \leadsto \omega_N,
  \qquad (N\in\NN)
\]
其中每一步均满足 $\mathsf{Rep}(\omega_k)$（即存在 $D_k$ 使 $\partial D_k=\omega_k$）。
则：
\begin{enumerate}
  \item （闭包）合成修复仍可证书化：存在合成证据链
  \[
    \widehat{D}:=(D_0,\dots,D_N)\in\widehat{\mathcal{D}}^{n+1}_N(X;G)
  \]
  使得
  \[
    \widehat{\partial}_N(\widehat{D})=(\omega_0,\dots,\omega_N),
  \]
  即“整体修复”成立（见定义~\ref{def:tex29-composed-evidence}）；
  \item （全集判定）“修复同伦类”的归属判定与障碍类型/修复形式无关，仅取决于谓词 $\mathsf{Rep}$：
  \[
    \omega\ \text{属于修复同伦类全集}\ \Longleftrightarrow\ \mathsf{Rep}(\omega).
  \]
\end{enumerate}
\end{theorem}

\begin{Certificate}[数学归纳法证书：按修复步数 $N$ 的闭包与判定稳定性]
\label{cert:tex29-axiom-downgraded}
令谓词 $P(N)$ 表示：
“对任意长度为 $N$ 的修复链 $\omega_0\leadsto\cdots\leadsto\omega_N$，若 $\forall k\le N,\ \mathsf{Rep}(\omega_k)$，
则存在合成证据链
\[
  \widehat{D}:=(D_0,\dots,D_N)\in\widehat{\mathcal{D}}^{n+1}_N(X;G)
\]
使得
\[
  \widehat{\partial}_N(\widehat{D})=(\omega_0,\dots,\omega_N),
\]
并且对链上任意 $\omega_k$，其是否属于修复同伦类全集只由 $\mathsf{Rep}(\omega_k)$ 判定。”
证书化要求为：
\[
  P(0)\ \text{成立},\qquad (\forall N\in\NN)\ \bigl(P(N)\Rightarrow P(N+1)\bigr).
\]
\end{Certificate}

\paragraph{证明（数学归纳法证书；谓词因果链）.}
由证书~\ref{cert:tex29-axiom-downgraded}：
\begin{itemize}
  \item 基步：$N=0$ 时链仅含 $\omega_0$。
  若 $\mathsf{Rep}(\omega_0)$，则存在 $D_0$ 使 $\partial D_0=\omega_0$。
  取合成证据链 $\widehat{D}:=(D_0)\in\widehat{\mathcal{D}}^{n+1}_0(X;G)$，
  则由定义~\ref{def:tex29-composed-evidence} 的序列边缘算子得
  \[
    \widehat{\partial}_0(\widehat{D})=(\partial D_0)=(\omega_0),
  \]
  闭包成立；且“$\omega_0$ 属于全集 $\Leftrightarrow \mathsf{Rep}(\omega_0)$”为定义性判定，故 $P(0)$ 成立；
  \item 归纳步：假设 $P(N)$ 成立。给定长度为 $N+1$ 的链 $\omega_0\leadsto\cdots\leadsto\omega_{N+1}$，
  且 $\forall k\le N+1,\ \mathsf{Rep}(\omega_k)$。
  由归纳假设 $P(N)$，存在
  \[
    \widehat{D}_N:=(D_0,\dots,D_N)\in\widehat{\mathcal{D}}^{n+1}_N(X;G)
    \quad\text{满足}\quad
    \widehat{\partial}_N(\widehat{D}_N)=(\omega_0,\dots,\omega_N),
  \]
  并且链上任意 $\omega_k$ 的全集归属由 $\mathsf{Rep}(\omega_k)$ 决定。
  又因 $\mathsf{Rep}(\omega_{N+1})$，存在 $D_{N+1}$ 使 $\partial D_{N+1}=\omega_{N+1}$。
  取
  \[
    \widehat{D}_{N+1}:=(D_0,\dots,D_N,D_{N+1})\in\widehat{\mathcal{D}}^{n+1}_{N+1}(X;G),
  \]
  则由定义~\ref{def:tex29-composed-evidence} 得
  \[
    \widehat{\partial}_{N+1}(\widehat{D}_{N+1})=(\partial D_0,\dots,\partial D_N,\partial D_{N+1})
    =(\omega_0,\dots,\omega_N,\omega_{N+1}),
  \]
  故闭包对 $N+1$ 成立；同时“归属判定仅依赖 $\mathsf{Rep}$”对新增节点仍成立，从而 $P(N+1)$ 成立。
\end{itemize}
由数学归纳法，$P(N)$ 对任意 $N$ 成立，定理得证。证毕.

\begin{remark}[备注：为何称为“全集范畴”】【剥离修复形式的依赖】]
定理~\ref{thm:tex29-axiom-downgraded} 的第二条表明：
修复同伦类的全集范畴只关心“是否可被边缘算子平凡化”，
而不关心“平凡化是一次性完成、递归完成、跨域映射完成、还是超限层完成”。
因此广义分形最多只能作为在 $\mathsf{Rep}$ 之上再叠加 $\mathsf{Fr}$ 的\emph{额外结构}。
\end{remark}

\subsection{命题：六类核心子类（按障碍拓扑本质与修复逻辑划分）}
\label{subsec:tex29-six-classes}

\begin{definition}[类型谓词与六类修复同伦子类]
\label{def:tex29-six-classes}
为把“覆盖全场景”的分类写成可判定对象，引入类型函数（或类型谓词族）
\[
  \tau(\omega)\in\{1,2,3,4,5,6\}.
\]
在本章\KouJing下，六类核心子类定义为：
\begin{enumerate}
  \item \textbf{0阶局域奇点修复类}（$\tau(\omega)=1$）：障碍为局域、单点、集中式奇点；修复为一次性局域拓扑手术；
  \item \textbf{有限阶全局拓扑修复类}（$\tau(\omega)=2$）：障碍由基本群/整体拓扑型不匹配导致；修复为有限层内的全局同胚式重构；
  \item \textbf{跨域同构修复类}（$\tau(\omega)=3$）：障碍来自异构系统间保结构映射断裂；修复为跨域同伦等价映射的构造；
  \item \textbf{多尺度递归修复类}（$\tau(\omega)=4$）：障碍为分布式、多尺度、层间连续；修复需逐层递归并保持层间同伦等价约束；
  \item \textbf{超限全序数修复类}（$\tau(\omega)=5$）：障碍沿超限序数层展开且层间类型可异构；修复仅需满足层间闭合校验（如 $\partial^2=0$）而不要求同伦等价自相似；
  \item \textbf{静态退化修复类}（$\tau(\omega)=6$）：障碍为静态闭系统中的有限已知缺陷；修复塔冻结为一次性边缘链平凡化。
\end{enumerate}
\end{definition}

\begin{proposition}[六类覆盖的“全集\KouJing”与互斥\KouJing]
\label{prop:tex29-six-classes-cover}
若类型函数 $\tau$ 的定义\KouJing满足：
\[
  \forall \omega,\ \mathsf{Rep}(\omega)\Rightarrow \tau(\omega)\in\{1,\dots,6\},
  \qquad
  \forall \omega,\ \tau(\omega)=i\wedge\tau(\omega)=j\Rightarrow i=j,
\]
则六类子类在“修复同伦类全集”上形成一个互斥覆盖：
\[
  \{\omega:\mathsf{Rep}(\omega)\}=\bigsqcup_{i=1}^{6}\{\omega:\mathsf{Rep}(\omega)\wedge \tau(\omega)=i\}.
\]
\end{proposition}

\begin{Certificate}[证书：类型函数的互斥性与覆盖性]
\label{cert:tex29-six-classes-cover}
证书登记谓词链：
\[
\begin{aligned}
  T_1 &: \mathsf{Rep}(\omega)\Rightarrow \tau(\omega)\in\{1,\dots,6\},\\
  T_2 &: \tau(\omega)=i\wedge\tau(\omega)=j\Rightarrow i=j,\\
  T_3 &: T_1\wedge T_2\Rightarrow \text{按 }\tau\text{ 形成互斥覆盖 } \bigsqcup_{i=1}^{6}.
\end{aligned}
\]
\end{Certificate}

\paragraph{证明（谓词因果链）.}
由 $T_1$，任取 $\omega$ 满足 $\mathsf{Rep}(\omega)$，必有某个 $i\in\{1,\dots,6\}$ 使 $\tau(\omega)=i$，
故 $\omega$ 落入并集右侧；由 $T_2$ 得不同 $i$ 的切片两两不交，从而并集为互斥并。
故等式成立。证毕。

\begin{remark}[备注：广义分形在六类中的位置]
第 $4$ 类（多尺度递归修复类）即是“广义分形”所对应的修复同伦子类；
其余 $1,2,3,6$ 类不具备分形定义所需的“层间递归自相似”，
第 $5$ 类在额外满足同伦等价递归约束时可退化为广义分形特例。
\end{remark}

\subsection{定理：广义分形是修复同伦类中的严格子集}
\label{subsec:tex29-fractal-boundary}

\begin{lemma}[分形约束蕴含可修复性]
\label{lem:tex29-fractal-implies-repairable}
若 $\mathsf{Fr}(\omega)$，则 $\mathsf{Rep}(\omega)$。
\end{lemma}

\begin{Certificate}[证书：$\mathsf{Fr}$ 的定义性约束 $\Rightarrow$ 存在逐层边缘链]
\label{cert:tex29-fractal-implies-repairable}
证书登记谓词链：
\[
\begin{aligned}
  G_1 &: \mathsf{Fr}(\omega)\Rightarrow \text{存在多尺度修复塔及障碍序列 }(\omega_0=\omega,\omega_1,\dots),\\
  G_2 &: \mathsf{Fr}(\omega)\Rightarrow (\forall k)\ \exists D_k\ \bigl(\partial D_k=\omega_k\bigr),\\
  G_3 &: G_1\wedge G_2\Rightarrow \exists D_0\ (\partial D_0=\omega)\Rightarrow \mathsf{Rep}(\omega).
\end{aligned}
\]
\end{Certificate}

\paragraph{证明（谓词因果链）.}
由 $G_1$ 得 $\omega$ 为修复塔第 $0$ 层障碍（$\omega_0=\omega$）；
由 $G_2$ 得存在 $D_0$ 使 $\partial D_0=\omega_0=\omega$；
于是由 $G_3$ 得 $\mathsf{Rep}(\omega)$。证毕。

\begin{theorem}[刚性边界：$\mathsf{Fr}$ 只能刻画第4类，且 $\mathsf{Fr}\subsetneq \mathsf{Rep}$]
\label{thm:tex29-fractal-strict-subset}
在本章定义\KouJing下：
\begin{enumerate}
  \item （定位）若 $\mathsf{Fr}(\omega)$，则 $\tau(\omega)=4$，从而 $\mathsf{Fr}(\omega)\Rightarrow \mathsf{Rep}(\omega)$；
  \item （严格性）存在 $\omega$ 使得 $\mathsf{Rep}(\omega)$ 成立但 $\mathsf{Fr}(\omega)$ 不成立，
  因而 $\{\omega:\mathsf{Fr}(\omega)\}\subsetneq \{\omega:\mathsf{Rep}(\omega)\}$。
\end{enumerate}
\end{theorem}

\begin{Certificate}[证书：分形约束的必要条件与反例证书]
\label{cert:tex29-fractal-strict-subset}
证书登记谓词链：
\[
\begin{aligned}
  F_1 &: \mathsf{Fr}(\omega)\Rightarrow \text{存在多尺度修复塔且层间保持同伦等价递归约束},\\
  F_2 &: F_1 \Rightarrow \tau(\omega)=4,\\
  F_3 &: \exists \omega_\star\ \bigl(\mathsf{Rep}(\omega_\star)\wedge \neg \mathsf{Fr}(\omega_\star)\bigr)\quad(\text{取 }
    \tau(\omega_\star)=1\ \text{的局域奇点障碍}).
\end{aligned}
\]
\end{Certificate}

\paragraph{证明（谓词因果链）.}
（1）由 $F_1$，$\mathsf{Fr}(\omega)$ 蕴含“逐层递归修复 + 层间同伦等价约束”，这正是第 $4$ 类的定义性特征，
故由 $F_2$ 得 $\tau(\omega)=4$。又由定义~\ref{def:tex29-six-classes} 中第 $4$ 类的定义性条件（逐层存在 $D_k$ 使 $\partial D_k=\omega_k$）
  可得每层均可由边缘算子平凡化，故 $\mathsf{Fr}(\omega)\Rightarrow \mathsf{Rep}(\omega)$。

（2）由 $F_3$ 取反例 $\omega_\star$：令 $\omega_\star$ 为一个局域奇点型障碍，满足一次性边缘链平凡化，
则 $\mathsf{Rep}(\omega_\star)$ 成立；但其不具备多尺度层级与层间同伦等价递归自相似结构，
因而 $\neg\mathsf{Fr}(\omega_\star)$。因此 $\mathsf{Fr}$ 不能覆盖 $\mathsf{Rep}$，严格包含成立。证毕。

\begin{corollary}[排除性推论：第1/2/3/6类与广义分形无关]
\label{cor:tex29-non-fractal-classes}
若 $\tau(\omega)\in\{1,2,3,6\}$，则必有 $\neg\mathsf{Fr}(\omega)$。
换言之，局域奇点修复、全局拓扑重构、跨域同构修复与静态退化修复均不需要、也不应被强制纳入分形结构。
\end{corollary}

\begin{Certificate}[证书：无层间递归自相似 $\Rightarrow$ 非分形]
\label{cert:tex29-non-fractal-classes}
证书登记谓词链：
\[
\begin{aligned}
  N_1 &: \tau(\omega)\in\{1,2,3,6\}\Rightarrow \text{修复为一次性或有限层完成，缺失层间递归自相似约束},\\
  N_2 &: \mathsf{Fr}(\omega)\Rightarrow \text{必须存在层间同伦等价递归自相似约束},\\
  N_3 &: N_1\wedge N_2\Rightarrow \neg\mathsf{Fr}(\omega).
\end{aligned}
\]
\end{Certificate}

\paragraph{证明（谓词因果链）.}
由 $N_1$ 得当 $\tau(\omega)\in\{1,2,3,6\}$ 时缺失 $\mathsf{Fr}$ 的必要结构；
由 $N_2$ 知该结构为 $\mathsf{Fr}$ 的必要条件；
故由 $N_1\wedge N_2$ 得 $\neg\mathsf{Fr}(\omega)$，推论成立。证毕。

\begin{corollary}[第5类的关系：广义分形至多是其极小特例]
\label{cor:tex29-transfinite-fractal-special}
若 $\tau(\omega)=5$（超限全序数修复类），则 $\mathsf{Fr}(\omega)$ 仅在额外满足“层间同伦等价递归约束”时成立；
在一般情形下（层间障碍异构），$\mathsf{Fr}(\omega)$ 不成立。
\end{corollary}

\begin{Certificate}[证书：第5类的充要分岔\KouJing]
\label{cert:tex29-transfinite-fractal-special}
证书登记谓词链：
\[
\begin{aligned}
  O_1 &: \tau(\omega)=5\Rightarrow \text{允许层间异构，仅要求层间闭合校验（如 }\partial^2=0\text{）},\\
  O_2 &: \mathsf{Fr}(\omega)\Rightarrow \text{要求层间同伦等价递归自相似},\\
  O_3 &: (\tau(\omega)=5\wedge \mathsf{Fr}(\omega))\Rightarrow \text{第5类退化为满足额外约束的特例}.
\end{aligned}
\]
\end{Certificate}

\paragraph{证明（谓词因果链）.}
由 $O_1$ 给出第5类的允许结构，由 $O_2$ 给出广义分形的必要结构。
两者合并得到 $O_3$：只有在额外加入层间同伦等价递归约束时，第5类成员才可能同时满足 $\mathsf{Fr}$。
因此广义分形至多是第5类中的极小特例。证毕。

\subsection{备注：为何“最优修复”通常排斥分形结构}
\label{subsec:tex29-remarks-optimality}

\begin{remark}[变分收敛\KouJing下的结构选择]
在许多场景（局域奇点、全局拓扑重构、跨域映射）中，
引入递归层级的分形结构往往会增加冗余拓扑结构并引入次生障碍；
从“使缺陷势最小/使泛函收敛到最小势能点”的变分观点看，
一次性定向修复更符合最小化原则。
因此“是否使用分形结构”应被视为\emph{策略选择}而非\emph{范畴判定条件}。
\end{remark}

\begin{remark}[最终闭环定论（形式系统\KouJing）]
修复同伦类的范畴广度由 $\mathsf{Rep}$ 给出，广义分形由 $\mathsf{Fr}$ 给出，且 $\mathsf{Fr}\subsetneq \mathsf{Rep}$。
换言之：任何修复形式都不能反向定义修复同伦类全集；
“边缘算子驱动的障碍平凡化”才是唯一原生内核。
\end{remark}

\clearpage

%%%%%%%%%%%%%%%%%%%%%%%%%%%%%%%%%%%%%%%%%%%%%%%%%%%%%%%%%%%%
\section{广义分形与超限递归同伦修复塔：同伦等价与严格从属边界}
\label{sec:tex29-ch2}
%%%%%%%%%%%%%%%%%%%%%%%%%%%%%%%%%%%%%%%%%%%%%%%%%%%%%%%%%%%%

\begin{remark}[边界陈述（公理级）]
在同伦修复体系\KouJing下，广义分形仅与\emph{超限递归同伦修复塔}同伦等价；
它无法覆盖\emph{有限阶}同伦修复的全场景。
因此广义分形严格地仅为同伦修复体系的一个子集/一个特例，而不等同于“全部修复逻辑”。
该边界句的价值在于：用于排除把“分形表象”过度泛化为“修复本质”的常见错误推断。
\end{remark}

\subsection{公理级术语对齐：把四个核心拓扑单元写成同一语言}
\label{subsec:tex29-term-alignment}

\begin{definition}[全量同伦修复塔（修复的全集层级载体）]
称\textbf{全量同伦修复塔}为一个序数索引族
\[
  \mathbb{T}:=\bigl\{(X_\alpha,\omega_\alpha,D_\alpha)\bigr\}_{\alpha\in\Lambda},
\]
其中 $\Lambda$ 为某个序数（可取有限序数或超限序数），并且对每个 $\alpha\in\Lambda$ 都满足修复谓词
\[
  \partial D_\alpha=\omega_\alpha,
  \qquad\text{即}\qquad
  \mathsf{Rep}(\omega_\alpha).
\]
该定义把“同伦修复塔”固定为\emph{层级载体}：层级可以是有限的，也可以延伸到超限序数层。
\end{definition}

\begin{definition}[有限阶同伦修复塔（有限序数层）]
若 $\Lambda=\{0,1,2,\dots,n\}$（某个有限 $n\in\NN$），则称 $\mathbb{T}$ 为\textbf{有限阶同伦修复塔}，记为 $\mathbb{T}_{<\omega}$。
其核心特征是：层数有限、可穷尽；修复动作可在有限步内终止（允许一次性完成），不要求层间递归自相似约束。
\end{definition}

\begin{definition}[超限递归同伦修复塔（$\omega$-递归层）]
若 $\Lambda=\NN$（理解为第一超限序数 $\omega$ 的离散索引），并且存在层间结构映射
\[
  h_n:X_{n+1}\to X_n,\qquad
  k_n:X_n\to X_{n+1}\qquad (n\in\NN),
\]
满足
\[
  h_n\circ k_n\simeq \mathrm{id}_{X_n},
  \qquad
  k_n\circ h_n\simeq \mathrm{id}_{X_{n+1}},
\]
并且障碍序列在层间保持同伦等价的递归一致性（作为定义性约束），
则称 $\mathbb{T}$ 为\textbf{超限递归同伦修复塔}，记为 $\mathbb{T}_{\omega}^{\mathrm{rec}}$。
\end{definition}

\begin{definition}[广义分形（弱同伦等价自相似）]
称一组层级对象 $\{X_n\}_{n\in\NN}$ 构成\textbf{广义分形}，若存在一组层间同伦等价
\[
  X_{n+1}\simeq X_n\qquad (n\in\NN),
\]
且该“自相似”是递归迭代的（无限层级作为分形底线），并允许“强同构自相似”放宽为“弱同伦等价自相似”。
\end{definition}

\begin{remark}[对齐表：定义边界一眼可见]
为避免\KouJing漂移，将四个单元的边界压缩为同一张对齐表（本章只引用其定义性要点）：
\begin{center}
\renewcommand{\arraystretch}{1.25}
\begin{tabular}{|p{0.22\linewidth}|p{0.50\linewidth}|p{0.20\linewidth}|}
\hline
\textbf{单元} & \textbf{公理级定义要点（贴合本笔记\KouJing）} & \textbf{拓扑特征关键词} \\
\hline
全量同伦修复塔 &
以边缘算子为唯一执行单元，以障碍平凡化为唯一目标的全场景层级载体；层级索引可为有限或超限序数。 &
“全集层级载体” \\
\hline
有限阶同伦修复塔 &
$\Lambda=\{0,\dots,n\}$；修复可有限步终止；不要求层间递归自相似约束。 &
有限层、非递归 \\
\hline
超限递归同伦修复塔 &
$\Lambda=\NN$；层间存在同伦等价 $(h_n,k_n)$；并把“层间同伦等价递归一致性”作为定义性约束。 &
无限递归、层间连续 \\
\hline
广义分形 &
层间弱同伦等价自相似 + 无限递归迭代（分形底线）。 &
弱自相似、递归 \\
\hline
\end{tabular}
\end{center}
其中，几何侧“雅可比场/共轭点/比安基恒等式”等标准概念可参见 \cite{Lee2018RiemannianManifolds,Petersen2016RiemannianGeometry}；
链侧“链复形/边缘算子闭合/同调不变量”等标准概念可参见 \cite{Hatcher2002AT,Weibel1994HomologicalAlgebra}。
\end{remark}

\begin{remark}[底层规则固定（本章不变）]
本章对“体系边界”的所有证明都只使用同一条底层规则：
\[
  \begin{aligned}
    &\text{边缘算子是修复核心；}\quad \text{障碍是原生驱动力；}\\
    &\text{同伦修复塔是层级载体；}\quad \text{修复本质是障碍平凡化。}
  \end{aligned}
\]
因此，“是否属于修复体系全集”仍由 $\mathsf{Rep}$ 判定；“是否属于广义分形/超限递归”由额外结构谓词判定。
\end{remark}

\subsection{引理与定理：广义分形仅与超限递归同伦修复塔同伦等价（公理降级）}
\label{subsec:tex29-fractal-omega-equivalence}

\begin{lemma}[层间同伦等价的递归闭包（尺度无关性）]
\label{lem:tex29-homotopy-closure-by-induction}
设 $\{X_n\}_{n\in\NN}$ 与映射对 $(h_n,k_n)$ 满足 $X_{n+1}\simeq X_n$（即 $h_n,k_n$ 为同伦逆）。
则对任意 $n\in\NN$，有
\[
  X_n\simeq X_0.
\]
\end{lemma}

\begin{Certificate}[数学归纳法证书：按层数 $n$ 归纳的同伦自相似闭包]
\label{cert:tex29-homotopy-closure-by-induction}
令谓词 $P(n)$ 表示：“存在映射对 $(H_n,K_n)$ 使得 $H_n:X_n\to X_0$、$K_n:X_0\to X_n$ 且
$H_n\circ K_n\simeq \mathrm{id}_{X_0}$、$K_n\circ H_n\simeq \mathrm{id}_{X_n}$。”
证书化要求为：
\[
  P(0)\ \text{成立},\qquad (\forall n\in\NN)\ \bigl(P(n)\Rightarrow P(n+1)\bigr).
\]
\end{Certificate}

\paragraph{证明（数学归纳法证书；谓词因果链）.}
由证书~\ref{cert:tex29-homotopy-closure-by-induction}：
\begin{itemize}
  \item 基步：$n=0$ 时取 $H_0=\mathrm{id}_{X_0}$、$K_0=\mathrm{id}_{X_0}$。
  则 $H_0\circ K_0=\mathrm{id}_{X_0}\simeq \mathrm{id}_{X_0}$ 且 $K_0\circ H_0=\mathrm{id}_{X_0}\simeq \mathrm{id}_{X_0}$，
  故 $P(0)$ 成立；
  \item 归纳步：假设 $P(n)$ 成立，即 $X_n\simeq X_0$ 由 $(H_n,K_n)$ 给出。
  又由假设 $X_{n+1}\simeq X_n$，存在 $(h_n,k_n)$ 使得 $h_n\circ k_n\simeq \mathrm{id}_{X_n}$、
  $k_n\circ h_n\simeq \mathrm{id}_{X_{n+1}}$。
  令
  \[
    H_{n+1}:=H_n\circ h_n,\qquad
    K_{n+1}:=k_n\circ K_n,
  \]
  则由同伦与复合的结合性可得 $H_{n+1}\circ K_{n+1}\simeq \mathrm{id}_{X_0}$ 且 $K_{n+1}\circ H_{n+1}\simeq \mathrm{id}_{X_{n+1}}$，
  从而 $P(n+1)$ 成立。
\end{itemize}
由数学归纳法，$P(n)$ 对任意 $n$ 成立，即 $X_n\simeq X_0$。引理得证。证毕。

\begin{theorem}[公理降级：广义分形 $\simeq$ 超限递归同伦修复塔（证书化\KouJing）]
\label{thm:tex29-gf-equiv-omega-rec-tower}
对任意 $\omega\in\Omega^n(X;G)$，以下两种\KouJing严格等价：
\begin{enumerate}
  \item \textbf{广义分形\KouJing：} $\mathsf{Fr}(\omega)$ 为真；
  \item \textbf{超限递归同伦修复塔\KouJing：} 存在一个 $\omega$-索引修复塔数据对象
  \[
    \mathbb{T}_{\omega}^{\mathrm{rec}}
    :=
    \Bigl(\mathcal{C},\{(X_k,\omega_k,D_k,h_k,k_k)\}_{k\in\NN}\Bigr),
  \]
  使得定义~\ref{def:tex29-rep-fr} 中 (F0)--(F4) 的全部条件逐项成立。
\end{enumerate}
并且在该\KouJing下，广义分形的“无限递归”是定义性约束：若仅给出有限层数据 $\{0,\dots,N\}$，则只能得到有限截断，
而不能推出 $\mathsf{Fr}(\omega)$（因为 (F1)--(F4) 需要对所有 $k\in\NN$ 成立）。
\end{theorem}

\begin{Certificate}[证书：谓词等价（定义展开）与“不可有限截断”]
\label{cert:tex29-gf-equiv-omega-rec-tower}
证书登记谓词链：
\[
\begin{aligned}
  G_1 &: \mathsf{Fr}(\omega)\ \Longleftrightarrow\ \exists \Bigl(\mathcal{C},\{(X_k,\omega_k,D_k,h_k,k_k)\}_{k\in\NN}
    \Bigr)\ \text{满足 (F0)--(F4)}\quad(\text{定义~\ref{def:tex29-rep-fr}}),\\
  G_2 &: \Bigl(\mathcal{C},\{(X_k,\omega_k,D_k,h_k,k_k)\}_{k\in\NN}\Bigr)\ \text{满足 (F0)--(F4)}
  \ \Longleftrightarrow\ \exists \mathbb{T}_{\omega}^{\mathrm{rec}}\ \text{如定理所述},\\
  G_3 &: \bigl[\exists N\in\NN\ \text{仅给出层集 }\{0,\dots,N\}\bigr]\ \Rightarrow\ \neg\mathsf{Fr}(\omega)\quad(\text{因为 (F1)
    --(F4) 含 }\forall k\in\NN).
\end{aligned}
\]
\end{Certificate}

\paragraph{证明（谓词因果链）.}
由证书~\ref{cert:tex29-gf-equiv-omega-rec-tower}：
\begin{itemize}
  \item 由 $G_1$，$\mathsf{Fr}(\omega)$ 的真值等价于一组见证数据满足 (F0)--(F4)；
  \item 由 $G_2$，将该见证数据整体记作 $\mathbb{T}_{\omega}^{\mathrm{rec}}$（或反向拆解）不改变任何条件，
  因而两种\KouJing严格等价；
  \item 由 $G_3$，若只给出有限截断层集，则无法满足 (F1)--(F4) 中对所有 $k\in\NN$ 的全称量词要求，
  因而不能推出 $\mathsf{Fr}(\omega)$。
\end{itemize}
定理得证。证毕.

\subsection{命题：动力内核与闭合校验的统一（\texorpdfstring{$\partial^2=0$}{d2=0}）}
\label{subsec:tex29-dynamics-and-closure}

\begin{definition}[层级缺陷势与收敛\KouJing（抽象变分骨架）]
为把“多尺度障碍的拓扑势能差”写成可检验对象，引入每层的缺陷势（defect potential）
\[
  \Phi_n:\Omega^n(X_n;G)\to \RR_{\ge 0},
\]
并称修复过程\emph{收敛}，若存在 $n\to\infty$ 的极限意义下 $\Phi_n(\omega_n)\to 0$（或达到稳定阈值）。
此外，称\emph{层间闭合校验}成立，若边缘算子满足
\[
  \partial^2=0,
\]
以避免“修复引入次生障碍”的断点。
\end{definition}

\begin{proposition}[广义分形/超限递归塔的动力一致性]
\label{prop:tex29-dynamics-consistency}
在 $\mathbb{T}_{\omega}^{\mathrm{rec}}$ \KouJing下，若每层修复满足 $\partial D_n=\omega_n$ 且层间同伦等价递归一致性成立，
则其“动力内核”（多尺度驱动 + 边缘算子逐层修复 + 变分收敛）与广义分形\KouJing一致；
并且 $\partial^2=0$ 可作为统一的闭合校验标准，保证层间无断点、无次生障碍。
\end{proposition}

\begin{Certificate}[证书：动力一致性与闭合校验归并]
\label{cert:tex29-dynamics-consistency}
证书登记谓词链：
\[
\begin{aligned}
  D_1 &: \text{多尺度/层间连续障碍}\Rightarrow \text{存在层级缺陷势}\ \Phi_n(\omega_n),\\
  D_2 &: \partial D_n=\omega_n\Rightarrow \text{可把修复视为对}\ \Phi_n\ \text{的定向下降步骤},\\
  D_3 &: \text{层间同伦等价递归一致性}\Rightarrow \text{跨尺度的“同一动力”保持},\\
  D_4 &: \partial^2=0\Rightarrow \text{修复链闭合}\Rightarrow \text{不引入次生障碍},\\
  D_5 &: D_1\wedge D_2\wedge D_3\wedge D_4\Rightarrow \text{动力内核与校验标准统一}.
\end{aligned}
\]
\end{Certificate}

\paragraph{证明（谓词因果链）.}
由 $D_1$ 可将“障碍强度”编码为缺陷势；
由 $D_2$ 将每步修复解释为使缺陷势下降的定向动作；
由 $D_3$ 保证该解释在层间递归下保持一致（因此与广义分形的弱自相似动力一致）；
由 $D_4$ 给出闭合校验以排除次生障碍。
合并 $D_1$--$D_4$ 得 $D_5$，命题成立。证毕。

\subsection{定理：广义分形无法覆盖有限阶同伦修复的全场景（刚性边界）}
\label{subsec:tex29-boundary-finite-not-covered}

\begin{theorem}[刚性边界：有限阶同伦修复塔 $\Rightarrow$ 非广义分形]
\label{thm:tex29-finite-tower-not-fractal}
若修复塔为有限阶（$\mathbb{T}_{<\omega}$），则它无法满足广义分形的定义性约束（无限递归的层间弱自相似），因此不属于广义分形范畴。
从而广义分形不可能覆盖有限阶同伦修复的全场景。
\end{theorem}

\begin{Certificate}[证书：有限层级与“无限递归底线”的矛盾证书]
\label{cert:tex29-finite-tower-not-fractal}
证书登记谓词链：
\[
\begin{aligned}
  B_1 &: \mathbb{T}_{<\omega}\Rightarrow \exists n\in\NN\ \text{使得层级仅到 }n\ \text{为止（无 }X_{n+1}\text{）},\\
  B_2 &: \text{广义分形}\Rightarrow (\forall n\in\NN)\ \exists X_{n+1}\ \text{且}\ X_{n+1}\simeq X_n,\\
  B_3 &: B_1\wedge B_2\Rightarrow \bot.
\end{aligned}
\]
\end{Certificate}

\paragraph{证明（谓词因果链）.}
由 $B_1$，有限阶修复塔存在最大层 $n$，其上不存在 $X_{n+1}$；
由 $B_2$，广义分形要求对任意 $n$ 都存在下一层 $X_{n+1}$ 并满足同伦等价自相似；
二者矛盾，得 $B_3$，故定理成立。证毕。

\begin{remark}[有限阶修复的典型形式为何构成“分形禁区”】【对齐第\ref{subsec:tex29-six-classes}小节的分类】]
有限阶修复塔常见的合法形式（例如：局域奇点一次性修复、有限阶全局同胚重构、有限阶跨域同构映射、静态退化一次性修复）
本质上都属于“可有限步终止/不需要层间递归自相似”的结构；
这与广义分形的定义性底线（无限递归 + 层间弱自相似）互斥。
因此，有限阶修复不仅\emph{不需要}分形结构，强行引入分形反而可能带来冗余结构与次生障碍。
\end{remark}

\subsection{推论：从属链条与体系闭环}
\label{subsec:tex29-subordination-chain}

\begin{corollary}[从属链条：广义分形的严格位置]
\label{cor:tex29-subordination-chain}
在本笔记\KouJing下，存在严格从属链条
\[
  \text{广义分形}
  \ \simeq\
  \mathbb{T}_{\omega}^{\mathrm{rec}}
  \ \subsetneq\
  \mathbb{T}\ (\text{全量同伦修复塔})
  \ =\
  \mathbb{T}_{<\omega}\ \cup\ \mathbb{T}_{\ge\omega},
\]
其中 $\mathbb{T}_{<\omega}$ 为有限阶修复塔族，$\mathbb{T}_{\ge\omega}$ 为超限修复塔族（允许异构层）。
因此“广义分形 = 全部修复逻辑”在形式系统内被严格否定。
\end{corollary}

\begin{Certificate}[证书：等价与严格包含的合并证书]
\label{cert:tex29-subordination-chain}
证书登记谓词链：
\[
\begin{aligned}
  C_1 &: \text{定理~\ref{thm:tex29-gf-equiv-omega-rec-tower}}\Rightarrow \text{广义分形}\ \simeq\ \mathbb{T}_{\omega}
    ^{\mathrm{rec}},\\
  C_2 &: \text{定理~\ref{thm:tex29-finite-tower-not-fractal}}\Rightarrow \mathbb{T}_{<\omega}\cap \text{广义分形}=\varnothing,
    \\
  C_3 &: \mathbb{T}=\mathbb{T}_{<\omega}\cup \mathbb{T}_{\ge\omega},\\
  C_4 &: C_1\wedge C_2\wedge C_3\Rightarrow \text{广义分形是 }\mathbb{T}\text{ 的严格子集}.
\end{aligned}
\]
\end{Certificate}

\paragraph{证明（谓词因果链）.}
由 $C_1$ 得广义分形与 $\mathbb{T}_{\omega}^{\mathrm{rec}}$ 等价；
由 $C_2$ 得有限阶修复塔与广义分形互斥；
由 $C_3$ 将全量修复塔分解为有限阶与超限两部分；
合并得 $C_4$，从属链条与严格包含成立。证毕。

\begin{remark}[最终公理级定论（边界句的形式化版本）]
广义分形是“超限递归 + 层间同伦等价自相似”这一附加结构的名称；
同伦修复体系的全集由 $\mathsf{Rep}$ 给出，包含大量有限阶修复场景。
因此，广义分形严格地仅为该体系的一个特例，而非全部。
\end{remark}

\clearpage

%%%%%%%%%%%%%%%%%%%%%%%%%%%%%%%%%%%%%%%%%%%%%%%%%%%%%%%%%%%%
\section{从传统分形到广义分形：强同构自相似的放宽、向下兼容与边界保持}
\label{sec:tex29-ch3}
%%%%%%%%%%%%%%%%%%%%%%%%%%%%%%%%%%%%%%%%%%%%%%%%%%%%%%%%%%%%

\begin{remark}[核心陈述（范式级拓展）]
将传统分形的“强同构自相似约束”放宽为“弱同伦等价自相似约束”，
本质上是在不改变修复内核（边缘算子驱动的障碍平凡化）的前提下，
把“分形结构”从\emph{强刚性表象}升级为\emph{弱拓扑功能保持}：
\begin{itemize}
  \item 在满足强同构约束的受限场景，广义分形退化为传统分形，因而传统结论可作为特例被完整回收；
  \item 在不满足强同构但满足同伦等价的更广场景，广义分形扩大了“分形修复”的适用边界；
  \item 同时仍严格守住边界：分形修复（无论传统/广义）只是在同伦修复体系中的一个特例，而非全部。
\end{itemize}
\end{remark}

\subsection{公理级约束拆解：强同构自相似与弱同伦自相似}
\label{subsec:tex29-strong-vs-weak-selfsimilarity}

\begin{definition}[强自相似谓词：同构约束（传统分形底线）]
对层级对象 $X_{n+1},X_n$，定义谓词
\[
  \mathsf{SS}_{\cong}(X_{n+1},X_n)
\]
为真，当且仅当存在结构同构（强自相似映射）
\[
  \phi_n:X_{n+1}\xrightarrow{\cong} X_n,
\]
并且 $\phi_n$ 与其逆映射均为连续（即为同胚意义下的同构），且保持传统分形所要求的全部“强结构”。
（在工程\KouJing下，可理解为：缩放/重写后\emph{严格重合}的 1:1 对应。）
\end{definition}

\begin{definition}[弱自相似谓词：同伦约束（广义分形底线）]
对层级对象 $X_{n+1},X_n$，定义谓词
\[
  \mathsf{SS}_{\simeq}(X_{n+1},X_n)
\]
为真，当且仅当存在同伦等价映射对
\[
  h_n:X_{n+1}\to X_n,\qquad k_n:X_n\to X_{n+1},
\]
使得 $h_n\circ k_n\simeq \mathrm{id}_{X_n}$ 且 $k_n\circ h_n\simeq \mathrm{id}_{X_{n+1}}$。
（在工程\KouJing下，可理解为：允许层间结构\emph{不必严格重合}，但必须保持拓扑连通性与修复闭合性。）
\end{definition}

\begin{remark}[对齐表：强/弱约束的边界差异（本章只依赖定义性部分）]
\begin{center}
\renewcommand{\arraystretch}{1.25}
\begin{tabular}{|p{0.12\linewidth}|p{0.39\linewidth}|p{0.39\linewidth}|}
\hline
\textbf{维度} & \textbf{强自相似：同构约束} & \textbf{弱自相似：同伦约束} \\
\hline
核心要求 &
层间必须严格同构（同胚/结构保持的 1:1 对应） &
层间只需同伦等价（存在双向连续映射且可连续形变互逆） \\
\hline
适配范围 &
更偏向“层间同构”的理想化场景（强结构可重复） &
覆盖“层间非同构但同伦等价”的更广场景（弱结构保持） \\
\hline
结构刚性 &
迭代规则静态、固定，通常需预设重写/缩放规则 &
允许规则随障碍驱动自适应变化，只保同伦等价的一致性 \\
\hline
闭合校验 &
常依赖强结构重合来保证闭合 &
可用统一的边缘算子闭合校验（如 $\partial^2=0$）来保证无断点 \\
\hline
\end{tabular}
\end{center}
\end{remark}

\subsection{引理：强同构约束蕴含弱同伦约束（向下兼容的最小证明）}
\label{subsec:tex29-strong-implies-weak}

\begin{lemma}[强自相似蕴含弱自相似]
\label{lem:tex29-ss-strong-implies-weak}
若 $\mathsf{SS}_{\cong}(X_{n+1},X_n)$，则 $\mathsf{SS}_{\simeq}(X_{n+1},X_n)$。
\end{lemma}

\begin{Certificate}[证书：同胚同构 $\Rightarrow$ 同伦等价]
\label{cert:tex29-ss-strong-implies-weak}
证书登记谓词链：
\[
\begin{aligned}
  S_1 &: \mathsf{SS}_{\cong}(X_{n+1},X_n)\Rightarrow \exists \phi_n:X_{n+1}\xrightarrow{\cong} X_n,\\
  S_2 &: \phi_n\ \text{为同胚}\Rightarrow \exists (\phi_n,\phi_n^{-1})\ \text{满足同伦逆的条件},\\
  S_3 &: S_2 \Rightarrow \mathsf{SS}_{\simeq}(X_{n+1},X_n).
\end{aligned}
\]
\end{Certificate}

\paragraph{证明（谓词因果链）.}
由 $S_1$ 得到同胚同构 $\phi_n$。
取 $h_n:=\phi_n$、$k_n:=\phi_n^{-1}$，则有 $h_n\circ k_n=\mathrm{id}_{X_n}$ 且 $k_n\circ h_n=\mathrm{id}_{X_{n+1}}$；
恒等映射等价于自身的同伦类，因而满足同伦逆条件，得到 $S_2$。
由 $S_2$ 得 $S_3$，引理成立。证毕。

\subsection{命题：放宽是严格扩展（弱同伦等价但非同构的层间对确实存在）}
\label{subsec:tex29-weak-strictly-larger}

\begin{proposition}[弱约束严格弱于强约束]
\label{prop:tex29-weak-strictly-larger}
存在层级对象对 $(X,Y)$ 使得 $X\simeq Y$ 但 $X\not\cong Y$。
因此 $\mathsf{SS}_{\simeq}$ 的允许范围严格大于 $\mathsf{SS}_{\cong}$：
\[
  \{(X,Y):\mathsf{SS}_{\cong}(X,Y)\}\subsetneq \{(X,Y):\mathsf{SS}_{\simeq}(X,Y)\}.
\]
\end{proposition}

\begin{Certificate}[证书：同伦等价的反例证书（典型 deformation retraction；参见 \cite{Hatcher2002AT}）]
\label{cert:tex29-weak-strictly-larger}
证书登记谓词链：
\[
\begin{aligned}
  E_1 &: X:=S^1\cup_{p}[0,1]\ \text{（在一点 }p\text{ 处附加线段）},\quad Y:=S^1,\\
  E_2 &: \exists r:X\to Y\ \text{为形变收缩（deformation retraction）},\\
  E_3 &: E_2 \Rightarrow X\simeq Y,\\
  E_4 &: X\ \text{含端点型割点而 }Y\ \text{不具备该局部结构}\Rightarrow X\not\cong Y.
\end{aligned}
\]
\end{Certificate}

\paragraph{证明（谓词因果链）.}
由 $E_1$ 给出构造。
由附加线段可沿线段方向收缩回附着点 $p$，从而得到到 $Y=S^1$ 的形变收缩 $r$，即 $E_2$。
形变收缩蕴含同伦等价，得 $E_3$。
另一方面，$X$ 在附加线段端点处存在“端点型”局部结构（可被割开），而 $S^1$ 的任一点邻域均为环状，
二者不可能同胚，得 $E_4$。
由 $E_3$ 与 $E_4$ 得到“同伦等价但非同构”的反例，从而严格包含结论成立。证毕。

\subsection{定理（公理降级）：向下兼容与边界保持（数学归纳法证书）}
\label{subsec:tex29-compatibility-and-boundary}

\begin{theorem}[公理降级：传统分形 $\Rightarrow$ 广义分形；并保持“仅为特例”的边界（数学归纳法证书）]
\label{thm:tex29-paradigm-extension}
设 $\{X_n\}_{n\in\NN}$ 为一条无限层级序列。
若对每个 $n$ 均满足强自相似谓词 $\mathsf{SS}_{\cong}(X_{n+1},X_n)$，则必满足弱自相似谓词 $\mathsf{SS}_{\simeq}(X_{n+1},X_n)$，
从而传统分形是广义分形的静态特例。
并且在同伦修复体系\KouJing下，广义分形仍只对应满足分形约束的修复子类，无法覆盖有限阶与非递归修复的全场景。
\end{theorem}

\begin{Certificate}[数学归纳法证书：按层数 $n$ 归纳的“强约束链 $\Rightarrow$ 弱约束链”】【谓词因果链】]
\label{cert:tex29-paradigm-extension}
令谓词 $P(n)$ 表示：“对所有 $k<n$，若 $\mathsf{SS}_{\cong}(X_{k+1},X_k)$，则 $\mathsf{SS}_{\simeq}(X_{k+1},X_k)$。”
证书化要求为：
\[
  P(0)\ \text{成立},\qquad (\forall n\in\NN)\ \bigl(P(n)\Rightarrow P(n+1)\bigr).
\]
\end{Certificate}

\paragraph{证明（数学归纳法证书；谓词因果链）.}
由证书~\ref{cert:tex29-paradigm-extension}：
\begin{itemize}
  \item 基步：$n=0$ 时，$k<n$ 无元素，$P(0)$ 平凡成立；
  \item 归纳步：假设 $P(n)$ 成立。
  对于 $k=n$ 的新增一步，由引理~\ref{lem:tex29-ss-strong-implies-weak}（证书~\ref{cert:tex29-ss-strong-implies-weak}）
  有
  \[
    \mathsf{SS}_{\cong}(X_{n+1},X_n)\Rightarrow \mathsf{SS}_{\simeq}(X_{n+1},X_n).
  \]
  与 $P(n)$ 合并得到 $P(n+1)$。
\end{itemize}
由数学归纳法，$P(n)$ 对任意 $n$ 成立，从而“强自相似链”逐层推出“弱自相似链”，即传统分形嵌入为广义分形特例。

关于“边界保持”的部分：由第\ref{subsec:tex29-boundary-finite-not-covered}小节定理~\ref{thm:tex29-finite-tower-not-fractal}
可知有限阶修复塔不属于广义分形；由第\ref{subsec:tex29-fractal-boundary}小节定理~\ref{thm:tex29-fractal-strict-subset}
可知广义分形严格包含于可修复障碍的全集。
因此广义分形仍只是同伦修复体系的一个子类，无法覆盖全部修复形式。证毕。

\subsection{推论：传统分形结论在强约束下可作为广义分形的特例回收}
\label{subsec:tex29-classical-results-recovered}

\begin{corollary}[向下兼容推论：经典分形定理的“特例继承”】【\KouJing说明】]
\label{cor:tex29-classical-results-recovered}
若一条广义分形链在每一层额外满足传统分形所要求的强约束（例如强同构自相似、相似映射与对应的度量条件等），
则任何仅依赖这些强约束给出的经典分形结论（例如自相似结构的递推性质、在满足经典附加条件时的豪斯多夫维数计算等）
都可作为该广义分形链的特例被直接引用。
\end{corollary}

\begin{Certificate}[证书：强约束特例化 $\Rightarrow$ 经典结论可调用]
\label{cert:tex29-classical-results-recovered}
证书登记谓词链：
\[
\begin{aligned}
  R_1 &: \text{强约束（同构自相似 + 相似映射/度量条件）成立}\Rightarrow \text{回到传统分形设定},\\
  R_2 &: \text{传统分形设定}\Rightarrow \text{经典分形定理适用},\\
  R_3 &: R_1\wedge R_2\Rightarrow \text{经典结论可作为广义分形的特例调用}.
\end{aligned}
\]
\end{Certificate}

\paragraph{证明（谓词因果链）.}
由 $R_1$ 得广义分形在强约束下退化为传统分形设定；
由 $R_2$ 得经典分形定理在该设定下成立；
合并得到 $R_3$，推论成立。证毕。

\subsection{备注：动态同伦维数（接口化）与工程化适配}
\label{subsec:tex29-dynamic-homotopy-dimension}

\begin{remark}[动态同伦维数的角色定位]
传统分形中常用的豪斯多夫维数 $\dim_H$ 是“静态几何维数”，典型地在固定相似比与固定迭代规则下为常数。
在广义分形\KouJing下，层间允许“非同构但同伦等价”的弱自相似，
并且修复塔的有效递归深度与障碍序列可随开放扰动动态变化；
因此更自然的做法是把“维数”视为随修复深度变化的\emph{接口化度量}，而不是预设为固定常数。
\end{remark}

\begin{definition}[动态同伦维数（接口化度量）]
给定一条广义分形修复链（或其有限截断）
\[
  \bigl\{(X_k,\omega_k,D_k)\bigr\}_{k=0}^{n},
\]
称映射
\[
  \dim_{\mathrm{ho}}:\NN\to\RR_{\ge 0},\qquad
  n\mapsto \dim_{\mathrm{ho}}(n),
\]
为一类“动态同伦维数”，若它满足：
\begin{enumerate}
  \item （证书可调用性）$\dim_{\mathrm{ho}}(n)$ 可由该截断链的证书对象（例如 $\omega_k,D_k$ 与闭合校验）确定；
  \item （退化一致性）在强同构自相似且度量相似的经典场景下，$\dim_{\mathrm{ho}}(n)$ 退化为常数并与经典维数\KouJing一致；
  \item （工程指向）$\dim_{\mathrm{ho}}(n)$ 反映“有效递归层级/有效修复深度”的增长与稳定性，用于指导收敛与资源预算。
\end{enumerate}
\end{definition}

\begin{remark}[边界不变：维数是表象度量，修复内核仍是平凡化]
无论采用静态 $\dim_H$ 还是动态 $\dim_{\mathrm{ho}}$，
它们都只是修复过程的可观测表象度量；
同伦修复体系的原生内核始终不变：边缘算子驱动的障碍平凡化与闭合校验（如 $\partial^2=0$）。
\end{remark}

\clearpage

%%%%%%%%%%%%%%%%%%%%%%%%%%%%%%%%%%%%%%%%%%%%%%%%%%%%%%%%%%%%
\section{修复的唯一内核与分形的从属定位：边缘算子平凡化、形式矩阵与过度泛化纠偏}
\label{sec:tex29-ch4}
%%%%%%%%%%%%%%%%%%%%%%%%%%%%%%%%%%%%%%%%%%%%%%%%%%%%%%%%%%%%

\begin{remark}[核心陈述（边界条件：精确表述）]
在本笔记的同伦修复体系\KouJing下：
\begin{itemize}
  \item \textbf{修复的唯一原生内核}是边缘算子驱动的障碍平凡化：对非平凡障碍闭链 $\omega$，构造证据链 $D$ 使得 $\partial D=\omega$，
  从而把“不可闭合的障碍”转化为“可闭合的边缘链”，实现同伦映射的无断点延拓；
  \item \textbf{分形}只是额外满足“层间同伦等价（弱自相似）约束”的一种特定修复形式，
  仅适用于特定类型（多尺度、递归、且需要自相似约束）的障碍场景；
  \item 因而分形修复严格地仅为同伦修复体系的一个子集/一个特例，不等同于修复的全部方式，亦不等同于修复的内核。
\end{itemize}
该边界陈述用于排除传统分形理论中常见的错误推断：把“分形表象”过度泛化为“修复本质”；
并与离散统公理体系的基本要求一致：所有关键结论必须可证书化、可闭合、可由谓词链条核验。
\end{remark}

\subsection{公理级本质区分：修复的判据与分形的从属}
\label{subsec:tex29-core-vs-fractal-positioning}

\begin{definition}[修复内核判据（唯一入口）]
沿用第\ref{subsec:tex29-domain-predicates}小节的记号。
称障碍代表元 $\omega$ \textbf{可修复}，若修复谓词
\[
  \mathsf{Rep}(\omega)\ :\Longleftrightarrow\ \exists D\ \bigl(\partial D=\omega\bigr)
\]
成立。
该谓词是“是否进入同伦修复体系”的唯一入口判据；其真值不依赖修复外观（是否递归、是否自相似、是否分形）。
\end{definition}

\begin{definition}[分形修复谓词（从属约束）]
称 $\omega$ 具有\textbf{分形修复结构}，若在可修复性之上还满足额外结构谓词：
\[
  \mathsf{FRep}(\omega)\ :\Longleftrightarrow\ \mathsf{Rep}(\omega)\ \wedge\ \mathsf{Fr}(\omega),
\]
其中 $\mathsf{Fr}(\omega)$ 为第\ref{subsec:tex29-domain-predicates}小节引入的“层间同伦等价递归约束”谓词。
因此 $\mathsf{FRep}$ 是对 $\mathsf{Rep}$ 的\emph{从属强化}，而非对修复本质的替代。
\end{definition}

\begin{definition}[传统分形修复谓词（强同构自相似 + 同构性障碍）]
称 $\omega$ 具有\textbf{传统分形修复结构}，若存在一条分形修复塔
\[
  \bigl\{(X_n,\omega_n,D_n)\bigr\}_{n\in\NN},
  \qquad \omega_0=\omega,
  \qquad \partial D_n=\omega_n,
\]
并且对每个 $n\in\NN$ 存在同构 $\phi_n:X_{n+1}\xrightarrow{\cong} X_n$ 使得：
\begin{enumerate}
  \item （强自相似）$\mathsf{SS}_{\cong}(X_{n+1},X_n)$ 成立（见第\ref{subsec:tex29-strong-vs-weak-selfsimilarity}小节）；
  \item （障碍同构性对齐）在该同构诱导的障碍表达变换 $\phi_n^\ast$ 下满足
  \[
    \omega_{n+1}=\phi_n^\ast(\omega_n).
  \]
\end{enumerate}
记该谓词为 $\mathsf{FRep}_{\cong}(\omega)$。
\end{definition}

\begin{lemma}[传统分形修复是广义分形修复的特例]
\label{lem:tex29-classical-fractal-repair-is-special-case}
\[
  \mathsf{FRep}_{\cong}(\omega)\Rightarrow \mathsf{FRep}(\omega).
\]
\end{lemma}

\begin{Certificate}[证书：强同构自相似与同构性对齐 $\Rightarrow$ 弱同伦递归约束]
\label{cert:tex29-classical-fractal-repair-is-special-case}
证书登记谓词链：
\[
\begin{aligned}
  C_1 &: \mathsf{FRep}_{\cong}(\omega)\Rightarrow \mathsf{Rep}(\omega),\\
  C_2 &: \mathsf{FRep}_{\cong}(\omega)\Rightarrow (\forall n)\ \mathsf{SS}_{\cong}(X_{n+1},X_n),\\
  C_3 &: \text{引理~\ref{lem:tex29-ss-strong-implies-weak}}\Rightarrow \mathsf{SS}_{\cong}\Rightarrow \mathsf{SS}_{\simeq},
    \\
  C_4 &: C_2\wedge C_3\Rightarrow \text{存在层间同伦等价递归约束}\Rightarrow \mathsf{Fr}(\omega),\\
  C_5 &: C_1\wedge C_4\Rightarrow \mathsf{Rep}(\omega)\wedge \mathsf{Fr}(\omega)\Rightarrow \mathsf{FRep}(\omega).
\end{aligned}
\]
\end{Certificate}

\paragraph{证明（谓词因果链；符号推演）.}
由 $C_1$ 得 $\mathsf{Rep}(\omega)$；
由 $C_2$ 与 $C_3$ 得层间强同构自相似逐层推出弱同伦等价自相似，从而由 $C_4$ 得 $\mathsf{Fr}(\omega)$；
合并得到 $C_5$，即 $\mathsf{FRep}(\omega)$。证毕。

\begin{theorem}[公理降级：修复内核与分形从属（数学归纳法证书）]
\label{thm:tex29-core-vs-fractal-positioning}
在上述定义下成立：
\begin{enumerate}
  \item （内核唯一性）$\omega$ 是否属于“可修复障碍同伦类全集”仅由 $\mathsf{Rep}(\omega)$ 判定；
  \item （从属性）$\mathsf{FRep}(\omega)\Rightarrow \mathsf{Rep}(\omega)$；
  \item （严格性）存在 $\omega$ 使得 $\mathsf{Rep}(\omega)$ 成立但 $\mathsf{FRep}(\omega)$ 不成立，
  因而分形修复严格小于修复全集：$\{\omega:\mathsf{FRep}(\omega)\}\subsetneq \{\omega:\mathsf{Rep}(\omega)\}$。
\end{enumerate}
\end{theorem}

\begin{Certificate}[数学归纳法证书：按修复步数的内核判定不变性与反例证书]
\label{cert:tex29-core-vs-fractal-positioning}
证书登记谓词链：
\[
\begin{aligned}
  K_1 &: \text{由定理~\ref{thm:tex29-axiom-downgraded}（第\ref{subsec:tex29-axiom-downgraded}小节）得：全集归属}\Leftrightarrow
    \mathsf{Rep},\\
  K_2 &: \mathsf{FRep}(\omega)\Rightarrow \mathsf{Rep}(\omega)\quad(\text{由定义直接推出}),\\
  K_3 &: \exists \omega_\star\ \bigl(\mathsf{Rep}(\omega_\star)\wedge \neg\mathsf{Fr}(\omega_\star)\bigr)
    \quad(\text{由定理~\ref{thm:tex29-fractal-strict-subset}}),\\
  K_4 &: K_3\Rightarrow \mathsf{Rep}(\omega_\star)\wedge \neg\mathsf{FRep}(\omega_\star).
\end{aligned}
\]
\end{Certificate}

\paragraph{证明（数学归纳法证书；谓词因果链；符号推演）.}
由 $K_1$ 得（1）成立（修复全集的判定仅依赖 $\mathsf{Rep}$）；
由 $K_2$ 得（2）成立（从属性）；
由 $K_3$ 取反例 $\omega_\star$，并由 $K_4$ 得 $\mathsf{Rep}(\omega_\star)$ 但 $\neg\mathsf{FRep}(\omega_\star)$，
故（3）成立，得到严格包含。证毕。

\begin{remark}[备注：为何“分形不是本质”】【以谓词结构表达】]
定理~\ref{thm:tex29-core-vs-fractal-positioning} 把“分形只是特例”写成谓词包含关系：
\[
  \mathsf{FRep}=\mathsf{Rep}\wedge \mathsf{Fr},
  \qquad\Rightarrow\qquad
  \mathsf{FRep}\subsetneq \mathsf{Rep}.
\]
因此任何把 $\mathsf{Fr}$（分形约束）当作修复本质的说法，都等价于把“额外从属约束”误当作“唯一入口判据”，在形式系统中是可被直接否定的。
\end{remark}

\subsection{命题：修复形式矩阵（分形仅为其中一类）}
\label{subsec:tex29-repair-form-matrix}

\begin{definition}[六类修复形式（与第\ref{subsec:tex29-six-classes}小节对齐）]
为把“同伦修复塔的核心修复形式体系”写成可检验对象，定义修复形式标签
\[
  \kappa(\omega)\in\{\mathsf{FR},\mathsf{SP},\mathsf{GH},\mathsf{CD},\mathsf{TR},\mathsf{SD}\},
\]
分别对应：
\begin{enumerate}
  \item $\mathsf{FR}$：分形递归式修复（多尺度递归 + 层间同伦等价自相似约束）；
  \item $\mathsf{SP}$：单点爆破式修复（局域奇点一次性拓扑手术）；
  \item $\mathsf{GH}$：全局同胚式修复（全局拓扑型/基本群级别的重构）；
  \item $\mathsf{CD}$：跨域同构式修复（异构系统之间保结构同伦映射的构造）；
  \item $\mathsf{TR}$：超限递归式修复（沿序数层逐层消解，允许层间异构，仅需闭合校验如 $\partial^2=0$）；
  \item $\mathsf{SD}$：静态退化式修复（静态闭系统的有限已知障碍，一次性平凡化，修复塔冻结）。
\end{enumerate}
该标签与第\ref{subsec:tex29-six-classes}小节的 $\tau(\omega)\in\{1,\dots,6\}$ 可一一对应（仅换记号以强调“修复形式矩阵”的工程语义）。
\end{definition}

\begin{proposition}[形式矩阵覆盖与分形的严格从属]
\label{prop:tex29-repair-form-matrix}
若 $\kappa$ 的定义\KouJing满足：
\[
  \mathsf{Rep}(\omega)\Rightarrow \kappa(\omega)\ \text{可被确定},
  \qquad
  \kappa(\omega)=\mathsf{FR}\Rightarrow \mathsf{FRep}(\omega),
\]
则：
\begin{enumerate}
  \item （全场景覆盖）所有可修复障碍都落入上述六类形式之一；
  \item （分形从属）分形递归式修复 $\mathsf{FR}$ 只是矩阵中的一类形式，且严格从属于修复全集。
\end{enumerate}
\end{proposition}

\begin{Certificate}[证书：覆盖性与从属性的合并证书]
\label{cert:tex29-repair-form-matrix}
证书登记谓词链：
\[
\begin{aligned}
  M_1 &: \mathsf{Rep}(\omega)\Rightarrow \kappa(\omega)\in\{\mathsf{FR},\mathsf{SP},\mathsf{GH},\mathsf{CD},\mathsf{TR},
    \mathsf{SD}\},\\
  M_2 &: \kappa(\omega)=\mathsf{FR}\Rightarrow \mathsf{FRep}(\omega),\\
  M_3 &: \text{由定理~\ref{thm:tex29-core-vs-fractal-positioning} 得 }\mathsf{FRep}\subsetneq \mathsf{Rep},\\
  M_4 &: M_1\wedge M_2\wedge M_3\Rightarrow \text{覆盖成立且 }\mathsf{FR}\ \text{为严格从属特例}.
\end{aligned}
\]
\end{Certificate}

\paragraph{证明（谓词因果链；符号推演）.}
由 $M_1$ 得覆盖性；由 $M_2$ 得 $\mathsf{FR}$ 必为分形修复；由 $M_3$ 得分形修复严格小于修复全集；
三者合并得到 $M_4$，命题成立。证毕。

\begin{remark}[修复形式矩阵（对齐表；用于工程落地的“边界一眼可见”】【示意】]
为便于工程语境快速定位边界，给出对齐表（仅作参照，不作为本章证明前提）：
\begin{center}
\begingroup
\setlength{\tabcolsep}{4pt}
\small
\renewcommand{\arraystretch}{1.25}
\begin{tabular}{|p{0.17\linewidth}|p{0.21\linewidth}|p{0.30\linewidth}|p{0.20\linewidth}|}
\hline
\textbf{修复形式} & \textbf{障碍类型} & \textbf{公理级核心逻辑（$\partial D=\omega$）} & \textbf{与分形的区别} \\
\hline
$\mathsf{FR}$ 分形递归式 &
多尺度、分布式、需递归 &
沿修复塔逐层构造 $D_n$ 使 $\partial D_n=\omega_n$，并加入层间同伦等价自相似约束（传统分形进一步要求强同构与障碍同构性对齐） &
必须自相似；仅适用特定场景 \\
\hline
$\mathsf{SP}$ 单点爆破式 &
局域奇点、集中式 &
一次性局域拓扑手术构造 $D$ 使 $\partial D=\omega$ &
无递归、无自相似 \\
\hline
$\mathsf{GH}$ 全局同胚式 &
全局拓扑型断裂 &
全局同胚形变/重构后使障碍在新空间中平凡化 &
非分层递归；全局一次性 \\
\hline
$\mathsf{CD}$ 跨域同构式 &
跨系统映射断裂 &
构造保结构同伦等价映射，把已证修复规则迁移到目标域 &
跨域对齐；非单域自相似 \\
\hline
$\mathsf{TR}$ 超限递归式 &
超限层、层间异构 &
沿序数层逐层构造 $D_\alpha$，只要求闭合校验（如 $\partial^2=0$） &
允许异构；自相似非必要 \\
\hline
$\mathsf{SD}$ 静态退化式 &
静态已知障碍 &
一次性构造静态 $D$ 完成平凡化；修复塔冻结 &
无动态迭代；兼容既有体系 \\
\hline
\end{tabular}
\endgroup
\end{center}
\medskip

\noindent\textbf{典型应用场景（示意）：}
\begin{itemize}
  \item $\mathsf{FR}$：血管/水系网络优化、湍流控制、海岸线稳态、分布式电网潮流优化；
  \item $\mathsf{SP}$：单点材料缺陷/奇点断点修复、金融黑天鹅局域风险隔离、认知链路核心断点修补；
  \item $\mathsf{GH}$：系统底层拓扑结构重构、基本群级别全局对齐、规则体系的本体论级翻转；
  \item $\mathsf{CD}$：跨域模型迁移与保结构对齐、异构系统间修复规则的同伦映射迁移；
  \item $\mathsf{TR}$：超限层级开放系统的全维度风险消解、覆盖全序数层的异构障碍逐层修复；
  \item $\mathsf{SD}$：静态闭系统内常规定理证明、固定参数工程优化与既有体系的兼容求解。
\end{itemize}
\end{remark}

\subsection{推论：澄清“过度泛化”并得到体系完备性闭环}
\label{subsec:tex29-anti-overgeneralization}

\begin{corollary}[澄清推论：分形表象不构成修复本质]
\label{cor:tex29-anti-overgeneralization}
任何把“分形自相似结构”当作“所有复杂系统的普遍底层修复规律”的叙述，都与本形式系统的入口判据 $\mathsf{Rep}$ 冲突。
正确\KouJing应为：
\[
  \text{修复本质}=\mathsf{Rep}\quad\text{而}\quad \text{分形修复}=\mathsf{FRep}\subsetneq \mathsf{Rep}.
\]
\end{corollary}

\begin{Certificate}[证书：把表象当本质 $\Rightarrow$ 与入口判据冲突]
\label{cert:tex29-anti-overgeneralization}
证书登记谓词链：
\[
\begin{aligned}
  P_1 &: \text{修复体系入口判据为 }\mathsf{Rep}\quad(\text{定理~\ref{thm:tex29-core-vs-fractal-positioning}}),\\
  P_2 &: \text{把分形当本质}\Rightarrow \mathsf{Rep}\ \text{被}\ \mathsf{Fr}\ \text{替代为入口判据},\\
  P_3 &: \mathsf{Fr}\ \text{为从属约束且 }\mathsf{FRep}\subsetneq \mathsf{Rep}\quad(\text{定理~\ref{thm:tex29-core-vs-fractal-positioning}}),\\
  P_4 &: P_2\wedge P_3\Rightarrow \text{排除大量可修复但非分形的场景}\Rightarrow \text{与 }P_1\text{ 冲突}.
\end{aligned}
\]
\end{Certificate}

\paragraph{证明（谓词因果链；符号推演）.}
由 $P_1$ 固定入口判据；由 $P_2$ 得“过度泛化”等价于把从属约束当作入口；
由 $P_3$ 得该替代会严格缩小可修复集合；于是由 $P_4$ 得与入口判据冲突。
故推论成立。证毕。

\begin{remark}[最终闭环定论（本章总结）]
本章用形式逻辑把一句话落实为三层结构：
\[
  \underbrace{\mathsf{Rep}}_{\text{修复内核（唯一入口）}}
  \quad\supsetneq\quad
  \underbrace{\mathsf{FRep}}_{\text{分形修复（从属特例）}}
  \quad\supseteq\quad
  \underbrace{\mathsf{FRep}_{\cong}}_{\text{传统分形修复（强约束特例）}}.
\]
因此，体系的终极完备性来自“边缘算子平凡化”的内核，而非来自任何特定表象结构；
分形（传统/广义）只能被放在\emph{从属定位}上理解与调用。
\end{remark}

\clearpage

%%%%%%%%%%%%%%%%%%%%%%%%%%%%%%%%%%%%%%%%%%%%%%%%%%%%%%%%%%%%
\section{离散统下的分形：\GZMonFull{}基元、障碍驱动与超限递归构造}
\label{sec:tex29-ch5}
%%%%%%%%%%%%%%%%%%%%%%%%%%%%%%%%%%%%%%%%%%%%%%%%%%%%%%%%%%%%

\begin{remark}[\KouJing：术语统一（\GZMonFull{}主名；莱布尼茨单子仅作索引预告）]
本笔记中“\GZMonFull{}”指离散统第一性对象域中的可构造基元与其可证书化对象化形态；
除非显式写作“莱布尼茨单子”，否则文中出现的“\GZMonization{}/\GZMonState{}/\GZMonMotherBundle{}”等术语均指 \GZMon{} 体系下的对应结构。
\end{remark}

\begin{remark}[核心陈述（离散统\KouJing下的公理级定义重写）]
在离散统公理体系\KouJing下，“分形”的原生本质不再被定义为“连续统平坦背景空间内的静态自相似子集”，
而被定义为如下动态构造对象：
\begin{quote}
由\GZMon{}基元出发，在障碍驱动下，通过边缘算子的定向同伦扭曲与障碍平凡化（$\partial D=\omega$），
并借助变分法（微观统计力学意义下的迭代优化）沿同伦修复塔进行超限递归构造，
所生成的满足层间同伦等价自相似约束的自相似层级连通流形（或其等价类）。
\end{quote}
传统连续统框架下的经典分形对象，可被视为上述动态构造过程在“静态冻结极限 + 强同构自相似 + 度量相似”的附加约束下的退化投影；
因此经典结论在强约束特例\KouJing下可被完整回收，而分形在修复体系中仍保持从属定位（见第\ref{sec:tex29-ch4}章）。
\end{remark}

\subsection{对象域与谓词：\GZMonFull{}基元、障碍、边缘算子与分形生成链}
\label{subsec:tex29-fractal-domain}

\begin{definition}[\GZMonFull{}基元与离散构型空间]
令 $\mathsf{Mon}$ 表示\GZMon{}基元的集合（离散统第一性对象域）。
令 $\mathsf{Cfg}(\mathsf{Mon})$ 表示由 $\mathsf{Mon}$ 生成的离散构型空间（例如：有限多重集、离散复形或可数图的同伦类等），
并称 $X\in\mathsf{Cfg}(\mathsf{Mon})$ 为系统的一个离散状态空间表示。
\end{definition}

\begin{definition}[连续统表示（静态投影/退化函子）]
为刻画“静态极限投影”的兼容\KouJing，引入一个表示函子（或投影算子）
\[
  \Pi:\mathsf{Cfg}(\mathsf{Mon})\longrightarrow \mathsf{Cont},
\]
其中 $\mathsf{Cont}$ 表示连续统对象范畴（例如嵌入到某个度量空间的子集、流形或测度空间对象）。
$\Pi(X)$ 仅作为“观测/描述层”对象参与讨论，不作为分形的第一性本体。
\end{definition}

\begin{definition}[分形生成链（障碍驱动的修复塔数据）]
\label{def:tex29-fractal-generation-chain}
称一组数据
\[
  \mathbb{F}:=\bigl\{(X_n,\omega_n,D_n)\bigr\}_{n\in\NN}
\]
为一条\textbf{分形生成链}，若满足：
\begin{enumerate}
  \item （离散统第一性）$X_n\in\mathsf{Cfg}(\mathsf{Mon})$；
  \item （障碍驱动）$\omega_n$ 为 $X_n$ 上的非平凡障碍表达（代表元），并可由边缘算子平凡化：
  \[
    \partial D_n=\omega_n;
  \]
  \item （层间弱自相似）存在同伦等价映射对 $(h_n,k_n)$ 使 $X_{n+1}\simeq X_n$（见第\ref{subsec:tex29-strong-vs-weak-selfsimilarity}
    小节的弱自相似\KouJing）；
  \item （闭合校验）$\partial^2=0$ 作为层间闭合校验成立，用于排除次生障碍断点。
\end{enumerate}
当上述链满足额外约束“静态冻结 + 强同构自相似 + 度量相似”时，可经 $\Pi$ 投影到经典分形对象（见第\ref{subsec:tex29-fractal-static-degeneration}小节）。
\end{definition}

\begin{remark}[传统分形理论的三类原生缺陷与本体论反转（对齐表）]
在本章\KouJing下，“澄清传统分形的过度泛化缺陷”可被压缩为三条可检验的本体论反转：
\begin{center}
\renewcommand{\arraystretch}{1.25}
\begin{tabular}{|p{0.40\linewidth}|p{0.52\linewidth}|}
\hline
\textbf{传统分形的原生缺陷（连续统第一性）} & \textbf{离散统\KouJing下的反转（本章采用）} \\
\hline
先验预设平坦连续统背景空间；分形为其内部子集（第二性） &
分形本体从离散\GZMon{}构型出发（第一性）；连续统仅为静态投影 $\Pi$ 的观测层 \\
\hline
自相似来自人为预设迭代规则（外生、静态） &
自相似由障碍驱动的同伦修复链内生生成；层间同伦等价为定义性约束（内生、可证书化） \\
\hline
偏向静态闭系统描述；难以表达开放动态演化 &
分形被定义为开放系统障碍修复的动态拓扑印记；变分迭代与闭合校验提供可核验闭环 \\
\hline
\end{tabular}
\end{center}
\end{remark}

\subsection{公理降级：分形生成的动力学闭环（数学归纳法证书）}
\label{subsec:tex29-fractal-dynamic-loop}

\begin{definition}[缺陷势与变分更新（微观统计力学\KouJing）]
为表达“变分法=微观统计力学的迭代优化”，对每层 $X_n$ 引入缺陷势（defect potential）
\[
  \Phi_n:\Omega^*(X_n;G)\to\RR_{\ge 0}.
\]
称一次更新为\textbf{变分可接受}，若它在满足修复约束 $\partial D_n=\omega_n$ 的前提下，使缺陷势不增：
\[
  \Phi_{n+1}(\omega_{n+1})\le \Phi_n(\omega_n).
\]
该定义只要求“可核验的单调性\KouJing”，不依赖任何具体的连续极限或外部拟合叙述。
\end{definition}

\begin{theorem}[公理降级：障碍驱动的分形生成链存在性与闭环保持（数学归纳法证书）]
\label{thm:tex29-fractal-dynamic-loop}
设给定初始离散构型 $X_0\in\mathsf{Cfg}(\mathsf{Mon})$ 与初始障碍代表元 $\omega_0$。
若对每一层 $n\in\NN$ 均存在一条更新规则，使得：
\begin{enumerate}
  \item （修复步）可构造 $D_n$ 使 $\partial D_n=\omega_n$；
  \item （弱自相似步）可构造 $X_{n+1}$ 与同伦等价映射对 $(h_n,k_n)$ 使 $X_{n+1}\simeq X_n$；
  \item （变分步）更新为变分可接受：$\Phi_{n+1}(\omega_{n+1})\le \Phi_n(\omega_n)$；
  \item （闭合步）层间闭合校验 $\partial^2=0$ 成立。
\end{enumerate}
则存在一条分形生成链 $\mathbb{F}=\{(X_n,\omega_n,D_n)\}_{n\in\NN}$（定义~\ref{def:tex29-fractal-generation-chain}），
并且该链对任意有限截断 $0\le k\le n$ 保持：
\[
  \mathsf{Rep}(\omega_k),\qquad X_{k+1}\simeq X_k,\qquad \Phi_{k+1}(\omega_{k+1})\le \Phi_k(\omega_k).
\]
\end{theorem}

\begin{Certificate}[数学归纳法证书：按层数 $n$ 归纳的“修复--自相似--变分--闭合”四元闭环]
\label{cert:tex29-fractal-dynamic-loop}
令谓词 $P(n)$ 表示：
“存在长度为 $n$ 的链 $\{(X_k,\omega_k,D_k)\}_{k=0}^{n}$，满足
$X_k\in\mathsf{Cfg}(\mathsf{Mon})$、$\partial D_k=\omega_k$（对 $k<n$）、$X_{k+1}\simeq X_k$、以及
$\Phi_{k+1}(\omega_{k+1})\le \Phi_k(\omega_k)$。”
证书化要求为：
\[
  P(0)\ \text{成立},\qquad (\forall n\in\NN)\ \bigl(P(n)\Rightarrow P(n+1)\bigr).
\]
\end{Certificate}

\paragraph{证明（数学归纳法证书；谓词因果链；符号推演）.}
由证书~\ref{cert:tex29-fractal-dynamic-loop}：
\begin{itemize}
  \item 基步：$n=0$ 时取给定的 $(X_0,\omega_0)$；此时 $P(0)$ 只要求对象域成立，而 $(X_0,\omega_0)$ 已作为输入给定并满足定义域条件，故 $P(0)$ 成立；
  \item 归纳步：假设 $P(n)$ 成立，即已构造到第 $n$ 层。
  由定理假设的四条更新规则：
  \begin{itemize}
    \item 由“修复步”得存在 $D_n$ 使 $\partial D_n=\omega_n$，从而 $\mathsf{Rep}(\omega_n)$ 成立；
    \item 由“弱自相似步”得存在 $X_{n+1}$ 与 $X_{n+1}\simeq X_n$；
    \item 由“变分步”得 $\Phi_{n+1}(\omega_{n+1})\le \Phi_n(\omega_n)$；
    \item 由“闭合步”得链条闭合校验可通过（不引入次生断点）。
  \end{itemize}
  将新增层并入既有长度为 $n$ 的链，即得到长度为 $n+1$ 的链，故 $P(n+1)$ 成立。
\end{itemize}
由数学归纳法，$P(n)$ 对任意 $n$ 成立，从而存在无限链 $\mathbb{F}$，并保持四元闭环不变性。证毕。

\begin{remark}[备注：边缘算子在本章的定位是“构造动作”】【与传统静态\KouJing区分】]
在本章\KouJing下，边缘算子 $\partial$ 的角色不是“事后计算一个既有集合的边界”，
而是“为障碍闭链 $\omega$ 主动构造一条证据链 $D$ 使 $\partial D=\omega$”，从而把修复动作内生到生成过程中。
因此，分形的层级分支与自相似结构被视为修复动作的可证书化痕迹，而非外部预设的迭代装饰。
\end{remark}

\subsection{引理：分形生成规则由障碍与变分闭环内生确定（无需外部预设迭代）}
\label{subsec:tex29-fractal-no-external-iteration}

\begin{lemma}[内生性：障碍场与变分\KouJing确定可接受更新族]
\label{lem:tex29-fractal-no-external-iteration}
在定理~\ref{thm:tex29-fractal-dynamic-loop}的设定下，
每一层允许的更新族可被约束为仅依赖于 $(X_n,\omega_n)$、修复约束 $\partial D_n=\omega_n$ 与变分可接受性，
因此“迭代规则”可视为由障碍驱动的内生产物，而非外部人为预设的先验公设。
\end{lemma}

\begin{Certificate}[证书：可接受更新族的生成谓词链]
\label{cert:tex29-fractal-no-external-iteration}
证书登记谓词链：
\[
\begin{aligned}
  I_1 &: \text{给定}(X_n,\omega_n),\\
  I_2 &: \text{修复约束：}\exists D_n\ (\partial D_n=\omega_n),\\
  I_3 &: \text{变分约束：}\Phi_{n+1}(\omega_{n+1})\le \Phi_n(\omega_n),\\
  I_4 &: I_1\wedge I_2\wedge I_3\Rightarrow \text{定义“可接受更新族”}\ \mathcal{U}_n\ \text{（无需外部迭代规则）}.
\end{aligned}
\]
\end{Certificate}

\paragraph{证明（谓词因果链；符号推演）.}
由 $I_1$ 固定当前层数据；由 $I_2$ 给出修复可行域；由 $I_3$ 给出优化可接受域；
因此由 $I_4$ 可将“下一层如何更新”的规则限制在 $\mathcal{U}_n$ 内，且该集合仅由当前层数据与可核验约束确定，
不需要额外引入外部预设迭代规则。证毕。

\subsection{命题：静态退化投影下回收经典分形对象与结论}
\label{subsec:tex29-fractal-static-degeneration}

\begin{definition}[静态冻结谓词与强约束退化]
称分形生成链 $\mathbb{F}=\{(X_n,\omega_n,D_n)\}$ 满足\textbf{静态冻结}，若存在 $N$ 使得对所有 $n\ge N$：
\[
  \omega_{n+1}=\omega_n,\qquad X_{n+1}\cong X_n,
\]
并且层间映射进一步满足传统分形所需的强同构自相似与度量相似条件（作为外加约束）。
\end{definition}

\begin{proposition}[经典分形作为静态退化投影]
\label{prop:tex29-fractal-static-degeneration}
若分形生成链 $\mathbb{F}$ 满足静态冻结，
则其连续统表示 $\Pi(X_n)$ 在 $n\to\infty$ 的极限意义下退化为一个经典分形对象（连续统内的静态自相似子集）；
并且在该强约束\KouJing下，经典分形理论中仅依赖强自相似与度量相似条件的结论可作为特例被调用（对齐推论~\ref{cor:tex29-classical-results-recovered}）。
\end{proposition}

\begin{Certificate}[证书：冻结 + 强自相似 + 投影 $\Rightarrow$ 经典对象可调用]
\label{cert:tex29-fractal-static-degeneration}
证书登记谓词链：
\[
\begin{aligned}
  S_1 &: \text{静态冻结}\Rightarrow (\exists N)\ (\forall n\ge N)\ X_{n+1}\cong X_n,\\
  S_2 &: X_{n+1}\cong X_n\Rightarrow \Pi(X_{n+1})\ \text{与}\ \Pi(X_n)\ \text{满足强自相似\KouJing},\\
  S_3 &: \text{强自相似 + 度量相似}\Rightarrow \Pi(X_n)\ \text{落入经典分形设定},\\
  S_4 &: S_3\Rightarrow \text{经典结论可按推论~\ref{cor:tex29-classical-results-recovered} 调用}.
\end{aligned}
\]
\end{Certificate}

\paragraph{证明（谓词因果链；符号推演）.}
由 $S_1$ 得后期层间同构恒定；由 $S_2$ 得其在投影下表现为强自相似；
再由外加的度量相似条件得到 $S_3$，即落入经典分形设定；
由 $S_4$ 得经典结论可作为强约束特例调用。证毕。

\begin{corollary}[兼容性推论：传统分形是离散统分形的静态特例]
\label{cor:tex29-fractal-static-special}
在静态冻结 + 强同构自相似 + 度量相似的附加条件下，
离散统\KouJing下的分形生成链经 $\Pi$ 退化投影后，与传统分形对象一致；
因此传统分形理论可被视为本体系的一个静态特例\KouJing。
\end{corollary}

\begin{Certificate}[证书：由命题~\ref{prop:tex29-fractal-static-degeneration} 直接推出]
\label{cert:tex29-fractal-static-special}
证书登记谓词链：
\[
\begin{aligned}
  Q_1 &: \text{命题~\ref{prop:tex29-fractal-static-degeneration}} \Rightarrow \Pi(X_n)\ \text{退化为经典分形对象},\\
  Q_2 &: Q_1\Rightarrow \text{传统分形为本体系静态特例\KouJing}.
\end{aligned}
\]
\end{Certificate}

\paragraph{证明（谓词因果链）.}
由 $Q_1$ 得退化结论；由 $Q_2$ 得特例定位。证毕。

\subsection{推论：分形的工程化定位（可构造、可优化、可核验）}
\label{subsec:tex29-fractal-engineering}

\begin{corollary}[工程化推论：分形从描述性对象提升为可执行修复引擎]
\label{cor:tex29-fractal-engineering}
在离散统\KouJing下，分形生成链的每一步都由三类可核验约束确定：修复约束（$\partial D=\omega$）、闭合校验（$\partial^2=0$）与变分可接受性（$\Phi$ 单调性）。
因此“分形结构的生成”可被接口化为一个可执行的优化过程：给定障碍表达与资源约束，输出一条可证书化的修复塔与其自相似层级结构，
从而把分形从“事后描述工具”提升为“可主动构造、可优化、可复算验证”的修复引擎组件。
\end{corollary}

\begin{Certificate}[证书：接口化闭环 $\Rightarrow$ 可执行性与可核验性]
\label{cert:tex29-fractal-engineering}
证书登记谓词链：
\[
\begin{aligned}
  G_1 &: \partial D_n=\omega_n\Rightarrow \text{每步修复可证书化},\\
  G_2 &: \partial^2=0\Rightarrow \text{层间闭合可核验（无次生断点）},\\
  G_3 &: \Phi_{n+1}(\omega_{n+1})\le \Phi_n(\omega_n)\Rightarrow \text{优化迭代可核验（单调性）},\\
  G_4 &: G_1\wedge G_2\wedge G_3\Rightarrow \text{分形生成过程可接口化为可执行优化引擎}.
\end{aligned}
\]
\end{Certificate}

\paragraph{证明（谓词因果链；符号推演）.}
由 $G_1$ 得每步有可独立核验的证据链；由 $G_2$ 得全链闭合无断点；由 $G_3$ 得迭代具备可核验的优化方向；
三者合并得 $G_4$，推论成立。证毕。

\begin{remark}[本章闭环总结]
离散统\KouJing下的“分形”被刻画为一条障碍驱动的同伦修复塔及其自相似层级结构：
\[
  \text{分形} \equiv \text{分形生成链 } \mathbb{F}\ \text{（可证书化/可闭合/可变分优化）}.
\]
传统分形对象只是在附加强约束与静态冻结条件下，经连续统表示 $\Pi$ 得到的退化投影。
因此，分形在本体系中既被提升为可构造的动力学对象，又被严格限制为修复体系中的从属形式之一（不构成修复本质与全集）。
\end{remark}

\clearpage

%%%%%%%%%%%%%%%%%%%%%%%%%%%%%%%%%%%%%%%%%%%%%%%%%%%%%%%%%%%%
\section{动力学元定理：障碍驱动、同伦扭曲、统计变分与全局修复闭环}
\label{sec:tex29-ch6}
%%%%%%%%%%%%%%%%%%%%%%%%%%%%%%%%%%%%%%%%%%%%%%%%%%%%%%%%%%%%

\begin{remark}[核心陈述（静态判定 $\Rightarrow$ 动力学闭环）]
前四章给出了“修复”的静态入口判据（$\mathsf{Rep}$）与边界结构（分形/修复塔/有限--超限等）。
本章在同一形式系统\KouJing下进一步补全\textbf{内生驱动闭环}：
\[
\begin{aligned}
  &\text{障碍驱动}\ \Longrightarrow\ \text{边缘算子定向同伦扭曲}\ \Longrightarrow\ \text{统计变分单调下降}\\
  &\Longrightarrow\ \text{全局稳态（最小势能点）}.
\end{aligned}
\]
其目标不是引入额外物理比喻，而是把“拓扑结构、微观统计演化、变分收敛与主动修复动作”
统一写成一条可证书化的谓词因果链，从而得到一个\textbf{自驱动、自收敛、自闭合}的通用修复引擎接口。
\end{remark}

\subsection{术语对齐：把“拓扑—统计—变分—同伦修复”写成同一谓词链}
\label{subsec:tex29-dynamics-meta-terminology}

\begin{remark}[对齐表：四个核心单元的形式化定位]
为避免“把静态工具当作动力核心”的\KouJing错位，本章将四个核心单元固定为如下对齐关系：
\begin{center}
\renewcommand{\arraystretch}{1.25}
\begin{tabular}{|p{0.22\linewidth}|p{0.50\linewidth}|p{0.20\linewidth}|}
\hline
\textbf{核心单元} & \textbf{公理级内涵（本章采用）} & \textbf{体系定位} \\
\hline
变分法（统计力学\KouJing） &
以\GZMon{}（或其他微观自由度）为自由度，对其集体演化的统计聚合给出缺陷势 $\Phi$；
变分更新不再是“外部给定路径上的静态极值”，而被写成“在可接受更新族内使 $\Phi$ 单调不增的选择规则” &
同伦修复的路径优化器 \\
\hline
同伦扭曲 &
由边缘算子约束（$\partial D=\omega$）所驱动的定向拓扑形变；
其输出不是“任意形变后的对象”，而是带证据链 $D$ 的可核验更新 &
同伦修复的核心执行动作 \\
\hline
障碍（驱动力源） &
障碍代表元 $\omega$ 通过缺陷势 $\Phi(\omega)$ 产生可核验的“驱动量”；
更新规则仅以 $\omega$（及其可证书化约束）为输入，从而实现内生驱动 &
构造起点/原生驱动力 \\
\hline
最小势能点（全局最优） &
当且仅当缺陷势达到零值或稳定阈值（$\Phi\to 0$ 或冻结），并且修复谓词 $\mathsf{Rep}$ 被满足；
此时闭合校验 $\partial^2=0$ 排除次生断点 &
收敛终点与闭环标志 \\
\hline
\end{tabular}
\end{center}
\end{remark}

\begin{definition}[统计变分缺陷势（变分法=微观统计力学的接口化定义）]
沿用第\ref{subsec:tex29-fractal-domain}小节的对象域：$\mathsf{Mon}$ 为\GZMon{}基元集合，$X_n\in\mathsf{Cfg}(\mathsf{Mon})$ 为离散构型表示。
对每一层 $n$，取系综分布 $\mu_n$（$\mathsf{Mon}$ 上的概率分布）与局域势能函数
\[
  \varepsilon_n:\mathsf{Mon}\times \Omega^*(X_n;G)\to \RR_{\ge 0}.
\]
定义统计变分意义下的缺陷势为
\[
  \Phi_n(\omega):=\mathbb{E}_{m\sim\mu_n}\bigl[\varepsilon_n(m,\omega)\bigr].
\]
称一次更新为\textbf{统计变分可接受}，若对相应障碍代表元满足单调性
\[
  \Phi_{n+1}(\omega_{n+1})\le \Phi_n(\omega_n).
\]
\end{definition}

\begin{definition}[定向同伦扭曲（同伦修复的主动执行数据）]
称从 $(X_n,\omega_n)$ 到 $(X_{n+1},\omega_{n+1})$ 的一次更新为\textbf{定向同伦扭曲}，
若存在证据链 $D_n\in\mathcal{D}^{*}(X_n;G)$ 与映射对 $(h_n,k_n)$ 使得：
\begin{enumerate}
  \item （修复证据）$\partial D_n=\omega_n$；
  \item （结构保持）$X_{n+1}\simeq X_n$，并由 $(h_n,k_n)$ 给出同伦等价证据；
  \item （闭合校验）$\partial^2=0$ 用于排除修复引入的次生断点（因此 $\partial\omega_n=0$ 自动成立）。
\end{enumerate}
\end{definition}

\begin{definition}[稳态判据（最小势能点）]
对一次动力学修复链，称第 $n$ 层达到\textbf{稳态}，若
\[
  \Phi_n(\omega_n)=0\quad(\text{或达到预设稳定阈值}).
\]
当稳态与修复谓词 $\mathsf{Rep}(\omega_n)$ 相容（见命题~\ref{prop:tex29-dynamics-meta-optimality}）时，
将该稳态称为\textbf{全局最优修复完成}的标志。
\end{definition}

\subsection{公理降级：障碍驱动的同伦修复动力学闭环（数学归纳法证书）}
\label{subsec:tex29-dynamics-meta-loop}

\begin{theorem}[公理降级：动力学元定理（自驱动/自收敛/自闭合）（数学归纳法证书）]
\label{thm:tex29-dynamics-meta-theorem}
设给定初始离散构型表示 $X_0\in\mathsf{Cfg}(\mathsf{Mon})$ 与初始障碍代表元 $\omega_0$。
若对每一层 $n\in\NN$，均存在一条仅由当前层数据 $(X_n,\omega_n)$ 决定的更新规则，
可构造下一层数据 $(X_{n+1},\omega_{n+1},D_n)$，并同时满足：
\begin{enumerate}
  \item （障碍驱动的修复步）$\partial D_n=\omega_n$（即 $\mathsf{Rep}(\omega_n)$ 成立）；
  \item （定向同伦扭曲步）存在同伦等价映射对 $(h_n,k_n)$ 使 $X_{n+1}\simeq X_n$；
  \item （统计变分步）更新为统计变分可接受：$\Phi_{n+1}(\omega_{n+1})\le \Phi_n(\omega_n)$；
  \item （闭合步）边缘算子满足 $\partial^2=0$，作为层间闭合校验。
\end{enumerate}
则存在一条动力学修复链
\[
  \mathbb{R}:=\bigl\{(X_n,\omega_n,D_n)\bigr\}_{n\in\NN},
\]
并对任意有限截断 $0\le k\le n$ 保持：
\[
  \mathsf{Rep}(\omega_k),\qquad X_{k+1}\simeq X_k,\qquad \Phi_{k+1}(\omega_{k+1})\le \Phi_k(\omega_k).
\]
此外，令 $a_n:=\Phi_n(\omega_n)$，则 $\{a_n\}$ 单调不增且下有界（$a_n\ge 0$），因此存在极限 $a_\infty$；
该极限为“自收敛”的形式化表达，而闭合步提供“自闭合”的可核验保证。
\end{theorem}

\begin{Certificate}[数学归纳法证书：按迭代步数 $n$ 归纳的“驱动--扭曲--变分--闭合”四元闭环]
\label{cert:tex29-dynamics-meta-theorem}
令谓词 $P(n)$ 表示：
“存在长度为 $n$ 的链 $\{(X_k,\omega_k,D_k)\}_{k=0}^{n}$，满足
$X_k\in\mathsf{Cfg}(\mathsf{Mon})$、$\partial D_k=\omega_k$（对 $k<n$）、$X_{k+1}\simeq X_k$、以及
$\Phi_{k+1}(\omega_{k+1})\le \Phi_k(\omega_k)$（对 $k<n$），并且闭合校验 $\partial^2=0$ 成立。”
证书化要求为：
\[
  P(0)\ \text{成立},\qquad (\forall n\in\NN)\ \bigl(P(n)\Rightarrow P(n+1)\bigr).
\]
\end{Certificate}

\paragraph{证明（数学归纳法证书；谓词因果链；符号推演）.}
由证书~\ref{cert:tex29-dynamics-meta-theorem}：
\begin{itemize}
  \item 基步：$n=0$ 时取给定的 $(X_0,\omega_0)$；此时 $P(0)$ 仅要求对象域成立，而 $(X_0,\omega_0)$ 已作为输入给定并满足定义域条件，故 $P(0)$ 成立；
  \item 归纳步：假设 $P(n)$ 成立，即已构造到第 $n$ 层。
  由定理假设的四条更新规则（对第 $n$ 层）：
  \begin{itemize}
    \item 由“修复步”得存在 $D_n$ 使 $\partial D_n=\omega_n$，从而 $\mathsf{Rep}(\omega_n)$ 成立；
    \item 由“同伦扭曲步”得存在 $X_{n+1}$ 与 $X_{n+1}\simeq X_n$；
    \item 由“统计变分步”得 $\Phi_{n+1}(\omega_{n+1})\le \Phi_n(\omega_n)$；
    \item 由“闭合步”得 $\partial^2=0$ 可作为层间校验通过（不引入次生断点）。
  \end{itemize}
  将新增层并入既有长度为 $n$ 的链，即得到长度为 $n+1$ 的链，故 $P(n+1)$ 成立。
\end{itemize}
由数学归纳法，$P(n)$ 对任意 $n$ 成立，从而存在无限链 $\mathbb{R}$，并保持四元闭环不变性。
最后令 $a_n:=\Phi_n(\omega_n)$，由单调性与下界 $0$ 得极限存在。证毕。

\begin{remark}[备注：与前述静态入口判据的关系]
定理~\ref{thm:tex29-axiom-downgraded} 给出“修复同伦类全集”的静态入口判据 $\mathsf{Rep}$；
而定理~\ref{thm:tex29-dynamics-meta-theorem} 进一步给出：当 $\mathsf{Rep}$ 可被逐层证书化地实现时，
如何以“障碍 $\Rightarrow$ 扭曲 $\Rightarrow$ 变分下降 $\Rightarrow$ 闭合校验”的四元闭环把静态判定提升为动力学引擎。
因此，本章结论是对前述静态框架的动力学补全，而非额外引入新的入口判据。
\end{remark}

\subsection{引理：障碍是内生驱动力（无外部先验迭代规则）}
\label{subsec:tex29-dynamics-meta-driver}

\begin{lemma}[驱动性：更新族由障碍表达与可核验约束内生确定]
\label{lem:tex29-dynamics-meta-driver}
在定理~\ref{thm:tex29-dynamics-meta-theorem} 的设定下，
把“可接受更新族”记为
\[
  \mathcal{U}_n:=\Bigl\{(X_{n+1},\omega_{n+1},D_n)\ :\
  \partial D_n=\omega_n,\ X_{n+1}\simeq X_n,\ \Phi_{n+1}(\omega_{n+1})\le \Phi_n(\omega_n),\ \partial^2=0\Bigr\}.
\]
则 $\mathcal{U}_n$ 的定义仅依赖于当前层的障碍表达 $\omega_n$ 与固定的可核验约束，
从而“动力来源于障碍本身”可被写成如下无歧义形式：
更新规则不需要额外外部先验迭代公设，即可由 $\omega_n$ 诱导出可接受更新域 $\mathcal{U}_n$。
\end{lemma}

\begin{Certificate}[证书：驱动内生性的谓词链]
\label{cert:tex29-dynamics-meta-driver}
证书登记谓词链：
\[
\begin{aligned}
  R_1 &: \text{给定}(X_n,\omega_n),\\
  R_2 &: \partial D_n=\omega_n\Rightarrow \text{修复可行域由 }\omega_n\text{ 确定},\\
  R_3 &: \Phi_{n+1}(\omega_{n+1})\le \Phi_n(\omega_n)\Rightarrow \text{变分可接受域由 }\omega_n\text{ 确定},\\
  R_4 &: X_{n+1}\simeq X_n\ \wedge\ \partial^2=0\Rightarrow \text{结构与闭合约束为固定规则},\\
  R_5 &: R_1\wedge R_2\wedge R_3\wedge R_4\Rightarrow \mathcal{U}_n\ \text{仅依赖 }\omega_n\text{ 与固定约束}.
\end{aligned}
\]
\end{Certificate}

\paragraph{证明（谓词因果链；符号推演）.}
由 $R_1$ 固定当前层输入；由 $R_2$ 与 $R_3$，修复与变分两类可行域均以 $\omega_n$ 为判据入口；
由 $R_4$，同伦保持与闭合校验为预先固定的结构规则；
合并得 $R_5$，从而 $\mathcal{U}_n$ 无需外部先验迭代公设即可由当前障碍内生确定。证毕。

\subsection{命题：最小势能点与全局修复完成的等价（无次生断点）}
\label{subsec:tex29-dynamics-meta-optimality}

\begin{proposition}[最小势能点 $\Leftrightarrow$ 障碍平凡化完成]
\label{prop:tex29-dynamics-meta-optimality}
若缺陷势与修复谓词满足相容性：
\[
  (\forall n)\ (\forall \omega)\qquad \Phi_n(\omega)=0\ \Longleftrightarrow\ \mathsf{Rep}(\omega),
\]
则在定理~\ref{thm:tex29-dynamics-meta-theorem} 的动力学修复链上，
达到稳态 $\Phi_n(\omega_n)=0$ 等价于障碍平凡化完成；并且由 $\partial^2=0$ 与 $\partial D_n=\omega_n$ 得
\[
  \partial\omega_n=\partial(\partial D_n)=0,
\]
从而修复链在稳态处通过闭合校验并排除次生断点。
\end{proposition}

\begin{Certificate}[证书：相容性 $\Rightarrow$ “最小势能点=全局修复”】【与闭合校验合并】]
\label{cert:tex29-dynamics-meta-optimality}
证书登记谓词链：
\[
\begin{aligned}
  O_1 &: \Phi_n(\omega_n)=0\Longleftrightarrow \mathsf{Rep}(\omega_n),\\
  O_2 &: \mathsf{Rep}(\omega_n)\Rightarrow \exists D_n\ (\partial D_n=\omega_n),\\
  O_3 &: \partial^2=0\Rightarrow \partial(\partial D_n)=0,\\
  O_4 &: O_2\wedge O_3\Rightarrow \partial\omega_n=0\ \text{（闭合校验通过）},\\
  O_5 &: O_1\wedge O_4\Rightarrow \text{稳态与全局修复完成等价，且无次生断点}.
\end{aligned}
\]
\end{Certificate}

\paragraph{证明（谓词因果链；符号推演）.}
由 $O_1$ 得稳态与 $\mathsf{Rep}$ 的等价；
由 $O_2$ 得修复证据链存在；
由 $O_3$ 得边缘算子闭合；
合并 $O_2,O_3$ 得 $O_4$，即 $\omega_n$ 在闭合意义下无断点；
再合并 $O_1$ 与 $O_4$ 得 $O_5$，命题成立。证毕。

\subsection{推论：规避“局部最优陷阱”的充分条件（停滞排除）}
\label{subsec:tex29-dynamics-meta-no-local-min}

\begin{corollary}[去停滞推论：不存在 $\Phi>0$ 的停点]
\label{cor:tex29-dynamics-meta-no-local-min}
在定理~\ref{thm:tex29-dynamics-meta-theorem} 的设定下，
若进一步满足“严格下降可得性”：
\[
  (\forall n)\ \bigl(\Phi_n(\omega_n)>0\Rightarrow \exists (X_{n+1},\omega_{n+1},D_n)\in\mathcal{U}_n\ \text{使}\
    \Phi_{n+1}(\omega_{n+1})<\Phi_n(\omega_n)\bigr),
\]
则动力学修复迭代不可能在 $\Phi>0$ 的状态上停滞；
从而“变分收敛”在形式系统\KouJing下具备排除非零停点的充分条件。
\end{corollary}

\begin{Certificate}[证书：严格下降可得 $\Rightarrow$ 排除 $\Phi>0$ 停点]
\label{cert:tex29-dynamics-meta-no-local-min}
证书登记谓词链：
\[
\begin{aligned}
  L_1 &: \Phi_n(\omega_n)>0\Rightarrow \exists \text{可接受更新使}\ \Phi_{n+1}(\omega_{n+1})<\Phi_n(\omega_n),\\
  L_2 &: \text{若在某层停滞则需 }\Phi_{n+1}(\omega_{n+1})=\Phi_n(\omega_n),\\
  L_3 &: L_1\wedge L_2\Rightarrow \text{停滞只能发生在 }\Phi_n(\omega_n)=0\ \text{（或阈值）}.
\end{aligned}
\]
\end{Certificate}

\paragraph{证明（谓词因果链；符号推演）.}
若 $\Phi_n(\omega_n)>0$，由 $L_1$ 得存在严格下降更新，故不可能满足“停滞所需的等号”（$L_2$）。
因此停滞仅可能发生在 $\Phi_n(\omega_n)=0$（或稳定阈值）处，即 $L_3$。
推论成立。证毕。

\begin{remark}[接口化适配（示意）：跨领域仅需提供 $\omega,\partial,\Phi$ 的可核验实现]
本章的动力学元定理只固定“修复引擎接口”所需的最小结构：
\[
  \text{障碍表达 }\omega,\quad \text{边缘算子 }\partial,\quad \text{缺陷势 }\Phi,\quad \text{以及同伦保持与闭合校验}.
\]
因此任何具体应用场景只需给出上述对象在对应领域的保结构表达（并通过跨域同伦等价进行对齐），即可将“可证书化修复链”迁移为可复算的工程流程。
典型映射对象可以包括：材料/网络中的局域奇点缺陷、磁场或流场拓扑断裂、电网潮流约束违例、认知链路断点等；
这些示例在此仅作为“对象域如何落地”的接口化提示，不在本形式系统中承诺具体工程结论。
\end{remark}

\clearpage

%%%%%%%%%%%%%%%%%%%%%%%%%%%%%%%%%%%%%%%%%%%%%%%%%%%%%%%%%%%%
\section{核心元定理：原生上同调障碍的系统性内生平凡化与“无空洞完成态”}
\label{sec:tex29-ch7}
%%%%%%%%%%%%%%%%%%%%%%%%%%%%%%%%%%%%%%%%%%%%%%%%%%%%%%%%%%%%

\begin{remark}[核心陈述（从“终止判据”到“构造步骤”）]
在传统静态\KouJing中，“出现非平凡障碍类 $[\omega]\neq 0$”常被用作映射延拓失败的终止判据；
而在离散统/同伦修复体系\KouJing中，\textbf{非平凡障碍类不是终止点，而是构造起点}：
\[
  [\omega]\neq 0\ \Longrightarrow\ \text{启动对应阶数的修复构造}\ \Longrightarrow\ \exists D\ (\partial D=\omega).
\]
本章将上述\KouJing写成可证书化的形式逻辑结论，并给出其直接推论：
当“系统性平凡化”在全阶上成立时，障碍群（作为“同调空洞的记账对象”）在修复完成态下平凡化，
从而得到一个无断点、可闭合、可标准化执行的“修复元算法”闭环。
\end{remark}

\subsection{术语对齐：上同调障碍、障碍群与平凡化并入}
\label{subsec:tex29-core-meta-terminology}

\begin{definition}[障碍闭链、可平凡化边界与障碍群（抽象上同调\KouJing）]
沿用第\ref{subsec:tex29-domain-predicates}小节的对象域与边缘算子记号，并补足其“降阶链复形”\KouJing：
对每个 $n\ge 1$，视 $\Omega^n(X;G)$ 组成一个链复形片段
\[
  \cdots \xrightarrow{\partial} \Omega^{n+1}(X;G)\xrightarrow{\partial}\Omega^{n}(X;G)\xrightarrow{\partial}\Omega^{n-1}
    (X;G)\xrightarrow{\partial}\cdots,
  \qquad \partial^2=0.
\]
定义：
\begin{enumerate}
  \item \textbf{障碍闭链集}（可作为“上同调障碍表达”）：
  \[
    Z^n(X;G):=\{\omega\in\Omega^n(X;G):\partial\omega=0\};
  \]
  \item \textbf{可平凡化边界集}（由修复证据链给出）：
  \[
    B^n(X;G):=\{\omega\in\Omega^n(X;G):\exists D\ \ (\partial D=\omega)\},
  \]
  其中 $D$ 为适当层级的证据链（见第\ref{subsec:tex29-domain-predicates}小节）；
  \item \textbf{障碍群}（把“同调空洞”写成可记账对象）：
  \[
    H^n_{\mathrm{obs}}(X;G):=Z^n(X;G)\big/ B^n(X;G).
  \]
\end{enumerate}
当 $[\omega]\in H^n_{\mathrm{obs}}(X;G)$ 为非零类时，可将其视为“第 $n$ 阶延拓障碍”的抽象表达；
而谓词 $\mathsf{Rep}(\omega)$（第\ref{subsec:tex29-domain-predicates}小节）等价于 $\omega\in B^n(X;G)$。
\end{definition}

\begin{definition}[平凡化并入（把障碍类转化为构造单元）]
称对障碍闭链 $\omega\in Z^n(X;G)$ 执行一次\textbf{平凡化并入}，若选择证据链 $D$ 使得
\[
  \partial D=\omega,
\]
并将该证据链作为新的构造数据并入修复塔/链复形的下一层数据中。
在该\KouJing下，“并入”不是删除障碍，而是把障碍闭链重写为可闭合的边界，从而消解其非平凡性来源。
\end{definition}

\begin{remark}[对齐：与经典上同调障碍类 $\omega\in H^n(X;G)$ 的关系]
本章使用 $H^n_{\mathrm{obs}}(X;G)=Z^n/B^n$ 作为“障碍群”的抽象记账对象；
在具体应用中，可将 $Z^n$ 选取为经典上同调的 cocycle 集合、$B^n$ 选取为 coboundary 集合，
从而使 $H^n_{\mathrm{obs}}$ 与经典上同调群 $H^n(X;G)$ \KouJing一致。
本笔记不依赖该等同的具体实现细节，仅要求其可证书化与可核验。
\end{remark}

\subsection{公理降级：全阶障碍的系统性平凡化与障碍群平凡（数学归纳法证书）}
\label{subsec:tex29-core-meta-trivialization}

\begin{theorem}[公理降级：全阶障碍可修复 $\Rightarrow$ 障碍群平凡（数学归纳法证书）]
\label{thm:tex29-core-meta-trivialization}
设边缘算子满足 $\partial^2=0$。
若对任意 $n\ge 1$ 与任意障碍闭链 $\omega\in Z^n(X;G)$，修复谓词成立：
\[
  \mathsf{Rep}(\omega)\quad\Longleftrightarrow\quad \omega\in B^n(X;G),
\]
则对任意 $n\ge 1$ 都有
\[
  H^n_{\mathrm{obs}}(X;G)=0.
\]
换言之，在“全阶系统性平凡化”成立的对象上，不存在非平凡的上同调障碍类（在该障碍群\KouJing下）。
\end{theorem}

\begin{Certificate}[数学归纳法证书：按阶数 $n$ 归纳的“闭链皆为边界”】【从而商群为零】]
\label{cert:tex29-core-meta-trivialization}
令谓词 $P(n)$ 表示：“对所有 $1\le k\le n$，有 $Z^k(X;G)=B^k(X;G)$，因此 $H^k_{\mathrm{obs}}(X;G)=0$。”
证书化要求为：
\[
  P(1)\ \text{成立},\qquad (\forall n\ge 1)\ \bigl(P(n)\Rightarrow P(n+1)\bigr).
\]
\end{Certificate}

\paragraph{证明（数学归纳法证书；谓词因果链；符号推演）.}
由证书~\ref{cert:tex29-core-meta-trivialization}：
\begin{itemize}
  \item 基步：$n=1$。
  任取 $\omega\in Z^1(X;G)$，由定理假设 $\omega\in Z^1\Rightarrow \mathsf{Rep}(\omega)$，从而 $\omega\in B^1$。
  因此 $Z^1\subseteq B^1$；反向包含 $B^1\subseteq Z^1$ 由 $\partial^2=0$ 得到：
  若 $\omega=\partial D$，则 $\partial\omega=\partial(\partial D)=0$，故 $\omega\in Z^1$。
  于是 $Z^1=B^1$，从而 $H^1_{\mathrm{obs}}=0$，即 $P(1)$ 成立。
  \item 归纳步：假设 $P(n)$ 成立。
  取任意 $\omega\in Z^{n+1}(X;G)$。
  由定理假设对任意闭链均可修复，得 $\mathsf{Rep}(\omega)$，故 $\omega\in B^{n+1}$，从而 $Z^{n+1}\subseteq B^{n+1}$。
  反向 $B^{n+1}\subseteq Z^{n+1}$ 同样由 $\partial^2=0$ 给出。
  因此 $Z^{n+1}=B^{n+1}$，并得 $H^{n+1}_{\mathrm{obs}}=0$。
  与归纳假设合并即得 $P(n+1)$。
\end{itemize}
由数学归纳法，$P(n)$ 对任意 $n\ge 1$ 成立，定理得证。证毕。

\begin{remark}[备注：本定理的逻辑边界]
定理~\ref{thm:tex29-core-meta-trivialization} 的结论是一个\emph{条件命题}：
它把“全阶系统性平凡化”作为前提，把“障碍群平凡化”作为必然后果。
这是一条可由代数拓扑基本规则（商群、闭链/边界、$\partial^2=0$）完全推出的恒真形式。
至于“在何种对象域/何种工程接口下可保证全阶系统性平凡化”属于模型设定问题，
需要在具体应用中给出可核验的更新规则与证据链构造。
\end{remark}

\subsection{命题：边缘算子是构造算子；闭合校验是推论（分工纠偏）}
\label{subsec:tex29-core-meta-edge-operator-positioning}

\begin{proposition}[构造优先：$\partial D=\omega$ 是原生动作，$\partial\omega=0$ 是派生校验]
\label{prop:tex29-core-meta-edge-operator-positioning}
在本形式系统\KouJing下：
\begin{enumerate}
  \item “修复动作”的原生定义为：对障碍表达 $\omega$ 构造证据链 $D$ 使 $\partial D=\omega$；
  \item “闭合校验”可由 $\partial^2=0$ 对修复结果直接推出：若 $\partial D=\omega$，则必有 $\partial\omega=0$。
\end{enumerate}
因此边缘算子 $\partial$ 的核心角色是\emph{构成性修复构造}；
而“校验”仅是对构造结果应用同一算子的派生一致性检查。
\end{proposition}

\begin{Certificate}[证书：构造 $\Rightarrow$ 校验 的谓词链]
\label{cert:tex29-core-meta-edge-operator-positioning}
证书登记谓词链：
\[
\begin{aligned}
  E_1 &: \text{修复定义：}\mathsf{Rep}(\omega)\Longleftrightarrow \exists D\ (\partial D=\omega),\\
  E_2 &: \partial^2=0,\\
  E_3 &: (\partial D=\omega)\wedge E_2\Rightarrow \partial\omega=\partial(\partial D)=0,\\
  E_4 &: E_1\wedge E_3\Rightarrow \text{校验是构造后的派生一致性检查}.
\end{aligned}
\]
\end{Certificate}

\paragraph{证明（谓词因果链；符号推演）.}
由 $E_1$ 给出修复的构造性定义；由 $E_2$ 与 $E_3$ 得一切通过构造得到的 $\omega$ 自动满足闭合条件；
故 $E_4$ 成立，命题得证。证毕。

\subsection{推论：开放系统难题的统一形式（目标映射的同伦延拓障碍）}
\label{subsec:tex29-core-meta-open-systems}

\begin{corollary}[统一推论：开放系统目标构造可表述为“延拓障碍的可证书化修复”】【接口化版本】]
\label{cor:tex29-core-meta-open-systems}
给定初始约束空间 $X$ 与目标稳态空间 $Y$，
把“求解一个可持续稳态”的任务抽象为构造一个无断点映射（或其同伦类）$f:X\to Y$。
若该构造在逐层延拓过程中产生的障碍表达 $\omega_n$ 皆满足：
\[
  \mathsf{Rep}(\omega_n)\quad\text{且}\quad \partial^2=0\ \text{作为闭合校验通过},
\]
并且在每层还满足变分可接受性（例如第\ref{subsec:tex29-dynamics-meta-loop}小节的 $\Phi$ 单调下降\KouJing），
则该任务可被标准化为一条可证书化的同伦修复链，因而具备可复算验证的“通用修复引擎”形式。
\end{corollary}

\begin{Certificate}[证书：任务 $\Rightarrow$ 障碍序列 $\Rightarrow$ 修复链的谓词链]
\label{cert:tex29-core-meta-open-systems}
证书登记谓词链：
\[
\begin{aligned}
  U_1 &: \text{目标构造任务}\Rightarrow \text{逐层延拓产生障碍表达序列 }\{\omega_n\},\\
  U_2 &: (\forall n)\ \mathsf{Rep}(\omega_n)\Rightarrow (\forall n)\ \exists D_n\ (\partial D_n=\omega_n),\\
  U_3 &: \partial^2=0\Rightarrow \text{每层闭合校验可核验},\\
  U_4 &: \Phi\text{ 单调性}\Rightarrow \text{迭代收敛方向可核验},\\
  U_5 &: U_1\wedge U_2\wedge U_3\wedge U_4\Rightarrow \text{得到可证书化的修复链（引擎接口化）}.
\end{aligned}
\]
\end{Certificate}

\paragraph{证明（谓词因果链；符号推演）.}
由 $U_1$ 得任务诱导障碍序列；由 $U_2$ 得每层存在证据链；
由 $U_3$ 得闭合校验可核验；由 $U_4$ 得迭代优化方向可核验；
合并得 $U_5$，推论成立。证毕。

\begin{remark}[工程示意（不作结论承诺）]
推论~\ref{cor:tex29-core-meta-open-systems} 只给出“接口化同构”的形式化壳层。
因此可以把不同领域的对象用跨域同伦等价对齐到同一壳层上（见第\ref{subsec:tex29-six-classes}小节的跨域同构修复类）：
例如材料稳态构造、潮流系统的稳态调度、磁场拓扑的稳定化、认知链路的断点修复等。

\medskip
\noindent\textbf{统一壳层（示意对齐表）：}
\begin{center}
\renewcommand{\arraystretch}{1.20}
\begin{tabular}{|p{0.14\linewidth}|p{0.36\linewidth}|p{0.40\linewidth}|}
\hline
\textbf{符号} & \textbf{形式系统角色} & \textbf{领域化示意（仅作对象化提示）} \\
\hline
$X$ & 初始约束空间/可调控自由度的状态空间 & 材料构型空间；电网潮流参数空间；磁场拓扑参数空间；认知链复形状态空间 \\
\hline
$Y$ & 目标稳态空间/希望到达的全局态 & 稳定超导态；无失控调度稳态；无灾难性重联的磁场稳态；自洽情绪稳态 \\
\hline
$f:X\to Y$ & 目标映射（或其同伦类） & “从约束到稳态”的构造过程（设计/调度/治理/修复） \\
\hline
$\omega_n$ & 第 $n$ 层延拓障碍表达 & 稳态破缺因子/约束违例/拓扑断裂与重联风险/认知断点等的可表达对象 \\
\hline
$\partial$ & 构造性修复算子（证据链生成） & 把障碍转写为边界：$\partial D_n=\omega_n$ 的可执行构造 \\
\hline
$\Phi$ & 统计变分缺陷势（单调性） & 能量/损耗/风险/应力等可度量目标的单调下降\KouJing \\
\hline
\end{tabular}
\end{center}

\noindent\textbf{跨域迁移\KouJing：}
若存在一个保结构同伦对齐把上述数据从领域 $\mathcal{A}$ 迁移到领域 $\mathcal{B}$（见第\ref{subsec:tex29-six-classes}小节的跨域同构修复类），
则可将 $\mathcal{A}$ 中已证书化的修复链/变分迭代策略迁移为 $\mathcal{B}$ 中的候选修复链（仍需在 $\mathcal{B}$ 的对象化与约束下复核证书与闭合校验）。
这些例子仅说明“如何落入同一形式系统”，并不等价于对具体物理/工程可行性的直接证明。
\end{remark}

\clearpage

%%%%%%%%%%%%%%%%%%%%%%%%%%%%%%%%%%%%%%%%%%%%%%%%%%%%%%%%%%%%
\section{保核收缩兼容：静态极限回收旧体系；动态生成要求\GZMonFull{}第一性}
\label{sec:tex29-ch8}
%%%%%%%%%%%%%%%%%%%%%%%%%%%%%%%%%%%%%%%%%%%%%%%%%%%%%%%%%%%%

\begin{remark}[核心陈述（兼容性的拓扑表达）]
本章给出“静态/动态完美兼容”的严格拓扑表述：离散统的动态总空间在静态冻结极限下，
可通过强变形收缩回到一个连续统意义下的收缩核；因此旧体系在其适用边界内保留全部合法性。
同时，动态生成（障碍修复/新结构构造）要求离开该收缩核并回到\GZMon{}第一性的对象域，
否则无法在核内完成对非平凡障碍的构造性平凡化。
\end{remark}

\subsection{定义：收缩核、强变形收缩与静态冻结谓词}
\label{subsec:tex29-core-contraction-defs}

\begin{definition}[收缩核与强变形收缩（保核收缩兼容）]
沿用第\ref{subsec:tex29-fractal-domain}小节的表示函子 $\Pi:\mathsf{Cfg}(\mathsf{Mon})\to\mathsf{Cont}$。
对任意 $X\in\mathsf{Cfg}(\mathsf{Mon})$，称 $A:=\Pi(X)\in\mathsf{Cont}$ 为 $X$ 的一个\textbf{收缩核}，
若存在映射
\[
  s:A\to X,\qquad r:X\to A
\]
满足：
\[
  r\circ s=\mathrm{id}_A,\qquad s\circ r\simeq \mathrm{id}_X.
\]
此时称 $X$ \textbf{强变形收缩到} $A$，并称该关系为\textbf{保核收缩兼容}。
\end{definition}

\begin{definition}[静态冻结谓词（全局冻结\KouJing）]
称 $X\in\mathsf{Cfg}(\mathsf{Mon})$ 满足\textbf{静态冻结}，若存在一个阈值层 $N$，
使得对任意 $n\ge N$ 的更新（若存在）都退化为恒等意义下的无效更新：
\[
  X_{n+1}\cong X_n,\qquad \omega_{n+1}=\omega_n,
\]
并且不再产生新的可证书化修复证据链（仅保留既有结构的静态描述）。
该谓词用于刻画“动态生成停止”的极限\KouJing。
\end{definition}

\subsection{公理降级：静态冻结 \texorpdfstring{$\Rightarrow$}{=>} 旧体系可回收（数学归纳法证书）}
\label{subsec:tex29-core-contraction-theorem}

\begin{theorem}[公理降级：静态冻结下的收缩核回收与同伦不变量一致（数学归纳法证书）]
\label{thm:tex29-core-contraction-theorem}
设 $X\in\mathsf{Cfg}(\mathsf{Mon})$ 满足静态冻结谓词。
则存在收缩核 $A=\Pi(X)\in\mathsf{Cont}$，使得 $X$ 强变形收缩到 $A$；
从而 $X$ 与 $A$ 具有相同的同伦不变量（例如同调/上同调群等）：
\[
  X\simeq A\quad\Longrightarrow\quad H_\ast(X;G)\cong H_\ast(A;G),\ \ H^\ast(X;G)\cong H^\ast(A;G).
\]
\end{theorem}

\begin{Certificate}[数学归纳法证书：按冻结层级 $n$ 归纳的“收缩同伦”构造]
\label{cert:tex29-core-contraction-theorem}
令谓词 $P(n)$ 表示：
“在第 $n$ 层之后冻结（$n\ge N$）的对象 $X_n$ 存在一个收缩核 $A_n:=\Pi(X_n)$，
并存在映射对 $(s_n,r_n)$ 使 $r_n\circ s_n=\mathrm{id}_{A_n}$ 且 $s_n\circ r_n\simeq \mathrm{id}_{X_n}$。”
证书化要求为：
\[
  P(N)\ \text{成立},\qquad (\forall n\ge N)\ \bigl(P(n)\Rightarrow P(n+1)\bigr).
\]
\end{Certificate}

\paragraph{证明（数学归纳法证书；谓词因果链）.}
由证书~\ref{cert:tex29-core-contraction-theorem}：
\begin{itemize}
  \item 基步：在冻结阈值 $N$ 处，静态冻结给出 $X_{N+1}\cong X_N$ 与“更新不再引入新结构”。
  取 $A_N:=\Pi(X_N)$，把 $\Pi$ 视为静态表示；令 $s_N$ 为选定的表示嵌入，$r_N$ 为表示投影，
  则在静态冻结\KouJing下满足 $r_N\circ s_N=\mathrm{id}$ 且 $s_N\circ r_N\simeq \mathrm{id}$（同伦可用恒等冻结路径给出），故 $P(N)$ 成立。
  \item 归纳步：若 $P(n)$ 成立且 $n\ge N$，由于冻结给出 $X_{n+1}\cong X_n$，
  则可取 $A_{n+1}\cong A_n$，并把 $(s_n,r_n)$ 通过同构传递到 $(s_{n+1},r_{n+1})$，
  从而 $P(n+1)$ 成立。
\end{itemize}
由数学归纳法，$P(n)$ 对所有 $n\ge N$ 成立；因此存在收缩核 $A=\Pi(X)$ 与强变形收缩映射对 $(s,r)$。
由强变形收缩得到 $X\simeq A$，同伦不变量一致性为代数拓扑基本结论，定理得证。证毕。

\begin{remark}[备注：兼容性的严格边界]
定理~\ref{thm:tex29-core-contraction-theorem} 只在“静态冻结”\KouJing下回收连续统核的合法性；
它不宣称旧体系在动态开放场景下仍能承担“障碍修复/新结构构造”的任务。
因此旧体系的合法性被\emph{边界化}为静态核内的可调用结论，而动态生成需要回到\GZMon{}第一性对象域（见下节）。
\end{remark}

\subsection{推论：动态生成的必然性（核内不可完成构造性平凡化）}
\label{subsec:tex29-core-contraction-dynamic-necessity}

\begin{corollary}[边界推论：若停留在静态核内，则非平凡障碍只能被识别，不能被构造性消解]
\label{cor:tex29-core-contraction-dynamic-necessity}
设在静态核 $A\in\mathsf{Cont}$ 的\KouJing下，允许的操作仅为“静态描述与既有结构内推演”，
即不允许引入新的证据链 $D$ 去改变 $B^n(A;G)$。
若存在非平凡障碍类 $[\omega]\in H^n_{\mathrm{obs}}(A;G)$（$n\ge 1$），
则在该静态核\KouJing内无法使其变为零类；
要实现 $\partial D=\omega$ 的构造性平凡化，必须离开核并进入动态对象域以生成新的修复证据链。
\end{corollary}

\begin{Certificate}[证书：禁止新增证据链 $\Rightarrow$ 障碍群类不变]
\label{cert:tex29-core-contraction-dynamic-necessity}
证书登记谓词链：
\[
\begin{aligned}
  C_1 &: \text{静态核\KouJing：不允许新增证据链}\Rightarrow B^n(A;G)\ \text{固定不变},\\
  C_2 &: [\omega]\neq 0\ \text{表示}\ \omega\notin B^n(A;G),\\
  C_3 &: C_1\wedge C_2\Rightarrow \omega\ \text{在静态核内仍不属于}\ B^n(A;G),\\
  C_4 &: C_3\Rightarrow [\omega]\ \text{在静态核\KouJing内无法被平凡化}.
\end{aligned}
\]
\end{Certificate}

\paragraph{证明（谓词因果链；符号推演）.}
由 $C_1$，静态核内 $B^n$ 不随操作改变；由 $C_2$，$[\omega]\neq 0$ 等价于 $\omega\notin B^n$；
因此由 $C_3$ 得该事实在核内保持；从而 $C_4$ 得无法在核内实现平凡化构造。
推论成立。证毕。

\begin{remark}[终极闭环\KouJing（本章总结）]
本章把“范式革命的兼容性”写成两条严格边界：
\[
\begin{aligned}
  &\text{静态冻结}\Rightarrow X\simeq \Pi(X)\Rightarrow \text{旧体系核内结论可回收},\\
  &\text{动态生成}\Rightarrow \text{需生成新的证据链}\Rightarrow \text{必须回到\GZMon{}第一性对象域}.
\end{aligned}
\]
因此新体系不是对旧体系的否定，而是“在静态核内保留全部可调用结论，同时在核外提供构造性修复能力”的严格升维覆盖。
\end{remark}

\clearpage

%%%%%%%%%%%%%%%%%%%%%%%%%%%%%%%%%%%%%%%%%%%%%%%%%%%%%%%%%%%%
\section{形式逻辑：莱布尼茨单子与本体系“\GZMonization{}”的严格对照}
\label{sec:tex29-leibniz-monad-vs-monadization}
\label{sec:tex29-ch9}
%%%%%%%%%%%%%%%%%%%%%%%%%%%%%%%%%%%%%%%%%%%%%%%%%%%%%%%%%%%%

\subsection{定义}
\label{subsec:tex29-leibniz-monad-vs-monadization-defs}

\begin{definition}[定义 D1（莱布尼茨单子：Leibniz Monad）]
\label{def:tex29-leibniz-monad}
称对象 $m$ 为一个（莱布尼茨意义下的）\textbf{单子}，若满足以下工作刻画：
\begin{enumerate}
  \item （简单实体；内生变化原则）存在其内部状态函数 $s_m(t)$ 及演化规则 $F_m$，使得
  \[
    s_m(t+1)=F_m\bigl(s_m(t)\bigr).
  \]
  \item （无窗性；因果封闭）不存在从外部到 $m$ 的因果输入通道；以形式语言记作
  \[
    \forall m,\ \neg\exists I\ \bigl(I\to m\bigr),
  \]
  其中 $I$ 表示外部接口/因果通道的抽象占位符。
  \item （预成和谐；全局一致性约束）存在一个全局一致性约束 $H$，使得对任意 $m_i,m_j$ 与任意时刻 $t$，
  \[
    \exists H\ \text{s.t.}\ \forall m_i,m_j,\forall t,\ R\bigl(s_{m_i}(t),s_{m_j}(t)\bigr)\ \text{由 }H\text{保证},
  \]
  其中 $R(\cdot,\cdot)$ 为给定的协调关系占位符。
\end{enumerate}
该定义的核心语义为：单子是因果封闭的简单实体；多单子的同步/一致性由全局约束承担，而非由可执行的交互算子实现。
\end{definition}

\begin{definition}[定义 D2（本笔记口径下的“\GZMonFull{}/\GZMonization{}”）]
\label{def:tex29-gframework-monadization}
结合本笔记中“\GZMonization{}总空间/构造性修复”的语义，取如下可核验的工作定义。
给定系统（空间/状态对象）$X$ 与其 $n$ 阶障碍表达 $\omega\in\Omega^n(X;G)$，称映射
\[
  \mathbb{M}:\ (X,\omega)\ \longmapsto\ (\widetilde{X},\widetilde{\omega})
\]
为一个\textbf{\GZMonization{}算子/构造}，若满足：
\begin{enumerate}
  \item （修复可证书化）存在证据链 $D\in\mathcal{D}^{n+1}(\widetilde{X};G)$ 使得
  \[
    \partial D=\widetilde{\omega};
  \]
  \item （同伦约束对齐）$(\widetilde{X},\widetilde{\omega})$ 与 $(X,\omega)$ 的关系由“同伦修复塔/层间同伦等价约束”控制：
  存在映射对 $(h,k)$ 使得 $hk\simeq \mathrm{id}$ 与 $kh\simeq \mathrm{id}$ 等可核验条件成立。
\end{enumerate}
直观上，$\mathbb{M}$ 的语义是\emph{为使障碍变为边界而进行的构造性扩展}（新增自由度/总空间/见证链），其目标是“可修复、可审计、可复用”，而非提供形而上解释框架。
\end{definition}

\begin{remark}[补充说明：若同时考虑范畴论意义的 monad 结构（仅作参照）]
若进一步把 $\mathbb{M}$ 视为某个自函子并配备单位与乘法自然变换 $(\eta,\mu)$，则可登记抽象恒等式：
\[
  \mathrm{Id}\xrightarrow{\ \eta\ }\mathbb{M},
  \qquad
  \mathbb{M}\circ\mathbb{M}\xrightarrow{\ \mu\ }\mathbb{M}.
  \qquad(\text{参照口径})
\]
但下文对比不依赖该额外结构，仅使用定义~\ref{def:tex29-leibniz-monad} 与定义~\ref{def:tex29-gframework-monadization} 的工作口径。
\end{remark}

\subsection{定理（证书+证明）}
\label{subsec:tex29-leibniz-monad-vs-monadization-thm}

\begin{theorem}[定理 T1（本质差异：因果封闭 vs 构造性可修复）]
\label{thm:tex29-leibniz-vs-monadization}
在定义~\ref{def:tex29-leibniz-monad} 与定义~\ref{def:tex29-gframework-monadization} 的口径下，
莱布尼茨单子在定义层面以“无窗/因果封闭”为核心限定；
而本笔记所称“\GZMonization{}”在定义层面以“通过构造引入证据链，使障碍满足 $\partial D=\widetilde{\omega}$”为核心接口。
因此二者属于不同的类型论层级：前者是形而上实体的刻画，后者是形式系统中的构造性算子/构造；二者不可混同。
\end{theorem}

\begin{Certificate}[证书：无窗谓词与见证链谓词的类型分离]
\label{cert:tex29-leibniz-vs-monadization}
证书登记谓词链：
\[
\begin{aligned}
  T_1 &: \text{(D1) }\forall m,\ \neg\exists I\ (I\to m)\quad(\text{无窗/外因不入}),\\
  T_2 &: \text{(D2) }\mathbb{M}(X,\omega)=(\widetilde{X},\widetilde{\omega})\ \wedge\ \exists D\ (\partial
    D=\widetilde{\omega}),\\
  T_3 &: T_1\Rightarrow \text{莱布尼茨单子的变化不以外部可执行构造为条件},\\
  T_4 &: T_2\Rightarrow \text{\GZMonization{}以可执行构造与证书化见证链为定义内核},\\
  T_5 &: T_3\wedge T_4\Rightarrow \text{二者的定义谓词与用途不同，故不可混同}.
\end{aligned}
\]
\end{Certificate}

\paragraph{证明（谓词因果链；符号推演）.}
由 $T_1$ 得莱布尼茨单子被刻画为无外因输入的因果封闭实体，因此其变化不依赖外部可执行构造，得 $T_3$；
由 $T_2$ 得“\GZMonization{}”被定义为一个输出证据链 $D$ 的构造性接口，因此以可执行构造与可复核证书为核心，得 $T_4$；
合并 $T_3,T_4$ 得 $T_5$，定理成立。证毕。

\subsection{引理（证书+证明）}
\label{subsec:tex29-leibniz-monad-vs-monadization-lem}

\begin{lemma}[引理 L1（预成和谐：协调是全局一致性约束而非交互算子）]
\label{lem:tex29-leibniz-harmony}
在定义~\ref{def:tex29-leibniz-monad} 的\KouJing下，
多单子的“协调”不等价于存在可执行耦合算子；
其最小可形式化表述为：存在一个全局一致性约束 $H$ 使得
\[
  \forall m_i,m_j,\forall t,\ R\bigl(s_{m_i}(t),s_{m_j}(t)\bigr)
\]
被 $H$ 保证。
\end{lemma}

\begin{Certificate}[证书：无窗性排除交互通道]
\label{cert:tex29-leibniz-harmony}
证书登记谓词链：
\[
\begin{aligned}
  L_1 &: \forall m,\ \neg\exists I\ (I\to m)\quad(\text{定义~\ref{def:tex29-leibniz-monad}}),\\
  L_2 &: \text{若存在耦合算子 }C:m_i\to m_j,\ \text{则 }m_j\text{ 存在外部输入通道},\\
  L_3 &: L_1\wedge L_2\Rightarrow \text{不存在可执行耦合算子实现“协调”},\\
  L_4 &: L_3\Rightarrow \text{“协调”只能被刻画为全局一致性约束（预成和谐）}.
\end{aligned}
\]
\end{Certificate}

\paragraph{证明（谓词因果链）.}
反证：若存在耦合算子 $C:m_i\to m_j$，则 $m_j$ 接受来自 $m_i$ 的外部输入通道，与 $L_1$ 的无窗性矛盾；
故由 $L_3$ 得“协调”不能来自交互算子，只能以全局约束方式登记，得 $L_4$。证毕。

\subsection{命题（证书+证明）}
\label{subsec:tex29-leibniz-monad-vs-monadization-prop}

\begin{proposition}[命题 P1（\GZMonization{}更接近可组合构造/函子，而非“终极原子实体”）]
\label{prop:tex29-monadization-compositional}
在定义~\ref{def:tex29-gframework-monadization} 的工作口径下，
“\GZMonization{}”更接近一种可组合、可迭代的构造接口：它允许在修复塔语义下进行迭代与复用，例如
\[
  (X,\omega)\xrightarrow{\ \mathbb{M}\ }(\widetilde{X}_1,\widetilde{\omega}_1)
  \xrightarrow{\ \mathbb{M}\ }(\widetilde{X}_2,\widetilde{\omega}_2)
  \xrightarrow{\ \mathbb{M}\ }\cdots.
\]
这与莱布尼茨单子作为“不可分的简单实体”的角色相反：前者强调构造的可组合性与证书链复用，后者强调本体论基元的不可分性与因果封闭。
\end{proposition}

\begin{Certificate}[证书：修复塔语义 $\Rightarrow$ 构造可迭代]
\label{cert:tex29-monadization-compositional}
证书登记谓词链：
\[
\begin{aligned}
  P_1 &: \mathbb{M}\ \text{是从输入数据 }(X,\omega)\ \text{到输出数据 }(\widetilde{X},\widetilde{\omega})\ \text{的构造接口},\\
  P_2 &: \text{输出数据可再次作为输入数据（同型接口）}\Rightarrow \mathbb{M}\ \text{可迭代},\\
  P_3 &: \exists D\ (\partial D=\widetilde{\omega})\Rightarrow \text{每次迭代均可登记证据链并拼接为证书链},\\
  P_4 &: P_1\wedge P_2\wedge P_3\Rightarrow \text{\GZMonization{}具有可组合/可复用的构造语义}.
\end{aligned}
\]
\end{Certificate}

\paragraph{证明（谓词因果链）.}
由 $P_1$，“\GZMonization{}”被刻画为一个可执行的构造接口；由 $P_2$，同型接口保证其可迭代；
由 $P_3$，每次迭代均可登记并拼接证据链；故由 $P_4$ 得命题成立。证毕。

\subsection{推论（证书+证明）}
\label{subsec:tex29-leibniz-monad-vs-monadization-cor}

\begin{corollary}[推论 C1（术语防混淆：四条对位差异口径）]
\label{cor:tex29-monad-disambiguation}
在定理~\ref{thm:tex29-leibniz-vs-monadization} 的\KouJing下，可将“莱布尼茨单子”与“本体系\GZMonization{}”之间的最小防混淆口径压缩为四条对位差异：
\begin{enumerate}
  \item \textbf{对象类型}：莱布尼茨单子是实体 $m$ 的刻画；本体系“\GZMonization{}”是构造/算子 $\mathbb{M}$ 或其输出对象 $(\widetilde{X},\widetilde{\omega}
    )$。
  \item \textbf{因果结构}：莱布尼茨口径为无窗 $\neg\exists(I\to m)$；本体系口径为显式见证链 $\exists D\ (\partial D=\widetilde{\omega})$。
  \item \textbf{一致性来源}：莱布尼茨以预成和谐 $H$ 的全局约束承担一致性；本体系以可核验的同伦等价数据与边缘算子证书链实现“局部构造$\to$全局闭合”的一致性。
  \item \textbf{可审计性}：莱布尼茨单子不是以证书化可复核为定义内核；本体系以证书/证明/可复用构造为对象层接口内核。
\end{enumerate}
\end{corollary}

\begin{Certificate}[证书：由 T1/L1/P1 抽取四条口径]
\label{cert:tex29-monad-disambiguation}
证书登记谓词链：
\[
\begin{aligned}
  C_1 &: \text{定理~\ref{thm:tex29-leibniz-vs-monadization}}\Rightarrow \text{对象类型差异与因果结构差异},\\
  C_2 &: \text{引理~\ref{lem:tex29-leibniz-harmony}}\Rightarrow \text{莱布尼茨一致性来源为全局约束},\\
  C_3 &: \text{定义~\ref{def:tex29-gframework-monadization}}\Rightarrow \text{本体系一致性来源为同伦数据+证书链拼接},\\
  C_4 &: \text{命题~\ref{prop:tex29-monadization-compositional}}\Rightarrow \text{本体系接口语义强调可复用与可审计},\\
  C_5 &: C_1\wedge C_2\wedge C_3\wedge C_4\Rightarrow \text{四条对位差异口径成立}.
\end{aligned}
\]
\end{Certificate}

\paragraph{证明（谓词因果链）.}
由 $C_1$ 得前两条（对象类型/因果结构）；
由 $C_2$ 与 $C_3$ 得第三条（一致性来源）；
由 $C_4$ 得第四条（可审计性）；合并得 $C_5$，推论成立。证毕。

\subsection{备注}
\label{subsec:tex29-leibniz-monad-vs-monadization-remark}

\begin{remark}[如何把两者翻译到同一形式语言（仅作参照，不作为证明前提）]
若希望给出“最小公共语义”的对照，可采用如下抽象化翻译：
\begin{itemize}
  \item 莱布尼茨单子可被抽象为\textbf{无输入的内部动力系统}：
  \[
    s_m(t+1)=F_m\bigl(s_m(t)\bigr).
  \]
  \item 本体系的\GZMonization{}可被抽象为\textbf{带见证输出的修复编译器}：
  \[
    (X,\omega)\ \longmapsto\ (\widetilde{X},\widetilde{\omega},D)
    \qquad\text{with}\qquad
    \partial D=\widetilde{\omega}.
  \]
\end{itemize}
两者都可以“自洽”：前者的自洽是世界解释层的自洽，后者的自洽是修复证书链的自洽。
若后续需要把定义~\ref{def:tex29-gframework-monadization} 上升为更“范畴论化”的正式定义（例如明确 $(\eta,\mu)$ 或与某个投影/代数引擎的关系），
则可在不改变本笔记符号体系的前提下补充对应接口，并据此给出更强的不可混同判据（例如不可同伦/不可等价的严格证书链）。
\end{remark}

\clearpage

%%%%%%%%%%%%%%%%%%%%%%%%%%%%%%%%%%%%%%%%%%%%%%%%%%%%%%%%%%%%
\section{形式逻辑：主纤维丛—Atiyah 序列—广义非交换代数口径下的“\GZMonFull{}”模型}
\label{sec:tex29-principal-bundle-atiyah-monad}
\label{sec:tex29-ch10}
%%%%%%%%%%%%%%%%%%%%%%%%%%%%%%%%%%%%%%%%%%%%%%%%%%%%%%%%%%%%

\subsection{定义}
\label{subsec:tex29-principal-bundle-atiyah-monad-defs}

\begin{definition}[定义 D1（主纤维丛与投影）]
\label{def:tex29-pb-projection}
令 $M$ 为“观测核/状态底空间”（可对齐为本笔记中的某个观测投影口径，例如 $\Pi(X)$ 或 $X$ 的可观测截面），
令
\[
  \pi:P\to M
\]
为一个以李群 $G$ 为结构群的\textbf{主纤维丛}。
其中投影 $\pi$ 被视为“从总空间到观测核”的严格映射（几何投影接口）。
\end{definition}

\begin{definition}[定义 D2（主纤维丛版：广义非交换李代数）]
\label{def:tex29-galgebra-jacobi-gap}
令 $\mathfrak{g}$ 为一组“可执行修复/变分/控制算子”的代数化载体，
给定一个双线性括号
\[
  [\cdot,\cdot]_{\star}:\mathfrak{g}\times\mathfrak{g}\to\mathfrak{g},
\]
但不强制其满足雅可比恒等式。
定义\textbf{雅可比子}（Jacobiator）
\[
  J_{\star}(x,y,z):=[x,[y,z]_{\star}]_{\star}+[y,[z,x]_{\star}]_{\star}+[z,[x,y]_{\star}]_{\star},
\]
并记
\[
  \mathsf{JacGap}(x,y,z):=J_{\star}(x,y,z)
\]
为“非李性缺口”的原始证书对象。
因此“主纤维丛版广义非交换李代数”在本章被严格落实为四元组
\[
  \bigl(\mathfrak{g},[\cdot,\cdot]_{\star},J_{\star},\mathsf{JacGap}\bigr).
\]
\end{definition}

\begin{definition}[定义 D3（伴随丛：把 $\mathfrak{g}$ 装到主纤维丛上）]
\label{def:tex29-adjoint-bundle}
在定义~\ref{def:tex29-pb-projection} 与定义~\ref{def:tex29-galgebra-jacobi-gap} 下，
定义伴随丛（内部算子纤维丛）
\[
  \mathrm{ad}(P):=P\times_G \mathfrak{g}.
\]
其纤维 $\mathrm{ad}(P)_x$（$x\in M$）可理解为“在观测点 $x$ 处可调用的内部算子代数态”。
\end{definition}

\begin{definition}[定义 D4（Atiyah 序列：投影的第二种严格形式）]
\label{def:tex29-atiyah-sequence}
定义 Atiyah 代数丛（规范 algebroid）
\[
  \mathsf{At}(P):=TP/G.
\]
并登记正合列（Atiyah 序列）
\[
  0\longrightarrow \mathrm{ad}(P)\xrightarrow{\ \iota\ }\mathsf{At}(P)\xrightarrow{\ a\ }TM\longrightarrow 0,
\]
其中 $a$ 称为\textbf{锚映射}（anchor），可被视为“把总空间的等价类速度投影到基底真实运动”的严格投影接口。
\end{definition}

\begin{definition}[定义 D5（投影：几何投影、动力投影与算子投影）]
\label{def:tex29-projection-multi}
在主纤维丛版表述中，“投影”通常至少包含两层：
\begin{enumerate}
  \item \textbf{几何投影}：$\pi:P\to M$（定义~\ref{def:tex29-pb-projection}）。
  \item \textbf{动力/可执行投影}：$a:\mathsf{At}(P)\to TM$（定义~\ref{def:tex29-atiyah-sequence}）。
\end{enumerate}
此外，若需要刻画“算子层面的可观测/隐含分解”，给定线性映射
\[
  \mathrm{pr}:\mathfrak{g}\to\mathfrak{g}_{\mathrm{obs}},
  \qquad
  \ker(\mathrm{pr})=\mathfrak{g}_{\mathrm{hid}},
\]
并记对应的伴随丛为
\[
  \mathrm{ad}_{\mathrm{obs}}(P):=P\times_G\mathfrak{g}_{\mathrm{obs}},
  \qquad
  \mathrm{ad}_{\mathrm{hid}}(P):=P\times_G\mathfrak{g}_{\mathrm{hid}}.
\]
若 $\mathrm{pr}$ 与结构群作用相容（等变），则它诱导出纤维态投影
\[
  \mathrm{pr}_x:\mathrm{ad}(P)_x\to \mathrm{ad}_{\mathrm{obs}}(P)_x.
\]
\end{definition}

\begin{definition}[定义 D6（\GZMonFull{}的主纤维丛模型）]
\label{def:tex29-monad-gz-atiyah}
在定义~\ref{def:tex29-pb-projection}--定义~\ref{def:tex29-projection-multi} 的口径下，
并记该模型为 $\mathsf{RDElem}_{GZ}$（仅作符号标记），
把“\GZMon{}”（非莱布尼茨意义的本体论原子，而是可证书化的工程—形式对象）定义为三元组
\[
  \mathbf{m}:=(x,\alpha_x,A),
\]
其中：
\begin{itemize}
  \item $x\in M$ 为观测核点；
  \item $\alpha_x\in \mathrm{ad}(P)_x$ 为该点处的内部算子态（来自伴随丛纤维）；
  \item $A$ 为连接（connection），等价地给出 Atiyah 序列的一个分裂
  \[
    s_A:TM\to\mathsf{At}(P),\qquad a\circ s_A=\mathrm{id}_{TM}.
  \]
\end{itemize}
为避免 Atiyah 序列的正合性导致 $a\circ\iota=0$ 从而使“可执行投影”恒为零，
本章额外固定一条\textbf{执行编码}（接口）
\[
  \rho_x:\mathrm{ad}(P)_x\to T_xM,
\]
并由此定义一个严格可计算的可执行投影接口：
\[
  \mathcal{V}_A(\alpha_x)
  :=a\bigl(\iota(\alpha_x)+s_A(\rho_x(\alpha_x))\bigr)
  \in T_xM.
\]
由 $a\circ\iota=0$ 与 $a\circ s_A=\mathrm{id}$ 立即得到
\[
  \mathcal{V}_A(\alpha_x)=\rho_x(\alpha_x),
\]
从而 $\mathcal{V}_A$ 把“内部算子态”映射成“基底可观察运动/控制方向”。
\end{definition}

\subsection{定理（证书+证明）}
\label{subsec:tex29-principal-bundle-atiyah-monad-thm}

\begin{theorem}[定理 T1（进一步区分核心：莱布尼茨单子不含可执行锚投影链）]
\label{thm:tex29-leibniz-vs-gz-atiyah}
若“\GZMon{}”对象具备定义~\ref{def:tex29-monad-gz-atiyah} 的结构，
且存在 $\alpha_x\in\mathrm{ad}(P)_x$ 使 $\mathcal{V}_A(\alpha_x)\neq 0$，
则它与定义~\ref{def:tex29-leibniz-monad} 的“无窗单子”在形式系统层面不等价：
莱布尼茨单子的无窗性排除任何“内部算子态 $\Rightarrow$ 外部可执行运动”的可证书化锚投影链条。
\end{theorem}

\begin{Certificate}[证书：由可执行投影推出“窗”，与无窗谓词冲突]
\label{cert:tex29-leibniz-vs-gz-atiyah}
证书登记谓词链：
\[
\begin{aligned}
  T_1 &: \text{(D6) }\exists \alpha_x\in\mathrm{ad}(P)_x\ \bigl(\mathcal{V}_A(\alpha_x)\neq 0\bigr),\\
  T_2 &: T_1\Rightarrow \text{存在“内部态}\ \alpha_x\ \Rightarrow\ \text{观测核可执行效应”}\ \mathcal{V}_A(\alpha_x)\ \text{的结构通道},\\
  T_3 &: \text{(莱布尼茨无窗) }\forall m,\ \neg\exists I\ (I\to m)\quad(\text{定义~\ref{def:tex29-leibniz-monad}}),\\
  T_4 &: T_2\Rightarrow \text{上述结构通道可被视为一类“窗/通道”的形式化实例},\\
  T_5 &: T_3\wedge T_4\Rightarrow \text{与无窗谓词冲突，故莱布尼茨单子与\GZMon{}不等价}.
\end{aligned}
\]
\end{Certificate}

\paragraph{证明（谓词因果链；符号推演）.}
由 $T_1$ 得存在非平凡可执行投影输出；
因此由 $T_2$ 得“内部态$\Rightarrow$外部可执行效应”的结构通道被形式化登记。
由 $T_3$，莱布尼茨无窗谓词排除任何此类通道；
结合 $T_4$ 即得冲突，从而由 $T_5$ 得不等价。证毕。

\subsection{引理（证书+证明）}
\label{subsec:tex29-principal-bundle-atiyah-monad-lem}

\begin{lemma}[引理 L1（广义非交换/非李性 $\Rightarrow$ 顺序与结合方式敏感）]
\label{lem:tex29-noncomm-jacgap-path-order}
在定义~\ref{def:tex29-galgebra-jacobi-gap} 的\KouJing下，
若括号 $[\cdot,\cdot]_{\star}$ 用于编码连接 $A$ 所诱导的“可执行复合动作/平行运输/修复路径”的代数语法，
则：
\begin{enumerate}
  \item 若存在 $x,y\in\mathfrak{g}$ 使 $[x,y]_{\star}\neq -[y,x]_{\star}$，则相应复合动作一般对顺序敏感；
  \item 若存在 $x,y,z\in\mathfrak{g}$ 使 $\mathsf{JacGap}(x,y,z)=J_{\star}(x,y,z)\neq 0$，则三步复合动作一般对结合方式敏感。
\end{enumerate}
\end{lemma}

\begin{Certificate}[证书：非交换给出顺序差异；Jacobiator 给出结合缺口]
\label{cert:tex29-noncomm-jacgap-path-order}
证书登记谓词链：
\[
\begin{aligned}
  L_1 &: \exists x,y\in\mathfrak{g}\ \bigl([x,y]_{\star}\neq -[y,x]_{\star}\bigr),\\
  L_2 &: L_1\Rightarrow \text{交换次序会改变括号表达的代数证书}\Rightarrow \text{复合动作对顺序敏感},\\
  L_3 &: \exists x,y,z\in\mathfrak{g}\ \bigl(J_{\star}(x,y,z)\neq 0\bigr)\ \Leftrightarrow\ \exists x,y,z\
    \bigl(\mathsf{JacGap}(x,y,z)\neq 0\bigr),\\
  L_4 &: L_3\Rightarrow \text{三步嵌套括号的代数证书依赖结合方式}\Rightarrow \text{复合动作对结合方式敏感}.
\end{aligned}
\]
\end{Certificate}

\paragraph{证明（谓词因果链）.}
由 $L_1$，存在交换次序导致括号证书不同的见证对；若括号用于编码复合动作语法，则得到顺序敏感性，得 $L_2$。
由 $L_3$，Jacobiator 非零给出三步嵌套括号的缺口见证，因此结合方式改变会改变证书，得 $L_4$。
引理成立。证毕。

\subsection{命题（证书+证明）}
\label{subsec:tex29-principal-bundle-atiyah-monad-prop}

\begin{proposition}[命题 P1（投影给出两层阴影：观测阴影与控制阴影）]
\label{prop:tex29-two-shadows}
对任意\GZMon{} $\mathbf{m}=(x,\alpha_x,A)$（定义~\ref{def:tex29-monad-gz-atiyah}），存在两种自然的“阴影/投影对象”：
\begin{enumerate}
  \item \textbf{观测阴影}：$x\in M$（由 $\pi$ 的口径对齐给定）；
  \item \textbf{控制阴影}：$\mathcal{V}_A(\alpha_x)\in T_xM$（由可执行投影接口给定）。
\end{enumerate}
此外，在定义~\ref{def:tex29-projection-multi} 的算子投影 $\mathrm{pr}$ 存在时，
还有\textbf{可观测算子阴影}：
\[
  \mathrm{pr}_x(\alpha_x)\in \mathrm{ad}_{\mathrm{obs}}(P)_x,
  \qquad
  \alpha_x=\mathrm{pr}_x(\alpha_x)+\alpha_x^{\mathrm{hid}},\ \alpha_x^{\mathrm{hid}}\in\ker(\mathrm{pr}_x).
\]
\end{proposition}

\begin{Certificate}[证书：像与核分解的标准结论]
\label{cert:tex29-two-shadows}
证书登记谓词链：
\[
\begin{aligned}
  P_1 &: \mathbf{m}=(x,\alpha_x,A)\Rightarrow x\in M\quad(\text{定义~\ref{def:tex29-monad-gz-atiyah}}),\\
  P_2 &: \mathbf{m}=(x,\alpha_x,A)\Rightarrow \mathcal{V}_A(\alpha_x)\in T_xM\quad(\text{定义~\ref{def:tex29-monad-gz-atiyah}}),\\
  P_3 &: \mathrm{pr}_x:\mathrm{ad}(P)_x\to \mathrm{ad}_{\mathrm{obs}}(P)_x\ \text{为线性映射}\Rightarrow \alpha_x=\mathrm{pr}
    _x(\alpha_x)+\alpha_x^{\mathrm{hid}},\\
  P_4 &: P_1\wedge P_2\wedge P_3\Rightarrow \text{三类阴影/投影对象存在且可核验}.
\end{aligned}
\]
\end{Certificate}

\paragraph{证明（谓词因果链）.}
由 $P_1$ 与 $P_2$ 直接得到观测阴影与控制阴影；
由 $P_3$，线性投影给出像/核的标准分解；
合并得 $P_4$，命题成立。证毕。

\subsection{推论（证书+证明）}
\label{subsec:tex29-principal-bundle-atiyah-monad-cor}

\begin{corollary}[推论 C1（核/像刻画隐含修复自由度：莱布尼茨单子不具备的区分量）]
\label{cor:tex29-hidden-repair-dof}
若定义~\ref{def:tex29-projection-multi} 中存在非平凡核
\[
  \mathfrak{g}_{\mathrm{hid}}=\ker(\mathrm{pr})\neq 0,
\]
则对任意 $\alpha_x\in\mathrm{ad}(P)_x$ 的分解
\[
  \alpha_x=\mathrm{pr}_x(\alpha_x)+\alpha_x^{\mathrm{hid}},\qquad \alpha_x^{\mathrm{hid}}\in\ker(\mathrm{pr}_x),
\]
给出一类“对观测不可见但对修复可用”的隐含自由度证书。
该证书在结构上依赖于伴随丛/投影核分解与（可选的）$\mathsf{JacGap}$ 等内部复杂性证书，
属于莱布尼茨无窗单子的形式系统中所没有的区分量。
\end{corollary}

\begin{Certificate}[证书：非平凡核 $\Rightarrow$ 存在不可观测但可执行的分量]
\label{cert:tex29-hidden-repair-dof}
证书登记谓词链：
\[
\begin{aligned}
  C_1 &: \ker(\mathrm{pr})\neq 0\Rightarrow \exists u\in\mathfrak{g}_{\mathrm{hid}}\ (u\neq 0),\\
  C_2 &: C_1\Rightarrow \exists x\in M,\ \exists \alpha_x\in\mathrm{ad}(P)_x\ \text{使得}\ \alpha_x^{\mathrm{hid}}\neq 0,
    \\
  C_3 &: \alpha_x^{\mathrm{hid}}\in\ker(\mathrm{pr}_x)\Rightarrow \text{在算子观测口径下不可见},\\
  C_4 &: \alpha_x^{\mathrm{hid}}\neq 0\Rightarrow \text{在内部算子语法中可参与复合/修复（可作为证书链输入）},\\
  C_5 &: C_2\wedge C_3\wedge C_4\Rightarrow \text{隐含修复自由度证书成立}.
\end{aligned}
\]
\end{Certificate}

\paragraph{证明（谓词因果链）.}
由 $C_1$，非平凡核给出非零隐含方向；
将其提升到伴随丛纤维即可得到 $C_2$ 的见证元。
由 $C_3$，它在投影像中不可见；由 $C_4$，它仍是内部算子代数中的有效输入，因此可进入证书链。
合并得 $C_5$，推论成立。证毕。

\subsection{备注}
\label{subsec:tex29-principal-bundle-atiyah-monad-remark}

\begin{remark}[把区分压缩成一句可直接引用的边界陈述（仅含结构条件）]
若需要一句不含形容词、只含结构条件的“区分句”，可取：
\[
  \boxed{
  \begin{aligned}
    \text{莱布尼茨单子}\ &\equiv\ \text{无窗谓词（定义~\ref{def:tex29-leibniz-monad}）下的内部态系统；}\\
    \text{GZ富结构离散元}\ &\equiv\ (x,\alpha_x,A)\ \text{且}\ \exists \alpha_x\ \bigl(\mathcal{V}_A(\alpha_x)\neq 0\bigr).
  \end{aligned}
  }
\]
以及进一步的“非交换/非李性”强化证书：
\[
  \boxed{
  \begin{aligned}
    &\exists x,y,z\in\mathfrak{g}:\ [x,y]_{\star}\neq -[y,x]_{\star}\ \text{或}\ J_{\star}(x,y,z)\neq 0\\
    &\Rightarrow\ \text{\GZMon{}的修复/控制是路径—顺序—结合方式敏感的可执行结构。}
  \end{aligned}
  }
\]
\end{remark}

\clearpage

%%%%%%%%%%%%%%%%%%%%%%%%%%%%%%%%%%%%%%%%%%%%%%%%%%%%%%%%%%%%
\section{形式逻辑：\GZMon{}—配置—退化投影骨架与主纤维丛增强\GZMonState{}的严格对齐}
\label{sec:tex29-mon-cfg-proj-bundle-alignment}
\label{sec:tex29-ch11}
%%%%%%%%%%%%%%%%%%%%%%%%%%%%%%%%%%%%%%%%%%%%%%%%%%%%%%%%%%%%

\subsection{定义}
\label{subsec:tex29-mon-cfg-proj-bundle-alignment-defs}

\begin{definition}[定义 D0（“\GZMon{}—配置—投影”骨架对齐）]
\label{def:tex29-align-d0-skeleton}
沿用本笔记的对象域与记号：
\[
  \begin{aligned}
    &\mathsf{Mon}\ (\text{\GZMon{}基元集合}),\\
    &X\in\mathsf{Cfg}(\mathsf{Mon})\ (\text{离散构型/状态}),\\
    &\Pi:\mathsf{Cfg}(\mathsf{Mon})\to\mathsf{Cont}\ (\text{连续统表示/退化投影}).
  \end{aligned}
\]
并记观测核（底空间）
\[
  M:=\Pi(X)\in\mathsf{Cont}.
\]
\end{definition}

\begin{definition}[定义 D1（主纤维丛版“算子层”的附加结构）]
\label{def:tex29-align-d1-principal-bundle}
在 $M=\Pi(X)$ 上赋予一条主纤维丛
\[
  \pi:P\to M,
  \qquad
  \text{结构群为 }G.
\]
该结构用于把“可执行算子层”作为观测核之上的可追踪附加数据（接口层）。
\end{definition}

\begin{definition}[定义 D2（广义非交换李代数作为“法/算子”纤维；JacGap 证书）]
\label{def:tex29-align-d2-law-jacgap}
给定一组“可执行算子”的线性空间 $\mathfrak{g}_{\star}$，以及双线性括号
\[
  [\cdot,\cdot]_{\star}:\mathfrak{g}_{\star}\times\mathfrak{g}_{\star}\to\mathfrak{g}_{\star},
\]
不强制其满足雅可比恒等式。
定义雅可比子（Jacobiator）
\[
  J_{\star}(x,y,z):=[x,[y,z]_{\star}]_{\star}+[y,[z,x]_{\star}]_{\star}+[z,[x,y]_{\star}]_{\star}.
\]
为与本笔记常用符号对齐，引入“法对象”（law-object）
\[
  L:=(\mathfrak{g}_{\star},[\cdot,\cdot]_{\star}),
\]
并定义
\[
  \mathsf{JacGap}(L):=J_{\star}
  \qquad
  (\text{作为缺口证书对象；若需数值证书，可另取 }|J_{\star}|).
\]
\end{definition}

\begin{definition}[定义 D3（伴随丛：把 $\mathfrak{g}_{\star}$ 装到 $P$ 上）]
\label{def:tex29-align-d3-adjoint}
在定义~\ref{def:tex29-align-d1-principal-bundle}--定义~\ref{def:tex29-align-d2-law-jacgap} 下，定义
\[
  \mathrm{ad}_{\star}(P):=P\times_G \mathfrak{g}_{\star}.
\]
其纤维元素与截面可理解为：在观测核 $M$ 上可调用的“内部算子态”（算子层状态）。
\end{definition}

\begin{definition}[定义 D4（Atiyah 序列与可执行投影驱动 $\mathcal{V}_A$）]
\label{def:tex29-align-d4-atiyah-drive}
令
\[
  \mathsf{At}(P):=TP/G,
\]
并登记 Atiyah 正合列
\[
  0\longrightarrow \mathrm{ad}_{\star}(P)\xrightarrow{\ \iota\ }\mathsf{At}(P)\xrightarrow{\ a\ }TM\longrightarrow 0.
\]
取连接 $A$（等价于给出分裂 $s_A:TM\to\mathsf{At}(P)$ 且 $a\circ s_A=\mathrm{id}_{TM}$）。

\textbf{注意}：由正合性有 $a\circ\iota=0$，因此若直接写 $a(\iota(\alpha_x))$ 则其值恒为零，无法承载“驱动场”语义。
为得到非平凡且可证书化的“可执行投影驱动”接口，本章额外固定一条执行编码（接口）
\[
  \rho_x:\mathrm{ad}_{\star}(P)_x\to T_xM,
\]
并对任意 $x\in M$ 与 $\alpha_x\in \mathrm{ad}_{\star}(P)_x$ 定义
\[
  \mathcal{V}_A(\alpha_x)
  :=a\bigl(\iota(\alpha_x)+s_A(\rho_x(\alpha_x))\bigr)
  \in T_xM.
\]
由 $a\circ\iota=0$ 与 $a\circ s_A=\mathrm{id}$ 立即得到 $\mathcal{V}_A(\alpha_x)=\rho_x(\alpha_x)$，
从而 $\mathcal{V}_A$ 可被解释为“变分/修复算子态投影到观测核的驱动场”的严格可计算接口。
\end{definition}

\begin{definition}[定义 D5（边缘算子证书链与缺陷势：对齐 $\partial D=\omega$ 与 $\Phi$）]
\label{def:tex29-align-d5-edge-cert-phi}
沿用本笔记的链复形口径：对每层 $n$ 给定
\[
  \omega_n\in\Omega^n(X_n;G),
  \qquad
  D_n\in\mathcal{D}^{n+1}(X_n;G),
  \qquad
  \partial D_n=\omega_n,
  \qquad
  \partial^2=0,
\]
以及缺陷势（defect potential）
\[
  \Phi_n:\Omega^n(X_n;G)\to\RR_{\ge 0}.
\]
\end{definition}

\subsection{定理（证书+证明）}
\label{subsec:tex29-mon-cfg-proj-bundle-alignment-thm}

\begin{theorem}[定理 T1（主纤维丛增强版\GZMonState{}对齐）]
\label{thm:tex29-monad-state-alignment}
在定义~\ref{def:tex29-align-d0-skeleton}--定义~\ref{def:tex29-align-d5-edge-cert-phi} 的对齐口径下，
本笔记中“\GZMonState{}”（可证书化工程—形式对象）可取为如下两级对象：

\paragraph{（I）局部\GZMonState{}（与 $(x,\alpha_x,A)$ 精确对齐）.}
对给定 $X\in\mathsf{Cfg}(\mathsf{Mon})$，令 $M=\Pi(X)$ 并给定 $(\pi:P\to M,A,L)$。
对任意 $x\in M$，定义局部\GZMonState{}
\[
  \mathbf{m}_x:=\bigl(x,\alpha_x,A;L\bigr),
  \qquad
  \alpha_x\in \mathrm{ad}_{\star}(P)_x,
\]
并把其“观测核可执行影子”定义为
\[
  \mathcal{V}_A(\alpha_x)\in T_xM.
\]

\paragraph{（II）全局\GZMonState{}（与“分形生成链/修复塔证书链”拼接）.}
定义全局\GZMonState{}（\GZMonization{}配置）为
\[
  \mathbf{M}:=\Bigl(X,\ M=\Pi(X),\ \pi:P\to M,\ A,\ L,\ \alpha,\ \omega,\ D,\ \Phi\Bigr),
\]
其中：
\begin{itemize}
  \item $\alpha\in\Gamma\bigl(\mathrm{ad}_{\star}(P)\bigr)$（给出每个 $x\in M$ 的 $\alpha_x$）；
  \item $(\omega,D)$ 满足 $\partial D=\omega$（边缘算子证书；见定义~\ref{def:tex29-align-d5-edge-cert-phi}）；
  \item $\Phi$ 为缺陷势（用于变分可接受性/单调性口径）；
  \item $L$ 给出 $\mathsf{JacGap}(L)$（法结构缺口证书；见定义~\ref{def:tex29-align-d2-law-jacgap}）。
\end{itemize}
从而对任意生成链
\[
  \mathbb{F}:=\bigl\{(X_n,\omega_n,D_n)\bigr\}_{n\in\NN},
\]
可自然增强为
\[
  \mathbb{F}^{\sharp}
  :=
  \Bigl\{\bigl(X_n,\Pi(X_n),P_n,A_n,L_n,\alpha_n,\omega_n,D_n,\Phi_n\bigr)\Bigr\}_{n\in\NN},
\]
其中 $(X_n,\omega_n,D_n)$ 保持原样，仅附加“主纤维丛—法对象—可执行投影驱动—JacGap 证书”数据。
\end{theorem}

\begin{Certificate}[证书：直积拼接不改变原证书链，仅附加算子层证书接口]
\label{cert:tex29-monad-state-alignment}
证书登记谓词链：
\[
\begin{aligned}
  A_1 &: X\in\mathsf{Cfg}(\mathsf{Mon})\Rightarrow M=\Pi(X)\quad(\text{定义~\ref{def:tex29-align-d0-skeleton}}),\\
  A_2 &: M=\Pi(X)\Rightarrow \exists(\pi:P\to M)\quad(\text{定义~\ref{def:tex29-align-d1-principal-bundle}}),\\
  A_3 &: L=(\mathfrak{g}_{\star},[\cdot,\cdot]_{\star})\Rightarrow \mathsf{JacGap}(L)=J_{\star}\quad(\text{定义~\ref{def:tex29-align-d2-law-jacgap}}),\\
  A_4 &: \alpha\in\Gamma(\mathrm{ad}_{\star}(P))\Rightarrow \forall x\in M,\ \alpha_x\in \mathrm{ad}_{\star}(P)
    _x\quad(\text{定义~\ref{def:tex29-align-d3-adjoint}}),\\
  A_5 &: A\Rightarrow \mathcal{V}_A:\mathrm{ad}_{\star}(P)_x\to T_xM\quad(\text{定义~\ref{def:tex29-align-d4-atiyah-drive}}
    ),\\
  A_6 &: (\partial D=\omega)\wedge(\partial^2=0)\wedge(\Phi:\Omega^n\to\RR_{\ge 0})\quad(\text{定义~\ref{def:tex29-align-d5-edge-cert-phi}}),\\
  A_7 &: A_1\wedge A_2\wedge A_3\wedge A_4\wedge A_5\wedge A_6\Rightarrow \mathbf{m}_x,\mathbf{M},\mathbb{F}^{\sharp}\
    \text{均为良定义的拼接对象}.
\end{aligned}
\]
\end{Certificate}

\paragraph{证明（谓词因果链；符号推演）.}
由 $A_1$ 得观测核 $M=\Pi(X)$；
由 $A_2,A_3,A_4,A_5$ 得算子层数据 $(P,A,L,\alpha)$ 及其在观测核上的可执行接口 $\mathcal{V}_A$；
由 $A_6$ 得修复证书链数据 $(\omega,D,\Phi)$，且其核心约束 $\partial D=\omega$ 与闭合校验 $\partial^2=0$ 不被改动。
因此把两类数据做直积拼接即可得到局部\GZMonState{} $\mathbf{m}_x$ 与全局\GZMonState{} $\mathbf{M}$；
对链 $\mathbb{F}$ 逐层同样附加算子层数据即可得到增强链 $\mathbb{F}^{\sharp}$。证毕。

\subsection{引理（证书+证明）}
\label{subsec:tex29-mon-cfg-proj-bundle-alignment-lem}

\begin{lemma}[引理 L1（$\mathcal{V}_A$ 的“驱动场”语义不需额外假设）]
\label{lem:tex29-drive-field-no-extra}
在定义~\ref{def:tex29-align-d4-atiyah-drive} 的口径下，
\[
  \mathcal{V}_A(\alpha_x)\in T_xM
\]
本身就是把“内部算子态”送到“观测核上的可执行方向”的映射；
因此把 $\mathcal{V}_A$ 解释为“驱动场/控制方向”的严格表示不需要额外假设。
\end{lemma}

\begin{Certificate}[证书：像空间落在 $T_xM$]
\label{cert:tex29-drive-field-no-extra}
证书登记谓词链：
\[
\begin{aligned}
  D_1 &: \mathcal{V}_A:\mathrm{ad}_{\star}(P)_x\to T_xM\quad(\text{定义~\ref{def:tex29-align-d4-atiyah-drive}}),\\
  D_2 &: D_1\Rightarrow \forall \alpha_x,\ \mathcal{V}_A(\alpha_x)\ \text{是观测核点 }x\text{ 处的切向量},\\
  D_3 &: D_2\Rightarrow \text{“驱动方向/驱动场”的最小形式化语义成立}.
\end{aligned}
\]
\end{Certificate}

\paragraph{证明（谓词因果链）.}
由 $D_1$ 得 $\mathcal{V}_A$ 的像空间为 $T_xM$；
因此任意 $\mathcal{V}_A(\alpha_x)$ 均是观测核上的可执行方向向量，得 $D_2$；
于是“驱动场”的最小语义无需额外假设即成立，得 $D_3$。证毕。

\subsection{命题（证书+证明）}
\label{subsec:tex29-mon-cfg-proj-bundle-alignment-prop}

\begin{proposition}[命题 P1（$\Phi$ 与 $\mathsf{JacGap}(L)$ 的一致性耦合：不改抽象定义，仅给出实例化分解）]
\label{prop:tex29-phi-jacgap-coupling}
在不改变 $\Phi_n:\Omega^n(X_n;G)\to\RR_{\ge 0}$ 的抽象地位前提下，
可给出一个“主纤维丛增强版”的实例化分解：
\[
  \Phi_n(\omega_n):=\Phi_n^{\mathrm{top}}(\omega_n)+\Phi_n^{\mathrm{law}}(L_n),
\]
其中：
\begin{itemize}
  \item $\Phi_n^{\mathrm{top}}$ 仅依赖障碍表达 $\omega_n$（纯拓扑/链复形层）；
  \item $\Phi_n^{\mathrm{law}}$ 仅依赖法对象 $L_n$，并可取为对缺口证书的标量化，例如
  \[
    \Phi_n^{\mathrm{law}}(L_n):=\rho_n\bigl(|\mathsf{JacGap}(L_n)|\bigr),
    \qquad
    \rho_n:\RR_{\ge 0}\to\RR_{\ge 0}\ \text{单调}.
  \]
\end{itemize}
\end{proposition}

\begin{Certificate}[证书：通过构造给出一类满足分解的 $\Phi_n$]
\label{cert:tex29-phi-jacgap-coupling}
证书登记谓词链：
\[
\begin{aligned}
  P_1 &: \rho_n:\RR_{\ge 0}\to\RR_{\ge 0}\ \text{单调}\Rightarrow \Phi_n^{\mathrm{law}}(L_n)\in\RR_{\ge 0},\\
  P_2 &: \Phi_n^{\mathrm{top}}(\omega_n)\in\RR_{\ge 0}\ \wedge\ P_1\Rightarrow \Phi_n(\omega_n)\in\RR_{\ge 0},\\
  P_3 &: P_2\Rightarrow \Phi_n\ \text{是一个与定义~\ref{def:tex29-align-d5-edge-cert-phi} 兼容的缺陷势实例}.
\end{aligned}
\]
\end{Certificate}

\paragraph{证明（谓词因果链）.}
由构造定义 $\Phi_n^{\mathrm{law}}(L_n):=\rho_n(|\mathsf{JacGap}(L_n)|)$，由 $P_1$ 得其非负；
取任意非负的 $\Phi_n^{\mathrm{top}}$，则由 $P_2$ 得 $\Phi_n$ 非负且满足势函数口径；
因此由 $P_3$ 得该分解给出一类可用实例。证毕。

\subsection{推论（证书+证明）}
\label{subsec:tex29-mon-cfg-proj-bundle-alignment-cor}

\begin{corollary}[推论 C1（投影接口 + JacGap 证书：对莱布尼茨单子的进一步区分量）]
\label{cor:tex29-extra-distinguish-vs-leibniz}
对齐后的主纤维丛增强版\GZMonState{}至少携带两类在莱布尼茨单子口径中不出现的可核验结构：
\begin{enumerate}
  \item \textbf{可执行投影接口}：对任意 $x\in M$，存在
  \[
    \mathcal{V}_A:\mathrm{ad}_{\star}(P)_x\to T_xM,
    \qquad
    \alpha_x\mapsto \mathcal{V}_A(\alpha_x),
  \]
  给出“内部算子态 $\Rightarrow$ 观测核可执行方向”的严格通道；
  \item \textbf{法结构缺口证书}：
  \[
    \mathsf{JacGap}(L)=J_{\star},
  \]
  把“可能非李”的结构缺口对象化、证书化，从而为后续的势函数耦合（命题~\ref{prop:tex29-phi-jacgap-coupling}）提供可计算输入。
\end{enumerate}
因此，“可执行性 + 缺口可证书化/可度量性”构成额外区分量；而莱布尼茨单子以无窗/非交互为本体论内核，并不提供上述接口层证书结构。
\end{corollary}

\begin{Certificate}[证书：由 D4 与 D2 直接抽取两类结构]
\label{cert:tex29-extra-distinguish-vs-leibniz}
证书登记谓词链：
\[
\begin{aligned}
  C_1 &: \text{定义~\ref{def:tex29-align-d4-atiyah-drive}}\Rightarrow \mathcal{V}_A:\mathrm{ad}_{\star}(P)_x\to T_xM,\\
  C_2 &: \text{定义~\ref{def:tex29-align-d2-law-jacgap}}\Rightarrow \mathsf{JacGap}(L)=J_{\star},\\
  C_3 &: C_1\wedge C_2\Rightarrow \text{存在“可执行投影接口 + JacGap 证书”两类可核验结构},\\
  C_4 &: \text{定义~\ref{def:tex29-leibniz-monad} 的无窗谓词口径}\Rightarrow \text{不提供上述接口层结构}.
\end{aligned}
\]
\end{Certificate}

\paragraph{证明（谓词因果链）.}
由 $C_1$ 与 $C_2$ 直接得到两类结构接口与证书对象；
由 $C_4$，莱布尼茨无窗口径不以“可执行投影接口/缺口证书对象”作为定义数据；
合并得 $C_3$ 所述区分量成立。证毕。

\subsection{备注}
\label{subsec:tex29-mon-cfg-proj-bundle-alignment-remark}

\begin{remark}[一句对齐版定义（最短且不含隐喻）]
若需要一句可直接写入文稿的对齐定义，可取：
\[
  \boxed{
  \begin{aligned}
    &\text{给定 }X\in\mathsf{Cfg}(\mathsf{Mon}),\ M=\Pi(X),\ \pi:P\to M,\ A,\\
    &\qquad L=(\mathfrak{g}_{\star},[\cdot,\cdot]_{\star});\\
    &\text{则局部\GZMonState{}为 }\mathbf{m}_x=(x,\alpha_x,A;L),\
    \mathcal{V}_A(\alpha_x)\in T_xM,\
    \mathsf{JacGap}(L)=J_{\star}.
  \end{aligned}
  }
\]
\end{remark}

\clearpage

%%%%%%%%%%%%%%%%%%%%%%%%%%%%%%%%%%%%%%%%%%%%%%%%%%%%%%%%%%%%
\section{形式逻辑：因果/可执行范畴口径下的莱布尼茨单子与\GZMonFull{}对比}
\label{sec:tex29-causal-category-contrast}
\label{sec:tex29-ch12}
%%%%%%%%%%%%%%%%%%%%%%%%%%%%%%%%%%%%%%%%%%%%%%%%%%%%%%%%%%%%

\subsection{定义：符号约定与两侧范畴}
\label{subsec:tex29-causal-category-contrast-defs}

\begin{definition}[定义 D0（符号约定与“因果态射”口径）]
\label{def:tex29-causal-notation}
记莱布尼茨单子体系为 $\mathbb{M}_L$，记本笔记体系的\GZMonState{}对象域为 $\mathbb{M}_G$。
下文所有“范畴结构/态射”均指\textbf{因果/可执行}意义上的态射结构：
态射仅代表可执行的影响、变分或控制操作；不把“表象对应/类比翻译/外部描述”混入因果态射之中。
\end{definition}

\subsubsection{莱布尼茨侧}

\begin{definition}[定义 D1（莱布尼茨单子作为内部态系统）]
\label{def:tex29-leibniz-internal-dynamics}
取指标集 $I$。
对每个 $i\in I$，单子 $m_i$ 具有内部状态空间 $\Sigma_i$ 与自治演化（离散或连续均可；此处写成连续口径）：
\[
  \dot{s}_i(t)=F_i\bigl(s_i(t)\bigr),
  \qquad
  s_i(t)\in \Sigma_i.
\]
\end{definition}

\begin{definition}[定义 D2（莱布尼茨“无窗”对应的因果范畴）]
\label{def:tex29-category-cl}
定义因果范畴 $\mathcal{C}_L$：
\begin{itemize}
  \item 对象：$\mathrm{Ob}(\mathcal{C}_L)=\{m_i\}_{i\in I}$；
  \item 态射：$\mathrm{Hom}_{\mathcal{C}_L}(m_i,m_j)$ 表示“从 $m_i$ 到 $m_j$ 的可执行因果影响”（若存在）。
\end{itemize}
\end{definition}

\subsubsection{本笔记体系侧（主纤维丛 + 广义非交换代数投影口径）}

\begin{definition}[定义 D3（主纤维丛与观测核）]
\label{def:tex29-category-cg-bundle}
给定构型 $X\in\mathsf{Cfg}(\mathsf{Mon})$，其观测核（退化投影）为
\[
  M:=\Pi(X)\in\mathsf{Cont}.
\]
在 $M$ 上取主纤维丛
\[
  \pi:P\to M,
  \qquad
  \text{结构群为 }G.
\]
\end{definition}

\begin{definition}[定义 D4（广义非交换代数与 Jacobiator-gap 证书）]
\label{def:tex29-category-cg-jacgap}
给定线性空间 $\mathfrak{g}_{\star}$ 及双线性括号 $[\cdot,\cdot]_{\star}$，
不要求满足雅可比恒等式。
定义雅可比子（Jacobiator）
\[
  J_{\star}(x,y,z):=[x,[y,z]_{\star}]_{\star}+[y,[z,x]_{\star}]_{\star}+[z,[x,y]_{\star}]_{\star},
\]
并引入法对象
\[
  L:=(\mathfrak{g}_{\star},[\cdot,\cdot]_{\star}),
  \qquad
  \mathsf{JacGap}(L):=J_{\star},
\]
作为“非李性缺口”的证书对象（若需数值证书，可另取 $|J_{\star}|$）。
\end{definition}

\begin{definition}[定义 D5（伴随丛与 Atiyah 序列：可执行投影链）]
\label{def:tex29-category-cg-atiyah}
定义伴随丛与 Atiyah 丛：
\[
  \mathrm{ad}_{\star}(P):=P\times_G \mathfrak{g}_{\star},
  \qquad
  \mathsf{At}(P):=TP/G.
\]
存在正合列（Atiyah 序列）
\[
  0\longrightarrow \mathrm{ad}_{\star}(P)\xrightarrow{\ \iota\ }\mathsf{At}(P)\xrightarrow{\ a\ }TM\longrightarrow 0.
\]
给定连接 $A$（等价于给出分裂 $s_A:TM\to\mathsf{At}(P)$ 且 $a\circ s_A=\mathrm{id}_{TM}$）。

\textbf{注意}：由正合性必有 $a\circ\iota=0$，因此若直接写 $a(\iota(\alpha_x))$ 将恒为零。
为得到非平凡且可证书化的“可执行投影驱动”接口，本章固定一条执行编码
\[
  \rho_x:\mathrm{ad}_{\star}(P)_x\to T_xM,
\]
并对任意 $\alpha_x\in\mathrm{ad}_{\star}(P)_x$ 定义
\[
  \mathcal{V}_A(\alpha_x)
  :=a\bigl(\iota(\alpha_x)+s_A(\rho_x(\alpha_x))\bigr)
  \in T_xM.
\]
\end{definition}

\begin{definition}[定义 D6（本笔记体系\GZMon{}对象：局部\GZMonState{}）]
\label{def:tex29-category-cg-monad}
定义（局部）\GZMonState{}为
\[
  m:=(x,\alpha_x,A;L),
  \qquad
  x\in M,\ \alpha_x\in \mathrm{ad}_{\star}(P)_x.
\]
并记 $\mathbb{M}_G$ 为全体此类局部\GZMonState{}的集合（或其带扇区扩充后的集合）。
若需把“扇区/对抗态/环境”纳入对象本身，可扩充为
\[
  m:=(x,\alpha_x,A;L,S),
\]
下文用 $S$ 表示扇区/环境/对手信息。
\end{definition}

\begin{definition}[定义 D7（缺陷势与可接受变分算子）]
\label{def:tex29-category-cg-phi}
给定缺陷势（或势能/目标）
\[
  \Phi:\mathbb{M}_G\times\mathcal{S}\to\RR_{\ge 0}.
\]
称变分算子（或可执行操作）$\delta$ 在 $(m,S)$ 处\textbf{可接受}，若满足
\[
  m'=\delta\cdot m
  \quad\wedge\quad
  \Phi(m',S)<\Phi(m,S)
  \quad\wedge\quad
  \mathsf{RuleCompat}(\delta;m,S),
\]
其中谓词 $\mathsf{RuleCompat}$ 表示“规则一致性证书”的口径占位符（可进一步展开为一组约束）。
\end{definition}

\begin{definition}[定义 D8（本笔记体系因果/可执行范畴）]
\label{def:tex29-category-cg}
为避免把 $S$ 作为外部参数混入 Hom 集，采用扇区扩充对象域：
\[
  \mathbb{M}_G^{\sharp}:=\mathbb{M}_G\times\mathcal{S},
  \qquad
  \bar{m}:=(m,S)\in\mathbb{M}_G^{\sharp}.
\]
定义范畴 $\mathcal{C}_G$：
\begin{itemize}
  \item 对象：$\mathrm{Ob}(\mathcal{C}_G)=\mathbb{M}_G^{\sharp}$；
  \item 态射：对 $\bar{m}=(m,S)$ 与 $\bar{m}'=(m',S')$，
  先定义“可接受”谓词
  \[
  \begin{aligned}
    \mathsf{Acc}(\delta;m,S,m'):\Longleftrightarrow\
    &m'=\delta\cdot m\ \wedge\\
    &\bigl(\delta=\mathrm{id}_m\ \vee\ (\Phi(m',S)<\Phi(m,S)\wedge \mathsf{RuleCompat}(\delta;m,S))\bigr).
  \end{aligned}
  \]
  并据此令
  \[
    \mathrm{Hom}_{\mathcal{C}_G}(\bar{m},\bar{m}')
    :=
    \begin{cases}
      \bigl\{\delta\ \big|\ \mathsf{Acc}(\delta;m,S,m')\bigr\}, & S'=S,\\[2pt]
      \varnothing, & S'\neq S;
    \end{cases}
  \]
  \item 合成：若 $\delta_1:\bar{m}\to \bar{m}_1$ 且 $\delta_2:\bar{m}_1\to \bar{m}_2$，则
  \[
    \delta_2\circ \delta_1:\bar{m}\to \bar{m}_2,
    \qquad
    (\delta_2\circ\delta_1)\cdot m:=\delta_2\cdot(\delta_1\cdot m);
  \]
  \item 恒等：$\mathrm{id}_{\bar{m}}:=\mathrm{id}_m:\bar{m}\to\bar{m}$，满足 $\mathrm{id}_m\cdot m=m$。
\end{itemize}
\end{definition}

\subsection{公理降级为定理（证书+证明）}
\label{subsec:tex29-causal-category-contrast-axioms}

\begin{theorem}[公理降级：无窗公理 $\Rightarrow$ $\mathcal{C}_L$ 为离散范畴（证书化\KouJing）]
\label{thm:tex29-axL-downgraded-discrete}
登记“无窗公理”（因果态射为空）：
\[
  \forall i\neq j,\ \mathrm{Hom}_{\mathcal{C}_L}(m_i,m_j)=\varnothing,
  \qquad
  \forall i,\ \mathrm{Hom}_{\mathcal{C}_L}(m_i,m_i)=\{\mathrm{id}_{m_i}\}.
\]
则 $\mathcal{C}_L$ 为离散范畴，且有范畴同构
\[
  \mathcal{C}_L\cong \bigsqcup_{i\in I}\mathbf{1},
\]
其中 $\mathbf{1}$ 表示仅含一个对象与一个态射的终对象范畴。
\end{theorem}

\begin{Certificate}[证书：离散范畴的定义性判据]
\label{cert:tex29-axL-downgraded-discrete}
证书登记谓词链：
\[
\begin{aligned}
  L_1 &: (\forall i\neq j)\ \mathrm{Hom}(m_i,m_j)=\varnothing,\\
  L_2 &: (\forall i)\ \mathrm{Hom}(m_i,m_i)=\{\mathrm{id}_{m_i}\},\\
  L_3 &: L_1\wedge L_2\Rightarrow \mathcal{C}_L\ \text{仅含恒等态射}\Rightarrow \mathcal{C}_L\ \text{离散},\\
  L_4 &: L_3\Rightarrow \mathcal{C}_L\cong \bigsqcup_{i\in I}\mathbf{1}.
\end{aligned}
\]
\end{Certificate}

\paragraph{证明（谓词因果链）.}
由 $L_1,L_2$ 知：不同对象间无态射，同一对象的态射仅为恒等态射；
这正是离散范畴的定义性判据，得 $L_3$；
而“离散范畴 $\cong$ 若干个 $\mathbf{1}$ 的不交并”亦为标准等价刻画，得 $L_4$。证毕。

\begin{theorem}[公理降级：可接受变分闭包的证书化充分条件（数学归纳法证书）]
\label{thm:tex29-axG-downgraded-closure}
固定扇区 $S$。
若对任意可接受变分 $\delta:m\to m'$（定义~\ref{def:tex29-category-cg-phi}）均有严格下降链
\[
  \Phi(m',S)<\Phi(m,S),
\]
且规则兼容谓词对合成封闭：
\[
  \mathsf{RuleCompat}(\delta_1;m,S)\wedge \mathsf{RuleCompat}(\delta_2;m_1,S)
  \Rightarrow
  \mathsf{RuleCompat}(\delta_2\circ\delta_1;m,S),
\]
则任意有限可接受变分链在合成下仍可接受；特别地，两步合成闭包公理成立：
若 $\delta_1:m\to m_1$ 与 $\delta_2:m_1\to m_2$ 均可接受，则 $\delta_2\circ\delta_1:m\to m_2$ 可接受。
\end{theorem}

\begin{Certificate}[数学归纳法证书：按态射链长度 $N$ 的严格下降与规则封闭]
\label{cert:tex29-axG-downgraded-closure}
令谓词 $P(N)$ 表示：
“对任意长度为 $N$ 的可接受变分链
\[
  m_0\xrightarrow{\delta_1} m_1\xrightarrow{\delta_2}\cdots\xrightarrow{\delta_N} m_N
  \qquad(\text{均在同一扇区 }S),
\]
其合成 $\delta_N\circ\cdots\circ\delta_1$ 仍为可接受变分，即
\[
  \Phi(m_N,S)<\Phi(m_0,S)\ \wedge\ \mathsf{RuleCompat}(\delta_N\circ\cdots\circ\delta_1;m_0,S).
\]
”
证书化要求为：
\[
  P(1)\ \text{成立},\qquad (\forall N\in\NN)\ \bigl(P(N)\Rightarrow P(N+1)\bigr).
\]
\end{Certificate}

\paragraph{证明（数学归纳法证书；谓词因果链）.}
由证书~\ref{cert:tex29-axG-downgraded-closure}：
\begin{itemize}
  \item 基步：$N=1$ 时链为 $m_0\xrightarrow{\delta_1}m_1$。
  由“可接受”的定义直接得 $\Phi(m_1,S)<\Phi(m_0,S)$ 且 $\mathsf{RuleCompat}(\delta_1;m_0,S)$，故 $P(1)$ 成立；
  \item 归纳步：假设 $P(N)$ 成立。
  给定长度 $N+1$ 的可接受链 $m_0\xrightarrow{\delta_1}\cdots\xrightarrow{\delta_{N+1}}m_{N+1}$。
  由归纳假设得
  \[
    \Phi(m_N,S)<\Phi(m_0,S),\qquad \mathsf{RuleCompat}(\delta_N\circ\cdots\circ\delta_1;m_0,S).
  \]
  又由 $\delta_{N+1}$ 可接受得 $\Phi(m_{N+1},S)<\Phi(m_N,S)$ 与 $\mathsf{RuleCompat}(\delta_{N+1};m_N,S)$。
  合并不等式链得
  \[
    \Phi(m_{N+1},S)<\Phi(m_N,S)<\Phi(m_0,S),
  \]
  且由规则兼容对合成封闭得
  \[
    \mathsf{RuleCompat}\bigl(\delta_{N+1}\circ(\delta_N\circ\cdots\circ\delta_1);m_0,S\bigr),
  \]
  即 $P(N+1)$ 成立。
\end{itemize}
因此 $P(N)$ 对全体 $N$ 成立，有限合成闭包成立；取 $N=2$ 即得到两步闭包公理。证毕。

\subsection{定理（证书+证明）}
\label{subsec:tex29-causal-category-contrast-theorems}

\begin{theorem}[定理 T1（范畴结构严格对比：离散 vs 富结构）]
\label{thm:tex29-category-discrete-vs-rich}
在定义~\ref{def:tex29-category-cl} 与定义~\ref{def:tex29-category-cg} 的语义下：
\[
  \mathcal{C}_L\cong \bigsqcup_{i\in I}\mathbf{1},
  \qquad
  \text{而}\qquad
  \mathcal{C}_G\ \text{一般非离散（存在非恒等可执行态射）}.
\]
\end{theorem}

\begin{Certificate}[证书：离散性来自无窗；非离散性来自存在严格下降的可接受变分]
\label{cert:tex29-category-discrete-vs-rich}
证书登记谓词链：
\[
\begin{aligned}
  T_1 &: \text{定理~\ref{thm:tex29-axL-downgraded-discrete}}\Rightarrow \mathcal{C}_L\ \text{离散},\\
  T_2 &: \exists (m,S),\ \exists \delta\neq \mathrm{id}_m\ \text{使}\ m'=\delta\cdot m\ \wedge\ \Phi(m',S)<\Phi(m,S)\
    \wedge\ \mathsf{RuleCompat}(\delta;m,S),\\
  T_3 &: T_2\Rightarrow \mathrm{Hom}_{\mathcal{C}_G}\bigl((m,S),(m',S)\bigr)\neq\varnothing\ \text{且含非恒等态射},\\
  T_4 &: T_3\Rightarrow \mathcal{C}_G\ \text{非离散}.
\end{aligned}
\]
\end{Certificate}

\paragraph{证明（谓词因果链）.}
由 $T_1$ 得 $\mathcal{C}_L$ 离散。
对 $\mathcal{C}_G$，若存在 $T_2$ 所述的可接受变分，则由定义~\ref{def:tex29-category-cg} 得存在非恒等态射，
因此由 $T_4$ 得 $\mathcal{C}_G$ 非离散。证毕。

\begin{theorem}[定理 T2（“$\mathcal{C}_L\subset \mathcal{C}_G$”的严格化口径：等价于一个离散全子范畴）]
\label{thm:tex29-category-fully-faithful-embed}
“$\mathcal{C}_L\subset \mathcal{C}_G$”的严格表述不是集合包含，而是存在全嵌入函子（fully faithful）
\[
  J:\mathcal{C}_L\hookrightarrow \mathcal{C}_G
\]
使得 $\mathrm{Im}(J)$ 是 $\mathcal{C}_G$ 的一个\textbf{离散全子范畴}（仅含恒等态射）。
\end{theorem}

\begin{Certificate}[证书（构造）：把每个 $m_i$ 映到“无可接受变分”的冻结对象]
\label{cert:tex29-category-fully-faithful-embed}
证书登记谓词链：
\[
\begin{aligned}
  E_1 &: \text{对每个 }m_i\in\mathrm{Ob}(\mathcal{C}_L)\ \text{选择 }J(m_i)=(m_i^G,S_0)\in\mathrm{Ob}(\mathcal{C}_G),\\
  E_2 &: \forall i,\ \forall \delta\neq \mathrm{id}_{m_i^G},\ \neg\bigl(\Phi(\delta\cdot m_i^G,S_0)<\Phi(m_i^G,S_0)\bigr)
    ,\\
  E_3 &: E_2\Rightarrow \mathrm{Hom}_{\mathcal{C}_G}\bigl(J(m_i),J(m_j)\bigr)=\varnothing\ (i\neq j),\quad
        \mathrm{Hom}_{\mathcal{C}_G}\bigl(J(m_i),J(m_i)\bigr)=\{\mathrm{id}\},\\
  E_4 &: \text{令 }J(\mathrm{id}_{m_i}):=\mathrm{id}_{J(m_i)}\Rightarrow J\ \text{保持恒等与合成},\\
  E_5 &: E_3\wedge E_4\Rightarrow J\ \text{在每个 Hom 集上给出双射}\Rightarrow J\ \text{fully faithful}.
\end{aligned}
\]
\end{Certificate}

\paragraph{证明（谓词因果链）.}
由 $E_1$ 定义 $J$ 在对象上的取值；由于 $\mathcal{C}_L$ 仅含恒等态射，
由 $E_4$ 定义 $J$ 在态射上的作用即可。
由 $E_2$，除恒等外不存在严格下降的可接受变分；结合定义~\ref{def:tex29-category-cg} 得 $E_3$，
从而 $\mathrm{Im}(J)$ 为离散且为全子范畴。
又由于两侧 Hom 集同为“不同对象为空、同一对象为单点”，故由 $E_5$ 得 $J$ fully faithful。证毕。

\subsection{引理（证书+证明）}
\label{subsec:tex29-causal-category-contrast-lemmas}

\begin{lemma}[引理 L1（动力系统严格对比：直积自治 vs 开放耦合）]
\label{lem:tex29-dynamics-product-vs-coupled}
在“把各单子状态拼成总状态向量”的形式化下：
\begin{enumerate}
  \item （莱布尼茨侧）取总态
  \[
    S(t):=(s_i(t))_{i\in I},
    \qquad
    \dot{s}_i=F_i(s_i),
  \]
  则总向量场为分块对角：对任意 $i\neq j$，
  \[
    \frac{\partial \dot{s}_i}{\partial s_j}=0;
  \]
  \item （本笔记体系侧；示例口径）若取总态 $m(t)\in\mathbb{M}_G$ 并采用开放系统梯度流模型
  \[
    \dot{m}(t)=-\nabla_m\Phi(m(t),S(t)),
  \]
  则一般存在 $i\neq j$ 使交叉影响非零（等价于 $\Phi$ 含耦合项）：
  \[
    \exists i\neq j:\quad
    \frac{\partial}{\partial m_j}\bigl(\nabla_{m_i}\Phi(m,S)\bigr)\neq 0.
  \]
\end{enumerate}
\end{lemma}

\begin{Certificate}[证书：自治系统右端仅依赖自身；耦合势函数导出交叉项]
\label{cert:tex29-dynamics-product-vs-coupled}
证书登记谓词链：
\[
\begin{aligned}
  D_1 &: \dot{s}_i=F_i(s_i)\Rightarrow \dot{s}_i\ \text{不显含}\ s_j\ (j\neq i),\\
  D_2 &: D_1\Rightarrow \frac{\partial \dot{s}_i}{\partial s_j}=0\ (i\neq j),\\
  D_3 &: \Phi(m,S)\ \text{含耦合项}\Rightarrow \exists i\neq j,\ \nabla_{m_i}\Phi\ \text{依赖}\ m_j,\\
  D_4 &: D_3\Rightarrow \exists i\neq j,\ \frac{\partial}{\partial m_j}(\nabla_{m_i}\Phi)\neq 0.
\end{aligned}
\]
\end{Certificate}

\paragraph{证明（谓词因果链）.}
由 $D_1$ 得自治系统的右端仅依赖自身状态，从而由 $D_2$ 得交叉导数为零。
若 $\Phi$ 含耦合项，则由 $D_3$ 得梯度分量依赖他变量，进而由 $D_4$ 得交叉项一般非零。证毕。

\subsection{命题（证书+证明）}
\label{subsec:tex29-causal-category-contrast-props}

\begin{proposition}[命题 P1（“广义分形/修复塔”差异的严格化：映射型一致性约束 vs 生成型递归算子）]
\label{prop:tex29-mapping-vs-generating}
为避免把“表象宇宙”误写为集合相等，采用两侧最小的可形式化结构：
\begin{enumerate}
  \item （莱布尼茨侧）存在表象映射族（仅表述一致性，不生成新单子）
  \[
    \mathrm{Per}_i:\mathcal{U}\to \Sigma_i,
  \]
  并存在全局一致性约束（预成和谐的形式化占位符）
  \[
    \exists H\ \text{s.t.}\ \forall i,j,t,\ \mathcal{R}\bigl(\mathrm{Per}_i(\mathcal{U};t),\mathrm{Per}_j(\mathcal{U};t)
      \bigr)\ \text{由 }H\text{保证};
  \]
  \item （本笔记体系侧）存在递归生成算子（修复塔/分形生成）
  \[
    \mathcal{C}:(X_k,\omega_k,D_k,\dots)\longmapsto (X_{k+1},\omega_{k+1},D_{k+1},\dots),
  \]
  且逐层携带证书条件（例如 $\partial D_k=\omega_k$ 与层间同伦等价约束），因此属于\textbf{生成型}结构。
\end{enumerate}
因此两侧差异可被严格化为：莱布尼茨侧为“映射型一致性约束”，本笔记体系侧为“生成型递归构造 + 证书链”。
\end{proposition}

\begin{Certificate}[证书：映射族不产生新对象；递归算子显式生成新层数据]
\label{cert:tex29-mapping-vs-generating}
证书登记谓词链：
\[
\begin{aligned}
  P_1 &: \mathrm{Per}_i:\mathcal{U}\to \Sigma_i\Rightarrow \text{给定输入仅得到表象值，不生成新的对象层数据},\\
  P_2 &: \exists H\ \text{s.t.}\ \forall i,j,t,\ \mathcal{R}(\cdot,\cdot)\ \text{由 }H\text{保证}\Rightarrow \text{一致性由约束承担}
    ,\\
  P_3 &: \mathcal{C}:(X_k,\omega_k,D_k,\dots)\mapsto (X_{k+1},\omega_{k+1},D_{k+1},\dots)\Rightarrow \text{显式生成新层数据对象},
    \\
  P_4 &: P_3\wedge (\partial D_k=\omega_k)\Rightarrow \text{生成过程逐层携带证书链},\\
  P_5 &: P_1\wedge P_2\wedge P_4\Rightarrow \text{两侧差异为“映射型约束” vs “生成型构造+证书链”}.
\end{aligned}
\]
\end{Certificate}

\paragraph{证明（谓词因果链）.}
由 $P_1,P_2$ 得莱布尼茨侧结构以映射与一致性约束为核心，不以“生成新层对象”为语义；
由 $P_3,P_4$ 得本笔记体系侧结构以递归生成与证书链为核心；
合并得 $P_5$，命题成立。证毕。

\subsection{推论（证书+证明）}
\label{subsec:tex29-causal-category-contrast-cor}

\begin{corollary}[推论 C1（可达集严格对比：离散范畴可达性平凡 vs 变分闭包下的可搜索）]
\label{cor:tex29-reachability-contrast}
对任意范畴 $\mathcal{C}$ 与对象 $x_0\in\mathrm{Ob}(\mathcal{C})$，定义可达集
\[
  \mathrm{Reach}_{\mathcal{C}}(x_0)
  :=\bigl\{x\in\mathrm{Ob}(\mathcal{C})\ \big|\ \exists n,\ \exists(\delta_1,\dots,\delta_n),\ x_0\xrightarrow{\delta_1}
    \cdots\xrightarrow{\delta_n}x\bigr\}.
\]
则：
\[
  \mathrm{Reach}_{\mathcal{C}_L}(m_0)=\{m_0\},
\]
并且在 $\mathcal{C}_G$ 中，从同一扇区内的初值 $\bar{m}_0=(m_0,S)$ 出发，
\[
  \mathrm{Reach}_{\mathcal{C}_G}(\bar{m}_0)
  =
  \bigl\{\bar{m}\ \big|\ \exists n,\ \exists(\delta_1,\dots,\delta_n),\
  \bar{m}_0\xrightarrow{\delta_1}\cdots\xrightarrow{\delta_n}\bar{m},\
  \Phi\ \text{沿链严格下降且规则兼容}\bigr\},
\]
从而 $\mathrm{Reach}_{\mathcal{C}_G}(\bar{m}_0)$ 可被理解为“由可接受变分链生成的可搜索对象集”（其与同伦类/轨道类的关系取决于对象层的等价定义）。
\end{corollary}

\begin{Certificate}[证书：离散范畴仅有恒等链；变分范畴存在严格下降链]
\label{cert:tex29-reachability-contrast}
证书登记谓词链：
\[
\begin{aligned}
  R_1 &: \mathcal{C}_L\ \text{离散}\Rightarrow \text{除恒等外无态射链},\\
  R_2 &: R_1\Rightarrow \mathrm{Reach}_{\mathcal{C}_L}(m_0)=\{m_0\},\\
  R_3 &: \text{定理~\ref{thm:tex29-axG-downgraded-closure}}\Rightarrow \text{可接受变分链对合成封闭},\\
  R_4 &: R_3\Rightarrow \text{任意严格下降且规则兼容的链给出一条可达路径},\\
  R_5 &: R_2\wedge R_4\Rightarrow \text{推论结论成立}.
\end{aligned}
\]
\end{Certificate}

\paragraph{证明（谓词因果链）.}
由 $R_1$ 得 $\mathcal{C}_L$ 中任意态射链只能由恒等组成，故可达集仅含自身，得 $R_2$。
对 $\mathcal{C}_G$，由 $R_3$ 得严格下降且规则兼容的可接受链可合成为可接受变分，
因此链终点对象可达，得 $R_4$；合并得 $R_5$。证毕。

\subsection{备注}
\label{subsec:tex29-causal-category-contrast-remark}

\begin{remark}[把原稿逐项修正为严格表达的三处关键点]
本章的严格化可压缩为三条写作边界：
\begin{enumerate}
  \item 把“$\mathcal{C}_L\subset\mathcal{C}_G$”改写为“存在 fully faithful 嵌入 $J$，其像为离散全子范畴”
  （定理~\ref{thm:tex29-category-fully-faithful-embed}）；
  \item 把 $\mathrm{Perception}(m_i)=\mathcal{U}$ 改写为 $\mathrm{Per}_i:\mathcal{U}\to\Sigma_i$，
  避免把“表象宇宙”误写成集合相等（命题~\ref{prop:tex29-mapping-vs-generating}）；
  \item 把态射条件显式写成“可接受 $\Leftrightarrow$ 势严格下降 $\wedge$ 规则兼容”，
  以便闭包、公理降级与可达集都能被证书化为谓词链（定义~\ref{def:tex29-category-cg-phi}--定理~\ref{thm:tex29-axG-downgraded-closure}）。
\end{enumerate}
\end{remark}

\clearpage

%%%%%%%%%%%%%%%%%%%%%%%%%%%%%%%%%%%%%%%%%%%%%%%%%%%%%%%%%%%%
\section{形式逻辑：\GZMonFull{}（TOP 记号）的公理化定义、ZF 超限递归证书与非交换闭包}
\label{sec:tex29-top}
\label{sec:tex29-ch13}
%%%%%%%%%%%%%%%%%%%%%%%%%%%%%%%%%%%%%%%%%%%%%%%%%%%%%%%%%%%%

\subsection{定义}
\label{subsec:tex29-top-defs}

\begin{definition}[定义 D0（序数、后继与极限）]
\label{def:tex29-top-ordinal}
令 $\mathrm{Ord}$ 表示序数类。对任意 $\alpha\in\mathrm{Ord}$，其后继序数定义为
\[
  \alpha^{+}:=\alpha\cup\{\alpha\}.
\]
称 $\lambda\in\mathrm{Ord}$ 为\textbf{极限序数}，若 $\lambda\neq 0$ 且
\[
  (\forall \beta<\lambda)\ \beta^{+}<\lambda.
\]
\end{definition}

\begin{definition}[定义 D1（\GZMonFull{}第一性对象域与构型空间）]
\label{def:tex29-top-mon-cfg}
令 $\mathsf{Mon}$ 为\GZMon{}基元的集合（本笔记中的“\GZMon{}基元域”）。
令 $\mathsf{Cfg}(\mathsf{Mon})$ 为由 $\mathsf{Mon}$ 生成的离散构型空间（例如有限多重集、离散复形、可数图之同伦类等；本章不固定具体实现，仅要求其可作为算子的作用域）。
\end{definition}

\begin{definition}[定义 D2（边缘算子与修复谓词）]
\label{def:tex29-top-boundary-rep}
对任意构型 $X\in\mathsf{Cfg}(\mathsf{Mon})$、系数对象 $G$ 与阶数 $n$，给定：
\begin{itemize}
  \item 障碍表达空间 $\Omega^n(X;G)$；
  \item 修复链空间 $\mathcal{D}^{n+1}(X;G)$；
  \item 边缘算子（证书接口）
  \[
    \partial:\mathcal{D}^{n+1}(X;G)\to \Omega^n(X;G).
  \]
\end{itemize}
定义修复谓词
\[
  \mathsf{Rep}(\omega)\ :\Longleftrightarrow\ \exists D\in\mathcal{D}^{n+1}(X;G)\ (\partial D=\omega).
\]
即给出 $D$ 便是修复成立的可核验证据。
\end{definition}

\begin{definition}[定义 D3（允许算子集、非交换性与交换子）]
\label{def:tex29-top-operators}
令 $\mathcal{O}\subseteq \mathrm{End}\bigl(\mathsf{Cfg}(\mathsf{Mon})\bigr)$ 为允许算子集合（可包含变分/修复/扇区切换等算子）。
若 $\mathcal{O}$ 对加法与数乘封闭，则将其视为算子代数。
在可作差的语境下，定义交换子（对易子）
\[
  [\delta_1,\delta_2]:=\delta_1\circ\delta_2-\delta_2\circ\delta_1.
\]
当存在 $\delta_1,\delta_2\in\mathcal{O}$ 使 $[\delta_1,\delta_2]\neq 0$ 时，称该结构\textbf{非交换}（非对易）。
\end{definition}

\begin{definition}[定义 D4（退化原子投影：原子化与交换化的可检验条件）]
\label{def:tex29-top-degeneration-projection}
引入一个“原子化（静态核）投影”
\[
  \Pi:\mathsf{Cfg}(\mathsf{Mon})\to \mathsf{Cont},
\]
其中 $\mathsf{Cont}$ 表示连续统/原子化点集等静态核对象域（用于承接静态冻结投影的像域）。
本章将“退化原子性”登记为两条可检验条件：
\begin{enumerate}
  \item \textbf{基元退化为原子：}对任意 $\mu\in\mathsf{Mon}$（将 $\mu$ 视为由其生成的最小构型，如 $\{\mu\}$），
  \[
    \Pi(\mu)\ \text{是原子（点式不可再分单位）};
  \]
  \item \textbf{非交换在投影下交换化：}对任意 $\delta_1,\delta_2\in\mathcal{O}$，
  \[
    \Pi\circ\delta_1\circ\delta_2=\Pi\circ\delta_2\circ\delta_1,
  \]
  等价地（在可作差语境下）
  \[
    \Pi\circ[\delta_1,\delta_2]=0.
  \]
\end{enumerate}
\end{definition}

\begin{definition}[定义 D5（\GZMonFull{}（TOP 记号）：Transfinite Operator Primitive）]
\label{def:tex29-top-primitive}
称 $\mu\in\mathsf{Mon}$ 为一个\textbf{\GZMon{}}（并记 $\mu\in\mathsf{TOP}$ 以兼容既有符号），
若存在一组见证数据 $(\mathcal{O},\partial,\Pi,F)$（其中 $F$ 为递归更新规则），使得：
\begin{enumerate}
  \item \textbf{(T1) 原基最小性（可操作意义下的不可再分）：}
  不存在真子对象域 $\mathsf{Mon}'\subsetneq\mathsf{Mon}$ 使 $\mu\in\mathsf{Mon}'$，
  且 $\mathsf{Cfg}(\mathsf{Mon}')$ 对 $\mathcal{O}$ 的作用封闭并能承载与 $\mathsf{Mon}$ 上同口径的修复谓词 $\mathsf{Rep}$（即缩小对象域会丢失可操作/可修复能力）。

  \item \textbf{(T2) 超限递归生成性（核心）：}
  对任意序数 $\theta$，存在一条长度为 $\theta$ 的“修复—生成链”
  \[
    \bigl\{(X_\alpha,\omega_\alpha,D_\alpha)\bigr\}_{\alpha<\theta}
  \]
  满足：
  \begin{itemize}
    \item 初始层 $X_0$ 由 $\mu$ 生成（例如 $X_0=\{\mu\}$ 或本章选定的初始构型）；
    \item 对每个 $\alpha<\theta$，有 $\omega_\alpha\in\Omega^n(X_\alpha;G)$ 且存在证据链 $D_\alpha$ 使 $\partial
      D_\alpha=\omega_\alpha$；
    \item 后继步：对任意 $\alpha+1<\theta$，
    \[
      (X_{\alpha^{+}},\omega_{\alpha^{+}},D_{\alpha^{+}})
      =F\bigl(\alpha, X_\alpha,\omega_\alpha,D_\alpha\bigr),
    \]
    其中 $F$ 被解释为“障碍驱动 + 变分/扇区”更新规则的形式接口；
    \item 极限步：若 $\lambda<\theta$ 为极限序数，则
    \[
      X_\lambda=\bigcup_{\beta<\lambda}X_\beta.
    \]
  \end{itemize}
  \item \textbf{(T3) 退化原子投影一致性：}
  $\Pi$ 满足定义~\ref{def:tex29-top-degeneration-projection} 的两条退化原子性条件。
\end{enumerate}
直观上，(T1)+(T2)+(T3) 把如下三点同时编码为可证书化条件：
\[
  \text{最小可操作单位}\ +\ \text{超限递归（沿序数）}\ +\ \text{算子非交换动力}\ \Rightarrow\ \Pi\text{ 下退化为原子并交换化}.
\]
\end{definition}

\subsection{公理（底座声明）}
\label{subsec:tex29-top-axioms}

\begin{theorem}[公理 A0（ZF 底座：用于超限递归的集合论可检验接口）]
\label{thm:tex29-top-zf-base}
本章在 ZF（Zermelo--Fraenkel）集合论底座上工作：
外延、配对、并集、幂集、无穷、分离、替换、正则（以及如需可加选择公理 AC）。
其中用于“超限递归”落地为集合对象的关键点是：\textbf{替换公理}与\textbf{并集公理}（并配合正则/分离等保证构造按序数良序推进且可截取所需谓词子集）。
\end{theorem}

\begin{Certificate}[基础底座证书：集合论公理输入接口]
\label{cert:tex29-top-zf-base}
本条用于固定元理论底座：后续所有“超限递归存在唯一性”推演均以 ZF（以及如需可加选择公理 AC）的公理模式为前提输入。
其可核验内容是：本章后续构造只调用如下接口：
外延、配对、并集、幂集、无穷、分离、替换、正则（以及如需 AC），
不引入超出该接口的额外集合论假设。
\end{Certificate}

\paragraph{证明（元理论接口核验；谓词因果链）.}
检查定理~\ref{thm:tex29-top-transfinite-recursion} 的构造步骤：
后继步把先前数据映射为下一阶段对象，使用替换公理；
极限步对阶段对象取并，使用并集公理；
对“所需谓词子集”的截取使用分离公理；
良基性与“按序数良序推进”的截断使用正则公理；
其余公理仅用于构造所需基本集合对象。
因此“后续推演所用接口 $\subseteq$ ZF（以及如需 AC）”成立。证毕。

\subsection{定理：ZF 下的超限递归存在唯一性（证书化版本）}
\label{subsec:tex29-top-transfinite-recursion}

\begin{theorem}[定理 T1（ZF 下的超限递归存在唯一性：证书化\KouJing）]
\label{thm:tex29-top-transfinite-recursion}
给定任意序数 $\theta$ 与一个确定性的递归规则 $F$（其输入为阶段 $\alpha<\theta$ 与先前数据 $f\upharpoonright \alpha$，输出为下一阶段的三元组），
则存在唯一序列
\[
  \bigl\{(X_\alpha,\omega_\alpha,D_\alpha)\bigr\}_{\alpha<\theta}
\]
满足：
\begin{enumerate}
  \item 初始层由规则给定；
  \item 后继步由 $F$ 给定；
  \item 极限步对 $X_\lambda$ 取并：
  \[
    (\lambda\ \text{极限})\Rightarrow X_\lambda=\bigcup_{\beta<\lambda}X_\beta,
  \]
  并且 $(\omega_\lambda,D_\lambda)$ 亦由同一规则在阶段 $\lambda$ 的输入确定。
\end{enumerate}
\end{theorem}

\begin{Certificate}[证书：替换 + 并集 + 正则 + 分离 + 超限归纳]
\label{cert:tex29-top-transfinite-recursion}
证书登记（使用的 ZF 公理点与推演链）：
\[
\begin{aligned}
  Z_1 &: \text{替换公理（Replacement）：收集“阶段}\ \alpha\mapsto f_\alpha\text{”的像为集合},\\
  Z_2 &: \text{并集公理（Union）：极限阶段取}\ \bigcup_{\beta<\lambda}f_\beta\ \text{仍为集合},\\
  Z_3 &: \text{正则公理（Foundation）：按序数良序推进，避免循环定义},\\
  Z_4 &: \text{分离公理（Separation）：从已构造集合中截取满足谓词的部分},\\
  Z_5 &: Z_1\wedge Z_2\wedge Z_3\wedge Z_4\Rightarrow \text{存在性（构造出 }f_\theta\text{）},\\
  Z_6 &: \text{超限归纳：若两序列在所有 }\beta<\alpha\text{ 相同，则由 }F\text{ 的确定性推出 }\alpha\text{ 相同},\\
  Z_7 &: Z_6\Rightarrow \text{唯一性}.
\end{aligned}
\]
\end{Certificate}

\paragraph{证明（超限归纳法证书；谓词因果链）.}
设 $\mathcal{U}$ 为容纳三元组 $(X,\omega,D)$ 的集合宇宙（例如用有序对嵌套编码三元组）。
对每个 $\alpha\le \theta$，构造“到阶段 $\alpha$ 为止”的部分函数
\[
  f_\alpha:\alpha\to \mathcal{U}.
\]
\begin{itemize}
  \item 起点：$\alpha=0$ 时取空函数 $f_0=\varnothing$。
  \item 后继步：若 $f_\alpha$ 已定义，则令
  \[
    f_{\alpha^{+}}
    :=
    f_\alpha\cup\Bigl\{\bigl(\alpha,\ F(\alpha,f_\alpha\upharpoonright \alpha)\bigr)\Bigr\}.
  \]
  该构造由配对与并集等基本公理给出（并可用分离截取函数图像性质）。
  \item 极限步：若 $\lambda\le\theta$ 为极限序数，定义
  \[
    f_\lambda:=\bigcup_{\beta<\lambda} f_\beta.
  \]
  由并集公理得右端为集合；又因递归构造保证 $\beta<\gamma\Rightarrow f_\beta\subseteq f_\gamma$，故并集仍为函数且 $\mathrm{dom}(f_\lambda)=\lambda$。
\end{itemize}
由替换公理可将族 $\{f_\beta\}_{\beta<\theta}$ 收集为集合并逐步推进，得到 $f_\theta$ 存在，即得所需序列的存在性（$Z_5$）。

唯一性：设 $g_\theta$ 亦满足相同递归规则。对任意 $\alpha<\theta$，作超限归纳。
若对所有 $\beta<\alpha$ 有 $f_\theta(\beta)=g_\theta(\beta)$，则由 $F$ 的确定性得
\[
  f_\theta(\alpha)=F(\alpha,f_\theta\upharpoonright \alpha)=F(\alpha,g_\theta\upharpoonright \alpha)=g_\theta(\alpha),
\]
故对全体 $\alpha$ 成立，从而 $f_\theta=g_\theta$，唯一性成立（$Z_7$）。证毕。

\subsection{定理：非交换算子代数的最小闭包（最小 *-子代数）}
\label{subsec:tex29-top-minimal-closure}

\begin{lemma}[引理 L1（任意 *-子代数族的交仍为 *-子代数）]
\label{lem:tex29-top-intersection-subalgebra}
设 $\mathcal{A}$ 为一个算子 $*$-代数（对 $+,\cdot,*$ 封闭）。
若 $\{\mathcal{B}_i\}_{i\in I}$ 为 $\mathcal{A}$ 的一族 $*$-子代数，则
\[
  \bigcap_{i\in I}\mathcal{B}_i
\]
仍为 $\mathcal{A}$ 的 $*$-子代数。
\end{lemma}

\begin{Certificate}[证书：逐条封闭性在交中保持]
\label{cert:tex29-top-intersection-subalgebra}
证书登记谓词链：
\[
\begin{aligned}
  I_1 &: x,y\in \bigcap_i\mathcal{B}_i\Rightarrow (\forall i)\ x,y\in\mathcal{B}_i,\\
  I_2 &: (\forall i)\ \mathcal{B}_i\ \text{对 }+,\cdot,*\ \text{封闭}\Rightarrow (\forall i)\ x+y,xy,x^*\in\mathcal{B}_i,\\
  I_3 &: I_1\wedge I_2\Rightarrow x+y,xy,x^*\in \bigcap_i\mathcal{B}_i.
\end{aligned}
\]
\end{Certificate}

\paragraph{证明（谓词因果链）.}
由 $I_1$，交中的元素属于每个 $\mathcal{B}_i$；
由 $I_2$，每个 $\mathcal{B}_i$ 对运算封闭；
故由 $I_3$ 得交对运算封闭，从而仍为 $*$-子代数。证毕。

\begin{theorem}[定理 T2（非交换算子代数的最小闭包存在性：最小 *-子代数）]
\label{thm:tex29-top-minimal-closure}
设 $\mathcal{A}$ 为一个算子 $*$-代数，取任意 $S\subseteq\mathcal{A}$。
定义
\[
  \langle S\rangle_{*}
  :=
  \bigcap\Bigl\{\mathcal{B}\subseteq\mathcal{A}\ \big|\ \mathcal{B}\ \text{是}\ *\text{-子代数且}\ S\subseteq\mathcal{B}
    \Bigr\}.
\]
则 $\langle S\rangle_{*}$ 是包含 $S$ 的\textbf{最小} $*$-子代数（即“最小闭包”）。
\end{theorem}

\begin{Certificate}[证书：非空性 + 交封闭性 + 交的极小性]
\label{cert:tex29-top-minimal-closure}
证书登记谓词链：
\[
\begin{aligned}
  M_1 &: \mathcal{A}\ \text{本身是}\ *\text{-子代数且}\ S\subseteq\mathcal{A}\Rightarrow \text{交族非空},\\
  M_2 &: \text{引理~\ref{lem:tex29-top-intersection-subalgebra}}\Rightarrow \langle S\rangle_{*}\ \text{为}\ *\text{-子代数},
    \\
  M_3 &: \text{由交定义}\Rightarrow S\subseteq \langle S\rangle_{*},\\
  M_4 &: \forall\mathcal{C}\ (\mathcal{C}\ \text{是}\ *\text{-子代数且}\ S\subseteq\mathcal{C})\Rightarrow \langle
    S\rangle_{*}\subseteq \mathcal{C},\\
  M_5 &: M_2\wedge M_3\wedge M_4\Rightarrow \langle S\rangle_{*}\ \text{为包含 }S\text{ 的最小}\ *\text{-子代数}.
\end{aligned}
\]
\end{Certificate}

\paragraph{证明（谓词因果链）.}
由 $M_1$ 得交族非空，故 $\langle S\rangle_{*}$ 良定义。
由 $M_2$ 得其为 $*$-子代数；由 $M_3$ 得其包含 $S$。
任取任意 $*$-子代数 $\mathcal{C}$ 且 $S\subseteq\mathcal{C}$，则 $\mathcal{C}$ 属于交族，故由交定义得 $M_4$。
合并得 $M_5$，定理成立。证毕。

\subsection{命题与推论：退化原子性、交换化与命名对应}
\label{subsec:tex29-top-prop-cor}

\begin{proposition}[命题 P1（退化原子投影必然交换化：投影消去交换子）]
\label{prop:tex29-top-degeneration-commutative}
若投影 $\Pi$ 满足定义~\ref{def:tex29-top-degeneration-projection}(2)，则对任意 $\delta_1,\delta_2\in\mathcal{O}$ 与任意
  $X\in\mathsf{Cfg}(\mathsf{Mon})$，
投影后的可观测序列不区分算子顺序：
\[
  \Pi\bigl(\delta_1(\delta_2(X))\bigr)=\Pi\bigl(\delta_2(\delta_1(X))\bigr),
\]
等价地，
\[
  (\Pi\circ[\delta_1,\delta_2])(X)=0.
\]
\end{proposition}

\begin{Certificate}[证书：对所有 $X$ 的函数等式即为交换化条件]
\label{cert:tex29-top-degeneration-commutative}
证书登记谓词链：
\[
\begin{aligned}
  P_1 &: \Pi\circ\delta_1\circ\delta_2=\Pi\circ\delta_2\circ\delta_1\quad(\text{定义~\ref{def:tex29-top-degeneration-projection}(2)}),\\
  P_2 &: P_1\Rightarrow (\forall X)\ \Pi(\delta_1(\delta_2(X)))=\Pi(\delta_2(\delta_1(X))),\\
  P_3 &: P_1\Rightarrow \Pi\circ(\delta_1\circ\delta_2-\delta_2\circ\delta_1)=0\Rightarrow \Pi\circ[\delta_1,\delta_2]
    =0.
\end{aligned}
\]
\end{Certificate}

\paragraph{证明（谓词因果链）.}
由 $P_1$ 直接得到 $P_2$；在可作差语境下同理得到 $P_3$。证毕。

\begin{corollary}[推论 C1（命名对应：为何“\GZMon{}”在投影下退化为“原子”）]
\label{cor:tex29-top-name}
对任意 $\mu\in\mathsf{TOP}$：
\begin{itemize}
  \item 在生成/修复的真实对象域中，$\mu$ 是“原基”（最小可操作单位），并通过 $\mathcal{O}$ 支持非交换与路径依赖；
  \item 在静态冻结投影 $\Pi$ 下，$\Pi(\mu)$ 退化为“原子”（点式不可分单位），且由命题~\ref{prop:tex29-top-degeneration-commutative} 知交换化消去非交换结构。
\end{itemize}
因此“\GZMon{}”一名可被理解为：
\[
\begin{aligned}
  &\text{富结构离散元（可构造闭合单元）}+\text{超限（沿序数递归）}\\
  &\qquad +\text{算子（非交换动力）}
  \Rightarrow\ \Pi\text{ 下退化为原子并交换化}.
\end{aligned}
\]
\end{corollary}

\begin{Certificate}[证书：由定义 D5 的三条条件抽取]
\label{cert:tex29-top-name}
证书登记谓词链：
\[
\begin{aligned}
  N_1 &: \mu\in\mathsf{TOP}\Rightarrow \text{满足 (T1) 最小性与 (T2) 超限递归生成性}\quad(\text{定义~\ref{def:tex29-top-primitive}}),\\
  N_2 &: \mu\in\mathsf{TOP}\Rightarrow \Pi(\mu)\ \text{为原子且}\ \Pi\ \text{交换化}\quad(\text{定义~\ref{def:tex29-top-primitive}
    的 (T3)}),\\
  N_3 &: \text{命题~\ref{prop:tex29-top-degeneration-commutative}}\Rightarrow \text{交换化等价于投影消去交换子},\\
  N_4 &: N_1\wedge N_2\wedge N_3\Rightarrow \text{推论陈述成立}.
\end{aligned}
\]
\end{Certificate}

\paragraph{证明（谓词因果链）.}
由 $N_1$ 得“原基 + 超限递归 + 算子动力”的对象层语义；
由 $N_2$ 得“投影下原子化 + 交换化”的退化语义；
由 $N_3$ 将交换化改写为可检验的交换子消去条件；合并得 $N_4$，推论成立。证毕。

\subsection{备注：命名方案（可落地）}
\label{subsec:tex29-top-remark}

\begin{remark}[命名建议：把退化原子性写成条件而非修辞]
为保证“证书化”与可核验性，本章建议：
\begin{enumerate}
  \item 主名优先使用\textbf{\GZMonFull{}}：其定义性约束为定义~\ref{def:tex29-top-primitive} 的 (T1)--(T3)；
  \item 为兼容既有符号与旧稿引用，本章仍使用 $\mu\in\mathsf{TOP}$ 作为“满足定义~\ref{def:tex29-top-primitive} 的\GZMon{}”的记号；
  \item “超限算子原基（TOP）”仅作为技术同义词保留，用于强调算子/非交换数据来源；正文默认不再以该旧名作为主名。
  \item “退化原子性”应固定为可检验条件（原子化 + 交换化），即定义~\ref{def:tex29-top-degeneration-projection}(1)(2)；
  其中 (2) 可由命题~\ref{prop:tex29-top-degeneration-commutative} 等价写成“投影消去交换子”。
\end{enumerate}
\end{remark}

\clearpage

%%%%%%%%%%%%%%%%%%%%%%%%%%%%%%%%%%%%%%%%%%%%%%%%%%%%%%%%%%%%
\section{本体论级拓扑重构：原子化平坦流形（带空洞）\texorpdfstring{$\to$}{->} \GZMonization{}同伦塔流形（无空洞且非平坦）}
\label{sec:tex29-ch14}
%%%%%%%%%%%%%%%%%%%%%%%%%%%%%%%%%%%%%%%%%%%%%%%%%%%%%%%%%%%%

\begin{remark}[核心陈述（从旧预设到新对象域）]
本章把一句最凝练的拓扑刻画落到可证书化的形式系统中，并给出无断点的严格对应拆解：
\begin{quote}
旧体系以“连续统平坦背景 + 原子化点集”为预设，其同调空洞由非平凡上同调障碍类（或障碍群的非零类）表达且无法在静态核内内生消解；
新体系以“\GZMon{}第一性 + 超限递归同伦修复塔”为对象域，通过边缘算子构造性平凡化（$\partial D=\omega$）把障碍从终止判据转写为构造步骤，
从而得到无全局空洞但带分层内生曲率的非平坦连通流形（或其等价类）。
\end{quote}
其中：
\begin{itemize}
  \item “无空洞”用障碍群平凡化表达（第\ref{sec:tex29-ch7}章）；
  \item “旧体系可回收”用保核收缩兼容表达（第\ref{sec:tex29-ch8}章）；
  \item “不平坦/内生曲率”用“核内非零类在总空间中可平凡化”的证书表达（本章第\ref{subsec:tex29-ch14-curvature}小节）。
\end{itemize}
本章不在旧体系内部证明诸如 CH/不完备性等命题；
而是把它们在本笔记中的角色固定为“可编码为障碍表达的原生空洞”，并给出在新对象域中对这些障碍进行构造性平凡化的形式化闭环。
\end{remark}

\subsection{旧体系：原子化平坦核与原生空洞（不可内生修复）}
\label{subsec:tex29-ch14-old}

\begin{definition}[原子化平坦核（连续统预设）]
令 $A\in\mathsf{Cont}$ 表示旧体系的连续统背景空间（平坦流形/平坦连续统的抽象化对象），
令 $\mathsf{Atm}$ 表示“原子”点集。
称 $(A,\mathsf{Atm})$ 为\textbf{原子化平坦核}，若满足：
\begin{enumerate}
  \item （背景依赖/第二性）$\mathsf{Atm}$ 仅作为 $A$ 内的点集（离散切片/采样），不携带内生闭合证据链数据（不含 $D$）；
  \item （静态性）允许的操作不改变边界像 $B^n(A;G)$：不允许新增证据链 $D$ 去扩展 $B^n(A;G)$（与推论~\ref{cor:tex29-core-contraction-dynamic-necessity}
    同\KouJing）。
\end{enumerate}
\end{definition}

\subsubsection{“原子”的本体论缺陷：背景依赖的第二性离散切片}
\label{subsubsec:tex29-ch14-old-atom}

\begin{remark}[严格对应：原子不携带内生证据链]
在本笔记的形式化\KouJing下，“原子”的关键特征不是“不可再分”，而是\emph{第二性}：
\begin{enumerate}
  \item 原子仅作为 $A$ 内点集 $\mathsf{Atm}\subseteq A$ 出现，其存在语义依赖于先验背景 $A$；
  \item 原子点集不携带任何用于修复的证据链数据（不含 $D$），因此无法在对象层面产生新的边界像 $B^n(A;G)$；
  \item 因此“修复/消洞”不可能来自原子本体的内生闭合结构，只能寄生于背景核外的新增构造（这正是旧体系需要“跳出纸面”才能补洞的形式化表达）。
\end{enumerate}
这与\GZMon{}第一性形成互斥：\GZMon{}对象域中，证据链 $D$ 与闭合校验 $\partial^2=0$ 被视为对象层数据而非外加规则。
\end{remark}

\subsubsection{“平坦流形”的预设陷阱：先验闭系统与动态生成禁令}
\label{subsubsec:tex29-ch14-old-flat}

\begin{remark}[严格对应：平坦=核内无新增边界能力]
本笔记不把“平坦”当作度量几何的曲率数值，而把其抽象为一个可核验的操作边界：
\[
  \text{平坦核\KouJing}\ \equiv\ \text{静态核\KouJing}\ \equiv\ \text{不允许新增证据链 }D\ \equiv\ B^n(A;G)\ \text{固定不变}.
\]
在该\KouJing下，核内推演最多能\emph{识别}障碍类是否为零（即判断 $[\omega]=0$ 与否），但无法\emph{生成}使其变零的构造。
这也解释了“交换性/理想闭系统”在本笔记中的对应物：核内运算的封闭性由 $B^n$ 的不变性承担，而“开放动态生成”被等价地刻画为必须引入新的 $D$ 去扩展边界像。
\end{remark}

\subsubsection{“含孔洞”的原生绝症：非平凡障碍群与不可内生修复}
\label{subsubsec:tex29-ch14-old-hole}

\begin{lemma}[空洞表述：非平凡障碍类 $\Leftrightarrow$ 障碍群非零]
\label{lem:tex29-ch14-hole-as-obstruction-group}
对任意 $n\ge 1$，以下两种表述等价：
\begin{enumerate}
  \item 存在闭链 $\omega\in Z^n(A;G)$ 使得 $\omega\notin B^n(A;G)$；
  \item 障碍群 $H^n_{\mathrm{obs}}(A;G)=Z^n(A;G)/B^n(A;G)$ 非零。
\end{enumerate}
\end{lemma}

\begin{Certificate}[证书：商群非零的充要条件]
\label{cert:tex29-ch14-hole-as-obstruction-group}
证书登记谓词链：
\[
\begin{aligned}
  H_1 &: (\exists \omega\in Z^n)\ (\omega\notin B^n)\Rightarrow Z^n/B^n\neq 0,\\
  H_2 &: Z^n/B^n\neq 0\Rightarrow (\exists [\omega]\neq 0)\Rightarrow (\exists \omega\in Z^n)\ (\omega\notin B^n),\\
  H_3 &: H_1\wedge H_2\Rightarrow \text{两种表述等价}.
\end{aligned}
\]
\end{Certificate}

\paragraph{证明（谓词因果链；符号推演）.}
由 $H_1$：若存在 $\omega\in Z^n$ 且 $\omega\notin B^n$，则其商类 $[\omega]\in Z^n/B^n$ 非零，故商群非零；
由 $H_2$：若商群非零，则存在某个非零类 $[\omega]\neq 0$，取代表元即可得到 $\omega\in Z^n$ 且 $\omega\notin B^n$。
合并得 $H_3$，引理成立。证毕。

\begin{proposition}[旧核死局：静态平坦核内不能构造性消解非平凡障碍]
\label{prop:tex29-ch14-old-no-repair}
设 $(A,\mathsf{Atm})$ 为原子化平坦核。
若存在 $n\ge 1$ 与非平凡障碍类 $[\omega]\neq 0\in H^n_{\mathrm{obs}}(A;G)$，
则在该核的静态操作\KouJing内不可能构造证据链 $D$ 使得 $\partial D=\omega$，
即非平凡障碍只能被识别，不能被内生消解。
\end{proposition}

\begin{Certificate}[证书：由边界推论~\ref{cor:tex29-core-contraction-dynamic-necessity} 直接推出]
\label{cert:tex29-ch14-old-no-repair}
证书登记谓词链：
\[
\begin{aligned}
  S_1 &: (A,\mathsf{Atm})\ \text{为原子化平坦核}\Rightarrow \text{不允许新增证据链}\Rightarrow B^n(A;G)\ \text{固定},\\
  S_2 &: [\omega]\neq 0\Rightarrow \omega\notin B^n(A;G),\\
  S_3 &: S_1\wedge S_2\Rightarrow \text{核内无法使 }\omega\text{ 进入 }B^n(A;G)\Rightarrow \neg\exists D\ (\partial D=\omega).
\end{aligned}
\]
\end{Certificate}

\paragraph{证明（谓词因果链；符号推演）.}
由 $S_1$ 得静态核内 $B^n$ 不随操作扩展；由 $S_2$ 得非零类代表元不属于 $B^n$；
因此由 $S_3$ 得无法在核内构造 $\partial D=\omega$。命题成立。证毕。

\begin{remark}[备注：旧体系“不可判定/不可延拓”在本笔记中的位置]
在本笔记的形式化\KouJing下：
“不可判定/不可延拓”被统一视为“在静态核内出现的非零障碍类 $[\omega]\neq 0$”；
而“能否突破”不在旧核内讨论，而被转移到新对象域中讨论：
是否存在可证书化的证据链 $D$ 使 $\partial D=\omega$，以及是否能在修复塔中保持闭合校验与变分可接受性。
\end{remark}

\subsection{新体系：\GZMonFull{}第一性、超限递归同伦塔与无空洞完成态}
\label{subsec:tex29-ch14-new}

\begin{definition}[\GZMonization{}同伦塔总空间（离散统第一性）]
沿用第\ref{subsec:tex29-fractal-domain}小节的对象域：$\mathsf{Mon}$ 为\GZMon{}基元集合，$X\in\mathsf{Cfg}(\mathsf{Mon})$ 为离散构型表示。
称 $X$ 为\textbf{\GZMonization{}同伦塔总空间}，若存在一条可证书化的修复塔/修复链数据（例如第\ref{subsec:tex29-dynamics-meta-loop}小节的动力学修复链），
使得每一层均可构造证据链 $D$ 满足 $\partial D=\omega$，并通过闭合校验 $\partial^2=0$。
当该链的层级索引被理解为第一超限序数 $\omega$（或更高序数）时，称其为\textbf{超限递归}修复塔\KouJing。
\end{definition}

\subsubsection{“\GZMonFull{}”的本体论反转：背景独立的第一性原生基元}
\label{subsubsec:tex29-ch14-new-monad}

\begin{remark}[严格对应：\GZMon{}是“可构造的闭合单元”】【与原子互斥】]
与“原子仅是背景 $A$ 的点集切片”不同：
\begin{enumerate}
  \item \GZMon{}对象域 $\mathsf{Mon}$ 被视为第一性；连续统对象仅通过投影/表示函子 $\Pi$ 作为观测层出现；
  \item \GZMonization{}总空间 $X\in\mathsf{Cfg}(\mathsf{Mon})$ 的定义性数据包含修复证据链 $D$ 与闭合校验 $\partial^2=0$，
  这使得“修复”成为对象层可生成、可登记、可验算的构造，而非对外部背景的被动依赖；
  \item 因此\GZMon{}并非“连续统的离散采样”，而是“可递归生成流形的构造基元”。
\end{enumerate}
这正是离散统第一性的形式化落点：不是先有平坦背景再切片出离散单元，而是先有可构造的离散基元，再以修复塔生成可观测结构。
\end{remark}

\begin{remark}[补充：主纤维丛视角的结构对齐（总空间/投影/结构群；仅作参照，不作为本章证明前提）]
为把“总空间 $X$、观测核 $\Pi(X)$、以及可执行修复操作族”压缩为一张可复用结构图，
可引入主纤维丛（principal bundle）的启发式对齐\KouJing：
\begin{itemize}
  \item 总空间 $P$：\GZMonization{}同伦塔总空间 $X\in\mathsf{Cfg}(\mathsf{Mon})$（携带证据链 $D$ 与闭合校验 $\partial^2=0$）；
  \item 底空间 $M$：观测/描述层对象 $\Pi(X)\in\mathsf{Cont}$（或更一般的“现实表象域”）；
  \item 丛投影 $\pi:P\to M$：观测/退化函子 $\Pi$；
  \item 结构群 $G$：由允许的修复变换（同伦变换、跨域保结构函子、链同构诱导等）生成的操作群（或操作半群）。
\end{itemize}
由于这些操作通常不可交换（不同修复顺序可能产生不同的中间塔结构），
可把其代数化为一种“广义非交换操作代数/广义李代数”的描述\KouJing。
该对齐不引入额外定理，只作为跨域沟通的结构参照点。
（主纤维丛基础参见 \cite{Husemoller1994FibreBundles}；与本仓库语境相关的“主纤维丛版广义非交换李代数”\KouJing可参照 \cite{GaoZheng2025GFramework}；
  李代数/非交换代数的常见教材\KouJing参见 \cite{Hall2015LieGroups}。）
\end{remark}

\subsubsection{超限递归同伦修复塔：连通性是构造结果而非背景预设}
\label{subsubsec:tex29-ch14-new-tower}

\begin{lemma}[同伦塔连通：层间同伦等价链 $\Rightarrow$ 全链同伦等价（数学归纳法证书）]
\label{lem:tex29-ch14-tower-connected}
设 $\{X_n\}_{n\in\NN}\subseteq\mathsf{Cfg}(\mathsf{Mon})$ 满足对任意 $n\in\NN$ 都有 $X_{n+1}\simeq X_n$。
则对任意 $n\in\NN$，均有 $X_n\simeq X_0$。
\end{lemma}

\begin{Certificate}[数学归纳法证书：同伦等价链的可传递合成]
\label{cert:tex29-ch14-tower-connected}
令谓词 $P(n)$ 表示“$X_n\simeq X_0$”。证书化要求为：
\[
  P(0)\ \text{成立},\qquad (\forall n\in\NN)\ \bigl(P(n)\Rightarrow P(n+1)\bigr).
\]
\end{Certificate}

\paragraph{证明（数学归纳法证书；谓词因果链）.}
由证书~\ref{cert:tex29-ch14-tower-connected}：
\begin{itemize}
  \item 基步：$n=0$ 时 $X_0\simeq X_0$ 恒成立；
  \item 归纳步：若 $X_n\simeq X_0$，且已知 $X_{n+1}\simeq X_n$，则由同伦等价的可合成性得 $X_{n+1}\simeq X_0$。
\end{itemize}
故对任意 $n$ 均成立。证毕。

\begin{remark}[严格对应：同伦塔把“全局连通”写成可执行的层级算法]
在该\KouJing下，“连通”不是来自先验平坦背景的假设，而是来自：
\[
  X_0\ \xrightarrow[\simeq]{}\ X_1\ \xrightarrow[\simeq]{}\ \cdots
\]
这一条可证书化的层间同伦等价链；它把“连通性”转写为“可由更新规则逐层构造并保持的同伦类连通”。
\end{remark}

\subsubsection{无空洞完成态：把障碍从终止判据转写为构造单元}
\label{subsubsec:tex29-ch14-new-nohole}

\begin{definition}[无空洞完成态谓词]
称 $X$ 达到\textbf{无空洞完成态}，若
\[
  \mathsf{NoHole}(X)\ :\Longleftrightarrow\ (\forall n\ge 1)\ H^n_{\mathrm{obs}}(X;G)=0.
\]
\end{definition}

\begin{corollary}[无空洞完成态：全阶系统性平凡化 $\Rightarrow$ $\mathsf{NoHole}(X)$]
\label{cor:tex29-ch14-no-hole}
若 $X$ 满足第\ref{subsec:tex29-core-meta-trivialization}小节定理~\ref{thm:tex29-core-meta-trivialization} 的前提（对任意闭链障碍均可构造性平凡化），
则 $X$ 达到无空洞完成态，即 $\mathsf{NoHole}(X)$ 成立。
\end{corollary}

\begin{Certificate}[证书：由定理~\ref{thm:tex29-core-meta-trivialization} 直接推出]
\label{cert:tex29-ch14-no-hole}
证书登记谓词链：
\[
\begin{aligned}
  N_1 &: (\forall n\ge 1)\ (\forall \omega\in Z^n(X;G))\ \mathsf{Rep}(\omega),\\
  N_2 &: N_1\wedge \partial^2=0\Rightarrow (\forall n\ge 1)\ H^n_{\mathrm{obs}}(X;G)=0\quad(\text{定理~\ref{thm:tex29-core-meta-trivialization}}),\\
  N_3 &: N_2\Rightarrow \mathsf{NoHole}(X).
\end{aligned}
\]
\end{Certificate}

\paragraph{证明（谓词因果链）.}
由 $N_2$ 得对任意 $n\ge 1$ 有 $H^n_{\mathrm{obs}}(X;G)=0$；
由定义即得 $\mathsf{NoHole}(X)$。证毕。

\begin{remark}[严格对应：障碍不是终点，而是构造的起点]
无空洞完成态并非“强行把洞填上”，而是把每一个可能出现的非零障碍类都纳入修复塔的某一层作为构造目标：
以 $\partial D=\omega$ 把“障碍”转写为“生成边界所需的证据链”，再由 $\partial^2=0$ 排除次生断点，
从而把“洞”从否定性终点转写为正向构造步骤。
\end{remark}

\subsection{公理降级：拓扑重构链的存在性与兼容闭环（数学归纳法证书）}
\label{subsec:tex29-ch14-reconstruction}

\begin{theorem}[公理降级：从平坦核到同伦塔总空间的重构链（数学归纳法证书）]
\label{thm:tex29-ch14-reconstruction}
给定旧核 $A\in\mathsf{Cont}$。
若存在初始离散构型 $X_0\in\mathsf{Cfg}(\mathsf{Mon})$ 使得 $\Pi(X_0)=A$，
并且对每一层 $n\in\NN$ 都存在一条更新规则，能从 $X_n$ 构造 $X_{n+1}$ 与数据 $(\omega_n,D_n)$，满足：
\begin{enumerate}
  \item （核保持）$\Pi(X_{n+1})=A$；
  \item （修复步）$\partial D_n=\omega_n$；
  \item （闭合步）$\partial^2=0$ 作为层间闭合校验成立；
  \item （层间连通）存在同伦等价映射对 $(h_n,k_n)$ 使 $X_{n+1}\simeq X_n$。
\end{enumerate}
则存在一条以 $A$ 为观测核的拓扑重构链
\[
  \mathbb{R}_A:=\bigl\{(X_n,\omega_n,D_n)\bigr\}_{n\in\NN},
  \qquad \Pi(X_n)=A\ \ (\forall n),
\]
从而把“旧核 + 静态描述”提升为“新对象域 + 可证书化修复链”的动力学构造载体。
\end{theorem}

\begin{Certificate}[数学归纳法证书：按层数 $n$ 归纳的“核保持--修复--闭合--连通”构造]
\label{cert:tex29-ch14-reconstruction}
令谓词 $P(n)$ 表示：
“存在长度为 $n$ 的链 $\{(X_k,\omega_k,D_k)\}_{k=0}^{n}$，满足
$\Pi(X_k)=A$、$\partial D_k=\omega_k$（对 $k<n$）、$X_{k+1}\simeq X_k$（对 $k<n$），并且 $\partial^2=0$ 作为闭合校验成立。”
证书化要求为：
\[
  P(0)\ \text{成立},\qquad (\forall n\in\NN)\ \bigl(P(n)\Rightarrow P(n+1)\bigr).
\]
\end{Certificate}

\paragraph{证明（数学归纳法证书；谓词因果链；符号推演）.}
由证书~\ref{cert:tex29-ch14-reconstruction}：
\begin{itemize}
  \item 基步：$n=0$ 时取给定的 $X_0$，且 $\Pi(X_0)=A$；因此 $P(0)$ 成立；
  \item 归纳步：假设 $P(n)$ 成立，即已构造到第 $n$ 层。
  由定理假设的更新规则可构造 $X_{n+1}$ 与 $(\omega_n,D_n)$，并满足核保持、修复、闭合与层间连通四条条件；
  将新增层并入既有链，即得长度为 $n+1$ 的链，故 $P(n+1)$ 成立。
\end{itemize}
由数学归纳法，$P(n)$ 对任意 $n$ 成立，从而存在无限链 $\mathbb{R}_A$。证毕。

\begin{remark}[备注：重构闭环的两端]
定理~\ref{thm:tex29-ch14-reconstruction} 给出“从旧核出发构造新修复链”的存在性；
若进一步满足全阶系统性平凡化，则由推论~\ref{cor:tex29-ch14-no-hole} 得到无空洞完成态；
若进一步满足静态冻结，则由定理~\ref{thm:tex29-core-contraction-theorem} 得到保核收缩回收旧体系。
因此，“旧体系可回收”与“新体系可修复”可在同一形式系统中闭环并存。
\end{remark}

\subsection{命题：非平坦性（内生曲率）作为“空洞可消解”的必然代价}
\label{subsec:tex29-ch14-curvature}

\begin{definition}[内生曲率谓词（相对收缩核）]
设 $A\in\mathsf{Cont}$ 为静态核，$X\in\mathsf{Cfg}(\mathsf{Mon})$ 为其\GZMonization{}总空间表示，并存在一个障碍表达的提升映射
\[
  \iota:\Omega^\ast(A;G)\longrightarrow \Omega^\ast(X;G),
\]
满足 $\iota\bigl(Z^n(A;G)\bigr)\subseteq Z^n(X;G)$（闭链在提升后仍为闭链）。
定义\textbf{内生曲率谓词}：
\[
  \mathsf{Curv}(X;A)\ :\Longleftrightarrow\ \exists n\ge 1,\ \exists \omega\in Z^n(A;G)\ \text{使}\ [\omega]\neq 0\in
    H^n_{\mathrm{obs}}(A;G)\ \wedge\ \iota(\omega)\in B^n(X;G).
\]
其语义解释（不作为证明步）为：在“静态核”中不可消解的空洞类，经由“\GZMonization{}总空间”的构造性修复而被转写为边界；
该“核外新增边界能力”即是本笔记\KouJing下的不平坦性/分层内生曲率证书。
\end{definition}

\begin{proposition}[非平坦性证书：核内非零类在总空间中可平凡化]
\label{prop:tex29-ch14-curvature}
若存在 $n\ge 1$ 使 $H^n_{\mathrm{obs}}(A;G)\neq 0$，并且 $X$ 达到无空洞完成态 $\mathsf{NoHole}(X)$，
则在上述提升映射 $\iota$ 存在的前提下，必有 $\mathsf{Curv}(X;A)$ 成立。
\end{proposition}

\begin{Certificate}[证书：核内非零 $\wedge$ 总空间无空洞 $\Rightarrow$ 存在“新增边界”的见证元]
\label{cert:tex29-ch14-curvature}
证书登记谓词链：
\[
\begin{aligned}
  K_1 &: H^n_{\mathrm{obs}}(A;G)\neq 0\Rightarrow \exists \omega\in Z^n(A;G)\ (\omega\notin B^n(A;G)),\\
  K_2 &: \iota\bigl(Z^n(A;G)\bigr)\subseteq Z^n(X;G)\Rightarrow \iota(\omega)\in Z^n(X;G),\\
  K_3 &: \mathsf{NoHole}(X)\Rightarrow Z^n(X;G)=B^n(X;G)\Rightarrow \iota(\omega)\in B^n(X;G),\\
  K_4 &: K_1\wedge K_2\wedge K_3\Rightarrow \mathsf{Curv}(X;A).
\end{aligned}
\]
\end{Certificate}

\paragraph{证明（谓词因果链；符号推演）.}
由 $K_1$ 取核内非零类的代表元 $\omega\in Z^n(A;G)$；
由 $K_2$ 得其提升 $\iota(\omega)$ 仍为闭链；
由 $K_3$（无空洞完成态）得总空间中闭链皆为边界，从而 $\iota(\omega)\in B^n(X;G)$；
合并得到 $K_4$，即 $\mathsf{Curv}(X;A)$ 成立。证毕。

\subsection{最终的公理级闭环定论：从原生空洞到无空洞完成态（并必然非平坦）}
\label{subsec:tex29-ch14-final}

\begin{theorem}[本体论级拓扑重构的形式化闭环]
\label{thm:tex29-ch14-final-loop}
设 $(A,\mathsf{Atm})$ 为原子化平坦核，$A\in\mathsf{Cont}$。
若存在\GZMonization{}总空间 $X\in\mathsf{Cfg}(\mathsf{Mon})$ 与提升映射 $\iota$，满足：
\begin{enumerate}
  \item （观测核保持）$\Pi(X)=A$；
  \item （无空洞完成态）$\mathsf{NoHole}(X)$；
  \item （旧核有洞）存在 $n\ge 1$ 使 $H^n_{\mathrm{obs}}(A;G)\neq 0$；
  \item （提升闭链保持）$\iota\bigl(Z^n(A;G)\bigr)\subseteq Z^n(X;G)$（同第\ref{subsec:tex29-ch14-curvature}小节的前提\KouJing）。
\end{enumerate}
则：
\begin{enumerate}
  \item 旧核内的非零障碍类无法被内生构造性消解（命题~\ref{prop:tex29-ch14-old-no-repair}）；
  \item 该非零类在 $X$ 中必可转写为边界，从而“空洞”在总空间中被构造性消解；
  \item 总空间必然满足 $\mathsf{Curv}(X;A)$，即具备核外新增边界能力（命题~\ref{prop:tex29-ch14-curvature}），这正是“不平坦/分层内生曲率”的形式化证书。
\end{enumerate}
\end{theorem}

\begin{Certificate}[证书：由“旧核不可修复 + 总空间无空洞 + 提升保持”拼接推出]
\label{cert:tex29-ch14-final-loop}
证书登记谓词链：
\[
\begin{aligned}
  F_1 &: H^n_{\mathrm{obs}}(A;G)\neq 0\Rightarrow \exists \omega\in Z^n(A;G)\ (\omega\notin B^n(A;G))
    \quad(\text{引理~\ref{lem:tex29-ch14-hole-as-obstruction-group}}),\\
  F_2 &: (A,\mathsf{Atm})\ \text{为原子化平坦核}\Rightarrow \neg\exists D\ (\partial D=\omega)\ \text{在核内成立}
    \quad(\text{命题~\ref{prop:tex29-ch14-old-no-repair}}),\\
  F_3 &: \iota\bigl(Z^n(A;G)\bigr)\subseteq Z^n(X;G)\Rightarrow \iota(\omega)\in Z^n(X;G),\\
  F_4 &: \mathsf{NoHole}(X)\Rightarrow Z^n(X;G)=B^n(X;G)\Rightarrow \iota(\omega)\in B^n(X;G),\\
  F_5 &: F_1\wedge F_3\wedge F_4\Rightarrow \mathsf{Curv}(X;A)\quad(\text{命题~\ref{prop:tex29-ch14-curvature}}).
\end{aligned}
\]
\end{Certificate}

\paragraph{证明（谓词因果链；符号推演）.}
由 $F_1$ 取核内非零类的代表元 $\omega\in Z^n(A;G)$ 且 $\omega\notin B^n(A;G)$；
由 $F_2$ 得其在旧核内无法通过新增证据链实现 $\partial D=\omega$ 的构造性平凡化；
由 $F_3$ 与 $F_4$ 得 $\iota(\omega)$ 在总空间中必为边界，从而核内“空洞类”在总空间中被构造性消解；
再由 $F_5$ 得 $\mathsf{Curv}(X;A)$ 成立。定理得证。证毕。

\begin{remark}[对齐表达：把自然语言叙述转写为三段可核验链条]
本章的形式化闭环可被压缩为：
\[
  \underbrace{\text{旧核内 }H^n_{\mathrm{obs}}(A;G)\neq 0}_{\text{原生空洞}}
  \ \Longrightarrow\
  \underbrace{\text{核内禁止新增 }D}_{\text{不可内生修复}}
  \ \Longrightarrow\
  \underbrace{\text{总空间中 }\iota(\omega)\in B^n(X;G)}_{\text{构造性消解 + 非平坦证书}}.
\]
其中，“不可判定/不完备”等仅作为“核内非零障碍类”的编码角色进入本笔记，不作为旧体系内部的可证明命题参与推演。
\end{remark}

\clearpage

%%%%%%%%%%%%%%%%%%%%%%%%%%%%%%%%%%%%%%%%%%%%%%%%%%%%%%%%%%%%
\section{范式重构的量级与文明跃迁：分支级公理拓展 vs 基础级地基重构}
\label{sec:tex29-ch15}
%%%%%%%%%%%%%%%%%%%%%%%%%%%%%%%%%%%%%%%%%%%%%%%%%%%%%%%%%%%%

\begin{remark}[核心陈述（历史定位的形式系统\KouJing）]
本章不在外部数学史叙事中争论“谁更伟大”，而仅在本笔记的对象域/谓词/证书\KouJing下给出一个可检验判据：
\begin{enumerate}
  \item 若一项变更仅改变某个分支理论的公设集合，但不改变第一性对象域，也不引入新的证据链 $D$ 去扩展边界像 $B^n$，
  则它属于\textbf{分支级公理拓展}：可产生并行分支，但不会触及“核内原生空洞/不可判定裂缝”的修复问题；
  \item 若一项变更将对象域从“静态连续统核”提升为“\GZMon{}第一性总空间”，并允许在修复塔中构造证据链 $D$ 使 $\partial D=\omega$，
  从而把核内非零障碍类提升后转写为边界（并最终达到 $\mathsf{NoHole}$），
  则它属于\textbf{基础级地基重构}：其效果不是增加分支，而是改变“障碍的存在论位置”（由终止判据转为构造步骤）。
\end{enumerate}
在该判据下，“第五公设的独立性”是分支级拓展的典型样例；
而“离散统第一性 + 超限递归同伦修复塔”是基础级地基重构：它直接对应第\ref{sec:tex29-ch7}章与第\ref{sec:tex29-ch14}章的无空洞闭环与核外新增边界能力证书。
\end{remark}

\subsection{定义：分支级拓展与地基级重构的可检验谓词}
\label{subsec:tex29-ch15-defs}

\begin{definition}[分支级公理拓展谓词（底层不变）]
设 $A\in\mathsf{Cont}$ 为某个静态核，并沿用第\ref{subsec:tex29-fractal-domain}小节引入的投影/表示函子 $\Pi$ 作为观测层\KouJing。
称一次理论变更 $\mathcal{E}$ 为\textbf{分支级公理拓展}，若存在两个分支理论 $\mathcal{T}_1,\mathcal{T}_2$ 满足：
\begin{enumerate}
  \item 两者共享同一底层对象域（同一静态核）$A$，且允许的核内操作不新增证据链 $D$（即 $B^n(A;G)$ 固定不变）；
  \item 两者的差异仅体现在对 $\Pi$ 的观测层描述/公设选择上，而不改变修复入口判据 $\mathsf{Rep}$ 的定义\KouJing；
  \item 两者可并行成立：不存在强制把 $\mathcal{T}_1$ 归约为 $\mathcal{T}_2$（或相反）的核内必然性。
\end{enumerate}
\end{definition}

\begin{definition}[基础级地基重构谓词（对象域升维）]
称一次理论变更 $\mathcal{R}$ 为\textbf{基础级地基重构}，若对给定核 $A\in\mathsf{Cont}$，存在\GZMonization{}总空间 $X\in\mathsf{Cfg}(\mathsf{Mon})$
  使得：
\begin{enumerate}
  \item （观测核保持）$\Pi(X)=A$；
  \item （修复链可生成）存在可证书化的修复塔/重构链（定理~\ref{thm:tex29-ch14-reconstruction} 的\KouJing）；
  \item （障碍可转写）对核内障碍表达的提升映射 $\iota$（第\ref{subsec:tex29-ch14-curvature}小节）成立，使得某些核内非零类在 $X$ 中变为边界；
  \item （闭环目标）最终可导出无空洞完成态 $\mathsf{NoHole}(X)$（推论~\ref{cor:tex29-ch14-no-hole} 的\KouJing）。
\end{enumerate}
\end{definition}

\subsection{公理降级：分支拓展不触及裂缝，地基重构给出消解通道（数学归纳法证书）}
\label{subsec:tex29-ch15-axiom}

\begin{theorem}[公理降级：底层不变的拓展不改障碍群；地基重构逐层生成“新增边界”能力（数学归纳法证书）]
\label{thm:tex29-ch15-branch-vs-foundation}
设 $A\in\mathsf{Cont}$ 为静态核，且存在某个 $n\ge 1$ 使 $H^n_{\mathrm{obs}}(A;G)\neq 0$。
\begin{enumerate}
  \item 若 $\mathcal{E}$ 为分支级公理拓展，则在 $\mathcal{E}$ 的核内\KouJing下无法构造证据链 $D$ 使 $\partial D=\omega$（对任何代表核内非零类的 $\omega$），因此
    $H^n_{\mathrm{obs}}(A;G)$ 的“非零裂缝”不被触及；
  \item 若 $\mathcal{R}$ 为基础级地基重构，则存在一条以 $A$ 为观测核的重构链 $\mathbb{R}_A$（定理~\ref{thm:tex29-ch14-reconstruction}），
    并可在某些层上生成见证元使得核内非零类在总空间中转写为边界，从而给出“裂缝可消解”的构造通道（命题~\ref{prop:tex29-ch14-curvature}）。
\end{enumerate}
\end{theorem}

\begin{Certificate}[数学归纳法证书：按修复层数登记“新增边界能力”的构造链]
\label{cert:tex29-ch15-branch-vs-foundation}
令谓词 $P(k)$ 表示：
“存在一条长度为 $k$ 的重构链 $\{(X_i,\omega_i,D_i)\}_{i=0}^{k}$，满足 $\Pi(X_i)=A$（对所有 $i$），并且对某个核内非零类代表元 $\omega$，
其提升 $\iota(\omega)$ 在某一层 $X_j$（$j\le k$）中落入 $B^n(X_j;G)$（即出现一次‘新增边界’见证）。”
证书化要求为：
\[
  P(0)\ \text{成立},\qquad (\forall k\in\NN)\ \bigl(P(k)\Rightarrow P(k+1)\bigr).
\]
并登记分支拓展侧的谓词链：
\[
\begin{aligned}
  B_1 &: \mathcal{E}\ \text{为分支级拓展}\Rightarrow \text{核内不允许新增证据链}\Rightarrow B^n(A;G)\ \text{固定},\\
  B_2 &: H^n_{\mathrm{obs}}(A;G)\neq 0\Rightarrow \exists \omega\in Z^n(A;G)\ (\omega\notin B^n(A;G)),\\
  B_3 &: B_1\wedge B_2\Rightarrow \neg\exists D\ (\partial D=\omega)\ \text{在核内成立}.
\end{aligned}
\]
\end{Certificate}

\paragraph{证明（数学归纳法证书；谓词因果链；符号推演）.}
\begin{itemize}
  \item 分支拓展侧：由 $B_1$ 得核内 $B^n(A;G)$ 不随操作扩展；由 $B_2$ 取 $\omega\notin B^n(A;G)$；
  从而由 $B_3$ 得核内无法构造 $\partial D=\omega$。对应结论即为命题~\ref{prop:tex29-ch14-old-no-repair} 的抽象复述。
  \item 地基重构侧（归纳链）：由定理~\ref{thm:tex29-ch14-reconstruction} 的基步可取 $X_0$ 使 $\Pi(X_0)=A$，故 $P(0)$ 成立；
  若 $P(k)$ 成立，则由更新规则可延拓到 $k+1$ 层并保持核保持/闭合校验，
  从而“已出现的新增边界见证”在扩展链中仍被保留，故 $P(k+1)$ 成立。
\end{itemize}
由归纳法得对任意 $k$ 可构造到对应层；再结合命题~\ref{prop:tex29-ch14-curvature} 的证书链，可得到核内非零类在总空间中转写为边界的必然性通道。
定理得证。证毕。

\begin{lemma}[裂缝保持：核内边界像固定 $\Rightarrow$ 非零类不可内生消解]
\label{lem:tex29-ch15-crack-persist}
设 $A\in\mathsf{Cont}$ 为静态核，且核内操作不允许新增证据链 $D$ 以扩展 $B^n(A;G)$。
若存在 $\omega\in Z^n(A;G)$ 使 $\omega\notin B^n(A;G)$，
则在该核内\KouJing下不存在 $D$ 使 $\partial D=\omega$；从而 $[\omega]\neq 0\in H^n_{\mathrm{obs}}(A;G)$ 在核内保持为非零裂缝。
\end{lemma}

\begin{Certificate}[证书：固定 $B^n$ 的直接否定证书]
\label{cert:tex29-ch15-crack-persist}
证书登记谓词链：
\[
\begin{aligned}
  R_1 &: \text{禁止新增证据链}\Rightarrow B^n(A;G)\ \text{固定不变},\\
  R_2 &: \omega\notin B^n(A;G)\Rightarrow \neg\exists D\ (\partial D=\omega),\\
  R_3 &: R_2\Rightarrow [\omega]\neq 0\in Z^n/B^n=H^n_{\mathrm{obs}}(A;G).
\end{aligned}
\]
\end{Certificate}

\paragraph{证明（谓词因果链；符号推演）.}
由 $R_1$ 得 $B^n(A;G)$ 不被扩展；若 $\omega\notin B^n(A;G)$，则由定义 $B^n=\{\partial D\}$ 立得 $R_2$；
从而由 $R_3$ 得其商类非零。证毕。

\begin{remark}[备注：为何说二者“不在同一量级”】【刚性判据】]
本章的判据不是“是否能产生新分支”，而是“是否能在体系内部提供构造性消洞能力”：
\[
\begin{aligned}
  \text{分支拓展}\ &\Rightarrow\ \text{核内 }B^n\text{ 固定}\Rightarrow\ \text{裂缝只可识别不可消解};\\
  \text{地基重构}\ &\Rightarrow\ \exists X\ \text{生成新 }B^n(X)\Rightarrow\ \text{裂缝可被转写并消解}.
\end{aligned}
\]
因此，分支级拓展（如几何公设的独立选择）与地基级重构（对象域升维 + 修复塔）在本笔记\KouJing下并非同一层级的问题。
\end{remark}

\subsection{对齐：第五公设的突破为何属于“分支级公理拓展”}
\label{subsec:tex29-ch15-parallel}

\begin{proposition}[并行分支：同一核上可同时承载不同几何公设选择]
\label{prop:tex29-ch15-parallel-branch}
设 $A\in\mathsf{Cont}$ 为静态核。
若两套几何公设系统 $\mathcal{T}_{\mathrm{E}},\mathcal{T}_{\mathrm{NE}}$ 仅在观测层（例如对 $\Pi(X)=A$ 的度量/平行性描述）上不同，
且均不允许在核内新增证据链 $D$ 去扩展 $B^n(A;G)$，
则它们在本笔记\KouJing下构成分支级公理拓展：可并行成立但不改变核内障碍群的可修复性结论。
\end{proposition}

\begin{Certificate}[证书：观测层差异 $\Rightarrow$ 底层不变 $\Rightarrow$ 分支级拓展]
\label{cert:tex29-ch15-parallel-branch}
证书登记谓词链：
\[
\begin{aligned}
  P_1 &: \mathcal{T}_{\mathrm{E}},\mathcal{T}_{\mathrm{NE}}\ \text{共享同一核 }A\Rightarrow \text{对象域不变},\\
  P_2 &: \text{两者均禁止新增 }D\Rightarrow B^n(A;G)\ \text{固定},\\
  P_3 &: P_1\wedge P_2\Rightarrow \text{差异仅为观测层公设选择}\Rightarrow \mathcal{E}\ \text{为分支级拓展}.
\end{aligned}
\]
\end{Certificate}

\paragraph{证明（谓词因果链）.}
由 $P_1$ 得底层对象域不变；由 $P_2$ 得核内修复能力不变；
因此由 $P_3$ 得该变更只能产生并行分支而不触及“裂缝可否消解”的问题。证毕。

\begin{remark}[备注：本命题的作用边界]
本命题只说明“第五公设的独立选择”属于观测层分支，不涉及任何具体几何学结论；
其唯一功能是为后文“基础级地基重构”的比较提供一个严格参照：分支选择不改变 $\mathsf{Rep}$ 的入口与 $B^n$ 的生成能力。
\end{remark}

\subsection{集合论从“原子”到“\GZMonFull{}”：空壳基元 vs 内禀良序与闭合基元}
\label{subsec:tex29-ch15-atom-to-monad}

\begin{definition}[空壳原子（urelement \KouJing的抽象化）]
在本笔记中，称 $\mathsf{Atm}$ 为\textbf{空壳原子}集合，若其元素仅作为静态核 $A$ 中的点集切片出现，
且不携带任何内生良序数据与证据链数据（不含 $D$）。
该\KouJing与第\ref{subsec:tex29-ch14-old}小节的“原子化平坦核”一致。
\end{definition}

\begin{definition}[\GZMon{}基元的两条刚性属性：内禀良序与闭合校验]
称\GZMon{}对象域 $\mathsf{Mon}$ 满足两条刚性属性，若：
\begin{enumerate}
  \item （内禀良序）存在一个关系 $\prec$ 使得 $(\mathsf{Mon},\prec)$ 为良序（或至少对可生成构型所需部分为良序），
  该序结构作为对象层数据存在，而非由外加选择规则事后赋予；
  \item （闭合校验）在由 $\mathsf{Mon}$ 生成的修复链中，闭合校验 $\partial^2=0$ 作为刚性一致性标准可被执行（见第\ref{sec:tex29-ch7}章）。
\end{enumerate}
\end{definition}

\begin{proposition}[良序内禀性 $\Rightarrow$ “排序入口”不再依赖外加选择]
\label{prop:tex29-ch15-wellorder-intrinsic}
若\GZMon{}对象域 $\mathsf{Mon}$ 具有内禀良序 $\prec$，
则对任意可生成的离散构型 $X\in\mathsf{Cfg}(\mathsf{Mon})$，存在一个由 $\prec$ 诱导的规范枚举/递归构造顺序；
因此“把构造过程良序化”在本笔记\KouJing下可作为对象层数据直接给出，而不必作为额外外加的选择式入口假设。
\end{proposition}

\begin{Certificate}[证书：良序诱导规范递归次序]
\label{cert:tex29-ch15-wellorder-intrinsic}
证书登记谓词链：
\[
\begin{aligned}
  W_1 &: (\mathsf{Mon},\prec)\ \text{良序}\Rightarrow \text{任意非空子集存在 }\prec\text{-最小元},\\
  W_2 &: W_1\Rightarrow \text{可对构造步骤/生成元进行规范选取}\Rightarrow \text{得到递归次序},\\
  W_3 &: W_2\Rightarrow \text{对 }X\in\mathsf{Cfg}(\mathsf{Mon})\text{ 存在规范枚举/构造顺序}.
\end{aligned}
\]
\end{Certificate}

\paragraph{证明（谓词因果链）.}
由 $W_1$，任取构造所需的候选生成元集合，均可选取 $\prec$-最小元作为“下一步”；
重复该操作即可得到一个规范递归次序，从而得到 $W_2,W_3$。证毕。

\begin{remark}[备注：与“连续统假设不可判定性”的关系（本笔记\KouJing）]
本笔记不在旧核内证明 CH 的真假；
而是在形式系统中把“不可判定裂缝”统一编码为核内非零障碍类（见第\ref{subsubsec:tex29-ch14-old-hole}小节与第\ref{subsec:tex29-ch14-old}小节），
再把“能否突破”转移到总空间中是否存在证据链 $D$ 使 $\partial D=\omega$（第\ref{sec:tex29-ch14}章）。
因此“从原子到\GZMon{}”的意义在本笔记中被压缩为：对象域是否允许把裂缝转写为可闭合的构造步骤。
\end{remark}

\begin{corollary}[推论：不可判定裂缝的对象域迁移（编码为障碍类 $\Rightarrow$ 转写为可修复对象）]
\label{cor:tex29-ch15-undecidability-transport}
设把某个旧体系命题（例如“核内不可判定/不可延拓”的代表）在本笔记中编码为：
存在 $n\ge 1$ 与 $\omega\in Z^n(A;G)$ 使 $[\omega]\neq 0\in H^n_{\mathrm{obs}}(A;G)$。
若存在基础级地基重构到 $X\in\mathsf{Cfg}(\mathsf{Mon})$，满足：
\[
  \Pi(X)=A,\qquad \mathsf{NoHole}(X),\qquad \iota\bigl(Z^n(A;G)\bigr)\subseteq Z^n(X;G),
\]
则必有 $\iota(\omega)\in B^n(X;G)$，从而该“不可判定裂缝”在新对象域中被转写为可构造性平凡化的对象。
\end{corollary}

\begin{Certificate}[证书：核内非零 $\wedge$ 总空间无空洞 $\Rightarrow$ 提升后为边界]
\label{cert:tex29-ch15-undecidability-transport}
证书登记谓词链：
\[
\begin{aligned}
  U_1 &: [\omega]\neq 0\in H^n_{\mathrm{obs}}(A;G)\Rightarrow \omega\notin B^n(A;G),\\
  U_2 &: \iota\bigl(Z^n(A;G)\bigr)\subseteq Z^n(X;G)\Rightarrow \iota(\omega)\in Z^n(X;G),\\
  U_3 &: \mathsf{NoHole}(X)\Rightarrow Z^n(X;G)=B^n(X;G)\Rightarrow \iota(\omega)\in B^n(X;G),\\
  U_4 &: U_2\wedge U_3\Rightarrow \iota(\omega)\in B^n(X;G).
\end{aligned}
\]
\end{Certificate}

\paragraph{证明（谓词因果链；符号推演）.}
由 $U_2$ 得 $\iota(\omega)$ 为总空间闭链；由 $U_3$（无空洞完成态）得总空间闭链皆为边界，
从而 $\iota(\omega)\in B^n(X;G)$。证毕。

\subsection{四层直觉的无断点拆解：本体论反转\texorpdfstring{$\to$}{->}构造闭环\texorpdfstring{$\to$}{->}跨域同构\texorpdfstring{$\to$}{->}工程收敛}
\label{subsec:tex29-ch15-four-levels}

\subsubsection{第一层：本体论级第一性反转（连续统第一性\texorpdfstring{$\to$}{->}离散统第一性）}
\label{subsubsec:tex29-ch15-level1}

\begin{proposition}[第一性反转的判据：存在核保持的总空间表示且裂缝可在总空间中被转写]
\label{prop:tex29-ch15-level1}
若存在 $X\in\mathsf{Cfg}(\mathsf{Mon})$ 使 $\Pi(X)=A$，并且存在提升映射 $\iota$ 使 $\mathsf{Curv}(X;A)$ 成立，
则在本笔记\KouJing下，连续统对象 $A$ 仅作为观测核出现，而“可构造的第一性对象域”由 $X$ 承担；
此时可将该情形记为“离散统第一性成立于对象层”。
\end{proposition}

\begin{Certificate}[证书：由非平坦性证书导出对象层第一性]
\label{cert:tex29-ch15-level1}
证书登记谓词链：
\[
\begin{aligned}
  L_1 &: \Pi(X)=A\Rightarrow A\ \text{为观测核},\\
  L_2 &: \mathsf{Curv}(X;A)\Rightarrow \text{核内非零类在总空间中可转写为边界},\\
  L_3 &: L_1\wedge L_2\Rightarrow \text{修复/构造能力来自 }X\text{ 而非 }A.
\end{aligned}
\]
\end{Certificate}

\paragraph{证明（谓词因果链）.}
由 $L_1$ 得 $A$ 仅承担观测核角色；由 $L_2$ 得“新增边界能力”发生在 $X$ 中；
因此由 $L_3$ 得第一性对象域应归于 $X$。证毕。

\subsubsection{第二层：结构构造直觉（同伦修复塔从静态分类工具升维为动态构造算法）}
\label{subsubsec:tex29-ch15-level2}

\begin{theorem}[公理降级：修复塔的“可执行构造性”来自核保持--修复--闭合--连通四件套（数学归纳法证书）]
\label{thm:tex29-ch15-level2}
若满足定理~\ref{thm:tex29-ch14-reconstruction} 的前提（核保持、修复步、闭合步、层间连通），
则同伦修复塔不仅可作为对既有对象的静态分类工具，
还可作为一个可证书化的动态构造算法：对每一层都能登记 $(\omega_n,D_n)$ 并通过 $\partial^2=0$ 排除断点。
\end{theorem}

\begin{Certificate}[证书：直接复用重构链存在性证书]
\label{cert:tex29-ch15-level2}
证书登记谓词链：
\[
\begin{aligned}
  A_1 &: \text{核保持}\Rightarrow \Pi(X_n)=A\ (\forall n),\\
  A_2 &: \text{修复步}\Rightarrow \partial D_n=\omega_n,\\
  A_3 &: \text{闭合步}\Rightarrow \partial^2=0\ \text{可核验},\\
  A_4 &: \text{层间连通}\Rightarrow X_{n+1}\simeq X_n,\\
  A_5 &: A_1\wedge A_2\wedge A_3\wedge A_4\Rightarrow \text{形成可执行的证据链数据流 }(X_n,\omega_n,D_n).
\end{aligned}
\]
\end{Certificate}

\paragraph{证明（谓词因果链）.}
由 $A_1$--$A_4$，每一层均可登记“障碍表达 + 修复证据链”，并以 $\partial^2=0$ 校验闭合；
因此由 $A_5$ 得修复塔具备可执行的构造性。证毕。

\subsubsection{第三层：跨域同构直觉（修复谓词作为跨域不变量）}
\label{subsubsec:tex29-ch15-level3}

\begin{definition}[跨域保结构函子（修复结构的同构搬运）]
设 $\mathsf{Dom}_1,\mathsf{Dom}_2$ 为两个异构系统对象域（可理解为“认知域/博弈域/工程域”等的抽象范畴），
它们各自携带障碍表达 $\Omega^\ast(-;G)$、边缘算子 $\partial$ 与修复谓词 $\mathsf{Rep}$。
称一个映射（函子）$F:\mathsf{Dom}_1\to\mathsf{Dom}_2$ 为\textbf{跨域保结构函子}，若满足：
\[
  F_\Omega(\omega)\ \text{保持闭链性且}\ \mathsf{Rep}(\omega)\Rightarrow \mathsf{Rep}\bigl(F_\Omega(\omega)\bigr),
  \qquad
  F_\partial\circ \partial=\partial\circ F_D.
\]
其中 $F_\Omega,F_D,F_\partial$ 表示在障碍表达、证据链与边缘算子层面的相容提升。
\end{definition}

\begin{proposition}[跨域同构：修复谓词与闭合校验在保结构函子下保持]
\label{prop:tex29-ch15-level3}
若 $F$ 为跨域保结构函子，则对任意对象 $X$ 与任意 $\omega\in\Omega^n(X;G)$：
\[
  \mathsf{Rep}(\omega)\Rightarrow \mathsf{Rep}\bigl(F_\Omega(\omega)\bigr),
  \qquad
  \partial^2=0\Rightarrow \partial^2=0\ \text{在像域同样成立}.
\]
因此“修复内核（唯一入口）”可作为跨域不变量被复制与调用。
\end{proposition}

\begin{Certificate}[证书：由函子相容性直接推出]
\label{cert:tex29-ch15-level3}
证书登记谓词链：
\[
\begin{aligned}
  D_1 &: \mathsf{Rep}(\omega)\Rightarrow \exists D\ (\partial D=\omega),\\
  D_2 &: F_\partial\circ \partial=\partial\circ F_D\Rightarrow \partial F_D(D)=F_\Omega(\omega),\\
  D_3 &: D_1\wedge D_2\Rightarrow \mathsf{Rep}\bigl(F_\Omega(\omega)\bigr),\\
  D_4 &: \partial^2=0\Rightarrow F_\partial(\partial^2)=0\Rightarrow \partial^2=0\ \text{在像域成立}.
\end{aligned}
\]
\end{Certificate}

\paragraph{证明（谓词因果链；符号推演）.}
由 $D_1$ 取见证链 $D$；由 $D_2$ 得 $\partial F_D(D)=F_\Omega(\omega)$，从而由 $D_3$ 得修复谓词保持；
闭合性由 $D_4$ 直接推出。证毕。

\subsubsection{第四层：工程收敛直觉（超限兜底 + 有限截断可落地）}
\label{subsubsec:tex29-ch15-level4}

\begin{theorem}[公理降级（ZFC-TR）：超限递归兜底保证全局修复映射存在唯一（超限递归证书）]
\label{thm:tex29-ch15-transfinite-recursion}
设 $\Lambda$ 为某个序数索引（可取全体序数类 $\mathrm{On}$ 的初段），
设 $\mathcal{S}$ 为“单层修复状态”的集合（例如四元组 $(X_\alpha,\omega_\alpha,D_\alpha,\partial^{(\alpha)})$）。
给定一个规则函数 $G$，满足对任意 $\alpha\in\Lambda$：
只要前缀函数 $F\!\upharpoonright\!\alpha$ 已定义，则 $G(F\!\upharpoonright\!\alpha)$ 唯一确定 $F(\alpha)\in\mathcal{S}$，
并且该层状态包含闭合校验通过的候选算子（例如 $(\partial^{(\alpha)})^2=0$）。
则在 ZFC 的超限递归原理下（参见 \cite{Jech2003SetTheory,Kunen2011SetTheory}），存在唯一函数 $F:\Lambda\to\mathcal{S}$ 使得
\[
  F(\alpha)=G(F\!\upharpoonright\!\alpha)\qquad(\forall \alpha\in\Lambda).
\]
因此，一旦把“定位--修复--闭合校验”的规则函数 $G$ 固定为某个确定的选择策略，
超限递归同伦修复塔（作为全序数层的证据链数据流）在形式系统内可被视为\emph{存在且唯一}，从而提供完备性兜底。
\end{theorem}

\begin{Certificate}[证书：规则函数的前缀决定性 + 超限递归存在唯一性]
\label{cert:tex29-ch15-transfinite-recursion}
证书登记谓词链：
\[
\begin{aligned}
  TR_1 &: \forall \alpha\in\Lambda,\ \text{$F\!\upharpoonright\!\alpha$ 已定义}\Rightarrow
    \text{$G(F\!\upharpoonright\!\alpha)$ 唯一确定 $F(\alpha)$},\\
  TR_2 &: \text{ZFC-TR（超限递归原理）}\Rightarrow \exists!\,F:\Lambda\to\mathcal{S}\ \text{使}\ F(\alpha)
    =G(F\!\upharpoonright\!\alpha),\\
  TR_3 &: \text{取 $G$ 为“定位--修复--闭合校验”的确定性规则}\Rightarrow \text{得到一条超限递归修复塔},\\
  TR_4 &: TR_2\wedge TR_3\Rightarrow \text{完备性兜底：只要规则可继续执行，塔即可继续延伸}.
\end{aligned}
\]
\end{Certificate}

\paragraph{证明（谓词因果链）.}
由 $TR_1$ 可把“下一层如何生成”写成前缀决定的规则函数；
由 $TR_2$（ZFC-TR）直接得到存在唯一的全局 $F$；
将 $F(\alpha)$ 解释为第 $\alpha$ 层的修复状态，并以 $TR_3$ 固定规则函数为修复闭环规则，即得 $TR_4$。证毕。

\begin{remark}[备注：与波斯尼科夫塔（Postnikov tower）的结构对齐（仅作参照，不作为本章证明前提）]
在经典代数拓扑中，Postnikov 分层逼近把“同伦类型”拆为逐层的障碍类与 $k$-不变量；
其“已知低层数据 $\Rightarrow$ 生成下一层数据”的构造方式，可被视为本定理中“规则函数 $G$”的典型实例。
本笔记不复述外部教材证明，只固定其在此处的角色：为“超限兜底”的存在性提供一个熟知的同伦论参照系（参见 \cite{Whitehead1978Elements,May1999Concise,
  Switzer1975HomotopyHomology}）。
\end{remark}

\begin{assumption}[有限障碍阶上界（工程/有限资源\KouJing）]
\label{asm:tex29-ch15-finite-bound}
假设对某类工程系统（有限资源可观测系统）存在一个上界 $N\in\NN$，
使得对所有 $n>N$，障碍表达空间在该系统\KouJing下退化为零：
\[
  \Omega^n(X;G)=0\qquad (n>N).
\]
等价地，可假设仅需处理有限阶的障碍闭链。
\end{assumption}

\begin{Certificate}[工程截断证书：有限资源 $\Rightarrow$ 有限障碍阶]
\label{cert:tex29-ch15-finite-bound}
本条假设的可核验内容是：对“有限资源可观测系统”，我们选择一个有限阶截断 $N$ 作为观测/执行接口上界，
并将障碍表达空间定义为该接口可表达的共链空间：
\[
  \Omega^n(X;G):=\mathrm{Span}\{\text{可观测接口可构造的 $n$-阶障碍表达式}\}.
\]
在此\KouJing下，当 $n>N$ 时不存在可构造的 $n$-阶表达式，因此 $\Omega^n(X;G)=0$。
\end{Certificate}

\paragraph{证明（构造性截断；谓词因果链）.}
令“有限资源可观测系统”的接口集只允许：
（i）有限数目的原子观测量；
（ii）有限次合成（例如有限阶的链合成/楔积/括号等）；
（iii）有限阶的证据链存储与校验。
则存在某个最大合成深度 $N$，使得任何高于 $N$ 阶的表达式都无法在该接口内被构造、存储或校验。
按上式定义 $\Omega^n(X;G)$ 为“接口可构造”的线性张成，则对 $n>N$，生成元集合为空，从而 $\Omega^n(X;G)=0$。
证毕。

\begin{remark}[备注：有限上界是“可观测/可执行截断”，不等价于真实同伦群消失断言]
假设~\ref{asm:tex29-ch15-finite-bound} 只固定工程语义：模型在有限资源下只追踪到有限阶的障碍表达，
从而把“无限层的完备性兜底”截断为“有限步可落地”。
它并不声称严格拓扑意义下的高阶同伦群/上同调群必然为零；
它只声称在该工程\KouJing下，$n>N$ 的障碍不再进入可证书化的数据流。
\end{remark}

\begin{theorem}[公理降级：有限阶上界 $\Rightarrow$ 有限步修复收敛（数学归纳法证书）]
\label{thm:tex29-ch15-finite-converge}
在假设~\ref{asm:tex29-ch15-finite-bound} 下，
若对每个 $1\le n\le N$ 的障碍闭链都能在某一层生成证据链 $D$ 使 $\partial D=\omega$ 并通过闭合校验 $\partial^2=0$，
则存在一个有限层数 $K$ 使得在第 $K$ 层之后不再出现新的有效障碍阶，
从而工程\KouJing下的修复过程可在有限步内收敛。
\end{theorem}

\begin{Certificate}[数学归纳法证书：按阶数 $n\le N$ 归纳消解障碍]
\label{cert:tex29-ch15-finite-converge}
令谓词 $Q(n)$ 表示：
“对所有阶数 $\le n$ 的闭链障碍，均已在某层被转写为边界并通过闭合校验。”
证书化要求为：
\[
  Q(1)\ \text{成立},\qquad (\forall n\in\{1,\dots,N-1\})\ \bigl(Q(n)\Rightarrow Q(n+1)\bigr).
\]
并登记有限上界谓词：
\[
  E_1 : n>N\Rightarrow \Omega^n(X;G)=0.
\]
\end{Certificate}

\paragraph{证明（数学归纳法证书；谓词因果链）.}
由假设给出 $Q(1)$；若 $Q(n)$ 成立，则对阶 $n+1$ 的任意闭链障碍可构造证据链实现平凡化并通过闭合校验，
故 $Q(n+1)$ 成立。由归纳法得 $Q(N)$ 成立。
再由 $E_1$ 得 $n>N$ 无有效障碍表达，因此在有限阶意义下不存在继续上升的修复需求；
取 $K$ 为覆盖 $\le N$ 阶修复所需的最大层即可。证毕。

\begin{remark}[备注：超限兜底与有限截断的关系]
超限递归修复塔提供“完备性兜底”（理论上覆盖任意阶数的障碍）；
而有限上界假设把现实系统的可执行路径截断为有限步。
这对应“无限的公理完备性为有限工程落地提供兜底”的结构化表述。
\end{remark}

\subsection{文明跃迁的形式化判据：从连续统垄断到离散统平权}
\label{subsec:tex29-ch15-civilization}

\begin{definition}[文明台阶谓词（破垄断 + 新框架）]
称某个文明运行范式 $\mathcal{C}$ 达到\textbf{台阶跃迁}，若同时满足两条谓词：
\begin{enumerate}
  \item （破垄断）存在一类核心能力在旧范式中由少数节点垄断，而在新范式中可被标准化、证书化并复制到任意节点；
  \item （新框架）新范式提供一套无内生断点的底层校验框架，使得“扩展/修复”可在其内部闭环执行（例如由 $\partial^2=0$ 与证书链约束）。
\end{enumerate}
\end{definition}

\begin{proposition}[离散统第一性 + 修复塔标准化 $\Rightarrow$ 满足台阶跃迁判据]
\label{prop:tex29-ch15-civilization}
若对象域采用离散统第一性（存在\GZMonization{}总空间 $X\in\mathsf{Cfg}(\mathsf{Mon})$），
并且修复过程可被证书化为标准化链数据流（定理~\ref{thm:tex29-ch15-level2}），
则在本笔记\KouJing下：
\begin{enumerate}
  \item “从障碍到修复”的构造能力不再依赖少数外部权威节点，而可被复制到任意可执行证书链的节点（破垄断）；
  \item $\partial^2=0$ 与证书链提供闭合校验，构成无断点的新框架（新框架）。
\end{enumerate}
从而文明台阶谓词成立。
\end{proposition}

\begin{Certificate}[证书：标准化证书链的可复制性 + 闭合校验]
\label{cert:tex29-ch15-civilization}
证书登记谓词链：
\[
\begin{aligned}
  C_1 &: \text{修复可证书化}\Rightarrow \text{存在可公开传递的构造步骤与校验步骤},\\
  C_2 &: C_1\Rightarrow \text{能力可复制到任意执行同一证书链的节点}\Rightarrow \text{破垄断谓词},\\
  C_3 &: \partial^2=0\ \text{与证书链}\Rightarrow \text{闭环校验}\Rightarrow \text{新框架谓词},\\
  C_4 &: C_2\wedge C_3\Rightarrow \text{文明台阶谓词成立}.
\end{aligned}
\]
\end{Certificate}

\paragraph{证明（谓词因果链）.}
由 $C_1$ 得修复步骤被编码为可复现的证书链；因此由 $C_2$ 得其不再是少数节点的不可复制“灵感”，满足破垄断谓词；
由 $C_3$ 得体系提供内生闭环校验框架；合并得 $C_4$。证毕。

\begin{remark}[备注：关于“当下为何可能无感”】【形式系统的最小解释】]
“无感”在本笔记\KouJing下可被理解为两类延迟叠加：
\begin{enumerate}
  \item （障碍阶延迟）大量节点仍停留在静态核 $A$ 的操作\KouJing内（推论~\ref{cor:tex29-core-contraction-dynamic-necessity}），
  因而只能识别障碍而不能构造性消解；
  \item （扩散延迟）证书链的可复制性并不等同于“即时扩散完成”，扩散过程本身可被建模为跨域函子的可达性问题（命题~\ref{prop:tex29-ch15-level3}）。
\end{enumerate}
本备注只给出最小的形式化解释框架，不对现实社会过程做经验断言。
\end{remark}

\clearpage

%%%%%%%%%%%%%%%%%%%%%%%%%%%%%%%%%%%%%%%%%%%%%%%%%%%%%%%%%%%%
\section{微分几何--代数拓扑跨分支同构：雅可比障碍、比安基恒等式与边缘算子闭合}
\label{sec:tex29-ch16}
%%%%%%%%%%%%%%%%%%%%%%%%%%%%%%%%%%%%%%%%%%%%%%%%%%%%%%%%%%%%

\begin{remark}[核心陈述（同构不是比喻，而是保结构映射\KouJing）]
本章把“微分几何（黎曼流形）中的测地线变分/雅可比场/共轭点/比安基恒等式”
与“代数拓扑（链复形）中的延拓路径/同伦变换/边缘障碍/闭合公理 $\partial^2=0$”
固定为一套\emph{可证书化的保结构对应}，并给出其在本笔记形式系统内的自洽闭环：
\begin{enumerate}
  \item 变分极值路径 $\leftrightarrow$ 约束下的可接受延拓路径（缺陷势/能量泛函\KouJing）；
  \item 雅可比场（变分方向）$\leftrightarrow$ 同伦变换（连续形变）；
  \item 共轭点导致的延拓断点 $\leftrightarrow$ 待闭合链复形中的边缘障碍（候选算子闭合失败）；
  \item 比安基恒等式作为“良定义前提”$\leftrightarrow$ 闭合公理 $\partial^2=0$ 作为“链复形良定义前提”。
\end{enumerate}
该对应在本笔记中被理解为：把“几何侧的延拓失败判据”转写为“链侧的可检验障碍表达”，
从而把静态判定升级为可执行的“定位--修复--闭合校验”流程（与第\ref{sec:tex29-ch7}章--第\ref{sec:tex29-ch14}章的修复内核一致）。
与本仓库更早的“边缘算子闭合/变分修复塔”推演链可互相参照（如 \cite{RepoNB17EdgeOperator}）。
\end{remark}

\begin{remark}[对齐表：几何侧--链侧的公理级对应（本章\KouJing不变）]
为把“严格同构项”压缩为一眼可见的参照点，给出对齐表（其中“雅可比判别”覆盖教材中对共轭点的各类可检验判据，
有时亦被非正式称作“雅可比不等式/判别式”；本笔记不固定其分量公式，仅将其作为 $\mathsf{JacObs}$ 的证书化代理谓词）：
\[
  \mathsf{JacCrit}(\gamma)\Rightarrow \mathsf{JacObs}(\gamma).
\]
\begin{center}
\renewcommand{\arraystretch}{1.25}
\begin{tabular}{|p{0.29\linewidth}|p{0.29\linewidth}|p{0.34\linewidth}|}
\hline
\textbf{微分几何（黎曼流形）} & \textbf{链/代数拓扑（链复形\KouJing）} & \textbf{公理级对应规则（本笔记）} \\
\hline
测地线 $\gamma$（$\mathcal{E}$ 的极值路径） &
可接受延拓路径（缺陷势 $\Phi$ 的下降/不增路径） &
以 $\mathcal{E}$ 或 $\Phi$ 作为入口泛函，极值/单调性给出“可接受更新”判据 \\
\hline
雅可比场 $J$（测地线的变分方向） &
同伦变换见证（连续形变数据） &
变分方向对应映射的连续形变，不改变同伦类不变量 \\
\hline
共轭点 / 雅可比障碍 $\mathsf{JacObs}$ &
边缘障碍 $\mathsf{EdgeObs}(\widetilde{\partial})\equiv(\widetilde{\partial}^{\,2}\neq 0)$ &
延拓失败信号 $\leftrightarrow$ 闭合失败信号（在第\ref{subsec:tex29-ch16-isomorphism}小节固定为工作\KouJing） \\
\hline
（可检验）雅可比判别 $\mathsf{JacCrit}$（共轭点判据代理谓词） &
非闭合性判定 $\widetilde{\partial}^{\,2}\neq 0$ &
“判别触发/违背”被统一编码为闭合失败谓词，从而可进入“定位--修复--闭合”流程 \\
\hline
比安基一致性谓词 $\mathsf{Bianchi}(M,g)$ &
闭合公理 $\partial^2=0$ &
良定义前提的对应：只有在闭合校验成立时，修复塔层数据才可被证书化并拼接 \\
\hline
\end{tabular}
\end{center}
\end{remark}

\subsection{对象域：测地线变分骨架与待闭合链复形骨架}
\label{subsec:tex29-ch16-domains}

\begin{definition}[测地线：能量泛函的极值路径（几何侧；参见 \cite{Lee2018RiemannianManifolds,doCarmo1992Riemannian,
  Petersen2016RiemannianGeometry}）]
设 $(M,g)$ 为黎曼流形，$p,q\in M$。
令 $\mathcal{P}(p,q)$ 为端点固定的分段光滑路径集合。
定义能量泛函
\[
  \mathcal{E}(\gamma):=\frac12\int_{0}^{1}\bigl\langle \dot\gamma(t),\dot\gamma(t)\bigr\rangle_g\,dt.
\]
称 $\gamma\in\mathcal{P}(p,q)$ 为\textbf{测地线}，若它是 $\mathcal{E}$ 的临界点（对所有端点固定变分一阶变分为零）。
\end{definition}

\begin{definition}[雅可比场、共轭点与雅可比障碍（几何侧；参见 \cite{Lee2018RiemannianManifolds,doCarmo1992Riemannian,
  Petersen2016RiemannianGeometry}）]
设 $\gamma$ 为测地线。
称沿 $\gamma$ 的向量场 $J$ 为\textbf{雅可比场}，若存在一族端点固定的测地线变分 $\gamma_s$ 使得
\[
  J(t)=\left.\frac{\partial}{\partial s}\right|_{s=0}\gamma_s(t).
\]
称 $t_0\in(0,1]$ 为\textbf{共轭点}，若存在非零雅可比场 $J$ 满足 $J(0)=J(t_0)=0$。
定义\textbf{雅可比障碍谓词}：
\[
  \mathsf{JacObs}(\gamma)\ :\Longleftrightarrow\ \exists t_0\in(0,1]\ \ (\gamma(t_0)\ \text{与}\ \gamma(0)\ \text{共轭}).
\]
\end{definition}

\begin{definition}[待闭合链复形与边缘障碍（链侧；参见 \cite{Hatcher2002AT,Weibel1994HomologicalAlgebra}）]
设 $C_*=\{C_n,\partial_n\}$ 为分次模（或向量空间）族，$\partial_n:C_n\to C_{n-1}$ 为线性算子族。
称 $C_*$ 为\textbf{良定义链复形}，若满足闭合公理
\[
  \partial_{n-1}\circ \partial_{n}=0\qquad (\forall n),
\]
即 $\partial^2=0$（见第\ref{sec:tex29-ch7}章的闭合校验\KouJing）。
若某候选算子族 $\widetilde{\partial}$ 不满足 $\widetilde{\partial}^2=0$，则称其为\textbf{待闭合链复形}并定义
\[
  \mathsf{EdgeObs}(\widetilde{\partial})\ :\Longleftrightarrow\ \widetilde{\partial}^{\,2}\neq 0,
\]
称其为\textbf{边缘障碍}（闭合失败的可检验断点信号）。
\end{definition}

\begin{remark}[术语对齐：变分缺陷势与能量泛函的统一接口]
在更一般的动力学表述方式中，本笔记以缺陷势 $\Phi$ 表达“可接受更新的单调性”（见第\ref{subsec:tex29-fractal-dynamic-loop}小节与第\ref{subsec:tex29-dynamics-meta-loop}小节）。
在本章几何侧，$\mathcal{E}$ 扮演同一角色：它给出“极值路径/可接受路径”的判据入口。
因此本章把“测地线”作为“变分极值路径”的一个具体化载体，并不改变修复体系的内核入口（仍以 $\mathsf{Rep}$ 与闭合校验为准）。
\end{remark}

\subsection{公理降级：雅可比障碍\texorpdfstring{$\Leftrightarrow$}{<->}边缘障碍；比安基\texorpdfstring{$\Leftrightarrow$}{<->}闭合公理}
\label{subsec:tex29-ch16-isomorphism}

\begin{definition}[跨分支同构数据与对应谓词]
称一组数据
\[
  \mathbb{I}:=\bigl((M,g),(\widetilde{C}_*,\widetilde{\partial}),\mathcal{T}\bigr)
\]
为一份\textbf{微分几何--代数拓扑跨分支同构数据}，若存在一个“对应翻译” $\mathcal{T}$，使得对几何侧候选延拓路径 $\gamma$ 与链侧候选算子 $\widetilde{\partial}$：
\begin{enumerate}
  \item （极值--可接受）$\gamma$ 为能量泛函极值路径 $\Rightarrow$ $\mathcal{T}(\gamma)$ 为修复塔中的可接受延拓路径（以 $\Phi/\mathcal{E}$ 作为入口泛函）；
  \item （变分--同伦）雅可比场 $J$ 的存在对应于链侧的同伦变换数据（连续形变见证）；
  \item （障碍对应）$\mathsf{JacObs}(\gamma)\ \Leftrightarrow\ \mathsf{EdgeObs}(\widetilde{\partial})$；
  \item （良定义前提对应）几何侧“比安基一致性谓词” $\mathsf{Bianchi}(M,g)$ 被编码为链侧闭合公理 $\widetilde{\partial}^{\,2}=0$。
\end{enumerate}
其中 $\mathsf{Bianchi}(M,g)$ 在本章仅作为“几何侧良定义一致性”的谓词符号使用，不展开其张量分量表达。
\end{definition}

\begin{theorem}[公理降级：跨分支同构的一致性与无矛盾性（数学归纳法证书）]
\label{thm:tex29-ch16-isomorphism-consistency}
设 $\mathbb{I}=((M,g),(\widetilde{C}_*,\widetilde{\partial}),\mathcal{T})$ 为跨分支同构数据。
对任意以层数 $n\in\NN$ 索引的候选延拓过程，若每层都执行闭合校验并在必要时更新候选算子使其闭合：
\[
  \widetilde{\partial}^{(n)}\ \leadsto\ \partial^{(n)}\quad\text{且}\quad (\partial^{(n)})^2=0,
\]
则：
\begin{enumerate}
  \item （无矛盾）在通过闭合校验的层上，链侧不存在边缘障碍；因此在同构\KouJing下，几何侧不存在对应的雅可比障碍；
  \item （可修复性入口保持）若某层出现 $\mathsf{JacObs}$（等价于链侧闭合失败信号），则它被登记为障碍表达并触发“定位--修复--再闭合”的标准流程；
  \item （层级稳定性）若对所有 $k\le n$ 均完成“修复并闭合”，则第 $n$ 层之前的对应关系保持一致，不会产生互相矛盾的断点判据。
\end{enumerate}
\end{theorem}

\begin{Certificate}[数学归纳法证书：按层数 $n$ 归纳的“候选--障碍--闭合”一致性]
\label{cert:tex29-ch16-isomorphism-consistency}
令谓词 $P(n)$ 表示：
“对前 $n$ 层候选数据，若每层最终都获得闭合算子 $\partial^{(k)}$ 满足 $(\partial^{(k)})^2=0$，
则链侧无边缘障碍；并且在同构\KouJing下，几何侧无对应雅可比障碍；且对应翻译 $\mathcal{T}$ 在前 $n$ 层保持一致。”
证书化要求为：
\[
  P(0)\ \text{成立},\qquad (\forall n\in\NN)\ \bigl(P(n)\Rightarrow P(n+1)\bigr).
\]
\end{Certificate}

\paragraph{证明（数学归纳法证书；谓词因果链）.}
由证书~\ref{cert:tex29-ch16-isomorphism-consistency}：
\begin{itemize}
  \item 基步：$n=0$ 时，闭合校验给出 $(\partial^{(0)})^2=0$，故链侧无 $\mathsf{EdgeObs}$；由同构数据第(3)(4)条可得几何侧无对应 $\mathsf{JacObs}$；因此 $P(0)
    $ 成立。
  \item 归纳步：假设 $P(n)$ 成立。对第 $n+1$ 层，若候选算子 $\widetilde{\partial}^{(n+1)}$ 闭合失败，则由同构数据第(3)条视为出现 $\mathsf{JacObs}$；
  执行“定位--修复--再闭合”后得到闭合算子 $\partial^{(n+1)}$ 使 $(\partial^{(n+1)})^2=0$，从而链侧不再有边缘障碍信号；
  因此几何侧也不再有对应雅可比障碍信号，并且既有层的对应关系由假设 $P(n)$ 保持一致，故 $P(n+1)$ 成立。
\end{itemize}
由归纳法，$P(n)$ 对任意 $n$ 成立，定理得证。证毕。

\begin{remark}[备注：本章“等价”是工作\KouJing，不是外部教材命题]
本章的“$\mathsf{JacObs}\Leftrightarrow \mathsf{EdgeObs}$”并非对外部标准教材结论做复述；
它在本笔记中被固定为\emph{跨分支翻译的工作\KouJing}：用于把几何侧的延拓失败信号编码为链侧的可检验闭合失败信号，
从而把“终止判据”转写为“可修复对象”的入口。
\end{remark}

\subsection{定理：同伦修复塔与变分法修复雅可比障碍的标准化流程}
\label{subsec:tex29-ch16-variational-tower}

\begin{theorem}[公理降级：变分定位\texorpdfstring{$\to$}{->}分层修复\texorpdfstring{$\to$}{->}闭合校验的四步闭环（数学归纳法证书）]
\label{thm:tex29-ch16-variational-tower}
设给定一个待延拓的认知映射（或系统态映射）$f$，其每一层延拓均可被几何侧的“极值路径 + 变分”定位为候选路径 $\gamma_n$，
并可在链侧被编码为障碍表达 $\omega_n$ 与候选算子 $\widetilde{\partial}^{(n)}$。
若对每一层 $n\in\NN$，存在证据链 $D_n$ 使得
\[
  \partial^{(n)} D_n=\omega_n,
  \qquad
  (\partial^{(n)})^2=0,
\]
并且候选路径更新满足变分可接受性（例如缺陷势/能量不增），
则存在一条可证书化的分层修复流程，使 $f$ 可逐层完成延拓并在每层通过闭合校验，从而不产生次生断点。
\end{theorem}

\begin{Certificate}[数学归纳法证书：按层数归纳的“极值定位--障碍登记--修复--闭合”闭环]
\label{cert:tex29-ch16-variational-tower}
令谓词 $R(n)$ 表示：
“已在前 $n$ 层完成：极值路径定位 $\gamma_k$、障碍登记 $\omega_k$、修复证据链 $D_k$（$\partial^{(k)}D_k=\omega_k$）、以及闭合校验 $(\partial^{(k)})
  ^2=0$（对所有 $k\le n$），
并且变分可接受性在层间保持。”
证书化要求为：
\[
  R(0)\ \text{成立},\qquad (\forall n\in\NN)\ \bigl(R(n)\Rightarrow R(n+1)\bigr).
\]
\end{Certificate}

\paragraph{证明（数学归纳法证书；谓词因果链）.}
由证书~\ref{cert:tex29-ch16-variational-tower}：
\begin{itemize}
  \item 基步：$n=0$ 时，由前提给出 $\gamma_0$、$\omega_0$、$D_0$ 与闭合校验成立，故 $R(0)$ 成立；
  \item 归纳步：假设 $R(n)$ 成立。由前提对第 $n+1$ 层同样可构造 $\gamma_{n+1}$、$\omega_{n+1}$ 与证据链 $D_{n+1}$，
  并通过闭合校验得到 $(\partial^{(n+1)})^2=0$；同时变分可接受性由假设保持；故 $R(n+1)$ 成立。
\end{itemize}
由归纳法，流程可对任意层继续执行并保持无断点闭环。证毕。

\begin{corollary}[推论：雅可比障碍被转写为“可修复障碍表达”】【与第\ref{sec:tex29-ch7}章一致]
\label{cor:tex29-ch16-jacobi-to-rep}
在定理~\ref{thm:tex29-ch16-variational-tower} 的前提下，
每一层的雅可比障碍信号（若出现）都被编码为障碍表达 $\omega_n$，
并通过证据链 $D_n$ 满足 $\partial^{(n)}D_n=\omega_n$ 被构造性平凡化；
因此“雅可比障碍”在本笔记\KouJing下不再是终止判据，而是修复流程的输入对象。
\end{corollary}

\begin{Certificate}[证书：由修复约束的直接翻译]
\label{cert:tex29-ch16-jacobi-to-rep}
证书登记谓词链：
\[
\begin{aligned}
  J_1 &: \mathsf{JacObs}\ \text{出现}\Rightarrow \text{登记为障碍表达 }\omega_n,\\
  J_2 &: \exists D_n\ (\partial^{(n)}D_n=\omega_n)\Rightarrow \mathsf{Rep}(\omega_n),\\
  J_3 &: J_1\wedge J_2\Rightarrow \text{雅可比障碍被转写为可修复对象}.
\end{aligned}
\]
\end{Certificate}

\paragraph{证明（谓词因果链）.}
由 $J_1$ 给出障碍登记；由 $J_2$ 得修复谓词成立；合并得 $J_3$。证毕。

\subsection{命题：边缘算子作为链映射的“诱导”与跨同伦类障碍的平凡化}
\label{subsec:tex29-ch16-chainmap}

\begin{definition}[核心链复形、待闭合链复形与修复型链同构]
设核心链复形为
\[
  C_*=\{C_n,\partial_n\},\qquad \partial^2=0,
\]
其同调群记为 $H_*(C)$。
设待闭合链复形为
\[
  D_*=\{D_n,d_n\},\qquad d^2\neq 0\ \ \text{（候选闭合失败信号）}.
\]
称一个分次同构 $\phi:D_*\to C_*$ 为\textbf{修复型链同构}，若以
\[
  d'_n:=\phi_n^{-1}\circ \partial_n\circ \phi_n
\]
诱导出的算子族 $d'$ 使得 $(D_*,d')$ 成为良定义链复形（即 $(d')^2=0$）。
该\KouJing可理解为：把核心边缘算子 $\partial$ 的“可复制形式”通过 $\phi$ 迁移到 $D$ 上，完成待闭合对象的闭合化。
\end{definition}

\begin{remark}[工程化解释：把“讲解/录音/代码”视为可复制的链同构载体]
在认知/工程语境中，$\phi$ 可被具体化为一条可公开传递的“规则载体”（例如一段讲解、录音或可执行代码）：
它不要求把两个节点在对象层面合并，也不要求节点互通；
它只提供一个把核心闭合算子 $\partial$ 迁移到待闭合对象 $D$ 的可执行接口，
从而在 $D$ 上诱导出满足 $(d')^2=0$ 的闭合算子 $d'$，并使其同调不变量与核心对象在该\KouJing下对齐。
\end{remark}

\begin{lemma}[诱导闭合：$\partial^2=0\Rightarrow (d')^2=0$]
\label{lem:tex29-ch16-induced-closure}
在上述定义下，若 $\partial^2=0$ 且 $\phi$ 为分次同构，则诱导算子 $d'$ 满足 $(d')^2=0$。
\end{lemma}

\begin{Certificate}[证书：共轭变换保持幂零性]
\label{cert:tex29-ch16-induced-closure}
证书登记谓词链：
\[
\begin{aligned}
  I_1 &: d'=\phi^{-1}\circ \partial\circ \phi,\\
  I_2 &: \partial^2=0,\\
  I_3 &: I_1\wedge I_2\Rightarrow (d')^2=\phi^{-1}\circ \partial^2\circ \phi=0.
\end{aligned}
\]
\end{Certificate}

\paragraph{证明（谓词因果链；符号推演）.}
由 $I_1$ 与算子合成得 $(d')^2=\phi^{-1}\partial\phi\phi^{-1}\partial\phi=\phi^{-1}\partial^2\phi$；
由 $I_2$ 得 $\partial^2=0$，故 $(d')^2=0$。证毕。

\begin{proposition}[同调并入：修复型链同构 $\Rightarrow$ 同调群同构]
\label{prop:tex29-ch16-homology-embed}
若 $\phi$ 为修复型链同构，则 $(D_*,d')$ 与 $(C_*,\partial)$ 链同构，从而
\[
  H_*(D,d')\cong H_*(C,\partial).
\]
\end{proposition}

\begin{Certificate}[证书：链同构诱导同调同构]
\label{cert:tex29-ch16-homology-embed}
证书登记谓词链：
\[
\begin{aligned}
  H_1 &: d'=\phi^{-1}\partial\phi\Rightarrow \phi\circ d'=\partial\circ \phi,\\
  H_2 &: H_1\Rightarrow \phi\ \text{为链映射且可逆},\\
  H_3 &: H_2\Rightarrow \phi_*\ \text{诱导同调群同构 }H_*(D,d')\cong H_*(C,\partial).
\end{aligned}
\]
\end{Certificate}

\paragraph{证明（谓词因果链）.}
由 $H_1$ 得链映射相容性；由 $\phi$ 可逆得 $H_2$；链同构诱导同调同构即得 $H_3$。证毕。

\begin{corollary}[跨同伦类上同调障碍的平凡化：闭合化 $\Rightarrow$ 延拓障碍可消解]
\label{cor:tex29-ch16-obstruction-trivialized}
把“映射 $f$ 无法从 $n$-骨架延拓到 $n{+}1$-骨架”的信息编码为障碍类
\[
  c_f\in H^{n+1}_{\mathrm{obs}}(D;\pi_n(C))
\]
（作为抽象记账对象，见第\ref{sec:tex29-ch7}章；外部障碍论/延拓障碍的标准\KouJing参见 \cite{Whitehead1978Elements,Switzer1975HomotopyHomology}。）
若存在修复型链同构 $\phi$ 使 $(D_*,d')$ 闭合良定义，
则该障碍类在闭合\KouJing下可被平凡化：其对应的“不可延拓断点”被转写为可修复对象，并可在修复塔中被消解。
\end{corollary}

\begin{Certificate}[证书：闭合良定义 $\Rightarrow$ 障碍群\KouJing可用 $\Rightarrow$ 可进入修复流程]
\label{cert:tex29-ch16-obstruction-trivialized}
证书登记谓词链：
\[
\begin{aligned}
  O_1 &: (d')^2=0\Rightarrow (D_*,d')\ \text{良定义}\Rightarrow H^{n+1}_{\mathrm{obs}}(D;\pi_n(C))\ \text{可登记},\\
  O_2 &: \text{可登记}\Rightarrow \text{可对 }c_f\text{ 进行 }\mathsf{Rep}\text{ 判定与证据链构造},\\
  O_3 &: O_1\wedge O_2\Rightarrow \text{“不可延拓断点”被转写为可修复对象并可被消解}.
\end{aligned}
\]
\end{Certificate}

\paragraph{证明（谓词因果链）.}
由 $O_1$ 得障碍群\KouJing成立；由 $O_2$ 得进入修复谓词入口；合并得 $O_3$。证毕。

\begin{remark}[备注：节点独立性与“同伦等价”的区别]
本章的“并入/同构”只在结构层面发生：它把待闭合对象 $D$ 的边缘算子闭合化，并使同调不变量与核心对象一致；
这并不等价于“节点互通”或“节点合并”：
节点间的刚性独立性可被表述为“对象仍为不同的生成元/不同的基元”，
而“同伦等价”仅意味着它们在所选不变量\KouJing下处于同一同伦类、满足同一闭合校验。
\end{remark}

\clearpage

%%%%%%%%%%%%%%%%%%%%%%%%%%%%%%%%%%%%%%%%%%%%%%%%%%%%%%%%%%%%
\section{障碍阶数与修复门槛：有限递归刚性约束下的层级谱}
\label{sec:tex29-ch17}
%%%%%%%%%%%%%%%%%%%%%%%%%%%%%%%%%%%%%%%%%%%%%%%%%%%%%%%%%%%%

\begin{remark}[核心陈述（“高阶上同调障碍”的形式化定义）]
本章把“高阶上同调障碍的阶数 $n$”固定为一个可计算的最小断点指标：
\[
  n=\min\Bigl\{k\in\NN_0:\ \text{映射无法从 $k$-维骨架延拓到 $k{+}1$-维骨架}\Bigr\},
\]
并把“修复难度必须逐层推进”固定为有限递归刚性约束：
每一层修复都必须通过闭合校验（$\partial^2=0$）才允许进入下一层；跳步将导致“塔稳定性”断裂，修复失败。
本章只在形式系统内给出\emph{结构性结论}（单调性、不可跳步性、有限上界），
不对现实世界中的人群比例/历史评价等经验性量作数学定理断言（这些只可作为外部刻度）。
\end{remark}

\subsection{定义：骨架延拓谓词、障碍类与最小障碍阶}
\label{subsec:tex29-ch17-order}

\begin{definition}[骨架延拓谓词与障碍类记账]
设 $X$ 为源对象、$Y$ 为目标对象（可理解为两套认知链复形/系统态空间），
记 $X^{(k)}$ 为其 $k$-维骨架（或 $k$ 层可观测子结构）。
设 $f^{(k)}:X^{(k)}\to Y$ 为在 $k$-骨架上已定义的映射。
定义延拓谓词：
\[
  \mathsf{Ext}(f^{(k)};k{+}1)\ :\Longleftrightarrow\ \exists f^{(k+1)}:X^{(k+1)}\to Y\ \ \text{使}\ \ f^{(k+1)}|_{X^{(k)}}
    =f^{(k)}.
\]
若 $\neg\mathsf{Ext}(f^{(k)};k{+}1)$，则把该失败在本笔记中记账为一个障碍类
\[
  c_{f,k}\in H^{k+1}_{\mathrm{obs}}\bigl(X;\pi_k(Y)\bigr)
\]
（作为抽象记账对象；其“可修复性”由谓词 $\mathsf{Rep}$ 入口判定，见第\ref{sec:tex29-ch7}章）。
\end{definition}

\begin{definition}[最小障碍阶（阶数 $n$ 的严格定义）]
在上述\KouJing下，定义映射 $f$ 的\textbf{最小障碍阶}为
\[
  \mathrm{ord}(f):=\min\bigl\{k\in\NN_0:\neg\mathsf{Ext}(f^{(k)};k{+}1)\bigr\},
\]
若对所有 $k$ 都满足 $\mathsf{Ext}(f^{(k)};k{+}1)$，则记 $\mathrm{ord}(f)=+\infty$（表示在该骨架序列上无有限阶断点）。
当 $\mathrm{ord}(f)\ge 4$ 时，称其为\textbf{高阶}障碍（本章仅作为术语约定，不引入外部价值判断）。
\end{definition}

\begin{assumption}[有限递归刚性约束（逐层修复 + 闭合校验）]
\label{asm:tex29-ch17-finite-recursion}
假设修复过程必须遵循以下刚性约束：
\begin{enumerate}
  \item （逐层）只能在 $k$ 层完成修复并通过闭合校验后，才允许进入 $k{+}1$ 层；
  \item （闭合）每层修复后的候选算子必须满足 $\partial^2=0$（或等价闭合谓词）；
  \item （跳步禁令）任何跳过某层闭合校验的“直接升阶”操作，均视为修复塔坍塌：高层结构不可再被证书化。
\end{enumerate}
\end{assumption}

\begin{Certificate}[过程不变量证书：逐层 + 闭合 + 跳步禁令可被实现为可核验约束]
\label{cert:tex29-ch17-finite-recursion}
本条假设被用作后续“塔稳定单调性”推演的前提输入。
其可核验内容是：修复算法（或修复规程）可以被实现为带不变量检查的有限状态机，使得：
\[
  \text{进入层 }(k{+}1)\ \Longrightarrow\ \mathsf{Stable}(k)\ \text{已被证书化，且当前层闭合校验通过。}
\]
\end{Certificate}

\paragraph{证明（不变量归纳；谓词因果链）.}
以层级索引 $k$ 做归纳，维护不变量：
\[
  I(k):=\bigl(\text{系统处于第 $k$ 层}\bigr)\Rightarrow\bigl(\mathsf{Stable}(k{-}1)\ \text{已证书化}\bigr),
\]
其中 $k=0$ 时不变量真（无前一层）。
算法规则规定：从第 $k$ 层迁移到第 $k{+}1$ 层之前，必须先执行闭合校验并生成证书，因而 $I(k{+}1)$ 成立。
同时，任何绕过该检查的迁移在规程中被定义为非法转移（跳步禁令），因此不会出现在可证书化轨迹中。
证毕。

\subsection{公理降级：不可跳步性与塔稳定单调性（数学归纳法证书）}
\label{subsec:tex29-ch17-no-skip}

\begin{theorem}[公理降级：塔稳定单调性（低层不稳 $\Rightarrow$ 高层不可定义）（数学归纳法证书）]
\label{thm:tex29-ch17-stability-monotone}
在假设~\ref{asm:tex29-ch17-finite-recursion} 下，定义第 $k$ 层的稳定谓词为
\[
  \mathsf{Stable}(k):=\bigl(\mathsf{Ext}(f^{(k)};k{+}1)\bigr)\ \wedge\ \bigl(\partial^{(k)}\bigr)^2=0.
\]
则对任意 $k\in\NN_0$，若 $\neg\mathsf{Stable}(k)$，则 $\neg\mathsf{Stable}(k{+}1)$。
因此，任何“跳过第 $k$ 层稳定性而直接处理第 $k{+}1$ 层”的方案都无法被证书化。
\end{theorem}

\begin{Certificate}[数学归纳法证书：稳定谓词的前缀闭包]
\label{cert:tex29-ch17-stability-monotone}
令谓词 $P(k)$ 表示：“若对某层 $k$ 有 $\neg\mathsf{Stable}(k)$，则对所有 $m\ge k$ 都有 $\neg\mathsf{Stable}(m)$。”
证书化要求为：
\[
  P(0)\ \text{成立},\qquad (\forall k\in\NN)\ \bigl(P(k)\Rightarrow P(k+1)\bigr).
\]
\end{Certificate}

\paragraph{证明（数学归纳法证书；谓词因果链）.}
由假设~\ref{asm:tex29-ch17-finite-recursion} 的“逐层/跳步禁令”可知：
一旦某层缺失延拓或闭合校验失败，则更高层所依赖的骨架数据与闭合算子均不可用；
因此“低层不稳 $\Rightarrow$ 高层不可定义”的蕴含成立。
将其登记为 $P(k)$ 的归纳步即可得到对任意 $m\ge k$ 的不可稳定性结论。证毕。

\begin{corollary}[推论：修复难度的阶数单调性（结构结论）]
\label{cor:tex29-ch17-difficulty-monotone}
在同一修复\KouJing下，若 $\mathrm{ord}(f)=n$，则修复到稳定态至少需要完成从 $0$ 到 $n$ 的全部层级稳定性证书化；
因此“需要通过的稳定层数”随 $n$ 单调不减。
该推论仅表达结构性门槛单调性，不对具体人群概率作定理断言。
\end{corollary}

\begin{Certificate}[结构推论证书：由塔稳定单调性推出“所需稳定层数”随阶数单调不减]
\label{cert:tex29-ch17-difficulty-monotone}
证书化目标谓词为：
\[
  Q(n):=\Bigl(\mathrm{ord}(f)=n\Bigr)\Rightarrow
  \Bigl(\forall k\in\{0,1,\dots,n\},\ \mathsf{Stable}(k)\Bigr).
\]
\end{Certificate}

\paragraph{证明（谓词因果链；符号推演）.}
取任意满足 $\mathrm{ord}(f)=n$ 的修复对象，并假设修复到稳定态意味着 $\mathsf{Stable}(n)$ 为真。
反设存在 $k\le n$ 使 $\neg\mathsf{Stable}(k)$。
由定理~\ref{thm:tex29-ch17-stability-monotone}（低层不稳 $\Rightarrow$ 高层不可定义），得 $\neg\mathsf{Stable}(n)$，与 $\mathsf{Stable}(n)$
  矛盾。
故对所有 $k\in\{0,\dots,n\}$ 必有 $\mathsf{Stable}(k)$，从而达到稳定态至少需要完成从 $0$ 到 $n$ 的全部层级稳定性证书化。
因此“需要通过的稳定层数”随 $n$ 单调不减。证毕。

\subsection{备注：障碍阶的解释表（0--5 阶与超限兜底；启发式，不纳入定理）}
\label{subsec:tex29-ch17-table}

\begin{remark}[对齐表：阶数的语义与典型修复动作（示意）]
下表给出一种常用解释\KouJing，把“最小障碍阶”映射为更便于沟通的语义层级；它用于沟通与工程分工，不作为形式系统内可证明命题：
\begin{center}
\renewcommand{\arraystretch}{1.25}
\begin{tabular}{|c|p{0.30\linewidth}|p{0.43\linewidth}|}
\hline
\textbf{阶数} & \textbf{数学\KouJing（最小断点）} & \textbf{解释\KouJing（示意）} \\
\hline
$0$ & 无法进入同一连通分支（$\pi_0$ 断裂） & 0 阶准入：拒绝进入对象域/拒绝承认裂缝可编码为障碍类 \\
\hline
$1$ & 基本群\KouJing不对齐导致延拓失败 & 1 阶闭环：核心规则未闭合，仍以旧核\KouJing套用新对象域 \\
\hline
$2$ & 无法完成单阶修复塔构造与证书化 & 2 阶工具：会读叙述但不会搭建/核验修复链（$\partial D=\omega$） \\
\hline
$3$ & 跨域翻译/保结构函子不可建立 & 3 阶穿透：纯数学闭系统到开放系统落地的同构迁移失败 \\
\hline
$4$ & 第一性预设无法反转（对象域升维失败） & 4 阶反转：连续统第一性假设仍在，无法进入离散统第一性对象层 \\
\hline
$5$ & 无法完成“公理完备\texorpdfstring{$\to$}{->}工程截断\texorpdfstring{$\to$}{->}全链条闭环” & 5 阶闭环：能局部证明但无法全链条闭环与边界管理 \\
\hline
$\omega$ & 超限递归兜底层（理论完备性） & 超限兜底：用于公理级存在性保证，通常不要求工程可执行 \\
\hline
\end{tabular}
\end{center}
\end{remark}

\subsection{备注：多主体分阶修复模型（避免把经验断言写成定理）}
\label{subsec:tex29-ch17-multiagent}

\begin{remark}[分工\KouJing：团队分阶修复 vs 个体全阶修复]
本笔记的形式系统允许把“修复任务”按阶数拆分为多个子任务：
不同角色分别负责不同层级的证书化（例如低阶连通/中阶工具/高阶对象域反转/工程截断），
再通过证书链拼接得到整体可用结论。
该模型解释了一个常见现象：集体可以在不完全掌握全体系的情况下识别其战略价值，
但单一个体要完成从 $0$ 阶到高阶的全链条闭环，必须跨越更高的综合门槛。
该备注只提供结构化分工框架，不对现实机构或概率作经验断言。
\end{remark}

\clearpage

%%%%%%%%%%%%%%%%%%%%%%%%%%%%%%%%%%%%%%%%%%%%%%%%%%%%%%%%%%%%
\section{终极元结构闭环：\GZMonMotherBundle{}版广义非交换李代数与轨道商/核商二分完备性}
\label{sec:tex29-meta-structure}
\label{sec:tex29-ch18}
%%%%%%%%%%%%%%%%%%%%%%%%%%%%%%%%%%%%%%%%%%%%%%%%%%%%%%%%%%%%

\begin{remark}[核心陈述（元结构闭环；轨道商/核商二分完备性）]
本章把“\GZMonMotherBundle{} + 广义非交换李代数”固定为离散统第一性对象域中的终极元结构：
它在同一套对象与操作语法下统一给出（i）\GZMon{}的拓扑载体，（ii）修复动作的代数引擎，
（iii）静态退化与动态生成的兼容接口。
在该元结构\KouJing下，\GZMon{}之间的等价判据满足\textbf{轨道商/核商二分完备性}：它被压缩为两类且不存在第三种稳定判据：
\textbf{同构等价}（动态完整态：结构群可逆作用的同一轨道）
与 \textbf{退化等价}（静态降维态：退化投影后的同一静态核/同一退化类）。
\end{remark}

\subsection{公理级术语对齐：母体主纤维丛、结构群与\GZMon{}截面}
\label{subsec:tex29-meta-alignment}

\begin{definition}[\GZMonMotherBundle{}（元本体总空间）]
\label{def:tex29-meta-mother-bundle}
称一个四元组
\[
  \mathbb{B}:=(P,M,G,\pi)
\]
为\textbf{\GZMonMotherBundle{}}，若满足：
\begin{enumerate}
  \item $P$ 为总空间（\GZMon{}母体/元本体载体），$M$ 为底空间（观测层/表象域），$\pi:P\to M$ 为丛投影；
  \item $G$ 为结构群，作用于 $P$ 的纤维并保持投影：对任意 $p\in P,g\in G$，
  \[
    \pi(p\cdot g)=\pi(p);
  \]
  \item \GZMon{}集合（对象域中的“\GZMon{}实例”）固定为截面集
  \[
    \mathsf{Mon}:=\Gamma(P):=\{s:M\to P:\ \pi\circ s=\mathrm{id}_{M}\}.
  \]
\end{enumerate}
并把 $P$ 理解为“所有\GZMon{}的统一母体”，把 $s\in\Gamma(P)$ 理解为一个具体\GZMon{}。
\end{definition}

\begin{definition}[退化投影（静态极限接口）]
\label{def:tex29-meta-degeneration}
在定义~\ref{def:tex29-meta-mother-bundle} 下，引入一个\textbf{退化投影/观测函子}（接口）
\[
  \Pi:\Gamma(P)\to \mathsf{Cont},
\]
其语义为“遗忘动态修复数据，仅保留静态描述层”。
本章仅要求其满足最小兼容\KouJing：存在某个静态核 $A\in\mathsf{Cont}$ 使得
\[
  (\forall s\in\Gamma(P))\ \Pi(s)=A,
\]
即退化投影把全部\GZMon{}压缩到同一静态核（旧体系回收\KouJing的抽象版本）。
\end{definition}

\begin{remark}[外部参照（仅作参照，不作为本章证明前提）]
主纤维丛的标准教材\KouJing可参见 \cite{Husemoller1994FibreBundles}；
交换子括号与李代数的常见教材\KouJing可参见 \cite{Hall2015LieGroups}；
与本仓库语境相关的“主纤维丛版广义非交换李代数（G-Framework/G-Algebra）”叙述可参见 \cite{GaoZheng2025GFramework}。
本章的定理与证书仍以“接口+谓词链”的工作\KouJing自洽，不把外部教材结论当作前提直接调用。
\end{remark}

\begin{remark}[对齐表：四个核心单元的角色参照点（\KouJing不变）]
为避免“类比化”导致的歧义，把本章固定的四个核心单元压缩为对齐表：
\begin{center}
\begingroup
\setlength{\tabcolsep}{4pt}
\small
\renewcommand{\arraystretch}{1.25}
\begin{tabular}{|p{0.21\linewidth}|p{0.46\linewidth}|p{0.23\linewidth}|}
\hline
\textbf{核心单元} & \textbf{公理级内涵（本章\KouJing）} & \textbf{定位} \\
\hline
\GZMonMotherBundle{} & 总空间 $P$；\GZMon{}为截面 $s\in\Gamma(P)$；$\pi:P\to M$ 给出表象投影 & 元本体载体 \\
\hline
广义非交换李代数 & 由“边缘算子生成元”诱导的非交换操作代数（括号为交换子） & 修复动作引擎 \\
\hline
同构等价 & 结构群 $G$ 在截面上的可逆作用同一轨道：$\exists g\in G,\ t=g\cdot s$ & 动态完整态 \\
\hline
退化等价 & 退化投影后同一静态核：$\Pi(s)=\Pi(t)=A$ & 静态降维态 \\
\hline
\end{tabular}
\endgroup
\end{center}
\end{remark}

\subsection{定义：广义非交换李代数（边缘算子生成元与交换子括号）}
\label{subsec:tex29-meta-galgebra}

\begin{definition}[广义非交换李代数（修复动作元引擎）]
\label{def:tex29-meta-galgebra}
在定义~\ref{def:tex29-meta-mother-bundle} 下，取 $\mathrm{End}(\Gamma(P))$ 为截面集上的端态算子集合。
称一对
\[
  \mathfrak{g}_{\partial}:=(\mathcal{O}_{\partial},[\cdot,\cdot])
\]
为\textbf{广义非交换李代数}，若存在一族“主动修复型边缘算子生成元”
\[
  \{\partial_i\}_{i\in I}\subseteq \mathcal{O}_{\partial}\subseteq \mathrm{End}(\Gamma(P)),
\]
并且：
\begin{enumerate}
  \item （交换子括号）对任意 $A,B\in\mathcal{O}_{\partial}$，定义
  \[
    [A,B]:=A\circ B - B\circ A;
  \]
  \item （闭包）$\mathcal{O}_{\partial}$ 对复合 $\circ$ 与括号 $[\cdot,\cdot]$ 封闭；
  \item （非交换性见证）存在 $A,B\in\mathcal{O}_{\partial}$ 使 $[A,B]\neq 0$。
\end{enumerate}
其中 $\partial_i$ 的语义仅固定为“可触发修复/同伦扭曲/延拓更新的无穷小生成元”，不强行绑定分量公式。
\end{definition}

\begin{remark}[非交换性的系统语义（顺序敏感）]
交换子 $[A,B]\neq 0$ 的存在意味着：
修复动作的执行顺序会改变中间塔结构，从而与“开放系统/非对称约束/动态演化”的顺序敏感性相容。
本章不把该语义当作经验断言，只把它登记为操作代数的结构特征。
\end{remark}

\subsection{公理（降级为定理）：元结构闭包与可执行操作的闭环（数学归纳法证书）}
\label{subsec:tex29-meta-closure}

\begin{theorem}[公理（降级为定理）：元结构闭包（有限字长操作闭包与截面保持）（数学归纳法证书）]
\label{thm:tex29-meta-closure}
在定义~\ref{def:tex29-meta-galgebra} 下，令 $\mathcal{W}_n$ 为由索引集 $I$ 生成的“字长为 $n$ 的操作词”集合：
若 $w=(i_1,\dots,i_n)\in \mathcal{W}_n$，定义其对应操作为
\[
  T_w:=\partial_{i_1}\circ\cdots\circ \partial_{i_n}\ \in\ \mathrm{End}(\Gamma(P)).
\]
并约定 $n=0$ 时 $T_{\emptyset}:=\mathrm{id}_{\Gamma(P)}$。
则对任意 $n\in\NN_0$，对任意 $w\in\mathcal{W}_n$，$T_w$ 为良定义的截面端态算子，且保持对象域闭包：
\[
  (\forall s\in\Gamma(P))\quad T_w(s)\in\Gamma(P).
\]
因此“可执行修复/变换操作”在本章\KouJing下形成一个有限字长闭包的可复核操作系统。
\end{theorem}

\begin{Certificate}[数学归纳法证书：按字长 $n$ 归纳的“操作闭包 + 截面保持”]
\label{cert:tex29-meta-closure}
令谓词 $P(n)$ 表示：
“对任意 $w\in\mathcal{W}_n$，操作 $T_w$ 良定义且对任意 $s\in\Gamma(P)$ 有 $T_w(s)\in\Gamma(P)$。”
证书化要求为：
\[
  P(0)\ \text{成立},\qquad (\forall n\in\NN)\ (P(n)\Rightarrow P(n+1)).
\]
\end{Certificate}

\paragraph{证明（数学归纳法证书；谓词因果链）.}
由证书~\ref{cert:tex29-meta-closure}：
\begin{itemize}
  \item 基步：$n=0$ 时 $T_{\emptyset}=\mathrm{id}_{\Gamma(P)}$。
  对任意截面 $s\in\Gamma(P)$ 有 $T_{\emptyset}(s)=s\in\Gamma(P)$，故 $T_{\emptyset}$ 保持 $\Gamma(P)$，从而 $P(0)$ 成立；
  \item 归纳步：假设 $P(n)$ 成立。任取 $w'=(i_1,\dots,i_n,i_{n+1})\in\mathcal{W}_{n+1}$，
  则 $T_{w'}=T_w\circ \partial_{i_{n+1}}$（其中 $w=(i_1,\dots,i_n)$）。
  由定义~\ref{def:tex29-meta-galgebra} 的闭包性，$T_w$ 与 $\partial_{i_{n+1}}$ 均为 $\Gamma(P)$ 上端态算子且复合良定义；
  再由归纳假设 $P(n)$，$T_w$ 保持 $\Gamma(P)$，故 $T_{w'}$ 亦保持 $\Gamma(P)$，从而 $P(n+1)$ 成立。
\end{itemize}
由数学归纳法，$P(n)$ 对任意 $n$ 成立，定理得证。证毕。

\subsection{引理：截面二分（正则/退化）与无第三态}
\label{subsec:tex29-meta-dichotomy}

\begin{definition}[正则截面与退化截面（动态/静态二分谓词）]
\label{def:tex29-meta-regular-degenerate}
引入一个接口赋值
\[
  \mathcal{J}:\Gamma(P)\to \mathcal{O}_{\partial},\qquad s\mapsto \partial_s,
\]
把“\GZMon{}可调用的边缘算子生成元”记为 $\partial_s$。
定义二分谓词：
\[
  \mathsf{Reg}(s)\ :\Longleftrightarrow\ \partial_s\neq 0,
  \qquad
  \mathsf{Deg}(s)\ :\Longleftrightarrow\ \partial_s=0.
\]
分别称 $s$ 为正则（动态）截面与退化（静态）截面。
\end{definition}

\begin{lemma}[截面二分：$\mathsf{Reg}\ \vee\ \mathsf{Deg}$（无第三态）]
\label{lem:tex29-meta-dichotomy}
对任意\GZMon{} $s\in\Gamma(P)$，必有
\[
  \mathsf{Reg}(s)\ \vee\ \mathsf{Deg}(s).
\]
并且不存在“既非正则亦非退化”的第三态谓词：
\[
  \neg\bigl(\neg\mathsf{Reg}(s)\wedge \neg\mathsf{Deg}(s)\bigr).
\]
\end{lemma}

\begin{Certificate}[证书：由 $\partial_s=0$ 的二值性直接推出]
\label{cert:tex29-meta-dichotomy}
证书登记谓词链：
\[
\begin{aligned}
  D_1 &: \mathsf{Reg}(s)\Longleftrightarrow (\partial_s\neq 0),\\
  D_2 &: \mathsf{Deg}(s)\Longleftrightarrow (\partial_s=0),\\
  D_3 &: (\partial_s=0)\ \vee\ (\partial_s\neq 0),\\
  D_4 &: D_1\wedge D_2\wedge D_3\Rightarrow \mathsf{Reg}(s)\vee\mathsf{Deg}(s)\Rightarrow \neg(\neg\mathsf{Reg}(s)
    \wedge\neg\mathsf{Deg}(s)).
\end{aligned}
\]
\end{Certificate}

\paragraph{证明（谓词因果链）.}
由 $D_3$ 得 $\partial_s$ 的二值划分；结合 $D_1,D_2$ 即得 $\mathsf{Reg}(s)\vee\mathsf{Deg}(s)$；
从而第三态谓词被排除。证毕。

\subsection{定理：轨道商/核商二分完备性——同构等价与退化等价的完备判据}
\label{subsec:tex29-meta-two-equivalences}

\begin{definition}[同构等价（轨道等价）与退化等价（核等价）]
\label{def:tex29-meta-equivalences}
在定义~\ref{def:tex29-meta-mother-bundle} 下，定义两类等价关系（均为标准的“商/核”构造）：
\begin{enumerate}
  \item \textbf{同构等价}（结构群作用的轨道等价；动态完整态）：
  \[
    s\simeq_{\mathrm{iso}} t\ :\Longleftrightarrow\ \exists g\in G\ \ \text{使得}\ \ t=g\cdot s,
  \]
  其中 $g\cdot s$ 表示结构群对截面的自然右作用（逐点作用）：
  \[
    (g\cdot s)(m):=s(m)\cdot g\qquad(m\in M).
  \]
  \item \textbf{退化等价}（退化投影 $\Pi$ 的核等价；静态降维态）：
  \[
    s\simeq_{\mathrm{deg}} t\ :\Longleftrightarrow\ \Pi(s)=\Pi(t),
  \]
  其中 $\Pi$ 为定义~\ref{def:tex29-meta-degeneration} 的退化投影接口。
\end{enumerate}
\end{definition}

\begin{remark}[等价关系视角：轨道商与核商]
等价关系可视为“把对象压缩为可复核等价类”的\textbf{商结构}。
上面的两类等价关系分别对应两种标准商：
\[
  \Gamma(P)\big/\simeq_{\mathrm{iso}}\ \cong\ \Gamma(P)/G\quad(\text{轨道商}),
  \qquad
  \Gamma(P)\big/\simeq_{\mathrm{deg}}\ \cong\ \mathrm{Im}(\Pi)\quad(\text{核商}).
\]
本章将以下性质称为\textbf{轨道商/核商二分完备性}：
在只允许调用“结构群作用/退化投影/生成元赋值”的元结构接口时，
任意可证书化的等价判定都只能归约为\textbf{轨道商}或\textbf{核商}这两种标准商机制；
换言之，不存在独立于二者的第三种等价机制。
\end{remark}

\begin{definition}[元结构相容等价关系（可证书化的等价判定口径）]
\label{def:tex29-meta-compatible-equivalence}
称 $\approx$ 为 $\Gamma(P)$ 上的一个\textbf{元结构相容}等价关系，若满足：
\begin{enumerate}
  \item （分层保持）$s\approx t\Rightarrow \bigl(\mathsf{Reg}(s)\Leftrightarrow \mathsf{Reg}(t)\bigr)$；
  \item （结构群不变）$s\approx t\Rightarrow (g\cdot s)\approx(g\cdot t)\quad(\forall g\in G)$；
  \item （退化不变）$s\approx t\Rightarrow \Pi(s)=\Pi(t)$。
\end{enumerate}
该定义把“允许使用的判定信息”固定在三类接口之内：正则/退化二分谓词、结构群作用、与退化投影。
\end{definition}

\begin{theorem}[轨道商/核商二分完备性（严格表述）：元结构相容等价关系至多给出两类商（证书化\KouJing）]
\label{thm:tex29-meta-two-equivalences}
在定义~\ref{def:tex29-meta-regular-degenerate}、定义~\ref{def:tex29-meta-equivalences} 与定义~\ref{def:tex29-meta-compatible-equivalence} 的\KouJing下，
登记以下接口前提：
\begin{enumerate}
  \item （正则层轨道传递）对任意 $s,t\in\Gamma(P)$，若 $\mathsf{Reg}(s)\wedge\mathsf{Reg}(t)$，则存在 $g\in G$ 使 $t=g\cdot s$；
  \item （退化像为单点）$\Pi$ 把任意\GZMon{}压缩到同一静态核 $A$（定义~\ref{def:tex29-meta-degeneration}），即 $\mathrm{Im}(\Pi)=\{A\}$。
\end{enumerate}
则对任意元结构相容等价关系 $\approx$，其商集 $\Gamma(P)/\approx$ 至多包含两类等价类，
并且这两类必分别对应于正则层与退化层：
\begin{enumerate}
  \item （正则层）若 $\mathsf{Reg}(s)\wedge\mathsf{Reg}(t)$，则 $s\simeq_{\mathrm{iso}} t$；
  \item （退化层）若 $\mathsf{Deg}(s)\wedge\mathsf{Deg}(t)$，则 $s\simeq_{\mathrm{deg}} t$（且由 $\mathrm{Im}(\Pi)=\{A\}$
    可知退化层为单一退化类）；
  \item （无第三类）不存在第三个与上述两层均不相交、且能在元结构接口内被证书化维持的等价类；
  换言之，任意元结构相容等价关系的分类信息至多是二值的（正则/退化）。
\end{enumerate}
\end{theorem}

\begin{Certificate}[证书：分层保持 + 正则层轨道传递 + 退化像单点]
\label{cert:tex29-meta-two-equivalences}
证书登记谓词链：
\[
\begin{aligned}
  E_1 &: \forall s\in\Gamma(P),\ \mathsf{Reg}(s)\vee\mathsf{Deg}(s)\quad(\text{引理~\ref{lem:tex29-meta-dichotomy}}),\\
  E_2 &: \mathsf{Reg}(s)\wedge\mathsf{Reg}(t)\Rightarrow \exists g\in G\ (t=g\cdot s)\Rightarrow s\simeq_{\mathrm{iso}}
    t,\\
  E_3 &: \mathrm{Im}(\Pi)=\{A\}\Rightarrow \bigl(\mathsf{Deg}(s)\wedge\mathsf{Deg}(t)\bigr)\Rightarrow \Pi(s)=\Pi(t)
    \Rightarrow s\simeq_{\mathrm{deg}} t,\\
  E_4 &: s\approx t\Rightarrow \bigl(\mathsf{Reg}(s)\Leftrightarrow \mathsf{Reg}(t)\bigr)\Rightarrow \Gamma(P)/\approx\
    \text{至多两类},\\
  E_5 &: E_1\wedge E_2\wedge E_3\wedge E_4\Rightarrow \text{分类信息至多二值（正则/退化），无第三类}.
\end{aligned}
\]
\end{Certificate}

\paragraph{证明（谓词因果链）.}
任取元结构相容等价关系 $\approx$。
由定义~\ref{def:tex29-meta-compatible-equivalence} 的“分层保持”，任意等价类只能落在正则层或退化层之一，
因此商集 $\Gamma(P)/\approx$ 的等价类数至多为两类（对应 $\mathsf{Reg}$ 与 $\mathsf{Deg}$ 两层），这给出 $E_4$。

对正则层，若 $\mathsf{Reg}(s)\wedge\mathsf{Reg}(t)$，由接口前提“正则层轨道传递”得到 $\exists g\in G$ 使 $t=g\cdot s$，
从而由定义~\ref{def:tex29-meta-equivalences} 得 $s\simeq_{\mathrm{iso}}t$，即 $E_2$。

对退化层，由 $\mathrm{Im}(\Pi)=\{A\}$ 得任意退化截面均满足 $\Pi(\cdot)=A$；
因此若 $\mathsf{Deg}(s)\wedge\mathsf{Deg}(t)$，有 $\Pi(s)=\Pi(t)$，从而 $s\simeq_{\mathrm{deg}}t$，即 $E_3$。

合并 $E_1,E_2,E_3,E_4$ 得 $E_5$：在元结构接口允许的判定信息下，商分类至多二值（正则/退化），不存在第三类。
定理得证。证毕。

\subsection{命题：元结构的内生派生——修复塔、边缘算子与静态/动态兼容统一为同一母体}
\label{subsec:tex29-meta-derivation}

\begin{proposition}[内生派生：同伦修复塔、边缘算子与变分选择可视为元结构接口的不同切片]
\label{prop:tex29-meta-derivation}
在元结构 $\mathbb{B}=(P,M,G,\pi)$ 与代数引擎 $\mathfrak{g}_{\partial}$（定义~\ref{def:tex29-meta-galgebra}）\KouJing下，
本笔记此前出现的核心构件可在形式系统中被统一解释为该元结构的派生接口：
\begin{enumerate}
  \item （边缘算子）修复谓词 $\mathsf{Rep}(\omega)\Longleftrightarrow \exists D\ (\partial D=\omega)$（第\ref{subsec:tex29-domain-predicates}小节）可视为生成元 $\partial$ 的构造性接口；
  \item （同伦修复塔）“定位--修复--闭合校验”的分层修复链（如第\ref{subsec:tex29-ch16-variational-tower}小节）可视为结构群作用轨道上的逐层截断与闭包校验；
  \item （静态/动态兼容）退化投影 $\Pi$（定义~\ref{def:tex29-meta-degeneration}）把动态对象回收为静态核，从而与第\ref{sec:tex29-ch8}章的“保核收缩兼容”\KouJing同构；

  \item （广义分形）在 $\mathsf{Rep}$ 之上叠加“层间同伦等价递归约束”即可得到广义分形特例（第\ref{subsec:tex29-fractal-boundary}小节）；
  \item （变分法）若在轨道上引入缺陷势/能量泛函（第\ref{sec:tex29-ch16}章的 $\Phi/\mathcal{E}$ 对齐\KouJing），则“选择最优修复路径”对应在可达轨道内选取极小值代表元。
\end{enumerate}
因此该元结构可作为统一母体，把“修复塔、边缘算子、变分选择、静态/动态兼容”收束到同一套可证书化操作语法之内。
\end{proposition}

\begin{Certificate}[证书：把分支规则登记为同一母体下的接口映射]
\label{cert:tex29-meta-derivation}
证书登记谓词链：
\[
\begin{aligned}
  M_1 &: \mathsf{Rep}(\omega)\Longleftrightarrow \exists D\ (\partial D=\omega)\quad(\text{修复入口}),\\
  M_2 &: \text{逐层修复 + 闭合校验}\Rightarrow \text{可写成操作词 }T_w\ \text{的逐步应用}\quad(\text{定理~\ref{thm:tex29-meta-closure}}),\\
  M_3 &: \Pi\ \text{遗忘动态修复数据}\Rightarrow \text{回收到静态核}\quad(\text{定义~\ref{def:tex29-meta-degeneration}}),\\
  M_4 &: \mathsf{Fr}\ \text{为 }\mathsf{Rep}\ \text{之上的递归自相似约束}\quad(\text{定理~\ref{thm:tex29-fractal-strict-subset}}),\\
  M_5 &: \text{变分极值/单调性}\Rightarrow \text{在轨道内选取极小代表}\quad(\text{第\ref{sec:tex29-ch16}章\KouJing}),\\
  M_6 &: M_1\wedge M_2\wedge M_3\wedge M_4\wedge M_5\Rightarrow \text{各分支规则可被同一母体接口化统一表述}.
\end{aligned}
\]
\end{Certificate}

\paragraph{证明（谓词因果链）.}
由 $M_1$ 固定修复入口；由 $M_2$ 把逐层修复写成有限字长操作词的可复核应用；
由 $M_3$ 给出静态极限接口；由 $M_4$ 给出分形约束作为附加结构；
由 $M_5$ 给出变分选择作为轨道内的代表元选择规则；合并得 $M_6$。证毕。

\subsection{推论：静态退化类丢失非交换修复能力（“安全锁”\KouJing；仅作结构结论）}
\label{subsec:tex29-meta-safety}

\begin{corollary}[退化等价的冻结性：退化投影杀死交换子括号（结构结论）]
\label{cor:tex29-meta-safety}
在退化投影 $\Pi$（定义~\ref{def:tex29-meta-degeneration}）的\KouJing下，
若把 $\Pi$ 视为对操作代数的遗忘映射（接口）并登记其满足
\[
  \Pi([A,B])=0\qquad(\forall A,B\in\mathcal{O}_{\partial}),
\]
则在退化等价类中，非交换修复信息被系统性丢弃：交换子括号退化为零。
因此退化等价类只能承载静态描述层（可被旧体系吸收/复述），而不能承载顺序敏感的动态修复能力；
动态修复能力只能在正则截面的同构等价轨道上以非交换操作语法被完整表达。
\end{corollary}

\begin{Certificate}[证书：退化映射杀死交换子 $\Rightarrow$ 非交换信息不可在退化类中保持]
\label{cert:tex29-meta-safety}
证书登记谓词链：
\[
\begin{aligned}
  S_1 &: \Pi([A,B])=0\quad(\forall A,B\in\mathcal{O}_{\partial}),\\
  S_2 &: S_1\Rightarrow \text{在退化像中，括号结构恒为零（非交换信息丢失）},\\
  S_3 &: S_2\Rightarrow \text{退化类仅保留静态可复述结构，无法区分顺序敏感的动态修复语法}.
\end{aligned}
\]
\end{Certificate}

\paragraph{证明（谓词因果链）.}
由 $S_1$ 得退化像中所有交换子归零；因此由 $S_2$ 得非交换结构被杀死；
从而由 $S_3$ 得退化类只能保留静态描述层。证毕。

\begin{remark}[备注：本章的“安全锁”仅为结构结论]
推论~\ref{cor:tex29-meta-safety} 只描述一个形式系统内的结构事实：退化投影会丢弃非交换修复信息。
它不对现实世界中的扩散、权限或个体能力做经验断言；这些均应作为外部系统的治理问题另行建模。
\end{remark}

\begin{remark}[兼容性边界：连续统与静态闭系统仅作为退化极限出现]
在本笔记的对象域\KouJing下：
\begin{itemize}
  \item \textbf{连续统/平坦背景不是第一性}：它仅作为退化投影 $\Pi$ 的像域（静态核 $A$）出现；
  \item \textbf{静态闭系统不是原生对象}：它仅作为“杀死交换子括号”的退化极限出现（即顺序敏感修复语法被冻结/丢弃）。
\end{itemize}
因此，“离散统第一性 + 动态开放系统”为对象层，“连续统 + 静态闭系统”为退化投影层；
两者并非对立本体，而是同一母体在不同退化层级下的表现形式。
该备注只固定形式系统的对象层级关系，不把外部数学基础争论（如 CH 等）直接纳入本章定理。
\end{remark}

\begin{remark}[终极闭环定论（本笔记\KouJing）]
在本章固定的元结构\KouJing下：
\GZMonMotherBundle{}给出统一拓扑载体，广义非交换李代数给出统一修复动作语法；
同构等价对应动态完整态的可逆保结构映射，退化等价对应静态极限的回收投影；
并由定理~\ref{thm:tex29-meta-two-equivalences} 得到轨道商/核商二分完备性的刚性结论。
因此，本章提供的是一个把“边缘算子是修复核心、离散统第一性、障碍是构造起点”的底层规则
压缩为统一母体的形式化闭环：后续任何分支推演若要保持与本体系一致，必须可被还原为这两类判据与其证书链的组合。
\end{remark}

\clearpage

%%%%%%%%%%%%%%%%%%%%%%%%%%%%%%%%%%%%%%%%%%%%%%%%%%%%%%%%%%%%
\section{形式逻辑条目化总览：增益与跃迁（定义—公理降级—定理—引理—命题—推论—备注）}
\label{sec:tex29-formal-logic-gain-jump}
\label{sec:tex29-ch19}
%%%%%%%%%%%%%%%%%%%%%%%%%%%%%%%%%%%%%%%%%%%%%%%%%%%%%%%%%%%%

\begin{remark}[本章定位（条目化复述；便于审计与复核）]
本章把本笔记中用到的核心对象、谓词与关键结论按“定义—（公理降级）定理—引理—命题—推论—备注”的形式逻辑模板汇总成一章，
并在每个条目后附上“证书 + 证明（谓词因果链）”的最小可复核版本。
更详细的背景解释与推演细节分散在前文各章：本章只做\emph{条目化可审计摘要}。
\end{remark}

\subsection{定义}
\label{subsec:tex29-formal-definitions}

\begin{definition}[定义 D0（体系增益/跃迁的比较对象）]
\label{def:tex29-formal-system-gain-jump}
记您原有体系的形式系统为 $\mathcal{T}_0$，
将本文档中新增的对象域、谓词、分类函数与证书链并入后得到扩展系统 $\mathcal{T}_1$。

所谓\textbf{增益}，是 $\mathcal{T}_1$ 相对 $\mathcal{T}_0$ 在如下三项上新增的可证书化结构：
\begin{enumerate}
  \item \textbf{可判定边界}（把“是否属于全集范畴”写成谓词真值入口）；
  \item \textbf{可证书化闭包}（把闭包/递归/链合成写成可检查的证书链）；
  \item \textbf{可构造性}（把“存在性断言”落实为 $\partial D=\omega$ 的构造见证）。
\end{enumerate}

所谓\textbf{跃迁}，是 $\mathcal{T}_1$ 引入的某些谓词使得“障碍”从\emph{终止判据}转写为\emph{构造单元}，
并可在全阶上导出闭环性质（例如无空洞完成态 $\mathsf{NoHole}$；见推论~\ref{cor:tex29-formal-no-hole}）。
\end{definition}

\begin{definition}[定义 D1（修复谓词与分形谓词：本文档的核心入口）]
\label{def:tex29-formal-rep-fr-entry}
对任意障碍代表元 $\omega\in\Omega^n(X;G)$，定义修复谓词
\[
  \mathsf{Rep}(\omega)\ :\Longleftrightarrow\ \exists D\in\mathcal{D}^{n+1}(X;G)\ \bigl(\partial D=\omega\bigr).
\]
进一步定义广义分形谓词
\[
  \mathsf{Fr}(\omega)\ :\Longleftrightarrow\ \exists\Bigl(\mathcal{C},\{(X_k,\omega_k,D_k,h_k,k_k)\}_{k\in\NN}\Bigr)\
  \text{使 (F0)--(F4) 逐项成立，且 }X_{k+1}\simeq X_k.
\]
其中 (F0)--(F4) 与定义~\ref{def:tex29-rep-fr} 一致：它们把“无限递归 + 层间同伦等价 + 逐层修复证据链”固定为\emph{定义性约束}。
\end{definition}

\begin{remark}[要点：$\mathsf{Rep}$ 与 $\mathsf{Fr}$ 的边界]
$\mathsf{Rep}$ 只问“能否被边缘算子平凡化”（即是否存在 $D$ 使 $\partial D=\omega$）；
$\mathsf{Fr}$ 在此基础上额外强制“无限递归 + 层间同伦等价”这一\emph{定义性约束}。
因此 $\mathsf{Fr}$ 只能刻画修复全集中的特定子类，而不能反向定义修复全集本身。
\end{remark}

\subsection{公理（降级为定理：数学归纳法证书 + 证明）}
\label{subsec:tex29-formal-axiom-downgraded}

\begin{theorem}[定理 A1（公理降级：修复闭包与“全集范畴”判定）]
\label{thm:tex29-formal-axiom-closure}
沿用定义~\ref{def:tex29-composed-evidence} 的合成证据链与序列边缘算子记号。
若有限链 $\omega_0\leadsto\cdots\leadsto\omega_N$ 满足 $(\forall k\le N)\ \mathsf{Rep}(\omega_k)$，
则存在合成证据链 $\widehat{D}=(D_0,\dots,D_N)$ 使
\[
  \widehat{\partial}_N(\widehat{D})=(\omega_0,\dots,\omega_N).
\]
并且“属于修复同伦类全集”当且仅当 $\mathsf{Rep}(\omega)$ 为真：
\[
  \omega\ \text{属于修复同伦类全集}\ \Longleftrightarrow\ \mathsf{Rep}(\omega).
\]
\end{theorem}

\begin{Certificate}[证书（归纳谓词）：按链长 $N$ 的闭包证书]
令 $P(N)$ 表示：
“对任意长度为 $N$ 的链 $\omega_0\leadsto\cdots\leadsto\omega_N$，若 $(\forall k\le N)\ \mathsf{Rep}(\omega_k)$，
则存在合成证据链 $\widehat{D}=(D_0,\dots,D_N)$ 使 $\widehat{\partial}_N(\widehat{D})=(\omega_0,\dots,\omega_N)$，
且链上任意节点的全集归属仅由 $\mathsf{Rep}$ 判定。”
证书化要求为 $P(0)$ 成立且 $(\forall N)\bigl(P(N)\Rightarrow P(N+1)\bigr)$。
\end{Certificate}

\paragraph{证明（数学归纳法证书；谓词因果链）.}
基步 $N=0$：由 $\mathsf{Rep}(\omega_0)$ 取 $D_0$ 使 $\partial D_0=\omega_0$，则 $\widehat{\partial}_0(D_0)=(\omega_0)$。
归纳步：由 $P(N)$ 得到 $(D_0,\dots,D_N)$，再由 $\mathsf{Rep}(\omega_{N+1})$ 取 $D_{N+1}$ 并拼接为 $\widehat{D}$，
由 $\widehat{\partial}_{N+1}$ 的分量定义即得结论。证毕.

\begin{theorem}[定理 A2（公理降级：广义分形 $\simeq \omega$-递归同伦修复塔）]
\label{thm:tex29-formal-fractal-rec-tower}
\[
  \mathsf{Fr}(\omega)\quad\Longleftrightarrow\quad
  \exists\ \mathbb{T}^{\mathrm{rec}}_{\omega}
  =\Bigl(\mathcal{C},\{(X_k,\omega_k,D_k,h_k,k_k)\}_{k\in\NN}\Bigr)\ \text{满足 (F0)--(F4)}.
\]
并且仅给出任意有限截断 $\{0,\dots,N\}$ 不能推出 $\mathsf{Fr}(\omega)$（因为 (F1)--(F4) 含 $\forall k\in\NN$ 的全称约束）。
\end{theorem}

\begin{Certificate}[证书要点：定义展开与“不可有限截断”]
证书登记两条谓词链：
\[
\begin{aligned}
  R_1 &: \mathsf{Fr}(\omega)\ \text{按定义即等价于“存在满足 (F0)--(F4) 的无限递归塔”},\\
  R_2 &: \text{有限截断只提供 }(\forall k\le N)\text{ 的信息，无法推出 }(\forall k\in\NN)\text{ 的定义性约束}.
\end{aligned}
\]
\end{Certificate}

\paragraph{证明（谓词因果链）.}
第一条等价由定义~\ref{def:tex29-rep-fr} 直接展开得到（即 $R_1$）。
第二条由量词结构给出：$\mathsf{Fr}$ 的定义包含对所有 $k\in\NN$ 的约束，
因此任何仅包含有限多层数据的证书都不足以推出其成立（即 $R_2$）。证毕.

\begin{theorem}[定理 A3（公理降级：全阶障碍可修复 $\Rightarrow$ 障碍群平凡）]
\label{thm:tex29-formal-all-degree-trivial}
设 $\partial^2=0$，并沿用第\ref{subsec:tex29-core-meta-terminology}小节的记号 $Z^n,B^n,H^n_{\mathrm{obs}}$。
若对任意 $n\ge 1$ 与任意闭链障碍 $\omega\in Z^n(X;G)$ 都有 $\mathsf{Rep}(\omega)$（等价于 $\omega\in B^n(X;G)$），
则
\[
  H^n_{\mathrm{obs}}(X;G)=Z^n(X;G)/B^n(X;G)=0\qquad(\forall n\ge 1).
\]
\end{theorem}

\begin{Certificate}[证书（按阶数归纳）：闭链皆为边界 $\Rightarrow$ 商群为零]
令 $P(n)$ 表示：“对所有 $1\le k\le n$，$Z^k(X;G)=B^k(X;G)$。”
证书链为：
\[
  P(1)\ \text{成立},\qquad (\forall n\ge 1)\ \bigl(P(n)\Rightarrow P(n+1)\bigr).
\]
\end{Certificate}

\paragraph{证明（谓词因果链；符号推演）.}
由前提得 $Z^{n}\subseteq B^{n}$（闭链皆可修复）；又由 $\partial^2=0$ 得 $B^{n}\subseteq Z^{n}$（边界必为闭链）。
故 $Z^{n}=B^{n}$，从而 $H^n_{\mathrm{obs}}=Z^n/B^n=0$。证毕.

\subsection{定理（证书 + 证明）}
\label{subsec:tex29-formal-theorems}

\begin{theorem}[定理 T1（刚性边界：$\mathsf{Fr}\subsetneq \mathsf{Rep}$）]
\label{thm:tex29-formal-fr-strict}
\begin{enumerate}
  \item 若 $\mathsf{Fr}(\omega)$，则 $\mathsf{Rep}(\omega)$；
  \item 存在 $\omega_\star$ 使 $\mathsf{Rep}(\omega_\star)\wedge \neg\mathsf{Fr}(\omega_\star)$。
\end{enumerate}
因此 $\mathsf{Fr}$ 为 $\mathsf{Rep}$ 的严格子集。
\end{theorem}

\begin{Certificate}[证书链：从定义到包含与反例]
证书登记谓词链：
\[
\begin{aligned}
  S_1 &: \mathsf{Fr}(\omega)\Rightarrow (\forall k)\exists D_k\ (\partial D_k=\omega_k),\\
  S_2 &: S_1\Rightarrow \exists D_0\ (\partial D_0=\omega)\Rightarrow \mathsf{Rep}(\omega),\\
  S_3 &: \exists\omega_\star\ (\mathsf{Rep}(\omega_\star)\wedge \neg\mathsf{Fr}(\omega_\star)).
\end{aligned}
\]
\end{Certificate}

\paragraph{证明（谓词因果链）.}
由 $S_1$ 取 $k=0$ 得 $\exists D_0(\partial D_0=\omega)$，即由 $S_2$ 得 $\mathsf{Rep}(\omega)$。
另一方面，取“局域奇点一次性修复”型障碍 $\omega_\star$（例如 $\tau(\omega_\star)=1$ 的情形），
其可满足 $\mathsf{Rep}$ 但不具备无限递归与层间同伦等价约束，故满足 $S_3$。
合并得严格包含。证毕.

\subsection{引理（证书 + 证明）}
\label{subsec:tex29-formal-lemmas}

\begin{lemma}[引理 L1（分形约束 $\Rightarrow$ 可修复性）]
\label{lem:tex29-formal-fr-implies-rep}
若 $\mathsf{Fr}(\omega)$，则 $\mathsf{Rep}(\omega)$。
\end{lemma}

\begin{Certificate}[证书：取 $k=0$ 层]
\[
  \mathsf{Fr}(\omega)\Rightarrow \exists D_0\ (\partial D_0=\omega)\Rightarrow \mathsf{Rep}(\omega).
\]
\end{Certificate}

\paragraph{证明（谓词因果链）.}
由 $\mathsf{Fr}(\omega)$ 的定义性约束，逐层存在 $D_k$ 使 $\partial D_k=\omega_k$，取 $k=0$ 即得。证毕.

\subsection{命题（证书 + 证明）}
\label{subsec:tex29-formal-propositions}

\begin{proposition}[命题 P1（六类修复形式矩阵：分形只是其中一类）]
\label{prop:tex29-formal-repair-form-matrix}
引入修复形式标签
\[
  \kappa(\omega)\in\{\mathsf{FR},\mathsf{SP},\mathsf{GH},\mathsf{CD},\mathsf{TR},\mathsf{SD}\},
\]
其中 $\mathsf{FR}$ 为分形递归式修复，其余对应单点、全局同胚、跨域同构、超限递归、静态退化等修复形式。
若 $\mathsf{Rep}(\omega)\Rightarrow \kappa(\omega)$ 可确定，则所有可修复障碍被上述六类覆盖；
并且 $\kappa(\omega)=\mathsf{FR}$ 蕴含 $\mathsf{Fr}(\omega)$，从而分形修复严格从属于修复全集。
\end{proposition}

\begin{Certificate}[证书链：覆盖性 + 从属性]
\[
\begin{aligned}
  M_1 &: \mathsf{Rep}(\omega)\Rightarrow \kappa(\omega)\in\{\mathsf{FR},\mathsf{SP},\mathsf{GH},\mathsf{CD},\mathsf{TR},
    \mathsf{SD}\},\\
  M_2 &: \kappa(\omega)=\mathsf{FR}\Rightarrow \mathsf{Rep}(\omega)\wedge \mathsf{Fr}(\omega),\\
  M_3 &: \text{由定理~\ref{thm:tex29-formal-fr-strict} 得 }\mathsf{Fr}\subsetneq \mathsf{Rep}.
\end{aligned}
\]
\end{Certificate}

\paragraph{证明（谓词因果链）.}
由 $M_1$ 得覆盖性；由 $M_2$ 得 $\mathsf{FR}$ 必为分形递归修复；再由 $M_3$ 得分形严格从属于修复全集。证毕.

\subsection{推论（证书 + 证明）}
\label{subsec:tex29-formal-corollaries}

\begin{corollary}[推论 C1（无空洞完成态）]
\label{cor:tex29-formal-no-hole}
定义
\[
  \mathsf{NoHole}(X)\ :\Longleftrightarrow\ (\forall n\ge 1)\ H^n_{\mathrm{obs}}(X;G)=0.
\]
若定理~\ref{thm:tex29-formal-all-degree-trivial} 的“全阶系统性平凡化”前提成立，则 $\mathsf{NoHole}(X)$ 成立。
\end{corollary}

\begin{Certificate}[证书：由 $H^n_{\mathrm{obs}}=0$ 得 $\mathsf{NoHole}$]
\[
  (\forall n\ge 1)\ H^n_{\mathrm{obs}}(X;G)=0\ \Rightarrow\ \mathsf{NoHole}(X).
\]
\end{Certificate}

\paragraph{证明（谓词因果链）.}
由定理~\ref{thm:tex29-formal-all-degree-trivial} 得 $(\forall n\ge 1)\ H^n_{\mathrm{obs}}(X;G)=0$，代入定义即得。证毕.

\begin{corollary}[推论 C2（雅可比障碍 $\Rightarrow$ 可修复障碍表达）]
\label{cor:tex29-formal-jacobi-to-rep}
在“分层极值定位—障碍登记—修复证据链—闭合校验”的流程\KouJing下，
每一层的雅可比障碍信号被编码为 $\omega_n$，
并通过证据链 $D_n$ 满足 $\partial^{(n)}D_n=\omega_n$ 被构造性平凡化；
因此“雅可比障碍”由终止判据转写为修复输入对象。
\end{corollary}

\begin{Certificate}[证书：编码 $\Rightarrow$ 平凡化]
\[
\begin{aligned}
  J_1 &: \text{雅可比障碍信号在第 }n\text{ 层被编码为 }\omega_n,\\
  J_2 &: \exists D_n\ (\partial^{(n)}D_n=\omega_n)\Rightarrow \mathsf{Rep}(\omega_n),\\
  J_3 &: J_1\wedge J_2\Rightarrow \text{雅可比障碍可作为修复输入对象}.
\end{aligned}
\]
\end{Certificate}

\paragraph{证明（谓词因果链）.}
由 $J_1$ 完成“几何侧失败信号”的对象化登记；由 $J_2$ 得其进入 $\partial D=\omega$ 的可修复口径；
合并得 $J_3$，从而完成“终止判据 $\Rightarrow$ 构造输入”的转写。证毕.

\begin{corollary}[推论 C3（隐含修复自由度证书）]
\label{cor:tex29-formal-hidden-dof}
若存在投影接口 $\mathrm{pr}$ 使得 $\ker(\mathrm{pr})\neq 0$，
则对任意点 $x$ 可作分解
\[
  \alpha_x=\mathrm{pr}_x(\alpha_x)+\alpha_x^{\mathrm{hid}},\qquad \alpha_x^{\mathrm{hid}}\in\ker(\mathrm{pr}_x),
\]
从而得到“观测不可见但对修复可用”的隐含自由度证书；
并指出这类区分量在“莱布尼茨无窗单子”口径下不出现（或被投影消去）。
\end{corollary}

\begin{Certificate}[证书：核非零 $\Rightarrow$ 存在不可见分量]
\[
\begin{aligned}
  H_1 &: \ker(\mathrm{pr})\neq 0\Rightarrow \exists u\neq 0\in\ker(\mathrm{pr}),\\
  H_2 &: u\in\ker(\mathrm{pr})\Rightarrow \text{在观测口径下被投影为零，但在内部语法中可参与复合/修复},\\
  H_3 &: H_1\wedge H_2\Rightarrow \text{隐含修复自由度证书成立}.
\end{aligned}
\]
\end{Certificate}

\paragraph{证明（谓词因果链）.}
由 $H_1$ 取非零核向量作为隐含分量候选；
由 $H_2$ 固定其“不可见但可用”的双重属性；
合并得 $H_3$。证毕.

\subsection{备注（直接回答“增益 + 跃迁”）}
\label{subsec:tex29-formal-remarks-gain-jump}

\begin{remark}[对体系的“增益”：可压缩为 5 个新增结构]
在定义~\ref{def:tex29-formal-system-gain-jump} 的比较口径下，本文档对体系的增益可压缩为 5 个新增结构：
\begin{enumerate}
  \item \textbf{边界增益（去混同）}：把“修复同伦类全集”的唯一真值入口固定为 $\mathsf{Rep}$，
  并以定理~\ref{thm:tex29-formal-fr-strict} 把 $\mathsf{Fr}\subsetneq\mathsf{Rep}$ 写成可检验的刚性边界；
  \item \textbf{分类增益（全场景可定位）}：以 $\tau(\omega)$ 与 $\kappa(\omega)$ 给出“修复形式矩阵”，使“该用哪一类修复逻辑”成为可判定对象（见命题~\ref{prop:tex29-formal-repair-form-matrix}）；
  \item \textbf{证书增益（可审计闭包）}：以定理~\ref{thm:tex29-formal-axiom-closure} 与定理~\ref{thm:tex29-formal-all-degree-trivial}
    为代表的“公理降级为定理（数学归纳法证书）”，
  把闭包、递归与链合成写成可检查的 $P(n)\Rightarrow P(n+1)$ 结构；
  \item \textbf{动力学增益（静态判定 $\Rightarrow$ 内生闭环接口）}：以缺陷势单调性 + $\partial D=\omega$ 的修复约束 + $\partial^2=0$ 的闭合校验，
  把“拓扑—统计—变分—同伦修复”统一成同一条谓词因果链，使修复成为可执行引擎接口；
  \item \textbf{跨分支增益（几何障碍可翻译为修复对象）}：把雅可比/比安基等“几何侧失败信号”系统性转写为 $\omega$ 并进入 $\mathsf{Rep}$ 处理（见推论~\ref{cor:tex29-formal-jacobi-to-rep}），
  并以核分解给出隐含修复自由度证书（见推论~\ref{cor:tex29-formal-hidden-dof}）。
\end{enumerate}
\end{remark}

\begin{remark}[对体系的“跃迁”：可形式化判定的结构条件]
本章所说的跃迁不是“多一个分支公设”，而是满足如下可检验结构：
\[
  [\omega]\neq 0\ \text{不再作为终止判据}
  \quad\Longrightarrow\quad
  \exists D\ (\partial D=\omega)\ \text{作为构造步骤},
\]
并最终可导出全阶闭环目标（无空洞完成态）：
\[
  (\forall n\ge 1)\ H^n_{\mathrm{obs}}(X;G)=0\quad\Longleftrightarrow\quad \mathsf{NoHole}(X).
\]
换言之：把“核内空洞/裂缝”从不可消解的判定结果转写为可生成边界的证据链任务——
这就是本文档所强调的“地基级重构”意义上的跃迁；并且它在形式系统中被写成可核验的谓词与证书链，而非叙述性的未证明断言。
\end{remark}

%%%%%%%%%%%%%%%%%%%%%%%%%%%%%%%%%%%%%%%%%%%%%%%%%%%%%%%%%%%%
\section{\texorpdfstring{形式逻辑：「离散统$\to$卡丘胚$\to$修复塔$\to$四维时空」完整生成链的统一闭环}{形式逻辑：离散统->卡丘胚->修复塔->四维时空完整生成链的统一闭环}}
\label{sec:tex29-chain-discrete-kqemb-tower-4d}
%%%%%%%%%%%%%%%%%%%%%%%%%%%%%%%%%%%%%%%%%%%%%%%%%%%%%%%%%%%%

\subsection{定义}
\label{subsec:tex29-chain-discrete-kqemb-tower-4d-def}

\begin{definition}[五环对象、生成箭头与“卡丘胚/投影核主丛胚”接口]
\label{def:tex29-chain5-objects}
定义以下对象与结构：
\[
  \mathfrak{D}:=\bigl(P_{\mathfrak D}\xrightarrow{\pi_{\mathfrak D}}B_{\mathfrak D},\mathfrak{g}_{\mathrm{GZ}}\bigr),
  \qquad
  \mathfrak{K}:=\bigl(P_{\mathfrak K}\xrightarrow{\pi_{\mathfrak K}}B_{\mathfrak K},\mathfrak{g}_{\mathrm{GZ}}\bigr),
\]
\[
  \{\mathfrak{T}_{\alpha}\}_{\alpha<\Lambda},
  \qquad
  (\mathcal{M}^{4},g,\nabla),
  \qquad
  (\mathcal{A}_{\mathrm{law}},\mathcal{F}_{\mathrm{law}}).
\]
其中：
\begin{enumerate}
  \item $\mathfrak{D}$ 表示离散统（允许离散状态空间与离散时间演化；参照 \cite{GaoZheng2025GFramework} 中“heterogeneous iterative systems”的起点设定）。
  \item $\Pi_{\mathrm{nc}}$ 表示“卡丘完备化/胚化”算子：它把离散统的运算空间在给定度量下做完备化，并闭包于卡丘列极限（与 \cite{GaoZheng2025GFramework}
    中“GZ-admissible law-system”对完备性与 Cauchy 极限的要求对齐）。
        记
        \[
          \mathfrak{K}:=\Pi_{\mathrm{nc}}(\mathfrak{D})
        \]
        并称 $\mathfrak{K}$ 为“卡丘胚”。
  \item （扩展数据，若启用）在局部化片 $A_f$ 上选取幂等投影 $p\in M_n(A_f)$（$p^2=p$），则可把“卡丘胚”细化为“非交换投影核主丛胚”之胚等价类
        \[
          \mathfrak{K}_{(f)}^{\mathrm{ncPK}} := [(f,p,\mathcal{G}_f)],
        \]
        其中 $\mathcal{G}_f\subseteq GL_n(A_f)$ 的作用为 $p\mapsto gpg^{-1}$，并记
        \[
          E_f:=pA_f^n,\qquad K_f:=\ker(p).
        \]
        （该细化数据在两部专著中尚未显式写出；此处作为接口定义使用。）
  \item $\mathfrak{T}_{\alpha}$ 是第 $\alpha$ 层超限同伦修复塔对象。
  \item $\mathcal{P}_{\mathrm{4D}}$ 是低维物理投影/切片算子；在 \cite{GaoZhengAppVol1} 的 GR 章节中，“四维流形 + 洛伦兹度规”属于被选定的背景，而非由 OHU 逻辑自动推出。

\end{enumerate}
完整生成链记为
\[
  \mathfrak{D}
  \xrightarrow{\Pi_{\mathrm{nc}}}
  \mathfrak{K}
  \xrightarrow{\mathcal{R}}
  \{\mathfrak{T}_{\alpha}\}_{\alpha<\Lambda}
  \xrightarrow{\mathcal{P}_{\mathrm{4D}}}
  (\mathcal{M}^{4},g,\nabla).
\]
\end{definition}

\begin{definition}[连续统分层/崩溃谓词与证书量]
\label{def:tex29-chain5-predicates}
对每一层 $\alpha$，定义雅可比间隙证书量与连续统破裂因子：
\[
  \mathrm{JacGap}_{\alpha}:=\lVert \mathcal{J}_{\alpha} \rVert,
  \qquad
  \mathfrak{B}_{\mathrm{CH},\alpha}:=\lVert \Delta_{\alpha}(\aleph_0,\aleph_1) \rVert.
\]
定义稳定谓词与崩溃谓词：
\[
  \mathsf{Stable}_{\alpha}\ :\Longleftrightarrow\
  \bigl(\mathrm{JacGap}_{\alpha}=0\bigr)\wedge\bigl(\mathfrak{B}_{\mathrm{CH},\alpha}=0\bigr),
\]
\[
  \mathsf{Collapse}_{\alpha}\ :\Longleftrightarrow\ \neg \mathsf{Stable}_{\alpha}.
\]
定义物理表观奇点谓词：
\[
  \mathsf{Sing}_{\alpha}\ :\Longleftrightarrow\ \mathsf{Collapse}_{\alpha}.
\]
故“连续统分层是否自洽”与“时空是否奇点化”被统一为同一判定接口。
\end{definition}

\subsection{公理（降级为定理：数学归纳法证书 + 证明）}
\label{subsec:tex29-chain-discrete-kqemb-tower-4d-axiom}

\begin{theorem}[公理降级：五环生成链有限截断的可构造闭包（数学归纳法证书）]
\label{thm:tex29-chain5-axiom-downgraded}
在定义~\ref{def:tex29-chain5-objects}--\ref{def:tex29-chain5-predicates} 下，设：
\begin{enumerate}
  \item[(A1)] 基步：存在 $\omega_0$ 与 $D_0$ 使 $\partial D_0=\omega_0$，且 $\mathsf{Stable}_0$；
  \item[(A2)] 递推：若第 $N$ 层存在 $D_N$ 使 $\partial^{(N)}D_N=\omega_N$ 且 $\mathsf{Stable}_N$，
  则存在 $D_{N+1}$ 使 $\partial^{(N+1)}D_{N+1}=\omega_{N+1}$，并有 $\mathsf{Stable}_{N+1}$；
  \item[(A3)] 相容：每层修复与投影满足
  \[
    \mathcal{P}_{\mathrm{4D}}\circ \mathcal{R}_{N\to N+1}
    \simeq
    \mathcal{R}^{\mathrm{4D}}_{N\to N+1}\circ \mathcal{P}_{\mathrm{4D}}.
  \]
\end{enumerate}
则对任意 $N\in\NN$，存在有限修复塔截断
\[
  \mathfrak{D}\to\mathfrak{K}\to(\mathfrak{T}_0\to\cdots\to\mathfrak{T}_N)\to(\mathcal{M}^{4}_{N},g_N,\nabla_N),
\]
并存在合成证据链 $\widehat{D}_{\le N}:=(D_0,\dots,D_N)$ 使
\[
  \widehat{\partial}_N(\widehat{D}_{\le N})=(\omega_0,\dots,\omega_N),
\]
且每层均满足 $\mathsf{Stable}_k$（$0\le k\le N$）。
\end{theorem}

\begin{Certificate}[数学归纳法证书：按修复塔层数 $N$]
\label{cert:tex29-chain5-axiom-downgraded}
\[
\begin{aligned}
  I_0 &:\exists D_0\ \bigl(\partial D_0=\omega_0\bigr)\wedge \mathsf{Stable}_0,\\
  I_{N\Rightarrow N+1} &: \Bigl(\exists D_N\ \bigl(\partial^{(N)}D_N=\omega_N\bigr)\wedge \mathsf{Stable}_N\Bigr)\\
  &\qquad\Rightarrow \Bigl(\exists D_{N+1}\ \bigl(\partial^{(N+1)}D_{N+1}=\omega_{N+1}\bigr)\wedge \mathsf{Stable}_{N+1}
    \Bigr),\\
  I_{\mathrm{comp}} &: \mathcal{P}_{\mathrm{4D}}\circ \mathcal{R}_{N\to N+1}
  \simeq
  \mathcal{R}^{\mathrm{4D}}_{N\to N+1}\circ \mathcal{P}_{\mathrm{4D}}.
\end{aligned}
\]
\end{Certificate}

\paragraph{证明（谓词因果链）.}
由 (A1) 得基步 $I_0$。
设对某 $N$ 已有 $\exists D_N(\partial^{(N)}D_N=\omega_N)\wedge \mathsf{Stable}_N$，则由 (A2) 得递推步 $I_{N\Rightarrow N+1}$，
从而对任意 $N$ 可构造证据链 $\widehat{D}_{\le N}=(D_0,\dots,D_N)$ 并保证各层 $\mathsf{Stable}_k$。
再由 (A3) 得相容性 $I_{\mathrm{comp}}$，因此有限修复塔截断与其四维投影链可构造且自洽。证毕.

\subsection{定理（证书 + 证明）}
\label{subsec:tex29-chain-discrete-kqemb-tower-4d-theorem}

\begin{theorem}[定理 T1（GZ-OHU 保证下的超限完备：塔极限存在；四维表观需额外投影假设）]
\label{thm:tex29-chain5-ohu-4d}
若存在极限序数 $\Lambda$ 使得：
\begin{enumerate}
  \item[(U1)] $\forall \alpha<\Lambda,\ \exists D_{\alpha}\ \bigl(\partial^{(\alpha)}D_{\alpha}=\omega_{\alpha}\bigr)$；
  \item[(U2)] $\forall \alpha<\Lambda,\ \mathsf{Stable}_{\alpha}$；
  \item[(U3)] 在极限层满足 GZ-OHU 小对象完备条件（参照 \cite{GaoZheng2025GFramework}），故修复塔极限对象 $\mathfrak{T}_{\Lambda}$ 存在且一致；
\end{enumerate}
则存在无预设时空背景的四维投影
\[
  \mathcal{P}_{\mathrm{4D}}(\mathfrak{T}_{\Lambda})=(\mathcal{M}^{4},g,\nabla),
\]
并且低能几何量由定律联络/曲率投影给出：
\[
  g=\mathbf{G}\bigl(\mathcal{P}_{\mathrm{4D}}(\mathcal{A}_{\mathrm{law}})\bigr),\qquad
  R=\mathbf{R}\bigl(\mathcal{P}_{\mathrm{4D}}(\mathcal{F}_{\mathrm{law}})\bigr).
\]
\end{theorem}

\begin{Certificate}[证书链：完备化 $\Rightarrow$ 塔极限；四维表观需额外投影假设]
\label{cert:tex29-chain5-ohu-4d}
\[
\begin{aligned}
  T_1 &: \forall \alpha<\Lambda,\ \exists D_{\alpha}\ (\partial^{(\alpha)}D_{\alpha}=\omega_{\alpha}),\\
  T_2 &: \forall \alpha<\Lambda,\ \mathsf{Stable}_{\alpha},\\
  T_3 &: T_1\wedge T_2\wedge \text{OHU 完备条件}\Rightarrow \exists \mathfrak{T}_{\Lambda},\\
  T_4 &: \text{(附加假设)}\ \exists\,\mathcal{P}_{\mathrm{4D}}:\mathfrak{T}_{\Lambda}\to (\mathcal{M}^{4},g,\nabla).
\end{aligned}
\]
\end{Certificate}

\paragraph{证明（谓词因果链）.}
由 $T_1$ 得逐层可修复；由 $T_2$ 排除层间崩溃；
合并 $T_1,T_2$ 与 OHU 完备条件即得 $T_3$，从而超限极限对象 $\mathfrak{T}_{\Lambda}$ 存在且一致。
四维表观并非由 $T_3$ 逻辑必然推出；若满足附加假设 $T_4$（给定投影/切片算子 $\mathcal{P}_{\mathrm{4D}}$），则可定义
$\mathcal{P}_{\mathrm{4D}}(\mathfrak{T}_{\Lambda})=(\mathcal{M}^{4},g,\nabla)$。
度规与曲率表达式由投影后的定律联络与定律曲率的函数化定义给出，故定理成立。证毕.

\subsection{引理（证书 + 证明）}
\label{subsec:tex29-chain-discrete-kqemb-tower-4d-lemma}

\begin{lemma}[引理 L1（五环生成链与连续统分层/崩溃的一一对应）]
\label{lem:tex29-chain5-layer-collapse}
在定义~\ref{def:tex29-chain5-objects}--\ref{def:tex29-chain5-predicates} 下，
五环生成链可与连续统分层/崩溃状态作如下对应：
\[
\begin{aligned}
  &\text{离散统层}\ \Longleftrightarrow\ \aleph_0\ \text{基底层},\quad
  \mathsf{Collapse}\Longleftrightarrow \mathrm{JacGap}\neq 0;\\
  &\text{卡丘胚层}\ \Longleftrightarrow\ \aleph_0\to\aleph_1\ \text{临界提升层},\quad
  \mathsf{Collapse}\Longleftrightarrow \mathfrak{B}_{\mathrm{CH}}\neq 0;\\
  &\text{修复元层}\ \Longleftrightarrow\ \text{层间映射修复层};\\
  &\text{超限修复塔层}\ \Longleftrightarrow\ (\aleph_0,\aleph_1,\aleph_2,\dots)\ \text{全阶展开层};\\
  &\text{四维投影层}\ \Longleftrightarrow\ \text{连续统宏观物理表观层}.
\end{aligned}
\]
\end{lemma}

\begin{Certificate}[证书：分层谓词到崩溃谓词的逐环翻译]
\label{cert:tex29-chain5-layer-collapse}
\[
\begin{aligned}
  L_1 &: \mathrm{JacGap}_{\alpha}\neq 0\Rightarrow \mathsf{Collapse}_{\alpha},\\
  L_2 &: \mathfrak{B}_{\mathrm{CH},\alpha}\neq 0\Rightarrow \mathsf{Collapse}_{\alpha},\\
  L_3 &: \mathsf{Stable}_{\alpha}\Rightarrow \neg\mathsf{Collapse}_{\alpha}\Rightarrow \text{分层可延拓}.
\end{aligned}
\]
\end{Certificate}

\paragraph{证明（谓词因果链）.}
由定义~\ref{def:tex29-chain5-predicates}，$L_1,L_2$ 直接成立；
由 $\mathsf{Stable}_{\alpha}$ 的定义即得 $L_3$。
因此五环对象可逐环翻译为“分层状态 + 崩溃状态 + 修复状态”的统一描述。证毕.

\subsection{命题（证书 + 证明）}
\label{subsec:tex29-chain-discrete-kqemb-tower-4d-proposition}

\begin{proposition}[命题 P1（连续统假设状态的曲率化转写）]
\label{prop:tex29-chain5-ch-curvature}
在同一层 $\alpha$，定义层间曲率状态量
\[
  \chi_{\alpha}:=\mathcal{F}_{\mathrm{law},\alpha}\!\llcorner\!(\aleph_0,\aleph_1).
\]
若存在阈值函数 $\Theta$ 使
\[
  \mathfrak{B}_{\mathrm{CH},\alpha}=\Theta(\chi_{\alpha}),
\]
则连续统假设在该层的成立/失效可写成：
\[
  \Theta(\chi_{\alpha})=0\Rightarrow \text{CH-相容层},
  \qquad
  \Theta(\chi_{\alpha})\neq 0\Rightarrow \text{CH-崩溃层}.
\]
故“CH 的独立性表述”在本\KouJing 下被转写为“层间曲率状态”的结构判定。
\end{proposition}

\begin{Certificate}[证书：曲率状态 $\Rightarrow$ CH 层态]
\label{cert:tex29-chain5-ch-curvature}
\[
\begin{aligned}
  C_1 &: \chi_{\alpha}\mapsto \Theta(\chi_{\alpha})=\mathfrak{B}_{\mathrm{CH},\alpha},\\
  C_2 &: \mathfrak{B}_{\mathrm{CH},\alpha}=0\Rightarrow \neg\mathsf{Collapse}_{\alpha}\ \text{(CH-相容层)},\\
  C_3 &: \mathfrak{B}_{\mathrm{CH},\alpha}\neq 0\Rightarrow \mathsf{Collapse}_{\alpha}\ \text{(CH-崩溃层)}.
\end{aligned}
\]
\end{Certificate}

\paragraph{证明（谓词因果链）.}
由 $C_1$ 把 CH 状态写成可计算的曲率函数值；
再由定义~\ref{def:tex29-chain5-predicates}，$C_2,C_3$ 直接成立；
故命题成立。证毕.

\subsection{推论（证书 + 证明）}
\label{subsec:tex29-chain-discrete-kqemb-tower-4d-corollary}

\begin{corollary}[推论 C1（连续统崩溃与时空奇点同源，并可由修复塔消解）]
\label{cor:tex29-chain5-singularity-collapse}
若某物理区域对应层 $\alpha$ 满足 $\mathsf{Collapse}_{\alpha}$，则该区域出现奇点化表观：
\[
  \mathsf{Collapse}_{\alpha}\Rightarrow \mathsf{Sing}_{\alpha}.
\]
若进一步满足 GZ-OHU 完备条件与可延拓修复链，
则存在 $\beta\ge \alpha$ 使该奇点化被消解：
\[
  \exists \beta\ge \alpha,\ \neg\mathsf{Collapse}_{\beta}
  \Rightarrow
  \neg\mathsf{Sing}_{\beta}.
\]
\end{corollary}

\begin{Certificate}[证书：崩溃同源 + 可修复消解]
\label{cert:tex29-chain5-singularity-collapse}
\[
\begin{aligned}
  R_1 &: \mathsf{Collapse}_{\alpha}\Longleftrightarrow \mathsf{Sing}_{\alpha},\\
  R_2 &: \text{OHU 完备} \wedge \text{可延拓修复链}\Rightarrow \exists \beta\ge\alpha,\ \neg\mathsf{Collapse}_{\beta},\\
  R_3 &: \neg\mathsf{Collapse}_{\beta}\Rightarrow \neg\mathsf{Sing}_{\beta}.
\end{aligned}
\]
\end{Certificate}

\paragraph{证明（谓词因果链）.}
由定义~\ref{def:tex29-chain5-predicates} 得 $R_1$；
由定理~\ref{thm:tex29-chain5-ohu-4d} 得 $R_2$；
由 $R_1$ 的逆向使用得 $R_3$。
合并 $R_1,R_2,R_3$ 即得推论。证毕.

\subsection{备注}
\label{subsec:tex29-chain-discrete-kqemb-tower-4d-remarks}

\begin{remark}[本章闭环的三条结构结论]
“离散统$\to$卡丘胚$\to$修复塔$\to$四维时空”在本章被落实为三条可核验结论，并可进一步展开为可审计的接口清单：
\begin{enumerate}
  \item \textbf{生成性}：四维时空是修复塔完备后的低维投影表观，而非预设背景；
  \item \textbf{统一性}：连续统分层与时空生成共享同一谓词链（$\mathsf{Stable}/\mathsf{Collapse}/\mathsf{Sing}$）；
  \item \textbf{可修复性}：连续统崩溃与奇点化不是终止判据，而是可继续输入修复塔的构造任务。
\end{enumerate}
对应到本章的“最小可核验接口”，可压缩为 5 个入口：
\begin{enumerate}
  \item \textbf{对象入口}：$\mathfrak{D},\mathfrak{K},\mathfrak{T}_{\alpha},(\mathcal{M}^4,g,\nabla)$；
  \item \textbf{箭头入口}：$\Pi_{\mathrm{nc}},\mathcal{R},\mathcal{P}_{\mathrm{4D}}$；
  \item \textbf{证书入口}：$\mathrm{JacGap}_{\alpha}$ 与 $\mathfrak{B}_{\mathrm{CH},\alpha}$；
  \item \textbf{谓词入口}：$\mathsf{Stable}_{\alpha},\mathsf{Collapse}_{\alpha},\mathsf{Sing}_{\alpha}$；
  \item \textbf{修复入口}：回到全文统一的 $\mathsf{Rep}(\omega)\Leftrightarrow \exists D(\partial D=\omega)$，并把每层的“失败信号”对象化为
    $\omega_{\alpha}$ 进入证据链。
\end{enumerate}
\end{remark}

\begin{remark}[表格化：把叙述性链条压缩为可核验对照]
为便于复核与复用，下面表~\ref{tab:tex29-chain5-continuum-map} 把“生成链环节—连续统分层—崩溃信号—修复入口”写成同一张对照表。
其中出现的 PL-PI、GZ-CBT、GZ-OHU 等名称在本\KouJing下仅作为“命名锚点/术语登记”，
形式证明仍以本章的谓词（$\mathsf{Stable}/\mathsf{Collapse}/\mathsf{Sing}$）与证书链（$D_{\alpha}$）为准。
\end{remark}

\begin{table}[htbp]
\centering
{\footnotesize
\setlength{\tabcolsep}{2pt}
\renewcommand{\arraystretch}{1.25}
\begin{tabular}{|p{0.17\linewidth}|p{0.23\linewidth}|p{0.26\linewidth}|p{0.28\linewidth}|}
\hline
\textbf{生成链环节} & \textbf{连续统分层对应} & \textbf{连续统崩溃的对应表现（证书量）} & \textbf{GZ 体系修复方案（形式入口）} \\
\hline
离散统 $\mathfrak{D}$（\GZMonFull{}） & 可数基底层 $\aleph_0$ &
$\mathrm{JacGap}_{\alpha}\neq 0$（非交换法结构缺口证书） &
以 $\mathsf{Rep}$ 作为真值入口；PL-PI 规则/容许定律约束给出可构造修复链 \\
\hline
投影生成卡丘胚 $\mathfrak{K}=\Pi_{\mathrm{nc}}(\mathfrak{D})$ &
临界提升层 $\aleph_0\to\aleph_1$ &
$\mathfrak{B}_{\mathrm{CH},\alpha}\neq 0$（层间提升破裂因子） &
结构化非交换投影 + 核携带；用（命名锚定的）GZ-CBT 口径构造层间单射/双射修复 \\
\hline
卡丘胚作为修复元（算子载体） &
连续统修复核心层（层间映射算子层） &
$\mathsf{Collapse}_{\alpha}$：投影核破裂、层间映射不可构造 &
以 $\partial D=\omega$ 的证据链平凡化作为修复动作；GZ-OHU 提供“无失败/同伦完备”兜底 \\
\hline
超限递归同伦修复塔 $\{\mathfrak{T}_{\alpha}\}_{\alpha<\Lambda}$ &
全阶分层展开 $\aleph_0,\aleph_1,\aleph_2,\dots$ &
高阶上同调障碍残留/曲率发散（对应 $\neg\mathsf{Stable}_{\alpha}$） &
逐层证书化消解：$\forall \alpha<\Lambda,\exists D_{\alpha}$；极限层由 OHU 小对象条件封顶 \\
\hline
四维时空投影 $(\mathcal{M}^{4},g,\nabla)=\mathcal{P}_{\mathrm{4D}}(\mathfrak{T}_{\Lambda})$ &
宏观表观层（低维连续实现） &
奇点化表观 $\mathsf{Sing}_{\alpha}\Leftrightarrow \mathsf{Collapse}_{\alpha}$ &
由定律联络/曲率投影给出 $g,R$；稳定性由 $\mathsf{Stable}$ 判定保障 \\
\hline
\end{tabular}
}
\caption{\texorpdfstring{“离散统$\to$卡丘胚$\to$修复塔$\to$四维时空”}{“离散统->卡丘胚->修复塔->四维时空”}与连续统分层/崩溃的形式对照表（\KouJing）}
\label{tab:tex29-chain5-continuum-map}
\end{table}

\begin{remark}[符号表（接口最小集）]
为使后续引用保持无歧义，表~\ref{tab:tex29-chain5-symbols} 列出本章使用的最小符号/接口集合。
其中 $\mathsf{Rep}$ 的定义已在全文前置给出（定义~\ref{def:tex29-rep-fr}）。
\end{remark}

\begin{table}[htbp]
\centering
{\footnotesize
\setlength{\tabcolsep}{2pt}
\renewcommand{\arraystretch}{1.2}
\begin{tabular}{|p{0.22\linewidth}|p{0.22\linewidth}|p{0.50\linewidth}|}
\hline
\textbf{符号} & \textbf{类型/域} & \textbf{说明（本章口径）} \\
\hline
$\mathfrak{D}$ & 离散统对象 & \GZMonFull{} 第一性基元的结构载体（主丛 + 法/代数数据） \\
\hline
$\Pi_{\mathrm{nc}}$ & 非交换结构化投影 & 从离散统抽取预几何胚并保留投影核信息 \\
\hline
$\mathfrak{K}$ & 卡丘胚对象 & $\mathfrak{K}:=\Pi_{\mathrm{nc}}(\mathfrak{D})$；既是预几何胚也是修复元（算子载体） \\
\hline
$\mathcal{R}$ & 修复塔生成算子 & 逐层构造 $\mathfrak{T}_{\alpha}$ 并生成 $D_{\alpha}$ 证据链 \\
\hline
$\mathfrak{T}_{\alpha}$ & 修复塔对象 & 第 $\alpha$ 层同伦修复/完备化对象（$\alpha<\Lambda$） \\
\hline
$\mathcal{P}_{\mathrm{4D}}$ & 低维物理投影 & 输出四维表观几何 $(\mathcal{M}^4,g,\nabla)$ \\
\hline
$\mathrm{JacGap}_{\alpha}$ & 证书量 & 非交换法结构缺口的诊断量；为稳定性判定的输入 \\
\hline
$\mathfrak{B}_{\mathrm{CH},\alpha}$ & 证书量 & $\aleph_0\to\aleph_1$ 层间提升的破裂因子（CH 层态的曲率化指标） \\
\hline
$\mathsf{Rep}(\omega)$ & 谓词 & 修复谓词：$\exists D(\partial D=\omega)$（定义~\ref{def:tex29-rep-fr}） \\
\hline
$\mathsf{Stable}_{\alpha}$ & 谓词 & $\mathrm{JacGap}_{\alpha}=0$ 且 $\mathfrak{B}_{\mathrm{CH},\alpha}=0$ \\
\hline
$\mathsf{Collapse}_{\alpha}$ & 谓词 & $\neg\mathsf{Stable}_{\alpha}$ \\
\hline
$\mathsf{Sing}_{\alpha}$ & 谓词 & 物理表观奇点谓词；本章取 $\mathsf{Sing}_{\alpha}\Longleftrightarrow\mathsf{Collapse}_{\alpha}$ \\
\hline
\end{tabular}
}
\caption{本章五环生成链的最小符号/接口表}
\label{tab:tex29-chain5-symbols}
\end{table}

\begin{remark}[与既有口径的对齐]
本章新增内容与全文既有“边缘算子平凡化 + 证书链 + 公理降级”口径严格一致：
所有关键断言都被写成“可计算状态量 + 可核验证书 + 谓词因果链证明”，
从而将叙述性主张统一压缩为可审计的形式逻辑条目。
\par
特别地：
\begin{enumerate}
  \item “连续统分层/崩溃 $\leftrightarrow$ 时空生成/奇点化”并非额外假设，而是以定义~\ref{def:tex29-chain5-predicates} 的谓词同一化作为接口约定；
  \item “修复塔兜底”并非叙述性承诺，而是以定理~\ref{thm:tex29-chain5-ohu-4d} 的证书链条目化表述；
  \item “对照表”并不替代证明：它只是把证书量与谓词入口集中登记，便于检查每一步调用的前提与输出。
\end{enumerate}
\end{remark}

%%%%%%%%%%%%%%%%%%%%%%%%%%%%%%%%%%%%%%%%%%%%%%%%%%%%%%%%%%%%
\section{\texorpdfstring{形式逻辑：两部专著对“离散统$\to$卡丘胚$\to$修复塔$\to$四维时空”生成链的全覆盖与内核等价}{形式逻辑：两部专著对离散统->卡丘胚->修复塔->四维时空生成链的全覆盖与内核等价}
  }
\label{sec:tex29-books-cover-chain}
%%%%%%%%%%%%%%%%%%%%%%%%%%%%%%%%%%%%%%%%%%%%%%%%%%%%%%%%%%%%

\subsection{定义}
\label{subsec:tex29-books-cover-chain-def}

\begin{definition}[两部专著、五环生成链与“覆盖”谓词]
\label{def:tex29-books-cover-chain}
记两部专著（抽象文献对象）为
\[
  \mathcal{B}:=\{\mathcal{B}^{(1)},\mathcal{B}^{(2)}\}.
\]
记五环生成链的对象与箭头沿用定义~\ref{def:tex29-chain5-objects}，并按环编号 $k\in\{1,2,3,4,5\}$ 记为
\[
  L_1=\mathfrak{D},\quad L_2=\mathfrak{K},\quad L_3=\text{“}\mathfrak{K}\text{作为修复元”},\quad
  L_4=\{\mathfrak{T}_{\alpha}\}_{\alpha<\Lambda},\quad L_5=(\mathcal{M}^{4},g,\nabla).
\]
对任意对象符号 $X$，定义三类“覆盖证书谓词”：
\[
  \mathsf{DefCov}_{\mathcal{B}}(X):\Longleftrightarrow \text{在 }\mathcal{B}\text{ 中存在 }X\text{ 的定义性锚定条目},
\]
\[
  \mathsf{ThmCov}_{\mathcal{B}}(X):\Longleftrightarrow \text{在 }\mathcal{B}\text{ 中存在支撑 }X\text{ 的定理级保障条目},
\]
\[
  \mathsf{AlgCov}_{\mathcal{B}}(X):\Longleftrightarrow \text{在 }\mathcal{B}\text{ 中存在支撑 }X\text{ 的可执行算法/证书链接口}.
\]
并定义“全覆盖谓词”：
\[
  \mathsf{Cover}_{\mathcal{B}}(X)\ :\Longleftrightarrow\
  \mathsf{DefCov}_{\mathcal{B}}(X)\wedge \mathsf{ThmCov}_{\mathcal{B}}(X)\wedge \mathsf{AlgCov}_{\mathcal{B}}(X).
\]
\end{definition}

\begin{definition}[“无字面术语但内核完备”谓词（以卡丘胚为例）]
\label{def:tex29-books-kqemb-kernel}
令字面术语 $\mathbf{t}_{\mathrm{KQE}}$ 表示字符串“卡丘胚”。定义“术语缺省”谓词：
\[
  \mathsf{TermOmit}_{\mathcal{B}}(\mathbf{t}_{\mathrm{KQE}})\ :\Longleftrightarrow\ \mathbf{t}_{\mathrm{KQE}}\ \text{不作为}
    \ \mathcal{B}\ \text{中的显式命名条目出现}.
\]
定义“卡丘胚内核谓词” $\mathsf{Ker}_{\mathrm{KQE}}(X)$（仅取最小可核验集合）为：
\[
\begin{aligned}
  \mathsf{Ker}_{\mathrm{KQE}}(X)\ :\Longleftrightarrow\
  &\text{(K1) }X=\Pi_{\mathrm{nc}}(\mathfrak{D})\ \text{为非交换结构化投影所得的主丛胚（预几何对象）；}\\
  &\text{(K2) }X\ \text{携带投影核信息并保留离散统的非交换代数数据（核携带）；}\\
  &\text{(K3) }X\ \text{在}\ \mathcal{B}\ \text{中满足“由规则导出且无任意性”的唯一性/可回溯性证书；}\\
  &\text{(K4) }X\ \text{存在内蕴维度特征证书（例如 }d_{\mathrm{int}}(X)=6\text{ 的占位性陈述）；}\\
  &\text{(K5) }X\ \text{可作为连续统层间修复的最小算子载体（修复元角色）；}\\
  &\text{(K6) }X\ \text{可作为超限同伦修复塔的基本构成元（与证据链 }D_{\alpha}\text{ 相容）；}\\
  &\text{(K7) }X\ \text{在低能投影下可被用于生成 }(\mathcal{M}^{4},g,\nabla)\ \text{的输入对象}.
\end{aligned}
\]
据此定义“内核完备”谓词：
\[
  \mathsf{KernelComplete}_{\mathcal{B}}(\mathbf{t}_{\mathrm{KQE}})\ :\Longleftrightarrow\
  \exists X\ \bigl(\mathsf{Ker}_{\mathrm{KQE}}(X)\wedge \mathsf{Cover}_{\mathcal{B}}(X)\bigr).
\]
\end{definition}

\subsection{公理（降级为定理：数学归纳法证书 + 证明）}
\label{subsec:tex29-books-cover-chain-axiom}

\begin{theorem}[公理降级：五环生成链的“全覆盖”可由逐环证书归纳推出]
\label{thm:tex29-books-cover-chain-induction}
若对每一环 $k\in\{1,2,3,4,5\}$ 均满足 $\mathsf{Cover}_{\mathcal{B}}(L_k)$，
并且相邻两环之间的箭头（$\Pi_{\mathrm{nc}},\mathcal{R},\mathcal{P}_{\mathrm{4D}}$）在 $\mathcal{B}$ 中同样被证书化为可复核的“跨层映射规则”，
则五环生成链作为一条\emph{连续逻辑主线}在本\KouJing下可被完整提炼，并不存在“环节缺失/逻辑断层”。
\end{theorem}

\begin{Certificate}[数学归纳法证书：按环编号 $k$ 的覆盖闭包]
\label{cert:tex29-books-cover-chain-induction}
令谓词 $P(N)$ 表示：
“对前 $N$ 环 $L_1,\dots,L_N$，均已在 $\mathcal{B}$ 中登记覆盖证书
（$\mathsf{Cover}_{\mathcal{B}}(L_i)$，$1\le i\le N$），并且存在从 $L_i$ 到 $L_{i+1}$ 的跨层映射规则登记。”
证书化目标为：
\[
  P(1),\qquad (\forall N\in\{1,2,3,4\})\ \bigl(P(N)\Rightarrow P(N+1)\bigr).
\]
\end{Certificate}

\paragraph{证明（谓词因果链）.}
由前提 $\mathsf{Cover}_{\mathcal{B}}(L_1)$ 得 $P(1)$。
假设 $P(N)$ 成立（$1\le N\le 4$），则已完成前 $N$ 环的覆盖登记；
又由前提 $\mathsf{Cover}_{\mathcal{B}}(L_{N+1})$ 与相邻映射规则的证书化登记，
可把第 $N+1$ 环并入同一条主线索引，得到 $P(N+1)$。
因此 $P(5)$ 成立，即五环生成链可被完整提炼且无断层，定理成立。证毕.

\subsection{定理（证书 + 证明）}
\label{subsec:tex29-books-cover-chain-theorem}

\begin{theorem}[定理 T1（“纲--目”关系：生成链是落地路径；专著是公理/工具/保障体系）]
\label{thm:tex29-books-cover-chain-outline-detail}
在定理~\ref{thm:tex29-books-cover-chain-induction} 的前提下，
可把生成链视为两部专著核心成果的“物理落地路径压缩表示”，而把专著视为该路径的三层支撑体系：
\[
  \begin{aligned}
    &\text{元数学公理层（PL-PI / GZ-G-代数）：给出 }L_1,L_2\text{ 的结构根基};\\
    &\Rightarrow\ \text{算法工具层（GZ-CBT / JacGap 证书）：给出 }L_3,L_4\text{ 的可执行修复接口};\\
    &\Rightarrow\ \text{物理投影层（定律空间 }\to\text{ 四维时空）：给出 }L_5\text{ 的投影生成规则}.
  \end{aligned}
\]
并且上述三层支撑不引入外部预设时空背景：所有“时空结构/维数/曲率/因果”均通过链上箭头与证书链对象化产生。
\end{theorem}

\begin{Certificate}[证书链：三层支撑 $\Rightarrow$ 路径压缩成立]
\label{cert:tex29-books-cover-chain-outline-detail}
\[
\begin{aligned}
  O_1 &: \mathsf{Cover}_{\mathcal{B}}(L_1)\wedge \mathsf{Cover}_{\mathcal{B}}(L_2)\Rightarrow
    \text{元数学公理层提供离散本原与投影胚的结构锚定},\\
  O_2 &: \mathsf{Cover}_{\mathcal{B}}(L_3)\wedge \mathsf{Cover}_{\mathcal{B}}(L_4)\Rightarrow \text{算法工具层提供修复元与修复塔的可执行接口}
    ,\\
  O_3 &: \mathsf{Cover}_{\mathcal{B}}(L_5)\Rightarrow \text{物理投影层提供 }(\mathcal{M}^4,g,\nabla)\text{ 的生成规则},\\
  O_4 &: O_1\wedge O_2\wedge O_3\Rightarrow \text{生成链可作为专著核心成果的压缩落地路径}.
\end{aligned}
\]
\end{Certificate}

\paragraph{证明（谓词因果链）.}
由 $O_1$ 得离散本原与预几何胚的定义/定理/算法锚定；
由 $O_2$ 得“修复元 + 修复塔”的证书化可执行接口；
由 $O_3$ 得低维物理投影输出四维时空的规则；
合并 $O_1,O_2,O_3$ 得 $O_4$，从而定理成立。证毕.

\subsection{引理（证书 + 证明）}
\label{subsec:tex29-books-cover-chain-lemma}

\begin{lemma}[引理 L1（无字面术语 $\not\Rightarrow$ 内核缺失：卡丘胚可由内核谓词唯一指认）]
\label{lem:tex29-books-kqemb-term-vs-kernel}
若 $\mathsf{TermOmit}_{\mathcal{B}}(\mathbf{t}_{\mathrm{KQE}})$ 且 $\mathsf{KernelComplete}_{\mathcal{B}}(\mathbf{t}
  _{\mathrm{KQE}})$，
则可在不改变任何证明义务的情况下，将满足 $\mathsf{Ker}_{\mathrm{KQE}}$ 的对象 $X$ 具象化命名为“卡丘胚”（记作 $\mathfrak{K}$），
并与定义~\ref{def:tex29-chain5-objects} 中的 $\mathfrak{K}$ 对齐。
\end{lemma}

\begin{Certificate}[证书：存在性 + 命名等价]
\label{cert:tex29-books-kqemb-term-vs-kernel}
\[
\begin{aligned}
  K_1 &: \mathsf{KernelComplete}_{\mathcal{B}}(\mathbf{t}_{\mathrm{KQE}})\Rightarrow \exists X\ \mathsf{Ker}
    _{\mathrm{KQE}}(X),\\
  K_2 &: \mathsf{Ker}_{\mathrm{KQE}}(X)\Rightarrow X\ \text{在链上承担 }L_2\text{ 的功能角色},\\
  K_3 &: (K_1\wedge K_2)\Rightarrow \text{可把 }X\text{ 具象化命名为 }\mathfrak{K}\ \text{而不改变任何证书链断言}.
\end{aligned}
\]
\end{Certificate}

\paragraph{证明（谓词因果链）.}
由 $K_1$ 取到满足内核谓词的对象 $X$；
由 $\mathsf{Ker}_{\mathrm{KQE}}$ 的 (K1)--(K3) 得 $X$ 在生成链中对应 $L_2$ 的最小功能角色，故得 $K_2$；
把“卡丘胚”作为该功能角色的具象化命名即得 $K_3$。
因此“无字面术语”不影响内核与证明义务，命题成立。证毕.

\subsection{命题（证书 + 证明）}
\label{subsec:tex29-books-cover-chain-proposition}

\begin{proposition}[命题 P1（连续统分层/崩溃的量化指标与修复方案使生成链可计算、可验证）]
\label{prop:tex29-books-continuum-metrics}
若 $\mathcal{B}$ 对连续统分层/崩溃给出量化指标与修复方案：
\[
  \text{指标：}\ \mathcal{F}_{\mathrm{law}},\ \mathfrak{B}_{\mathrm{CH},\alpha};\qquad
  \text{诊断：}\ \mathrm{JacGap}_{\alpha};\qquad
  \text{映射修复：}\ \text{GZ-CBT 结构化重写},
\]
并且这些量进入本章稳定性谓词（定义~\ref{def:tex29-chain5-predicates}），
则生成链从“叙述性框架”提升为“可计算、可校验、可迭代调控”的构造流程：
\[
  \mathsf{Stable}_{\alpha}\ \text{成为可核验接口},\qquad
  \mathsf{Collapse}_{\alpha}\ \text{成为可继续修复的输入对象}.
\]
\end{proposition}

\begin{Certificate}[证书：指标/诊断/修复 $\Rightarrow$ 可计算接口]
\label{cert:tex29-books-continuum-metrics}
\[
\begin{aligned}
  Q_1 &: \text{给出 }\mathcal{F}_{\mathrm{law}},\mathfrak{B}_{\mathrm{CH},\alpha}\Rightarrow \text{分层/崩溃状态可量化},\\
  Q_2 &: \text{给出 }\mathrm{JacGap}_{\alpha}\Rightarrow \text{结构偏差可诊断并进入证书链},\\
  Q_3 &: \text{给出 CBT 结构化重写}\Rightarrow \text{层间映射/修复可算法化执行},\\
  Q_4 &: Q_1\wedge Q_2\wedge Q_3\Rightarrow \mathsf{Stable}/\mathsf{Collapse}\ \text{成为可计算接口}.
\end{aligned}
\]
\end{Certificate}

\paragraph{证明（谓词因果链）.}
由 $Q_1$ 得连续统状态可量化；
由 $Q_2$ 得偏差诊断进入证书链（使“失败信号”对象化）；
由 $Q_3$ 得层间修复可执行；
合并得 $Q_4$，故命题成立。证毕.

\subsection{推论（证书 + 证明）}
\label{subsec:tex29-books-cover-chain-corollary}

\begin{corollary}[推论 C1（全覆盖 + 内核等价 $\Rightarrow$ 生成链可作为统一索引与复核入口）]
\label{cor:tex29-books-cover-chain-index}
若定理~\ref{thm:tex29-books-cover-chain-induction} 成立并且引理~\ref{lem:tex29-books-kqemb-term-vs-kernel} 给出“术语缺省但内核完备”，
则生成链可被用作对两部专著的统一索引与复核入口：
\[
  \text{任取链上任一环 }L_k\ \Rightarrow\ \text{存在对应的定义锚定、算法接口与定理保障条目可回溯},
\]
从而“元数学$\to$数学工具$\to$数学物理$\to$物理落地”在同一索引下闭环。
\end{corollary}

\begin{Certificate}[证书：回溯性 + 闭环]
\label{cert:tex29-books-cover-chain-index}
\[
\begin{aligned}
  I_1 &: (\forall k)\ \mathsf{Cover}_{\mathcal{B}}(L_k)\Rightarrow \text{每环均具备可回溯条目},\\
  I_2 &: \text{跨层映射规则证书化}\Rightarrow \text{环与环之间可在同一索引链上复核},\\
  I_3 &: I_1\wedge I_2\Rightarrow \text{生成链成为统一索引入口并形成闭环}.
\end{aligned}
\]
\end{Certificate}

\paragraph{证明（谓词因果链）.}
由 $I_1$ 得逐环回溯性；由 $I_2$ 得跨环复核性；合并得 $I_3$。证毕.

\subsection{备注}
\label{subsec:tex29-books-cover-chain-remarks}

\begin{remark}[表格化索引：专著锚定条目与五环生成链的对照]
表~\ref{tab:tex29-books-cover-chain-map} 给出“生成链环节 $\leftrightarrow$ 专著锚定条目（命名锚点）$\leftrightarrow$ 本笔记证书接口”的压缩索引。
其中“命名锚点”仅用于对齐术语与索引位置，不把外部文本当作本\KouJing的前提。
\end{remark}

\begin{table}[htbp]
\centering
{\footnotesize
\setlength{\tabcolsep}{2pt}
\renewcommand{\arraystretch}{1.25}
\begin{tabular}{|p{0.16\linewidth}|p{0.26\linewidth}|p{0.26\linewidth}|p{0.26\linewidth}|}
\hline
\textbf{链环节} & \textbf{专著锚定（命名锚点）} & \textbf{本笔记形式入口} & \textbf{可复核输出} \\
\hline
$L_1:\mathfrak{D}$ &
PL-PI 元规则；GZ-G-代数（主丛版广义非交换李代数） &
$\mathsf{Rep}(\omega)$；$\partial D=\omega$ 证据链 &
离散本原的结构自洽与可执行操作语义 \\
\hline
$L_2:\mathfrak{K}$ &
非交换结构化投影；投影核携带（对应“主丛胚”原型） &
$\Pi_{\mathrm{nc}}$；$\mathsf{Ker}_{\mathrm{KQE}}$ &
离散到连续的临界预几何胚与核信息保留 \\
\hline
$L_3:\mathfrak{K}$ 作为修复元 &
GZ-OHU 无失败保障（修复元功能条目化） &
$\mathrm{JacGap}_{\alpha}$；$\mathsf{Stable}_{\alpha}$ &
奇点/崩溃信号对象化并转为修复输入 \\
\hline
$L_4:\{\mathfrak{T}_{\alpha}\}$ &
超限同伦完备化；GZ-CBT 重写；JacGap 证书监控 &
$D_{\alpha}$ 链；$\mathsf{Stable}/\mathsf{Collapse}$ &
逐层消解障碍并在极限层封顶一致性 \\
\hline
$L_5:(\mathcal{M}^4,g,\nabla)$ &
定律空间几何投影；$\mathcal{A}_{\mathrm{law}},\mathcal{F}_{\mathrm{law}}$ 的低能表观 &
$\mathcal{P}_{\mathrm{4D}}$；$g,R$ 的投影表达 &
四维时空作为低能投影表观而非预设背景 \\
\hline
\end{tabular}
}
\caption{\texorpdfstring{两部专著对五环生成链的“覆盖锚定”压缩索引表}{两部专著对五环生成链的覆盖锚定压缩索引表}}
\label{tab:tex29-books-cover-chain-map}
\end{table}

\begin{remark}[避免把叙述性价值判断写成定理]
“体系化/顶级/学术价值”等表述属于叙述性评价，不纳入本节定理体系；
本节形式化目标仅是把“覆盖性、锚定性、可回溯性、术语缺省但内核完备”等主张压缩为可核验谓词与证书链。
\end{remark}

\clearpage

%%%%%%%%%%%%%%%%%%%%%%%%%%%%%%%%%%%%%%%%%%%%%%%%%%%%%%%%%%%%
\section{参考文献}
\label{sec:tex29-bibliography}
\label{sec:tex29-ch20}
%%%%%%%%%%%%%%%%%%%%%%%%%%%%%%%%%%%%%%%%%%%%%%%%%%%%%%%%%%%%

\begin{thebibliography}{99}

\bibitem{GaoZheng2025GFramework}
Gao, Z. (2025).
\newblock Meta-Mathematical Theory based on Pan-Logic Analysis and
Pan-Iterative Analysis (GaoZheng G-Framework) and the
Principal-Bundle-Based Generalized Noncommutative Lie Algebra
(GaoZheng G-Algebra).
\newblock In \emph{Meta-Mathematical Theory based on Pan-Logic Analysis and
Pan-Iterative Analysis (GaoZheng G-Framework) and the
Principal-Bundle-Based Generalized Noncommutative Lie Algebra
(GaoZheng G-Algebra): An Integrated Construction (v1.0)}.
\newblock Zenodo.
\newblock \url{https://doi.org/10.5281/zenodo.17651584}.


\bibitem{GaoZhengAppVol1}
Gao, Z. (2025).
\newblock GaoZheng G-Framework and GaoZheng G-Algebra: Applied Mathematics Volume I — Law-Space Geometry, GRL, Quantum
  Computing, and Superconductivity.
\newblock In \emph{GaoZheng G-Framework and GaoZheng G-Algebra: Applied Mathematics Volume I — Law-Space Geometry, GRL,
  Quantum Computing, and Superconductivity (v1.2)}.
\newblock Zenodo.
\newblock \url{https://doi.org/10.5281/zenodo.17686133}.


\bibitem{RepoNB17EdgeOperator}
【仓内 TeX】边缘算子闭合、同伦闭包与变分修复塔（形式系统笔记）：
\code{NoteBook/notebook1/tex17_pub/gframework_edge_operator_homotopy_closure_variational_tower_formal_system.tex}。

\bibitem{Hatcher2002AT}
Hatcher, A. (2002).
\newblock \emph{Algebraic Topology}.
\newblock Cambridge University Press.

\bibitem{Switzer1975HomotopyHomology}
Switzer, R. M. (1975).
\newblock \emph{Algebraic Topology --- Homotopy and Homology}.
\newblock Springer.

\bibitem{Whitehead1978Elements}
Whitehead, G. W. (1978).
\newblock \emph{Elements of Homotopy Theory}.
\newblock Springer.

\bibitem{May1999Concise}
May, J. P. (1999).
\newblock \emph{A Concise Course in Algebraic Topology}.
\newblock University of Chicago Press.

\bibitem{Weibel1994HomologicalAlgebra}
Weibel, C. A. (1994).
\newblock \emph{An Introduction to Homological Algebra}.
\newblock Cambridge University Press.

\bibitem{Lee2018RiemannianManifolds}
Lee, J. M. (2018).
\newblock \emph{Riemannian Manifolds: An Introduction to Curvature} (2nd ed.).
\newblock Springer.

\bibitem{Petersen2016RiemannianGeometry}
Petersen, P. (2016).
\newblock \emph{Riemannian Geometry} (3rd ed.).
\newblock Springer.

\bibitem{doCarmo1992Riemannian}
do Carmo, M. (1992).
\newblock \emph{Riemannian Geometry}.
\newblock Birkh\"{a}user.

\bibitem{Husemoller1994FibreBundles}
Husemoller, D. (1994).
\newblock \emph{Fibre Bundles} (3rd ed.).
\newblock Springer.

\bibitem{Hall2015LieGroups}
Hall, B. C. (2015).
\newblock \emph{Lie Groups, Lie Algebras, and Representations: An Elementary Introduction} (2nd ed.).
\newblock Springer.

\bibitem{Jech2003SetTheory}
Jech, T. (2003).
\newblock \emph{Set Theory: The Third Millennium Edition, Revised and Expanded}.
\newblock Springer.

\bibitem{Kunen2011SetTheory}
Kunen, K. (2011).
\newblock \emph{Set Theory}.
\newblock College Publications.

\end{thebibliography}

\end{CJK*}
\end{document}
